\documentclass[11pt,a4paper]{article}
\usepackage[margin=2.4cm]{geometry}
\usepackage{amsmath,amsthm,mathtools}
\usepackage{fontspec}
\usepackage{unicode-math}
\setmainfont{texgyrepagella}[Path=fonts/, Extension=.otf,
  UprightFont=*-regular, ItalicFont=*-italic,
  BoldFont=*-bold, BoldItalicFont=*-bolditalic]
\setmathfont{texgyrepagella-math}[Path=fonts/, Extension=.otf]
% Pagella uses U+2216 for set difference.
\AtBeginDocument{\let\setminus\smallsetminus}
\usepackage[protrusion=true,expansion=false]{microtype}
\usepackage{booktabs,enumitem}
\usepackage{bm}
\usepackage{xcolor}
\usepackage[psdextra,hidelinks]{hyperref}
% Leave room for two-digit subsection numbers such as 10.10 in the contents.
\makeatletter
\renewcommand*{\l@subsection}{\@dottedtocline{2}{1.5em}{3em}}
\makeatother
\emergencystretch=3em
\hbadness=10000
\hfuzz=2pt
\newtheorem{theorem}{Theorem}[section]
\newtheorem{lemma}[theorem]{Lemma}
\newtheorem{proposition}[theorem]{Proposition}
\newtheorem{corollary}[theorem]{Corollary}
\theoremstyle{definition}
\newtheorem{definition}[theorem]{Definition}
\newtheorem{construction}[theorem]{Construction}
\newtheorem{example}[theorem]{Example}

% kind-coded theorem blocks, palette shared with the figures and the website:
% definitions/constructions blue (tierin), assumptions amber (tierfull),
% proved statements a neutral bar; remarks and examples run unframed
\usepackage{mdframed}
\providecolor{tierin}{RGB}{60,130,180}
\providecolor{tierfull}{RGB}{210,140,40}
\mdfdefinestyle{kindbar}{topline=false,bottomline=false,rightline=false,
  linewidth=1.5pt,innerleftmargin=9pt,innerrightmargin=4pt,
  innertopmargin=5pt,innerbottommargin=6pt,leftmargin=0pt,rightmargin=0pt,
  skipabove=10pt,skipbelow=10pt}
\surroundwithmdframed[style=kindbar,linecolor=tierin!70,backgroundcolor=tierin!4]{definition}
\surroundwithmdframed[style=kindbar,linecolor=tierin!70,backgroundcolor=tierin!4]{construction}
\surroundwithmdframed[style=kindbar,linecolor=tierfull!80,backgroundcolor=tierfull!5]{assumption}
\surroundwithmdframed[style=kindbar,linecolor=black!30]{theorem}
\surroundwithmdframed[style=kindbar,linecolor=black!30]{lemma}
\surroundwithmdframed[style=kindbar,linecolor=black!30]{proposition}
\surroundwithmdframed[style=kindbar,linecolor=black!30]{corollary}

\usepackage{fancyhdr}
\pagestyle{fancy}
\fancyhf{}
\fancyhead[L]{\footnotesize\scshape The Zcash Arboretum}
\fancyhead[R]{\footnotesize\itshape\nouppercase{\leftmark}}
\fancyfoot[C]{\footnotesize\thepage}
\renewcommand{\headrulewidth}{0.3pt}
\renewcommand{\sectionmark}[1]{\markboth{\thesection\ \ #1}{}}
\theoremstyle{remark}
\newtheorem{remark}[theorem]{Remark}
\newcommand{\N}{\ensuremath{\mathbb{N}}}
\newcommand{\Z}{\ensuremath{\mathbb{Z}}}
\newcommand{\Q}{\ensuremath{\mathbb{Q}}}
\newcommand{\R}{\ensuremath{\mathbb{R}}}
\providecommand{\C}{}\renewcommand{\C}{\ensuremath{\mathbb{C}}}
\newcommand{\F}{\ensuremath{\mathbb{F}}}
\newcommand{\G}{\ensuremath{\mathbb{G}}}
\newcommand{\Fp}{\ensuremath{\mathbb{F}_{p}}}
\newcommand{\Fq}{\ensuremath{\mathbb{F}_{q}}}
\newcommand{\E}{\ensuremath{\mathbb{E}}}
\newcommand{\ip}[2]{\langle #1,\, #2 \rangle}
\DeclareMathOperator{\ord}{ord}
\DeclareMathOperator{\lcm}{lcm}
\DeclareMathOperator{\negl}{negl}
\DeclareMathOperator{\poly}{poly}
\DeclareMathOperator{\Hom}{Hom}
\DeclareMathOperator{\img}{im}
\DeclareMathOperator{\Tr}{Tr}
\newcommand{\legendre}[2]{\left(\frac{#1}{#2}\right)}
\DeclareMathOperator{\sgn}{sgn}
\DeclareMathOperator{\rank}{rank}
\DeclareMathOperator{\End}{End}
\DeclareMathOperator{\Aut}{Aut}
\DeclareMathOperator{\id}{id}
\DeclareMathOperator{\Frob}{Frob}
\DeclareMathOperator{\disc}{disc}
\DeclareMathOperator{\charf}{char}
\providecommand{\abs}[1]{\lvert #1 \rvert}
\providecommand{\norm}[1]{\lVert #1 \rVert}
\providecommand{\ceil}[1]{\lceil #1 \rceil}
\providecommand{\floor}[1]{\lfloor #1 \rfloor}
\providecommand{\indicator}{\mathbf{1}}
\providecommand{\mathbbm}[1]{\mathbf{#1}}
\newcommand{\Fr}{\ensuremath{\mathbb{F}_{r}}}
\newcommand{\Adv}{\ensuremath{\mathcal{A}}}
\usepackage[nameinlink]{cleveref}
\renewenvironment{abstract}{%
  \small\begin{center}{\bfseries\abstractname}\end{center}\medskip\noindent\ignorespaces
}{\par\medskip}

\title{\textbf{\Huge Crypto Guide}\\[6pt]\large Cryptographic primitives from scratch}
\author{m@rek.onl}
\date{}
\begin{document}
\maketitle

\begin{abstract}
This volume of \emph{The Zcash Arboretum} assembles the cryptographic
toolbox: the primitives from which a shielded protocol is built, each
defined, constructed, and proved against a named enemy. It assumes the
\emph{Math Guide}, whose groups, finite fields, elliptic curves,
probability, and model of computation are used throughout without
re-derivation.

Every guarantee here is a claim about an adversary, and every security
definition is a game: a challenger, an efficient adversary, and an
advantage that must be negligible. A proof is a reduction that turns
any winning adversary into a solver for a long-attacked computational
problem, with the reduction's loss accounted for down to a concrete
bit-security budget. Under that discipline the volume denies, in turn:
the cryptanalyst who would recover the exponent behind a public group
element, and with it every key; the forger by coincidence, who would
collide two documents into one digest or replay a digest outside its
role; the peeker and the equivocator on either side of a sealed
envelope, who would read a commitment early or open it to a different
value; the oracle distinguisher, who would tell a short secret key
from blind chance; the owner of the channel, who would read, splice,
or counterfeit what passes over it; the woman in the middle, who would
terminate both ends of a ``secure'' conversation at her own desk; the
forger of signatures, who would exhibit an authorisation the signer
never made; the cheating prover and the curious verifier, who would
certify a falsehood or squeeze a transcript for the witness it hides;
and the teller of lies about sets, who would prove membership in a
collection that never recorded it. What survives their attentions---
hash functions, commitments, pseudorandom functions, authenticated
encryption, key agreement, signatures, zero-knowledge arguments, and
Merkle trees---is the toolbox the deployed protocol spends.
\end{abstract}

\clearpage
\tableofcontents
\clearpage

\section{The provable-security model}\label{sec:model}

Every claim this volume makes ends, once unwound, in a sentence about
an enemy. No forger produces a valid signature on a message the signer
never signed; no observer tells an encryption of \(m_0\) from an
encryption of \(m_1\); no prover opens one commitment to two different
values. Before any such sentence can be a theorem, the enemy must be a
mathematical object: something with a precisely delimited repertoire,
about which a universally quantified statement can be proved. The
preceding volume supplied the algebraic raw material---groups in which
discrete logarithms appear hard, finite fields, elliptic curves---but
an object is not yet a guarantee. This section builds the adversary,
and around it the framework in which ``the scheme is secure'' becomes
a claim with mathematical content.

The central methodological idea is the following. One almost never
knows how to prove unconditionally that a cryptographic scheme is
secure; such a proof would typically imply
\(\mathsf{P} \neq \mathsf{NP}\)---that problems whose solutions can be
checked in polynomial time are not all solvable in polynomial
time---and a great deal more. Cryptography therefore adopts a
\emph{relativised} standard of proof. One isolates a small number of
computational problems---factoring, discrete logarithm, decisional
Diffie--Hellman, and the like---that have resisted decades of attack,
and \emph{reduces} the security of the schemes to the conjectured
hardness of those problems. A proof of
security is then a proof of an implication: \emph{if} the underlying
problem is hard, \emph{then} the scheme is secure. Equivalently, in
its contrapositive and more operationally useful form: \emph{any}
adversary that breaks the scheme can be mechanically transformed into
an algorithm that solves the underlying hard problem. Since no such
algorithm is believed to exist, one concludes that no such adversary
exists.

To render any of this rigorous, we must pin down four things: what
counts as a ``scheme'', parametrised by a notion of input size; what
counts as an ``adversary''; what it means for an adversary to
``break'' a scheme, measured by a quantity called its
\emph{advantage}; and what it means for one problem to reduce to
another. We develop these in turn, followed by the two idealisations
under which proofs proceed (the standard model and the random oracle
model), and finally the translation between the clean asymptotic
language and the concrete ``\(n\)-bit security'' numbers that
practitioners actually quote.

\subsection{Adversaries and the security parameter}
\label{model:subsec:adversaries}

The vocabulary of efficiency and smallness is the \emph{Math Guide}'s,
fixed in its closing section; we recall it here in the form this
volume uses and extend it with the refinements that are specifically
cryptographic.

The \emph{security parameter} \(\lambda \in \N\) is a positive integer
that controls the sizes of all objects in a cryptographic scheme---key
lengths, group orders, output lengths---and is conventionally supplied
to every algorithm in unary, written \(1^\lambda\), the string of
\(\lambda\) ones, so that an algorithm running in time polynomial in
the length of its input automatically receives time polynomial in
\(\lambda\) (\emph{Math Guide}, \S``The security parameter''). Every
scheme in this monograph is accordingly not a single object but an
infinite \emph{family} of objects indexed by \(\lambda\).

Two questions immediately arise: why a family, and why unary.

The reason for a \emph{family} is that security is inherently
asymptotic. There is no such thing as an absolutely unbreakable scheme
with fixed finite key length: an adversary with enough time can always
enumerate a finite key space. What we can hope for is that as
\(\lambda\) grows, the resources required to break the scheme grow far
faster than the resources required to use it. The honest parties'
costs (key generation, encryption, signing) should grow \emph{slowly}
in \(\lambda\)---polynomially---while the best attack's cost should
grow \emph{quickly}---super-polynomially, ideally exponentially. The
security parameter is the dial that separates these two growth rates,
and security is a statement about the limit \(\lambda \to \infty\).

The reason for \emph{unary} encoding, \(1^\lambda\) rather than the
\(\lceil \log_2 \lambda \rceil\)-bit binary representation of
\(\lambda\), is a bookkeeping convenience that aligns two notions of
``polynomial time''. Complexity theory measures running time as a
function of input \emph{length}. If \(\lambda\) came in binary it
would have length \(\Theta(\log \lambda)\), and an algorithm
``polynomial in its input length'' would get only
\(\poly(\log \lambda)\) steps---not even enough to write down a
\(\lambda\)-bit key. Padding the parameter to length \(\lambda\)
ensures that ``polynomial in the input length'' coincides with the
cryptographically meaningful ``polynomial in \(\lambda\)''. This is
purely a normalisation; nothing of substance depends on it.

\begin{remark}\label{model:rem:everything-asymptotic}
Because everything is indexed by \(\lambda\), every quantity we
discuss is really a \emph{function} of \(\lambda\): the adversary's
success probability, its running time, the key length. A single
number---``the advantage is at most \(2^{-128}\)''---either fixes a
concrete \(\lambda\) or describes the function's behaviour for the
\(\lambda\) of interest. Keeping the dependence on \(\lambda\)
explicit, at least mentally, prevents nearly every confusion in this
subject.
\end{remark}

Having fixed how scheme sizes scale, we fix what an attacker may do.
An algorithm \(\mathcal{A}\) is \emph{probabilistic polynomial-time}
(PPT) if it is a probabilistic Turing machine---one equipped with a
read-only tape of independent uniformly random bits, its \emph{random
tape}---and there is a polynomial \(p\) such that on every input
\(x\), \emph{and for every setting of the random tape}, \(\mathcal{A}\)
halts within \(p(|x|)\) steps: the strict, worst-case notion, exactly
as fixed in the \emph{Math Guide} (\S``Algorithms, running time, and
PPT''). Write \(y \gets \mathcal{A}(x)\) for the random variable
obtained by running \(\mathcal{A}\) on \(x\) with freshly sampled
randomness, and \(y := \mathcal{A}(x; r)\) for the deterministic
output when the random tape is fixed to the string \(r\). The honest
algorithms and the adversary alike are modelled as PPT machines.

Two of the modelling choices folded into that sentence warrant
comment.

\emph{Why polynomial time?} Polynomial time is the standard
mathematical abstraction of ``feasible computation''. It is
robust---closed under composition, and invariant across all
reasonable deterministic computational models (the strong
Church--Turing thesis)---which makes it a clean class to quantify
over. Crucially, we require the \emph{honest} parties to run in
polynomial time too; a scheme is useful only if it can be used
efficiently. The security claim is then an asymmetry between a
polynomial-time honest world and a quantified statement that \emph{no}
polynomial-time attack succeeds.

\emph{Why probabilistic?} Randomness is not a luxury but a necessity
in cryptography. Deterministic encryption leaks whether two
ciphertexts encrypt the same message; key generation must sample
secrets unpredictably. We therefore grant the adversary randomness
too, both for fairness and because a deterministic adversary is the
special case where the random tape is ignored.

The adversary class itself admits a strengthening that the \emph{Math
Guide} mentions (\S``Algorithms, running time, and PPT'') and this
volume adopts as standard.

\begin{definition}[Non-uniform PPT]\label{model:def:nonuniform}
A \emph{non-uniform} PPT adversary is a family
\(\{\mathcal{A}_\lambda\}_{\lambda\in\N}\) of probabilistic circuits
(or, equivalently, Turing machines each supplied with an \emph{advice
string} \(z_\lambda\) of length \(\poly(\lambda)\)) such that a single
polynomial in \(\lambda\) bounds the size of \(\mathcal{A}_\lambda\).
The advice may depend arbitrarily on \(\lambda\) but not on the
particular inputs drawn at run time.
\end{definition}

\begin{remark}\label{model:rem:uniform-vs-nonuniform}
The distinction between uniform adversaries (a single machine that
reads \(1^\lambda\) and works for all \(\lambda\)) and non-uniform
ones (a separate circuit per \(\lambda\), or a uniform machine with
per-\(\lambda\) advice) is a genuine one. Non-uniformity grants the
adversary free precomputed advice tailored to each parameter size---
for example, a short hint about the group of the day. Security against
non-uniform adversaries is the stronger and now-standard requirement,
partly because non-uniform models compose more cleanly under reduction
(one may ``hard-wire'' a fixed good choice of randomness or auxiliary
input). All definitions below can be read in either flavour; we state
hardness assumptions against non-uniform adversaries when it matters
and otherwise leave the choice implicit.
\end{remark}

\begin{remark}[Expected versus strict polynomial time, and quantum
adversaries]\label{model:rem:variants}
Two refinements recur. First, some definitions allow \emph{expected}
polynomial-time adversaries, bounding the average rather than
worst-case running time; this matters for certain extraction arguments
in zero-knowledge, and we flag it where it arises. Second, for
post-quantum security one replaces PPT by \emph{quantum}
polynomial-time (QPT)---polynomial-size circuits of quantum gates
followed by a measurement, a model this volume does not formalise;
the consequences for the assumptions this volume relies on are taken
up under ``The quantum caveat: Shor's algorithm'' in the
next section. The structural definitions of this section (advantage,
indistinguishability, reduction) are agnostic to which class of
adversaries one quantifies over; only the underlying hardness
assumptions change.
\end{remark}

\subsection{Distribution ensembles and indistinguishability}
\label{model:subsec:indist}

The most basic thing an adversary can attempt is to \emph{tell two
situations apart}---a real protocol transcript from a simulated one,
an encryption of \(m_0\) from an encryption of \(m_1\), a pseudorandom
string from a truly random one. The most basic security goal is that
it cannot. To formalise ``look the same'' we first formalise the two
situations as \emph{ensembles} of probability distributions, then
define three grades of sameness.

\begin{definition}[Probability ensemble]\label{model:def:ensemble}
A \emph{probability ensemble} is a sequence
\(X = \{X_\lambda\}_{\lambda\in\N}\) of probability distributions
(equivalently, random variables), one for each value of the security
parameter, where \(X_\lambda\) is supported on binary strings whose
length is bounded by a polynomial in \(\lambda\). There is often an
additional index---an input \(a\)---written \(\{X_\lambda(a)\}\); for
clarity we suppress it unless needed.
\end{definition}

The natural metric on a single pair of distributions is the
statistical distance of the \emph{Math Guide} (\S``Statistical
distance''), which we recall in the slightly wider generality this
volume needs: supports there are finite, while ensembles as just
defined live on countable sets of strings.

\begin{definition}[Statistical distance]\label{model:def:statdist}
Let \(X, Y\) be discrete random variables over a common finite or
countable set \(S\). Their \emph{statistical distance} (or total
variation distance) is
\[
  \Delta(X, Y) \;:=\; \frac{1}{2} \sum_{s \in S}
    \bigl| \Pr[X = s] - \Pr[Y = s] \bigr| .
\]
\end{definition}

\begin{proposition}[Variational characterisation]\label{model:prop:statdist}
For \(X, Y\) over \(S\),
\[
  \Delta(X, Y) \;=\; \max_{T \subseteq S}
    \bigl( \Pr[X \in T] - \Pr[Y \in T] \bigr)
  \;=\; \max_{D} \bigl( \Pr[D(X) = 1] - \Pr[D(Y) = 1] \bigr),
\]
where the last maximum ranges over all (even computationally
unbounded, deterministic) decision functions
\(D : S \to \{0,1\}\). In particular, no procedure whatsoever---
efficient or not---can distinguish \(X\) from \(Y\) with advantage
exceeding \(\Delta(X,Y)\).
\end{proposition}

\begin{proof}
The set form is the variational characterisation in the
statistical-distance theorem of the \emph{Math Guide}
(\S``Statistical distance''), whose argument is unchanged over a
countable \(S\). Only the reformulation over decision functions is
new: a deterministic \(D\) is the indicator of the set
\(T = D^{-1}(1)\), so the maximum over \(D\) equals the maximum over
\(T\); a randomised \(D\) is a convex combination of deterministic
ones and cannot do better.
\end{proof}

\begin{proposition}[Properties of statistical distance]
\label{model:prop:statdist-props}
The statistical distance is a metric on probability distributions:
\(0 \leq \Delta(X,Y) \leq 1\), it is symmetric, \(\Delta(X,Y)=0\) iff
\(X\) and \(Y\) are identically distributed, and it satisfies the
triangle inequality
\(\Delta(X,Z) \leq \Delta(X,Y)+\Delta(Y,Z)\). Moreover it never
increases under (possibly randomised) post-processing: for any
function---indeed any randomised algorithm---\(f\),
\[
  \Delta\bigl(f(X), f(Y)\bigr) \;\leq\; \Delta(X, Y).
\]
\end{proposition}

\begin{proof}
These are the metric and data-processing parts of the
statistical-distance theorem of the \emph{Math Guide}
(\S``Statistical distance''), again with the argument unchanged over
a countable \(S\).
\end{proof}

The weakest of the three grades below needs the asymptotic vocabulary
of the \emph{Math Guide} (\S``Polynomial, exponential, and negligible
functions''), which we recall in one breath. A function
\(f \colon \N \to \R_{\ge 0}\) is \emph{negligible}, written
\(f = \negl(\lambda)\), if for every \(c \in \N\) there exists
\(\lambda_0\) such that \(f(\lambda) < \lambda^{-c}\) for all
\(\lambda \ge \lambda_0\)---it decays faster than the reciprocal of
every polynomial. A function is \emph{non-negligible} if it is not
negligible; it is \emph{noticeable} if there exists \(c\) with
\(f(\lambda) \ge \lambda^{-c}\) for all sufficiently large \(\lambda\)
(noticeable functions are non-negligible, but not conversely); it is
\emph{overwhelming} if \(1 - f\) is negligible. The class of
negligible functions is closed under addition and under multiplication
by any polynomial---the closure property, proved there, on which the
hybrid argument of the next subsection depends.

\begin{definition}[Perfect, statistical, and computational
indistinguishability]\label{model:def:indist}
Let \(X = \{X_\lambda\}\) and \(Y = \{Y_\lambda\}\) be probability
ensembles.
\begin{itemize}
\item The ensembles \(X\) and \(Y\) are \emph{perfectly
  indistinguishable}, written
  \(X \equiv Y\), if \(X_\lambda\) and \(Y_\lambda\) are
  \emph{identically distributed} for every \(\lambda\); equivalently
  \(\Delta(X_\lambda, Y_\lambda) = 0\) for all \(\lambda\).
\item The ensembles \(X\) and \(Y\) are \emph{statistically
  indistinguishable}, written \(X \approx_s Y\), if their
  statistical distance is negligible:
  \[
    \Delta(X_\lambda, Y_\lambda) \;=\; \negl(\lambda).
  \]
\item The ensembles \(X\) and \(Y\) are \emph{computationally
  indistinguishable}, written \(X \approx_c Y\), if for every PPT
  algorithm \(D\) (the \emph{distinguisher}), the function
  \[
    \mathrm{Adv}^{\mathrm{dist}}_{X,Y}(D, \lambda)
    \;:=\; \Bigl| \Pr\bigl[D(1^\lambda, X_\lambda) = 1\bigr]
                 - \Pr\bigl[D(1^\lambda, Y_\lambda) = 1\bigr] \Bigr|
  \]
  is negligible in \(\lambda\). The probabilities are over the draw
  from the ensemble and over \(D\)'s internal randomness.
\end{itemize}
\end{definition}

These three notions sit in a strict hierarchy. The implications are
immediate from the definitions and the variational characterisation;
standard examples witness the strictness. The second example needs
two primitives, which we define here in the form the volume uses
throughout; write \(U_n\) for the uniform distribution on
\(\{0,1\}^n\). A \emph{one-way function} is a polynomial-time
computable \(h \colon \{0,1\}^* \to \{0,1\}^*\) that no PPT algorithm
inverts on a random input: for every PPT \(\mathcal{A}\),
\[
  \Pr\bigl[h(\mathcal{A}(1^\lambda, h(x))) = h(x)\bigr]
  \;=\; \negl(\lambda)
  \qquad (x \gets U_\lambda).
\]
A \emph{pseudorandom generator} (PRG) is a polynomial-time computable
\(G \colon \{0,1\}^\lambda \to \{0,1\}^{2\lambda}\) whose output on a
uniform seed is computationally indistinguishable from uniform,
\(\{G(U_\lambda)\} \approx_c \{U_{2\lambda}\}\); the doubling stretch
is a normalisation. Whether one-way functions exist is open, which is
why the second strictness below is stated conditionally.

\begin{proposition}[Hierarchy]\label{model:prop:hierarchy}
For all ensembles \(X, Y\)\textup{:}
\[
  X \equiv Y \;\Longrightarrow\; X \approx_s Y
  \;\Longrightarrow\; X \approx_c Y .
\]
The first implication is strict: there exist ensembles that are
statistically but not perfectly indistinguishable. The second is
strict provided one-way functions exist: under that assumption there
exist ensembles that are computationally but not statistically
indistinguishable.
\end{proposition}

\begin{proof}
If \(X \equiv Y\) then \(\Delta(X_\lambda,Y_\lambda)=0\), which is
negligible, so \(X \approx_s Y\). If \(X \approx_s Y\), then for
any---in particular any PPT---distinguisher \(D\),
Proposition~\ref{model:prop:statdist} gives
\(\mathrm{Adv}^{\mathrm{dist}}_{X,Y}(D,\lambda) \leq
\Delta(X_\lambda,Y_\lambda) = \negl(\lambda)\), so
\(X \approx_c Y\).

For strictness of the first implication, let \(X_\lambda\) be uniform
on \(\{0,1\}^\lambda\) and let \(Y_\lambda\) be uniform on
\(\{0,1\}^\lambda \setminus \{0^\lambda\}\). They are not identically
distributed (the point \(0^\lambda\) has probabilities
\(2^{-\lambda}\) and \(0\)), so not perfectly indistinguishable, yet
\(\Delta(X_\lambda,Y_\lambda) = 2^{-\lambda} = \negl(\lambda)\), so
they are statistically indistinguishable.

For strictness of the second, assume one-way functions exist; then a
PRG exists---the H\aa stad--Impagliazzo--Levin--Luby theorem,
together with the routine lemma that extends any stretch to the
doubling one, both used without proof---and furnishes an ensemble:
let \(G\) be a PRG, let \(X_\lambda := G(U_\lambda)\), and let
\(Y_\lambda := U_{2\lambda}\). By definition of a PRG,
\(X \approx_c Y\). But \(X_\lambda\) is supported on at most
\(2^\lambda\) strings out of \(2^{2\lambda}\), so a set \(T\) equal to
the image of \(G\) has \(\Pr[X_\lambda \in T] = 1\) while
\(\Pr[Y_\lambda \in T] \leq 2^\lambda/2^{2\lambda} = 2^{-\lambda}\);
hence \(\Delta(X_\lambda,Y_\lambda) \geq 1 - 2^{-\lambda}\), which is
not negligible. Thus \(X \approx_c Y\) but \(X \not\approx_s Y\).
\end{proof}

\begin{remark}[The price and the prize of computational
indistinguishability]\label{model:rem:comp-price}
The strictness example is instructive. Statistically the PRG output
and a random string are almost as far apart as two distributions can
be---distance nearly \(1\)---because the PRG image is a vanishing
fraction of the codomain. An unbounded distinguisher simply checks
membership in the image and wins. The entire content of computational
indistinguishability is that \emph{finding} that image membership test
is infeasible. This is the bargain at the heart of cryptography: we
trade the impossible (information-theoretic security with short keys)
for the merely conjectured-hard (computational security), and in
exchange we obtain the ability to encrypt long messages under short
keys, build public-key primitives, and so on. Statistical and perfect
security, when attainable, are unconditional and need no hardness
assumption; computational security buys vastly more functionality at
the cost of resting on an assumption.
\end{remark}

\subsection{The hybrid argument}\label{model:subsec:hybrid}

Indistinguishability would be of little use if every application had
to be argued from scratch. Two lemmas make it composable: efficient
post-processing cannot create a distinguishing edge, and a
polynomially long chain of pairwise-indistinguishable steps cannot
accumulate one. The first is a reduction in miniature---the template
for every reduction in this volume; the second, the hybrid argument,
is the workhorse of the entire proof literature.

\begin{lemma}[Closure under efficient post-processing]
\label{model:lem:closure-postproc}
If \(X \approx_c Y\) and \(M\) is any PPT algorithm, then the
ensembles \(\{M(1^\lambda, X_\lambda)\}\) and
\(\{M(1^\lambda, Y_\lambda)\}\) satisfy
\(\{M(X_\lambda)\} \approx_c \{M(Y_\lambda)\}\).
\end{lemma}

\begin{proof}
Suppose some PPT distinguisher \(D\) separates
\(\{M(X_\lambda)\}\) from \(\{M(Y_\lambda)\}\) with advantage
\(\delta(\lambda)\). Define \(D' := D \circ M\), i.e.\
\(D'(1^\lambda, z) := D(1^\lambda, M(1^\lambda, z))\). Since \(M\)
and \(D\) are both PPT, so is \(D'\). Then
\[
  \mathrm{Adv}^{\mathrm{dist}}_{X,Y}(D',\lambda)
  = \bigl|\Pr[D(M(X_\lambda))=1] - \Pr[D(M(Y_\lambda))=1]\bigr|
  = \delta(\lambda).
\]
By \(X \approx_c Y\) the left side is negligible, so \(\delta\) is
negligible.
\end{proof}

Many proofs must compare two ensembles that are far apart but
connected by a chain of pairwise-close intermediate ensembles. The
hybrid argument turns such a chain into a single conclusion. It is
essentially the triangle inequality, but with a refinement---the
uniformity hypothesis---that explains why the number of hybrids must
be polynomial and that careless statements omit.

\begin{lemma}[Hybrid lemma]\label{model:lem:hybrid}
Let \(t = t(\lambda)\) be polynomially bounded, and let
\(H^0_\lambda, H^1_\lambda, \dots, H^{t}_\lambda\) be a sequence of
distributions (the \emph{hybrids}). Suppose that for every PPT
distinguisher \(D\) the neighbouring advantages
\[
  \alpha_i(D,\lambda) :=
   \bigl|\Pr[D(H^{i}_\lambda)=1]-\Pr[D(H^{i+1}_\lambda)=1]\bigr|
\]
are \emph{uniformly} negligible: a single negligible function \(\nu\)
(depending on \(D\)) satisfies
\(\alpha_i(D,\lambda) \le \nu(\lambda)\) for all
\(i \in \{0,\dots,t(\lambda)-1\}\) simultaneously. Equivalently,
\(\max_{0\le i < t(\lambda)} \alpha_i(D,\lambda) = \negl(\lambda)\),
i.e.\ \(\alpha_{i(\lambda)}(D,\lambda)\) is negligible along
\emph{every} index function \(i(\lambda)\)---arbitrary, not merely
efficiently computable, since the maximising index need not be
efficient. Then \(\{H^0_\lambda\}\) and \(\{H^t_\lambda\}\) are
computationally indistinguishable, i.e.\ for every PPT \(D\),
\(\bigl|\Pr[D(H^0_\lambda)=1]-\Pr[D(H^t_\lambda)=1]\bigr| =
\negl(\lambda)\).
\end{lemma}

\begin{proof}
Fix a PPT distinguisher \(D\) and let \(\nu\) be the uniform bound the
hypothesis supplies for \(D\). By the triangle inequality applied to
the telescoping sum,
\[
  \bigl|\Pr[D(H^0_\lambda)=1]-\Pr[D(H^t_\lambda)=1]\bigr|
  = \Bigl| \sum_{i=0}^{t-1}\bigl(\Pr[D(H^i_\lambda)=1]
    -\Pr[D(H^{i+1}_\lambda)=1]\bigr)\Bigr|
  \leq \sum_{i=0}^{t-1} \alpha_i(D,\lambda),
\]
and the hypothesis bounds the sum by
\(t(\lambda)\cdot\nu(\lambda)\), which is negligible since \(t\) is a
polynomial (closure of the negligible class under polynomial
multiplication). Hence \(\{H^0\}\approx_c\{H^t\}\).
\end{proof}

Uniformity in \(i\) is essential, and no argument derives it from the
weaker hypothesis that each \(\alpha_i\) is negligible for every
\emph{fixed} \(i\): there the threshold beyond which the \(i\)-th term
is small can depend on \(i\), and the dependence cannot be removed. A
countermodel makes this concrete: take \(t(\lambda) = \lambda + 1\)
and let \(H^i_\lambda\) be the point mass at \(0\) if
\(i \le \lambda\) and the point mass at \(1\) otherwise. For every
fixed \(i\) the advantage \(\alpha_i(D,\lambda)\) is nonzero only at
\(\lambda = i\), hence negligible; yet \(H^0_\lambda\) and
\(H^{t(\lambda)}_\lambda\) are the point masses at \(0\) and \(1\),
perfectly distinguishable.

Applications usually obtain the conclusion by a stronger, structural
route. Suppose the chain arises from a \emph{single} underlying
assumption \(X \approx_c Y\): neighbouring hybrids \(H^{i}_\lambda\)
and \(H^{i+1}_\lambda\) differ only in that one designated component
is drawn from \(X_\lambda\) in the former and from \(Y_\lambda\) in
the latter, and each hybrid is efficiently samplable given
\((1^\lambda, i)\) and a sample of that component. Then the
\emph{random-index} distinguisher \(D'\) for the pair \((X, Y)\), on
challenge \(z\), samples \(i\) uniformly from \(\{0,\dots,t-1\}\),
builds the hybrid with \(z\) in the designated position, and runs
\(D\); writing \(p_i := \Pr[D(H^i_\lambda)=1]\), its distinguishing
advantage equals
\(\bigl|\frac{1}{t}\sum_{i=0}^{t-1}(p_i - p_{i+1})\bigr| =
|p_0 - p_t|/t\). Since \(D'\) is a single PPT machine, the assumption
makes its advantage negligible, and
\(|p_0 - p_t| \le t(\lambda)\cdot\mathrm{Adv}(D') = \negl(\lambda)\).
Note that this argument bounds \(D\)'s total advantage directly,
bypassing the lemma's uniform per-step hypothesis rather than
discharging it; the lemma itself covers the general case in which no
such common structure is available.

\begin{remark}\label{model:rem:hybrid-poly}
The hybrid argument is precisely where the ``polynomially many''
constraint earns its keep. A negligible per-step advantage multiplied
by a polynomial number of steps remains negligible; a negligible
per-step advantage multiplied by, say, \(2^\lambda\) steps need not.
This is why constructions that would require exponentially many
hybrids do not yield security proofs, and why the closure of
negligibility under polynomial factors is not a technicality but the
load-bearing wall of the whole edifice.
\end{remark}

\subsection{Security as a game}\label{model:subsec:games}

We can now state what it means to ``break'' a scheme. The modern idiom
expresses every security definition as a \emph{game} (or experiment)
played between a \emph{challenger} and an adversary \(\mathcal{A}\).
The challenger, following a fixed protocol, sets up the scheme,
answers whatever queries the adversary may pose, and at the end
declares the adversary a winner or loser by outputting a single bit.
Security says that no PPT adversary wins by more than it could by
trivial guessing.

% FIGURE (planned, dedicated figures pass): fig:game --- the generic
% challenger/adversary interaction: a challenger column (setup; oracle
% query/response arrows crossing to the adversary column; challenge;
% final output arrow back) ending in the win predicate and the output
% bit, with the advantage read off the bit. Deliberately abstract ---
% no scheme instantiated; the caption may forward-point to the
% IND-CPA, PRF, and unforgeability instances of later sections.

\begin{definition}[Security game and advantage]\label{model:def:game}
A \emph{security game} (or \emph{experiment}) \(\mathsf{G}\) is a PPT
interactive algorithm, the challenger, that takes \(1^\lambda\),
interacts with an adversary \(\mathcal{A}\), and finally outputs a bit
\(\mathsf{G}^{\mathcal{A}}(\lambda) \in \{0,1\}\), where \(1\) denotes
``adversary wins''. The \emph{advantage} of \(\mathcal{A}\) is a
measure of how much better than trivial its winning probability is;
its precise form depends on the type of game:
\begin{itemize}
\item \emph{Distinguishing (indistinguishability) games}, where the
  challenger secretly chooses a uniform bit \(b\in\{0,1\}\) and
  \(\mathcal{A}\) outputs a guess \(b'\); the adversary wins iff
  \(b'=b\). Here
  \[
    \mathrm{Adv}_{\mathsf{G}}(\mathcal{A},\lambda)
    \;:=\; \Bigl| \Pr\bigl[\mathsf{G}^{\mathcal{A}}(\lambda)=1\bigr]
                  - \tfrac{1}{2} \Bigr|
    \;=\; \tfrac12\bigl|\Pr[b'=1\mid b=1]-\Pr[b'=1\mid b=0]\bigr|.
  \]
  Trivial guessing achieves winning probability \(1/2\), hence
  advantage \(0\).
\item \emph{Search (unforgeability/inversion) games}, where the
  adversary must produce a specific kind of object---a forgery, a
  preimage, a discrete logarithm---and wins iff it succeeds. Here the
  baseline winning probability is (essentially) \(0\), and one simply
  sets
  \[
    \mathrm{Adv}_{\mathsf{G}}(\mathcal{A},\lambda)
    \;:=\; \Pr\bigl[\mathsf{G}^{\mathcal{A}}(\lambda)=1\bigr].
  \]
\end{itemize}
\end{definition}

\begin{definition}[Security relative to a game]\label{model:def:secure}
A scheme is \emph{secure} with respect to the game \(\mathsf{G}\) if
for every PPT adversary \(\mathcal{A}\),
\[
  \mathrm{Adv}_{\mathsf{G}}(\mathcal{A},\lambda) \;=\; \negl(\lambda).
\]
\end{definition}

Two remarks on the shape of this definition are in order. First, the
quantifier order is essential: \emph{for every} PPT adversary, the
advantage is negligible. For each fixed adversary there may be a
different negligible bound; the claim is universal over the entire
class. Second,
the use of \(\negl\) on the right is what makes the definition
asymptotic and clean: it is robust to the polynomial slop that
reductions introduce, and it captures ``the adversary does no better
than the trivial strategy, up to an amount that vanishes faster than
any inverse polynomial''.

\begin{example}[Pseudorandomness, as a game]
\label{model:exp:prg-game}
To anchor the abstraction with machinery already on the table,
consider the distinguishing game for a pseudorandom generator
\(G : \{0,1\}^\lambda \to \{0,1\}^{2\lambda}\)---the primitive that
witnessed strictness in Proposition~\ref{model:prop:hierarchy}:
\begin{enumerate}
\item The challenger samples \(b \leftarrow \{0,1\}\) uniformly. If
  \(b = 0\) it samples a seed \(s \leftarrow \{0,1\}^\lambda\) and
  sends \(z := G(s)\); if \(b = 1\) it sends a uniform
  \(z \leftarrow \{0,1\}^{2\lambda}\).
\item The adversary \(\mathcal{A}\) outputs a guess \(b'\); the game
  outputs \(1\) iff \(b'=b\).
\end{enumerate}
The generator is secure iff
\(\mathrm{Adv}(\mathcal{A},\lambda)=\bigl|\Pr[b'=b]-\tfrac12\bigr|\)
is negligible for every PPT \(\mathcal{A}\). Here the game packaging
and the ensemble packaging coincide: the two situations are the fixed
ensembles \(\{G(U_\lambda)\}\) and \(\{U_{2\lambda}\}\), and the
displayed identity in Definition~\ref{model:def:game} shows the game
advantage of \(\mathcal{A}\) is exactly half its distinguishing
advantage in the sense of Definition~\ref{model:def:indist}---so the
game is secure iff the ensembles are computationally
indistinguishable. The game form earns its keep beyond this case:
when the adversary chooses part of the challenge or queries oracles
adaptively, the interaction no longer collapses to a comparison of
two fixed ensembles, and the game is the more general packaging---the
shape the security definitions of the later sections take.
\end{example}

\begin{remark}[Game-hopping]\label{model:rem:game-hopping}
Once we phrase security as a game, proofs become \emph{sequences of
games}. One starts from the real security game and rewrites it through
a series of games \(\mathsf{G}_0, \mathsf{G}_1, \dots, \mathsf{G}_k\),
each differing from the last by a small, justifiable change, until
reaching a final game in which the adversary provably has zero (or
trivially negligible) advantage. Each hop is justified by exactly one
of three moves: the two games are identical unless some
negligible-probability ``bad'' event occurs, so their winning
probabilities differ by at most \(\Pr[\text{bad}]\) (on the complement
of that event the two runs coincide); or distinguishing the two games
would break a hardness assumption (a reduction, below); or the change
is purely syntactic. Bounding the total advantage is then a
telescoping sum across hops---a structured cousin of the hybrid
argument (Lemma~\ref{model:lem:hybrid}). This discipline keeps
otherwise sprawling proofs honest and auditable.
\end{remark}

\subsection{Reductions}\label{model:subsec:reductions}

The engine of provable security is the reduction. A \emph{reduction}
is an efficient transformation that converts a successful adversary
against a scheme into a successful algorithm against a problem
believed to be hard. Its existence proves the implication ``\(\Pi\) is
hard \(\Rightarrow\) \(X\) is secure'' by establishing the
contrapositive ``\(X\) is broken \(\Rightarrow\) \(\Pi\) is solved''.

\begin{definition}[Computational problem and hardness]
\label{model:def:problem}
A \emph{computational problem} \(\Pi\) consists of a PPT instance
generator \(\mathrm{Inst}(1^\lambda)\) producing a problem instance
together with any side information, and a relation deciding when a
candidate output \emph{solves} the instance. The \emph{success
probability} of an algorithm \(\mathcal{B}\) is
\[
  \mathrm{Succ}_\Pi(\mathcal{B},\lambda)
  := \Pr\bigl[\mathcal{B}(1^\lambda, \mathrm{Inst}(1^\lambda))
     \text{ solves the instance}\bigr].
\]
The problem \(\Pi\) is \emph{hard} if
\(\mathrm{Succ}_\Pi(\mathcal{B},\lambda) = \negl(\lambda)\) for every
PPT \(\mathcal{B}\) (for a search problem), or if every PPT
\(\mathcal{B}\) decides the instance with at most negligible
advantage, \(\mathrm{Adv}(\mathcal{B},\lambda) = \negl(\lambda)\) in
the sense of Definition~\ref{model:def:game} (for a decision problem).
\end{definition}

\begin{definition}[Black-box reduction]\label{model:def:reduction}
Let \(X\) be a scheme with security game \(\mathsf{G}_X\) and \(\Pi\)
a computational problem. A \emph{(black-box) reduction} from breaking
\(X\) to solving \(\Pi\) is a PPT oracle algorithm \(\mathcal{R}\)
with the following per-adversary guarantee: for \emph{every}
adversary \(\mathcal{A}\), writing
\(\varepsilon_X(\lambda):=\mathrm{Adv}_{\mathsf{G}_X}(\mathcal{A},
\lambda)\) for its advantage, the algorithm
\(\mathcal{R}^{\mathcal{A}}\)---which runs \(\mathcal{A}\) as a
subroutine, simulating \(\mathcal{A}\)'s game internally---solves
\(\Pi\) with success probability
\[
  \mathrm{Succ}_\Pi\bigl(\mathcal{R}^{\mathcal{A}},\lambda\bigr)
  \;\geq\; f\bigl(\varepsilon_X(\lambda), \lambda\bigr),
\]
where \(f\) is a function fixed by the reduction (not by
\(\mathcal{A}\)) that maps every non-negligible \(\varepsilon_X\) to
a non-negligible bound. The running time of \(\mathcal{R}\) is
measured with each call to \(\mathcal{A}\) counted as one step, so
``PPT'' bounds \(\mathcal{R}\)'s own work, and hence its number of
calls to \(\mathcal{A}\), by a polynomial in \(\lambda\); the running
time of \(\mathcal{R}^{\mathcal{A}}\) is then polynomial whenever
\(\mathcal{A}\)'s is, however many times \(\mathcal{R}\) runs
\(\mathcal{A}\). How far \(f\) degrades \(\varepsilon_X\) is the
reduction's \emph{loss}, discussed below.
The reduction is
\emph{black-box} because \(\mathcal{R}\) uses \(\mathcal{A}\) only
through its input/output interface, making no use of
\(\mathcal{A}\)'s code.
\end{definition}

% FIGURE (planned, dedicated figures pass): fig:reduction --- dataflow
% of a black-box reduction: a random Pi-instance entering R from the
% left; inside R, the instance embedded into a simulated security game
% drawn as a wrapper around the (opaque) adversary A; A's winning
% output flowing through the extraction step and leaving R as the
% Pi-solution. The wrapper boundary marks what A sees (a faithful
% game); the R boundary marks black-box access.

The structure of a reduction is always the same and is worth stating
as a recipe. To prove ``if \(\Pi\) is hard then \(X\) is secure'':

\begin{enumerate}
\item Assume toward the contrapositive a PPT adversary \(\mathcal{A}\)
  that breaks \(X\) with non-negligible advantage \(\varepsilon_X\).
\item Construct a PPT algorithm \(\mathcal{R}\) that, given a random
  instance of \(\Pi\), \emph{embeds} that instance into a simulated
  security game for \(\mathcal{A}\). The simulation must be faithful:
  the view \(\mathcal{R}\) presents to \(\mathcal{A}\) must follow
  (perfectly, statistically, or computationally) the distribution of
  the real game, so that \(\mathcal{A}\)'s advantage carries over.
\item Run \(\mathcal{A}\). From whatever \(\mathcal{A}\) does to win
  its game, \(\mathcal{R}\) \emph{extracts} a solution to the
  \(\Pi\)-instance.
\item Conclude that \(\mathcal{R}\) solves \(\Pi\) with non-negligible
  probability, contradicting the hardness of \(\Pi\). Therefore no
  such \(\mathcal{A}\) exists, and \(X\) is secure.
\end{enumerate}

A complete, genuine reduction illustrates the point.

\begin{theorem}[Hard-core security of the Diffie--Hellman-style key,
schematically]\label{model:thm:reduction-example}
Let \(\Pi = \mathrm{CDH}\) be the computational Diffie--Hellman
problem in a cyclic group \(\G=\langle g\rangle\) of prime order
\(q\), the setting of the \emph{Math Guide} (\S``The discrete
logarithm problem''): given \((g, g^a, g^b)\) for uniform
\(a,b \in \Z/q\Z\), compute \(g^{ab}\). Consider the ``key-recovery''
game \(\mathsf{G}_X\) for a key-exchange scheme in which the
adversary, given the public transcript \((g^a, g^b)\), wins iff it
outputs the shared key \(g^{ab}\). If \(\mathrm{CDH}\) is hard, then
no PPT adversary wins \(\mathsf{G}_X\) with noticeable probability.
\end{theorem}

\begin{proof}
Suppose \(\mathcal{A}\) is a PPT adversary that, on input a transcript
\((g^a, g^b)\), outputs \(g^{ab}\) with probability
\(\varepsilon_X(\lambda)=\mathrm{Adv}_{\mathsf{G}_X}(\mathcal{A},
\lambda)\). Construct \(\mathcal{R}\) solving \(\mathrm{CDH}\). On
input a \(\mathrm{CDH}\) instance \((g, A, B) = (g, g^a, g^b)\):
\begin{enumerate}
\item The reduction \(\mathcal{R}\) forms the simulated transcript
  \((A, B)\) and
  feeds it to \(\mathcal{A}\) as the public transcript of
  \(\mathsf{G}_X\).
\item Because \((a,b)\) are uniform in the \(\mathrm{CDH}\) instance
  \emph{and} in the real game, the view of \(\mathcal{A}\) follows
  \emph{identically} the distribution of its view in \(\mathsf{G}_X\)
  (the simulation is perfect; no loss in faithfulness).
\item When \(\mathcal{A}\) outputs a candidate key \(K\),
  \(\mathcal{R}\) outputs \(K\) as its \(\mathrm{CDH}\) answer.
\end{enumerate}
Since the simulated and real views are identically distributed,
\(\mathcal{A}\) outputs \(K = g^{ab}\) with exactly probability
\(\varepsilon_X(\lambda)\); whenever it does, \(\mathcal{R}\) returns
the correct \(\mathrm{CDH}\) solution. Hence
\[
  \mathrm{Succ}_{\mathrm{CDH}}(\mathcal{R},\lambda)
  \;=\; \varepsilon_X(\lambda).
\]
The running time of \(\mathcal{R}\) is that of \(\mathcal{A}\) plus a
constant amount of group bookkeeping, hence polynomial. If
\(\varepsilon_X\) were noticeable, \(\mathcal{R}\) would solve
\(\mathrm{CDH}\) with noticeable probability, contradicting its
hardness. Therefore \(\varepsilon_X(\lambda)\) is not noticeable, as
claimed. In fact, since \(\mathcal{R}\) is PPT and \(\mathrm{CDH}\) is
hard (Definition~\ref{model:def:problem}, search branch),
\(\varepsilon_X(\lambda) =
\mathrm{Succ}_{\mathrm{CDH}}(\mathcal{R},\lambda) = \negl(\lambda)\).
\end{proof}

This example is deliberately \emph{tight}: the reduction loses
nothing, calls the adversary once, and converts advantage
\(\varepsilon_X\) into success probability exactly \(\varepsilon_X\).
Most reductions are not so fortunate, which leads to the quantitative
anatomy of a reduction.

\begin{definition}[Reduction loss and tightness]\label{model:def:tightness}
Consider a reduction \(\mathcal{R}\) that transforms an adversary
running in time \(t\) with advantage \(\varepsilon\) into a
\(\Pi\)-solver running in time \(t' = t + t_{\mathcal{R}}\) with
success probability \(\varepsilon'\). The \emph{security loss} is the
factor
\[
  L(\lambda) \;:=\; \frac{\varepsilon / t}{\varepsilon' / t'}
  \;=\; \frac{\varepsilon}{\varepsilon'}\cdot\frac{t'}{t},
\]
the ratio of the solver's work-per-success to the adversary's
work-per-success. A reduction is \emph{tight} if
\(L(\lambda) = O(1)\) (a small constant, ideally \(1\)), and
\emph{loose} if \(L(\lambda)\) grows with \(\lambda\) (e.g.\
\(L = q_H\), the number of random-oracle queries
(Definition~\ref{model:def:rom} below), or \(L = q_s\), the number of
signature queries, or \(L = \lambda\)). Two common sources
of loss are: an \emph{advantage gap} \(\varepsilon' \ll \varepsilon\)
(the solver succeeds far less often than the adversary, e.g.\ because
the reduction must \emph{guess} where to embed its instance among
\(q\) queries, succeeding only with probability \(1/q\)); and a
\emph{time blow-up} \(t' \gg t\) (the reduction runs the adversary
many times, as in rewinding/forking arguments).
\end{definition}

\begin{remark}[Why loss matters concretely]\label{model:rem:loss-matters}
Asymptotically, loss is invisible: a polynomial loss factor turns a
negligible advantage into a negligible advantage and a hard problem
into security, without qualification. Concretely, however, loss
carries a cost. Suppose a scheme reduces to a problem believed to
offer \(128\) bits of security, but the reduction loses a factor
\(L = 2^{40}\) (not unusual: \(q_H \approx 2^{60}\) random-oracle
queries times a forking-lemma square-root loss can be much worse).
Then the \emph{proof} guarantees only \(128 - 40 = 88\) bits of
security for the scheme. To restore a \(128\)-bit guarantee one must
enlarge the group so that the \emph{underlying} problem offers
\(168\) bits---larger keys, slower operations. A tight reduction lets
the scheme inherit the full strength of the assumption at the smallest
parameters; a loose one imposes a permanent cost on every user. This
is why tightness is a first-class design goal in modern cryptography,
not a theoretical nicety.
\end{remark}

\begin{remark}[Rewinding and the forking lemma, in brief]
\label{model:rem:forking}
A recurrent source of large loss is \emph{rewinding}: running the
adversary, recording its behaviour, then restarting it from a saved
state with fresh randomness on some branch to obtain a second,
correlated transcript. From two transcripts that share a prefix but
diverge afterward one can often extract a secret (e.g.\ two valid
signatures on the same commitment yield the signing key, or two
accepting proof transcripts yield a witness). The \emph{forking
lemma}, stated as the general forking lemma under ``Security in the
random oracle model and the forking lemma'', quantifies the success
probability of obtaining such a pair: if the adversary succeeds with
probability \(\varepsilon\) after at most \(q\) random-oracle queries
whose answers are drawn from a set of size \(h\), the forked run
succeeds with probability at least \(\varepsilon(\varepsilon/q - 1/h)\),
i.e.\ on the order of \(\varepsilon^2/q\). The square in
\(\varepsilon\) is a quadratic security loss: to keep
\(\varepsilon^2/q\) noticeable one needs \(\varepsilon\) noticeably
large, roughly halving the bit-security extracted. Such
arguments---central to the analyses of Schnorr-type and Fiat--Shamir
signatures that this volume builds toward---are the archetype of
loose but still valid reductions.
\end{remark}

\subsection{The standard model and the random oracle model}
\label{model:subsec:rom}

Reductions reduce security to assumptions. \emph{Which} assumptions
one is willing to make defines the \emph{model} in which a proof
resides. Two models predominate.

\begin{definition}[Standard model]\label{model:def:standardmodel}
A proof resides in the \emph{standard model} if it assumes only the
computational hardness of explicitly stated problems (factoring,
discrete logarithm, the decisional Diffie--Hellman problem of the next
section, and so on) and grants the adversary no idealised access to any
component of the scheme. Every primitive, including every hash
function, is a concrete algorithm whose code the adversary may inspect
and exploit.
\end{definition}

\begin{definition}[Random oracle model]\label{model:def:rom}
A proof resides in the \emph{random oracle model} (ROM) if it
idealises one or more hash functions as a \emph{random oracle}: a
publicly accessible function \(\mathcal{O}\) that, on each distinct
input, returns an \emph{independent, uniformly random} output of the
appropriate length, consistently answering repeated queries with the
same value. The oracle is available to the honest parties, to the
adversary, and---crucially---to the reduction. No party holds the
``code'' of \(\mathcal{O}\); it can only be queried.
\end{definition}

The random oracle model is an idealisation: real hash functions are
fixed, public algorithms, not truly random functions. The motivation
for adopting it is that the reduction gains two capabilities that the
standard model withholds and that render many otherwise-intractable
proofs tractable.

\begin{remark}[The powers of the random oracle: observability and
programmability]\label{model:rem:rom-powers}
In the ROM the reduction \emph{simulates} the oracle for the
adversary, and this yields two additional capabilities.
\emph{Observability}: the reduction sees every query the adversary
makes to \(\mathcal{O}\). Since the adversary, to use a hash value
meaningfully, must query it, the reduction learns the adversary's
intermediate computations---for instance the preimage the adversary is
about to hash. \emph{Programmability}: the reduction may choose each
oracle answer on the fly, subject only to the answers appearing
uniform and consistent. It can therefore \emph{plant} its hard-problem
instance inside an oracle response---e.g.\ set
\(\mathcal{O}(x) := B\) where \(B\) is part of a CDH challenge---so
that the adversary's success necessarily reveals the solution.
Neither capability exists for a concrete standard-model hash function,
which fixes its outputs in advance and hides its internal queries.
\end{remark}

\begin{example}[The benefit of the ROM: full-domain hash, sketched]
\label{model:exp:rom-fdh}
In the ``full domain hash'' signature, the signer signs a message
\(m\) as \(\sigma = H(m)^{d} \bmod N\), where \((N,d)\) is an RSA
private key---a modulus \(N = pq\) with secret prime factors and a
secret exponent \(d\) paired with a public \(e\) so that
\(x^{ed} \equiv x \pmod N\) for \(x\) coprime to \(N\), by Euler's
theorem (\emph{Math Guide}, \S``Fermat's little theorem and Euler's
theorem''); thus \(x \mapsto x^{e}\) is a permutation of the unit
group \((\Z/N\Z)^\times\) that only the holder of \(d\) can
invert---and \(H\) is a hash function mapping messages into
\((\Z/N\Z)^\times\). (Restricting to units costs nothing: a nonzero
non-unit of \(\Z/N\Z\) shares a factor with \(N\), so stumbling on
one factors \(N\). Hash functions appear here informally, as
this sketch needs them; their formal syntax and security games are the
business of the next-but-one section.) To prove unforgeability one
reduces to RSA inversion: given a challenge \(y\) (find \(y^{d}\)),
the reduction programs the oracle so that for a random subset of
messages \(H(m) := y \cdot r_m^{e}\) (embedding the challenge) and for
the rest \(H(m) := r_m^{e}\) (where it knows the ``signature''
\(r_m\)); with each \(r_m\) drawn uniformly from
\((\Z/N\Z)^\times\), the programmed values are distributed exactly as
honest oracle answers on \(H\)'s range. When the forger outputs a
forgery on a message where the
reduction embedded the challenge, the reduction extracts \(y^{d}\).
This proof has \emph{no known standard-model analogue} for the plain
scheme: it relies essentially on programming \(H\). The cost is a
security loss (the reduction must guess which message the forger
attacks, incurring a factor \(q_H\) unless cleverly optimised)---
connecting back to Definition~\ref{model:def:tightness}.
\end{example}

\begin{remark}[The status of the ROM: heuristic, not theorem]
\label{model:rem:rom-caveat}
A proof in the ROM is a proof about an \emph{idealised} world:
deploying the scheme instantiates \(\mathcal{O}\) with a concrete
hash function, and the proof does \emph{not} formally transfer---a ROM
proof rules out all attacks that treat the hash as a black box, while
offering no guarantee against attacks on the specific structure of the
chosen hash. The case for trusting it runs in three steps: prove the
scheme secure against a perfect random oracle; instantiate
\(\mathcal{O}\) with a hash believed to have no exploitable structure;
and observe that any remaining attack must exploit structure the
random oracle lacks, so would be a direct break of the presumed
structure-free hash. The gap is nonetheless a theorem: Canetti,
Goldreich, and Halevi exhibit contrived schemes provably secure in
the ROM yet insecure under \emph{every} efficiently computable
instantiation, so no concrete hash soundly instantiates the oracle
for all schemes. No natural deployed scheme is known to suffer this
fate, and the working stance of this volume is the practical one: a
ROM proof is strong positive evidence, not a guarantee, and a
standard-model proof is preferred when one is available at
comparable cost. One caveat looks forward. A quantum adversary
(Remark~\ref{model:rem:variants}) queries the oracle \emph{in
superposition}, a single query addressing a weighted combination of
all inputs at once; in the resulting quantum random oracle model
(QROM) the two powers of Remark~\ref{model:rem:rom-powers} become
subtler and classical ROM theorems do not automatically carry over,
so post-quantum claims for arguments made non-interactive by hashing
the transcript (the Fiat--Shamir transform, developed under ``The
Fiat--Shamir transform: from interactive to non-interactive'') must
be made in that model.
\end{remark}

\begin{remark}[Relatives: CRS, generic-group, and algebraic-group
models]\label{model:rem:other-models}
Between and beyond these poles lie further idealisations relevant to
the proof systems this volume targets. The \emph{common reference
string} (CRS) model posits a trusted public string sampled from a
prescribed distribution, used by many succinct proof systems. The
\emph{generic group model} (GGM) idealises the group dual to the
manner in which the ROM idealises a hash: the adversary may only
perform group operations through an oracle on opaque ``handles'',
never exploiting the bit representation of group elements; it is the
natural setting for lower bounds on discrete-log-type problems. The
\emph{algebraic group model} (AGM) is an intermediate, milder
idealisation requiring the adversary to ``explain'' every group
element it outputs as a known combination of its inputs. These models
trade realism for provability along the same axis as the ROM, and the
constructions ahead invoke them by name where they require them.
\end{remark}

\subsection{Concrete security and the bit-security budget}
\label{model:subsec:concrete}

The asymptotic language---``negligible'', ``polynomial'', ``for all
sufficiently large \(\lambda\)''---is mathematically clean but
operationally silent. It establishes that a scheme is eventually
secure; it does not establish whether \emph{this} deployment with
\emph{these} parameters resists an adversary with \emph{that} budget.
Concrete (or exact) security supplies the missing numbers.

\begin{definition}[Concrete security:
\((t,\varepsilon)\)-security]\label{model:def:concrete}
Fix the security parameter to a concrete value (so all sizes are
fixed). A scheme is \emph{\((t,\varepsilon)\)-secure} for a game
\(\mathsf{G}\) if every adversary running in time at most \(t\) has
advantage \(\mathrm{Adv}_{\mathsf{G}}(\mathcal{A}) \leq \varepsilon\).
More refined statements bound the advantage as an explicit function of
the adversary's resources, e.g.\ time \(t\), number of queries \(q\),
and memory \(s\):
\(\mathrm{Adv}_{\mathsf{G}}(\mathcal{A}) \leq f(t,q,s)\).
\end{definition}

\begin{definition}[\(n\)-bit security]\label{model:def:bitsec}
A scheme (or problem) is commonly said to offer \(n\) \emph{bits of
security} when the cheapest known attack achieving constant success
probability (say at least \(\tfrac12\)) costs on the order of
\(2^{n}\) elementary operations. A proof may
justify the stronger, explicit claim that every adversary running in
time \(t\) has advantage at most about \(t/2^{n}\), but that linear
advantage bound is not implied merely by the absence of a faster known
attack. Thus ``\(128\)-bit security'' is a concrete estimate, relative
to a specified attack model and state of knowledge, that breaking the
scheme costs roughly \(2^{128}\) work.
\end{definition}

The relationship to the asymptotic picture is the following: one
chooses the security parameter \(\lambda\) \emph{precisely so that}
the resulting concrete scheme delivers the desired number of bits of
security. In the cleanest case the two coincide, \(n = \lambda\), but
in general one must account for the actual cost of the best attack on
the underlying problem.

\begin{example}[Why \(\lambda \neq n\) for discrete logarithm]
\label{model:exp:bitsec-dlog}
For a generic-group or hash-based primitive, doubling the output
length doubles the bit security: a hash with \(2n\)-bit output offers
\(n\)-bit collision resistance---by the birthday bound (\emph{Math
Guide}, \S``The union bound and a birthday calculation'')---and
\(2n\)-bit preimage resistance, so the parameter and the security
level track simply. For discrete-logarithm-based primitives the
accounting differs and instructs. In a group of prime order
\(q \approx 2^{\ell}\), generic attacks (baby-step/giant-step,
Pollard's rho) solve discrete logarithm in time
\(\approx \sqrt{q} = 2^{\ell/2}\) (\emph{Math Guide}, \S``The discrete
logarithm problem''). Hence a prime-order elliptic-curve
group whose order is \(256\) bits offers only about \(128\) bits of
security---one halves the bit length of the group order to obtain the
generic security level. For discrete logarithm in a finite-field
multiplicative group, sub-exponential index-calculus attacks
additionally force the ambient field modulus into the thousands of
bits for \(128\)-bit security. This is precisely why elliptic-curve
groups, for which no comparable sub-exponential classical attack is
known, are preferred: they deliver \(n\)-bit security at
\(\ell \approx 2n\) rather than thousands of bits. The single number
\(n\) conceals a
problem-specific translation from all parameter sizes to the security
level.
The working groups of this volume are priced by exactly this
arithmetic: the Pasta curves assembled in the \emph{Math Guide}
(\S``Pallas and Vesta assembled'') have prime order close to
\(2^{254}\), so the square-root rule prices generic discrete
logarithm there at about \(2^{127}\) operations; the next section
audits them criterion by criterion under ``Instantiation on elliptic
curves; the Pasta curves''.
\end{example}

\begin{remark}[Reading a bit-security budget end to end]
\label{model:rem:budget}
Combining the preceding components yields the practitioner's
parameter-selection recipe. Decide the target security level \(n\)
(commonly \(128\)). Identify the underlying hard problem and the cost
of its best attack as a function of the parameter (e.g.\
\(2^{\ell/2}\) for elliptic-curve discrete log), and set the parameter
so the problem itself offers \(n\) bits (here \(\ell = 2n = 256\)).
Then subtract the reduction's security loss
(Definition~\ref{model:def:tightness}): if the proof loses a factor
\(L = 2^{k}\) (that is, \(k\) bits), the problem must offer
\(n + k\) bits for the
\emph{scheme} to offer \(n\), so one enlarges the parameter
accordingly---or, with a tight reduction, one need not. Finally,
sanity-check against the ROM caveat
(Remark~\ref{model:rem:rom-caveat}): a ROM proof bounds only
black-box attacks on the hash, so the concrete hash must additionally
exhibit no exploitable structure. The clean asymptotic theorem ``the
scheme is secure if the problem is hard'' thus unfolds, through the
quantitative machinery of this section, into a concrete choice of
bytes.
\end{remark}

\begin{remark}[A note on what security does and does not promise]
\label{model:rem:scope}
A caution that frames everything above is in order. A proof of
security is a proof \emph{relative to a model}: a definition of the
goal (the game), a class of adversaries (PPT, possibly non-uniform or
quantum), and a set of assumptions (the hard problems, and any
idealisation such as the ROM). It guarantees nothing outside that
model. Side-channel leakage (timing, power), faulty randomness,
implementation bugs, broken assumptions, and goals the game did not
capture all lie beyond its reach. The value of the framework is not
that it renders schemes unconditionally safe---nothing can---but that
it renders the \emph{conditions} of safety explicit, falsifiable, and
modular, so that progress on assumptions and attacks translates
directly into recalibrated, auditable guarantees. Everything this
monograph builds rests on this discipline.
\end{remark}

\section{Hardness assumptions}\label{sec:assumptions}

Every secret in this volume is an exponent. A spending key, the
trapdoor of a commitment, the nonce of a signature---each lives in
\(\Z/q\Z\) and faces the world only through a group element
\(g^{x}\). The adversary of this section is the patient cryptanalyst
who collects those elements from the public record and asks for
nothing more than the exponent behind one of them: an adversary who
computes discrete logarithms recovers signing keys from verification
keys, opens and re-opens commitments at will, and strips the blinding
from every rerandomised key the later sections construct. What she can
\emph{forge} is everything; what she must be \emph{denied} is a single
inversion. This section assembles the family of assumptions that deny
it: the discrete logarithm problem itself, the ladder of
Diffie--Hellman variants above it, the generic-attack bounds that
calibrate how hard the problem can possibly be, the composite-order
failure mode that squanders that hardness, the elliptic curves on
which the remainder of the volume instantiates the family, and the
quantum caveat under which the whole edifice stands. Throughout, the
lodestar is a single number: the security parameter \(\lambda\), read
as the bit-length of security demanded. The recurring conclusion is
that to obtain \(\lambda\) bits of security against the best generic
attacks the group must have \emph{prime} order \(q\) with
\(\log_2 q \ge 2\lambda\). Notation is fixed once:
\(G = \langle g\rangle\) denotes a cyclic group of prime order \(q\),
written multiplicatively until \S\ref{assumptions:subsec:pasta}
switches to the additive notation of elliptic curves.

\subsection{The discrete logarithm problem}
\label{assumptions:subsec:dlp}

The stage is a finite cyclic group. Recall from the \emph{Math Guide}
(\S``The discrete logarithm problem'') that a group \(G\) is
\emph{cyclic} if \(G = \langle g \rangle = \{\, g^{k} : k \in \Z \,\}\)
for some \emph{generator} \(g\); written additively, as we shall write
elliptic curves, this reads \(G = \{\, kg : k \in \Z \,\}\). If
\(\abs{G} = q\) then \(\Z/q\Z \cong G\) via \(k \mapsto g^{k}\), an
isomorphism trivial to evaluate forwards---repeated squaring costs
\(O(\log k)\) group operations---but conjecturally infeasible to
invert. That asymmetry \emph{is} the discrete logarithm problem.

\begin{definition}[Discrete logarithm]
\label{assumptions:def:dlog}
Let \(G = \langle g \rangle\) be a cyclic group of order \(q\),
written multiplicatively. For \(h \in G\), the \emph{discrete
logarithm} of \(h\) to the base \(g\), denoted \(\log_g h\), is the
unique residue \(x \in \Z/q\Z\) with \(g^{x} = h\). Existence follows
from \(G = \langle g\rangle\) and \(\ord(g) = q\); uniqueness because
\(g^{x} = g^{x'}\) forces \(g^{x - x'} = 1\), whence
\(q \mid x - x'\).
\end{definition}

\begin{definition}[Discrete logarithm problem, DLP]
\label{assumptions:def:dlp}
Fix a family \(\{G_\lambda\}\) of cyclic groups, where
\(G_\lambda = \langle g_\lambda\rangle\) has \emph{prime} order
\(q_\lambda\) with \(\log_2 q_\lambda = \Theta(\lambda)\), together
with a sampler that on input \(1^{\lambda}\) outputs a description of
\((G_\lambda, g_\lambda, q_\lambda)\). The game
\(\mathsf{DLOG}^{\mathcal A}(\lambda)\) runs as follows.
\begin{enumerate}
\item The challenger runs the sampler, draws
  \(x \xleftarrow{\$} \Z/q_\lambda\Z\) uniformly, and sends
  \((G_\lambda, g_\lambda, q_\lambda,\, h = g_\lambda^{\,x})\) to the
  adversary \(\mathcal A\).
\item The adversary outputs \(x' \in \Z/q_\lambda\Z\).
\item The game outputs \(1\) iff \(g_\lambda^{\,x'} = h\).
\end{enumerate}
The \emph{DLP assumption} for the family states that for every PPT
\(\mathcal A\),
\[
\Pr\bigl[\mathsf{DLOG}^{\mathcal A}(\lambda) = 1\bigr]
  \;\in\; \negl(\lambda).
\]
\end{definition}

The order is restricted to be prime for reasons that
\S\ref{assumptions:subsec:pohlig} makes emphatic: in prime order every
non-identity element is a generator, there are no proper nontrivial
subgroups through which information can leak, and the problem cannot
be split across factors of the order. For now we record the
convenience that prime order affords immediately.

\begin{proposition}[Random self-reducibility of DLP]
\label{assumptions:prop:rsr}
In a cyclic group \(G = \langle g\rangle\) of prime order \(q\), the
discrete logarithm problem is \emph{random self-reducible}: an
algorithm solving a fraction \(\varepsilon\) of instances, over
uniformly random \(h\), converts into one solving \emph{any} fixed
instance with probability \(\varepsilon\) per trial. Consequently
worst-case and average-case hardness coincide.
\end{proposition}

\begin{proof}
Let \(\mathcal A\) succeed on a uniform target with probability
\(\varepsilon\). Given a fixed challenge \(h = g^{x}\), draw
\(r \xleftarrow{\$} \Z/q\Z\) and form
\(h' = h \cdot g^{r} = g^{x + r}\). Addition modulo \(q\) is a
bijection and \(r\) is uniform, so \(h'\) is uniform in \(G\) and
independent of \(h\); hence \(\mathcal A(h')\) returns
\(x + r \bmod q\) with probability \(\varepsilon\), and subtracting
\(r\) recovers \(x\). Each candidate is checked by testing
\(g^{x} \stackrel{?}{=} h\), so failures are recognised. Fresh
randomness \(r\) per trial makes trials independent: \(t\) of them
succeed with probability \(1 - (1 - \varepsilon)^{t}\), overwhelming
once \(t = \omega(\varepsilon^{-1} \log \lambda)\).
\end{proof}

Random self-reducibility is the formal licence to speak of ``the''
hardness of DLP in a group: an average-case solver also solves every
fixed instance after rerandomisation. This does not literally make
every instance difficult---the logarithm of the identity, for one, is
immediate.
% Deviation from brief (y2-6): the brief's "no weak instances; every
% instance is as hard as a random one" is false as literally stated
% (the identity's logarithm is immediate); qualified to the
% average-to-worst conversion the proposition actually gives.

\subsection{Diffie--Hellman: computational and decisional}
\label{assumptions:subsec:dh}

The discrete logarithm problem asks to invert exponentiation. Many
protocols---the Diffie--Hellman key exchange and ElGamal encryption
are the historical motivations---need only a weaker-looking guarantee
about the \emph{product} of two exponents. Suppose Alice publishes
\(g^{a}\) and Bob publishes \(g^{b}\); their shared secret is
\(g^{ab}\), which each computes from the other's public value and
their own exponent. The eavesdropper sees \(g, g^{a}, g^{b}\) and
wants \(g^{ab}\).

\begin{definition}[Computational Diffie--Hellman, CDH]
\label{assumptions:def:cdh}
With notation as in Definition~\ref{assumptions:def:dlp}, the game
\(\mathsf{CDH}^{\mathcal A}(\lambda)\) runs as follows: the challenger
draws \(a, b \xleftarrow{\$} \Z/q\Z\) independently and sends
\((g, g^{a}, g^{b})\); the adversary outputs \(z \in G\); the game
outputs \(1\) iff \(z = g^{ab}\). The \emph{CDH assumption} states
that for every PPT \(\mathcal A\),
\[
\Pr_{a,b}\bigl[\mathsf{CDH}^{\mathcal A}(\lambda) = 1\bigr]
  \;\in\; \negl(\lambda).
\]
\end{definition}

The computational guarantee delivers the shared key but not its
secrecy: an adversary may be unable to \emph{produce} \(g^{ab}\) yet
perfectly able to \emph{distinguish} it from a random group element,
and that alone already breaks, say, the indistinguishability of an
ElGamal ciphertext. The decisional variant excludes it.

\begin{definition}[Decisional Diffie--Hellman, DDH]
\label{assumptions:def:ddh}
The \emph{decisional Diffie--Hellman problem} is to distinguish two
distributions over \(G^{3}\),
\[
\mathcal D_{\mathrm{dh}} = \bigl(g^{a}, g^{b}, g^{ab}\bigr),
\qquad
\mathcal D_{\mathrm{rand}} = \bigl(g^{a}, g^{b}, g^{c}\bigr),
\]
with \(a, b, c\) independent and uniform in \(\Z/q\Z\); this is the
distinguishing game of Definition~\ref{model:def:game}, the challenger
handing over a sample of one or the other according to its hidden bit.
The \emph{DDH assumption} states that every PPT distinguisher \(D\)
has negligible advantage
\[
\mathrm{Adv}_{\mathrm{ddh}}(D)
= \Bigl| \Pr_{a,b}\bigl[D(g^{a}, g^{b}, g^{ab}) = 1\bigr]
  - \Pr_{a,b,c}\bigl[D(g^{a}, g^{b}, g^{c}) = 1\bigr] \Bigr|
  \;\in\; \negl(\lambda).
\]
A triple \((g^{a}, g^{b}, g^{c})\) with \(c \equiv ab \pmod q\) is a
\emph{Diffie--Hellman triple}; the assumption says no efficient
distinguisher recognises such triples.
\end{definition}

% Deviation from brief (y2-10): the brief's three-randomiser map
% (u^r g^t, v^s, w^{rs} v^{ts}) is defective (the second component is
% fixed whenever b = 0, and its non-DH analysis is loose); replaced by
% the standard four-randomiser map, with the conclusion refined to
% statistical distance 1/q from D_rand.
\begin{remark}[DDH admits a random self-reduction]
\label{assumptions:rem:ddh-self-reducible}
Rerandomisation works for DDH as it did for DLP, though the correct
map takes some care. Given a target triple
\((u, v, w) = (g^{a}, g^{b}, g^{c})\), choose
\(\alpha, \gamma \xleftarrow{\$} (\Z/q\Z)^{\times}\) and
\(\beta, \delta \xleftarrow{\$} \Z/q\Z\), and form
\[
(u', v', w')
= \bigl(u^{\alpha} g^{\beta},\; v^{\gamma} g^{\delta},\;
  w^{\alpha\gamma}\, u^{\alpha\delta}\, v^{\beta\gamma}\,
  g^{\beta\delta}\bigr).
\]
Writing \(A = \alpha a + \beta\), \(B = \gamma b + \delta\), and
\(E = c - ab\), the new exponents are \(A\), \(B\), and
\(AB + \alpha\gamma E\), with \(A\) and \(B\) uniform and independent.
A DH triple (\(E = 0\)) therefore becomes a fresh uniform DH triple; a
non-DH triple becomes uniform conditional on being non-DH, because
\(\alpha\gamma E\) is uniform in \((\Z/q\Z)^{\times}\). The latter
distribution lies at statistical distance
(Definition~\ref{model:def:statdist}) \(1/q\) from
\(\mathcal D_{\mathrm{rand}}\), whose third component completes a DH
triple by accident with probability \(1/q\). A noticeable
distinguishing advantage on random instances thus transfers to every
fixed triple, up to this negligible term.
\end{remark}

\subsection{The assumption hierarchy}
\label{assumptions:subsec:hierarchy}

The three problems form a chain. Write \(A \le B\) for ``\(A\) reduces
to \(B\)'': an oracle solving \(B\) yields an efficient solver for
\(A\); equivalently, \(B\) is at least as hard as \(A\), so that if
\(B\) is easy then \(A\) is easy; equivalently again, the assumption
``\(A\) is hard'' is the \emph{stronger} commitment, since it implies
``\(B\) is hard''. The reductions of this subsection give
\[
\mathrm{DDH} \;\le\; \mathrm{CDH} \;\le\; \mathrm{DLP},
\]
ordering the problems by increasing hardness and the assumptions by
decreasing strength: assuming DDH hard is the strongest commitment,
assuming DLP hard the weakest.

% TikZ figure placeholder --- fig:dh-hierarchy: reduction ordering
% DLP >= CDH >= DDH (solver arrows pointing down the chain, assumption
% strength increasing the other way), with the pairing separation edge
% of Example~\ref{assumptions:exa:pairing} marking where DDH detaches
% from CDH in gap groups. Dedicated figures pass supplies the drawing.

\begin{theorem}[DLP is at least as hard as CDH]
\label{assumptions:thm:dlp-le-cdh}
If DLP is solvable in PPT, then CDH is solvable in PPT. Equivalently,
the CDH assumption implies the DLP assumption.
\end{theorem}

\begin{proof}
Suppose \(\mathcal A\) solves DLP. Given a CDH instance
\((g, g^{a}, g^{b})\), run \(\mathcal A\) on \((g, g^{a})\) to recover
\(a\), then compute \((g^{b})^{a} = g^{ab}\) by fast exponentiation.
The cost is one DLP call plus \(O(\log q)\) group operations, and only
one of the two exponents ever needed inverting.
\end{proof}

\begin{theorem}[CDH is at least as hard as DDH]
\label{assumptions:thm:cdh-le-ddh}
If CDH is solvable in PPT with non-negligible success probability,
then DDH is solvable in PPT with non-negligible advantage; if the
success probability is noticeable, rerandomisation drives the
advantage overwhelmingly close to \(1\). Equivalently, the DDH
assumption implies the CDH assumption.
\end{theorem}

% Deviation from brief (y2-13): the brief amplifies to advantage
% 1 - negl from bare CDH solvability; that needs the noticeable
% hypothesis (a merely non-negligible epsilon exceeds an inverse
% polynomial only infinitely often, so the amplified advantage is
% overwhelming only along that subsequence). Stated in two tiers.

\begin{proof}
Suppose \(\mathcal A\) solves CDH with probability \(\varepsilon\).
Build a distinguisher \(D\): on input
\((g^{a}, g^{b}, z)\), call \(\mathcal A(g, g^{a}, g^{b})\) to obtain
a candidate \(w\), and output \(1\) iff \(w = z\). On a
Diffie--Hellman triple, \(z = g^{ab}\), so \(D\) outputs \(1\)
whenever \(\mathcal A\) succeeds---probability \(\varepsilon\). On a
random triple, \(z = g^{c}\) with \(c\) uniform and independent of
\(a, b\), while \(w\) depends only on \((g^{a}, g^{b})\); the event
\(w = z\) therefore has probability exactly \(1/q\). Hence
\[
\mathrm{Adv}_{\mathrm{ddh}}(D) \;\ge\; \varepsilon - \frac{1}{q},
\]
non-negligible whenever \(\varepsilon\) is, since \(1/q\) is
negligible.

For the amplified form let \(\varepsilon\) be noticeable. The
distinguisher \(D'\) rerandomises its input \(t\) times
independently by the map of
Remark~\ref{assumptions:rem:ddh-self-reducible}, runs \(D\) with
fresh coins on each copy, and outputs \(1\) iff \emph{some} copy
answered \(1\). A DH input becomes \(t\) independent uniform DH
triples, each accepted with probability \(\varepsilon\), so \(D'\)
accepts with probability at least \(1 - (1 - \varepsilon)^{t}\). A
sample of \(\mathcal D_{\mathrm{rand}}\) fails to be a DH triple
except with probability \(1/q\), and then each copy's third
component is uniform over the \(q - 1\) non-completing values given
the first two, so each copy is accepted with probability at most
\(1/(q - 1)\) and \(D'\) accepts with probability at most
\(1/q + t/(q - 1)\). With
\(t = \omega(\varepsilon^{-1} \log \lambda)\), polynomial because
\(\varepsilon\) is noticeable, the advantage of \(D'\) is
\(1 - \negl(\lambda)\).
\end{proof}

\begin{example}[DDH can be easy while CDH and DLP are hard]
\label{assumptions:exa:pairing}
Some elliptic-curve groups carry an efficiently computable,
non-degenerate \emph{bilinear pairing}
\(e \colon G \times G \to G_T\) with
\(e(g^{a}, g^{b}) = e(g, g)^{ab}\). On such a group DDH is
\emph{easy}: to test whether \((g^{a}, g^{b}, z)\) is a DH triple,
check
\[
e(g^{a}, g^{b}) \stackrel{?}{=} e(g, z),
\]
which holds iff \(\log_g z \equiv ab \pmod q\). Yet CDH---and a
fortiori DLP---may remain hard, because the pairing reveals
\(g^{ab}\) only inside the \emph{target} group \(G_T\), never as an
element of \(G\). Such groups, called \emph{gap groups}, separate DDH
from CDH and found pairing-based cryptography. The moral is that DDH
is a genuinely stronger assumption: choosing a ``DDH group'' means
choosing a group that rules out pairings of this convenient kind on
\(G\) itself.
\end{example}

\begin{remark}[The converses are open]
\label{assumptions:rem:converse}
The reverse reductions are not known in general. Whether CDH implies
DLP---whether the ability to compute \(g^{ab}\) can be leveraged to
extract \(a\)---is the longstanding Diffie--Hellman versus discrete
log question; partial results of den Boer and Maurer establish
equivalence when a suitable auxiliary group of \emph{smooth}
order---an order that is a product of small primes---is available,
but no unconditional equivalence is known. In the other
direction DDH may be strictly easier than CDH, as
Example~\ref{assumptions:exa:pairing} shows. Practice therefore
assumes exactly the strength a protocol requires: signatures and key
exchange often need only CDH, while ElGamal-style encryption needs
DDH for its ciphertexts to hide the plaintext.
\end{remark}

\subsection{Generic algorithms and the square-root barrier}
\label{assumptions:subsec:generic}

How hard can the discrete logarithm problem possibly be? In a group
with no exploitable structure beyond the group law, the best known
algorithms are \emph{generic}: they manipulate group elements only
through multiplication, inversion, and equality testing, never
inspecting the bit-representation of an element. Shoup's idealisation
makes the class precise: a random injective \emph{labelling}
\(\sigma \colon \Z/q\Z \to S\) encodes the element \(g^{k}\) as the
opaque string \(\sigma(k)\), and the algorithm may only ask an oracle
for \(\sigma(i + j)\) given \(\sigma(i)\) and \(\sigma(j)\), and
compare labels for equality. Two generic algorithms achieve
\(O(\sqrt q)\) group operations; a matching lower bound then shows
that, generically, nothing does better.

\subsubsection*{Baby-step giant-step}

\begin{proposition}[Baby-step giant-step, BSGS]
\label{assumptions:prop:bsgs}
In a cyclic group of order \(q\), any discrete logarithm can be
computed in \(O(\sqrt q)\) group operations and \(O(\sqrt q)\) stored
group elements, deterministically and unconditionally.
\end{proposition}

\begin{proof}
Let \(m = \ceil{\sqrt q}\). Since \(m^{2} \ge q\), every
\(x \in \{0, \dots, q - 1\}\) decomposes as
\[
x = i \cdot m + j, \qquad 0 \le i < m, \quad 0 \le j < m,
\]
and the target relation \(g^{x} = h\) becomes
\(g^{j} = h \cdot (g^{-m})^{i}\). \emph{Baby steps:} compute and store
the table \(\{(j, g^{j}) : 0 \le j < m\}\), indexed by group element
in a hash map---\(m - 1\) multiplications. \emph{Giant steps:} set
\(u = g^{-m}\) (one inversion plus one exponentiation, \(O(\log q)\)
operations) and for \(i = 0, 1, 2, \dots\) compute \(h \cdot u^{i}\)
incrementally, looking each value up in the table. The decomposition
guarantees some pair \((i, j)\) matches, yielding \(x = im + j\) after
at most \(m\) giant steps. Total: \(O(\sqrt q)\) operations and
\(O(\sqrt q)\) memory.
\end{proof}

\begin{remark}[Memory is the bottleneck]
\label{assumptions:rem:bsgs-memory}
The \(O(\sqrt q)\) \emph{memory} of BSGS, not its time, is what bites
in practice: at the \(128\)-bit security level no attacker can store
\(\sqrt q \approx 2^{128}\) table entries. Pollard's rho method
matches the time bound with negligible memory, which is why rho,
rather than BSGS, is the realistic generic attack.
\end{remark}

\subsubsection*{Pollard's rho}

\begin{proposition}[Pollard's rho, sketch]
\label{assumptions:prop:rho}
In a cyclic group of prime order \(q\) there is a randomised algorithm
computing discrete logarithms in expected \(O(\sqrt q)\) group
operations and \(O(1)\) storage.
\end{proposition}

\begin{proof}
\emph{Sketch.} Walk pseudo-randomly on \(G\), keeping with each
visited element its known representation \(g^{\alpha} h^{\beta}\).
Partition \(G\) into three roughly equal sets \(S_1, S_2, S_3\) by an
easily computed function of the element's label, and from
\(P = g^{\alpha} h^{\beta}\) step to
\[
f(P) =
\begin{cases}
h \cdot P = g^{\alpha} h^{\beta + 1}, & P \in S_1,\\[2pt]
P^{2} = g^{2\alpha} h^{2\beta}, & P \in S_2,\\[2pt]
g \cdot P = g^{\alpha + 1} h^{\beta}, & P \in S_3,
\end{cases}
\]
each step updating \((\alpha, \beta) \in (\Z/q\Z)^{2}\) explicitly at
\(O(1)\) cost. The sequence eventually cycles---a tail leading into a
loop, the shape of the letter \(\rho\). Modelling \(f\) as a random
function on \(q\) elements, the birthday bound (\emph{Math Guide},
\S``The union bound and a birthday calculation'') produces a collision
\(P_i = P_j\), \(i \ne j\), within \(O(\sqrt q)\) steps in
expectation. A collision gives
\(g^{\alpha_i} h^{\beta_i} = g^{\alpha_j} h^{\beta_j}\), hence
\(g^{\alpha_i - \alpha_j} = h^{\beta_j - \beta_i}\); writing
\(h = g^{x}\),
\[
\alpha_i - \alpha_j \equiv x\,(\beta_j - \beta_i) \pmod q,
\]
and if \(\beta_j \not\equiv \beta_i\) the difference is invertible
modulo the \emph{prime} \(q\), so
\(x \equiv (\alpha_i - \alpha_j)(\beta_j - \beta_i)^{-1} \pmod q\).
Floyd's tortoise-and-hare cycle detection, or Brent's improvement,
finds the collision with \(O(1)\) memory, never storing the
trajectory. The exceptional case \(\beta_j \equiv \beta_i\) occurs
with probability \(O(1/q)\); re-randomising the start point handles
it.
\end{proof}

\begin{remark}[Cost, not time, and constant factors]
\label{assumptions:rem:rho-parallel}
Pollard rho parallelises almost perfectly via the van
Oorschot--Wiener method of \emph{distinguished points}: each of
\(P\) processors walks independently and reports to a shared table
only the points whose label satisfies a cheap predicate, say a run of
leading zero bits; two walks that ever meet coincide from then on and
are detected at the next distinguished point, so the processors
achieve \(\sqrt q / P\) wall-clock time with modest communication.
The security-relevant
measure is therefore the \emph{cost} of the attack---the total
operation count \(\sqrt q\), read as the attacker's budget---not the
time on any single machine. A further constant-factor saving applies
in groups admitting the \(\pm\) automorphism \(P \mapsto -P\), as
elliptic curves do: identifying \(P\) with \(-P\) halves the
effective search space and lowers the expected count from
\(\sqrt{\pi q / 2}\) to \(\sqrt{\pi q / 4}\), a factor \(\sqrt 2\)
that leaves the \(\Theta(\sqrt q)\) scaling untouched. The constant
is the birthday integral: a random walk on \(N\) points is still
collision-free after \(k\) steps with probability
\(\prod_{i < k}(1 - i/N) \approx e^{-k^{2}/2N}\) (the product in the
birthday calculation cited above), and summing over \(k\) gives an
expected first-collision index of
\(\int_{0}^{\infty} e^{-k^{2}/2N}\,dk = \sqrt{\pi N / 2}\).
\end{remark}

\subsubsection*{The generic lower bound}

The two attacks are not merely the best known; up to constants they
are the best \emph{possible} for any generic algorithm. This is
Shoup's theorem, the formal statement of the square-root barrier.

\begin{theorem}[Shoup's generic lower bound]
\label{assumptions:thm:shoup}
Let \(q\) be prime. Any generic algorithm that makes at most \(n\)
queries to the group oracle outputs the discrete logarithm of a
uniformly random challenge with probability at most
\[
\frac{\binom{n+2}{2} + 1}{q} \;=\; O\!\left(\frac{n^{2}}{q}\right).
\]
In particular, constant success probability requires
\(n = \Omega(\sqrt q)\) queries.
\end{theorem}

\begin{proof}
\emph{Sketch---the linear-polynomial technique.} In the idealisation,
the challenge logarithm is an indeterminate \(X\) drawn uniformly from
\(\Z/q\Z\). By induction on the queries, every element the algorithm
forms is \(\sigma\) of a \emph{known linear polynomial}
\(a_i + b_i X \bmod q\) whose coefficients the algorithm chose. Treat
\(X\) formally: as long as the polynomials produced are pairwise
distinct as functions of \(X\), the labels seen are consistent with
\(X\) taking any value, so the algorithm has learnt nothing and can
only guess, succeeding with probability \(1/q\). Information leaks
only through an accidental collision: two distinct polynomials
agreeing at the secret \(x\), that is,
\((a_i - a_j) + (b_i - b_j)\,x \equiv 0 \pmod q\). For a fixed pair
this nonzero linear equation has at most one root modulo the prime
\(q\)---the degree-one case of the Schwartz--Zippel lemma
(\emph{Math Guide}, \S``The Schwartz--Zippel lemma'')---so it holds
with probability at most \(1/q\) over uniform \(x\). With \(n\)
queries plus the two inputs \(g\) and \(h\) there are at most
\(\binom{n+2}{2}\) pairs, and the union bound (\emph{Math Guide},
\S``The union bound and a birthday calculation'') caps the collision
probability at \(\binom{n+2}{2}/q\). Adding the \(1/q\) guess yields
the bound; solving \(n^{2}/q = \Omega(1)\) gives
\(n = \Omega(\sqrt q)\).
\end{proof}

\begin{corollary}[Why \(\log_2 q \ge 2\lambda\)]
\label{assumptions:cor:bitlength}
To force every generic attacker to spend at least \(2^{\lambda}\)
group operations---\(\lambda\) bits of security against generic
attacks---the prime group order must satisfy
\[
\sqrt q \;\gtrsim\; 2^{\lambda}
\qquad\Longleftrightarrow\qquad
\log_2 q \;\ge\; 2\lambda.
\]
\end{corollary}

\begin{proof}
Propositions~\ref{assumptions:prop:bsgs}
and~\ref{assumptions:prop:rho} show that \(\Theta(\sqrt q)\)
operations suffice for a generic attacker, and
Theorem~\ref{assumptions:thm:shoup} shows no generic attacker does
asymptotically better. Equating the attacker's cost \(\sqrt q\) with
the target work factor \(2^{\lambda}\) gives
\(q \approx 2^{2\lambda}\); demanding at least this much is
\(\log_2 q \ge 2\lambda\).
\end{proof}

\begin{example}[The \(128\)-bit target]
\label{assumptions:exa:128}
For \(\lambda = 128\) the corollary demands a prime group order with
\(\log_2 q \ge 256\). This explains the standardised \(256\)-bit
elliptic curves: NIST P-256 (secp256r1) has a \(256\)-bit field and a
\(256\)-bit group order, so \(\sqrt q \approx 2^{128.0}\). The Pasta
curves of \S\ref{assumptions:subsec:pasta}, with \(255\)-bit orders
just above \(2^{254}\) and about \(126\) bits of security, sit just
below this class. Contrast the multiplicative group
\(\Fp^{\times}\): there the sub-exponential index calculus of
Remark~\ref{assumptions:rem:index-calculus} forces \(\log_2 p\) into
the thousands of bits for the same \(\lambda\). Elliptic curves earn
their keep precisely because no sub-exponential, generic-beating
algorithm is known for them.
\end{example}

\begin{remark}[Why the barrier is special to elliptic curves]
\label{assumptions:rem:index-calculus}
The square-root barrier governs \emph{generic} groups only; a
concrete representation can leak structure that beats it. The
multiplicative group \(\Fp^{\times}\) is the cautionary case:
\emph{index calculus} exploits the factorisation of integers into
small primes---smooth numbers---to solve DLP in sub-exponential time
\[
\exp\!\bigl(O\bigl((\log p)^{1/3} (\log\log p)^{2/3}\bigr)\bigr),
\]
far below \(\sqrt p\). No analogue is known for the points of a
well-chosen elliptic curve over \(\Fp\): there is no useful notion of
a ``small'' point over which to factor. This non-generic gap is why
elliptic-curve groups can live at \(256\) bits where \(\Fp^{\times}\)
must reach about \(3072\) bits for the same security. It also means
the rule \(\log_2 q \ge 2\lambda\) protects against generic attacks
\emph{only}; a curve must additionally be screened against the
special-purpose attacks of
Definition~\ref{assumptions:def:secure-curve}.
\end{remark}

\subsection{Pohlig--Hellman and the necessity of large prime order}
\label{assumptions:subsec:pohlig}

The square-root barrier priced the attack at \(\sqrt q\) for
\emph{prime} \(q\). If the group order factors, the discrete
logarithm splits into independent pieces the size of the factors, and
the difficulty collapses to that of the \emph{largest prime} dividing
the order. This is the Pohlig--Hellman reduction---the reason
Definition~\ref{assumptions:def:dlp} insisted on prime order.

\begin{theorem}[Pohlig--Hellman]
\label{assumptions:thm:pohlig}
Let \(G = \langle g\rangle\) have order
\(N = \prod_{i=1}^{r} p_i^{e_i}\) with the \(p_i\) distinct primes.
Then a discrete logarithm in \(G\) can be computed in
\[
O\!\left(\sum_{i=1}^{r} e_i \bigl(\log N + \sqrt{p_i}\,\bigr)\right)
\]
group operations. The term \(\sqrt{p_{\max}}\) dominates, where
\(p_{\max}\) is the largest prime dividing \(N\).
\end{theorem}

\begin{proof}
We compute \(x = \log_g h\) modulo each prime power \(p_i^{e_i}\) and
recombine by the Chinese Remainder Theorem, the moduli being pairwise
coprime with product \(N\).

\emph{Reduction to prime-power order.} Fix \(p = p_i\), \(e = e_i\),
and put \(g_i = g^{N/p^{e}}\), \(h_i = h^{N/p^{e}}\). Then \(g_i\) has
order exactly \(p^{e}\), and
\(\log_{g_i} h_i \equiv x \pmod{p^{e}}\): raising to \(N/p^{e}\) kills
the part of \(x\) coprime to \(p\)'s power and leaves
\(x \bmod p^{e}\). It therefore suffices to solve DLP in
\(\langle g_i\rangle\) of order \(p^{e}\).

\emph{Reduction to prime order, digit by digit.} Write the unknown in
base \(p\):
\[
x \equiv x_0 + x_1 p + \dots + x_{e-1} p^{\,e-1} \pmod{p^{e}},
\qquad 0 \le x_k < p.
\]
Let \(\gamma = g_i^{\,p^{e-1}}\), an element of order exactly \(p\).
Raising \(h_i\) to the power \(p^{e-1}\) gives
\[
h_i^{\,p^{e-1}} = g_i^{\,x p^{e-1}} = \gamma^{\,x_0},
\]
because every higher digit contributes a multiple of \(p^{e}\) to the
exponent and annihilates \(\gamma\). Thus \(x_0\) is a single
discrete logarithm in a group of prime order \(p\), solvable in
\(O(\sqrt p)\) operations by BSGS or rho
(Propositions~\ref{assumptions:prop:bsgs}
and~\ref{assumptions:prop:rho}). Having \(x_0\), set
\(h_i^{(1)} = h_i \cdot g_i^{-x_0}\); its logarithm is divisible by
\(p\), and the same step at exponent \(p^{e-2}\) yields \(x_1\).
Iterating recovers all \(e\) digits with \(e\) prime-order DLPs of
cost \(O(\sqrt p)\) each, plus \(O(e \log N)\) operations of
exponentiation. Summing over \(i\) and applying the CRT gives the
stated bound.
\end{proof}

\begin{example}[A composite-order trap]
\label{assumptions:exa:smooth}
Suppose, contrary to good practice, a group has order
\(N = 2^{256} - 1\). The order factors completely into
\begin{align*}
2^{256} - 1 ={}& 3 \cdot 5 \cdot 17 \cdot 257 \cdot 641 \cdot 65537
  \cdot 274177 \cdot 6700417 \cdot 67280421310721 \\
&\cdot 59649589127497217 \cdot 5704689200685129054721,
\end{align*}
a product of small primes---it is \emph{smooth}. Despite \(N\) having
\(256\) bits, the largest prime factor is only about \(2^{72.3}\),
far below \(2^{128}\), so Pohlig--Hellman with rho in the largest
subgroup solves DLP in roughly
\(\sqrt{p_{\max}} \approx 2^{36.1}\) operations---about \(2^{36}\),
nowhere near the apparent \(128\)-bit level. A \(256\)-bit order is
necessary but \emph{not} sufficient: it must be (almost) prime.
\end{example}

\begin{corollary}[Design rule for the group order]
\label{assumptions:cor:large-prime}
For DLP-based security at level \(\lambda\), the group order must be
divisible by a prime \(q\) with \(\log_2 q \ge 2\lambda\), and the
cryptographic group must be, or be restricted to, the subgroup of
order \(q\). When the ambient group has order \(N = h \cdot q\) with
small \emph{cofactor} \(h\) and large prime \(q\), one works in the
order-\(q\) subgroup; clearing the cofactor---multiplying incoming
elements by \(h\), or checking subgroup membership---prevents
small-subgroup attacks that would otherwise leak \(x \bmod h\) via
Pohlig--Hellman on the order-\(h\) part.
\end{corollary}

\begin{remark}[Prime order, revisited]
\label{assumptions:rem:cofactor}
Here is the deeper justification for ``prime order'' in
Definition~\ref{assumptions:def:dlp}. Prime order \(q\) makes
Pohlig--Hellman vacuous (the only factor is \(q\) itself), makes
every non-identity element a generator (no element lands in a weak
proper subgroup), and removes cofactor bookkeeping entirely. The
Pasta curves below are engineered to have prime order: each group is
its own prime-order subgroup, with cofactor \(1\).
\end{remark}

\subsection{Instantiation on elliptic curves; the Pasta curves}
\label{assumptions:subsec:pasta}

We now fix the concrete groups the remainder of the volume uses.
Recall from the \emph{Math Guide} (\S``Elliptic curves'', with the
group law in \S``The group structure of \(E(\Fp)\)'') that an elliptic
curve over \(\Fp\), for \(p > 3\) prime, in short Weierstrass form is
\[
E \colon y^{2} = x^{3} + ax + b,
\qquad a, b \in \Fp, \quad 4a^{3} + 27b^{2} \ne 0,
\]
and that its set of \(\Fp\)-points,
\[
E(\Fp) = \bigl\{ (x, y) \in \Fp^{2} : y^{2} = x^{3} + ax + b \bigr\}
  \cup \{\mathcal O\},
\]
equipped with chord-and-tangent addition and the point at infinity
\(\mathcal O\) as identity, forms a finite abelian group. The discrete
logarithm problem in this group---given \(P\) and \(Q = [x]P\), find
\(x\)---is the \emph{elliptic-curve discrete logarithm problem}
(ECDLP), stated in its own habitat in the \emph{Math Guide}
(\S``The elliptic-curve discrete logarithm problem''); CDH and DDH
transcribe verbatim with multiplicative notation replaced by the
additive scalar multiplication \(x \mapsto [x]P\), giving ECCDH and
ECDDH. The group size is pinned down by the Hasse bound, likewise
recalled.

\begin{theorem}[Hasse bound, recalled]
\label{assumptions:thm:hasse}
For \(E\) over \(\Fp\) (\emph{Math Guide}, \S``The group structure of
\(E(\Fp)\)''),
\[
\bigl| \#E(\Fp) - (p + 1) \bigr| \;\le\; 2\sqrt p,
\qquad\text{i.e.}\qquad
\#E(\Fp) = p + 1 - t, \quad \abs{t} \le 2\sqrt p,
\]
where \(t\) is the \emph{trace of Frobenius}. The group order is thus
\(p + 1\) up to a deviation of order \(\sqrt p\).
\end{theorem}

The parameter recipe for \(\lambda\)-bit security follows: choose
\(p\) of about \(2\lambda\) bits, so that
\(\#E(\Fp) \approx p \approx 2^{2\lambda}\), and insist by
Corollary~\ref{assumptions:cor:large-prime} that \(\#E(\Fp)\) be
prime---cofactor \(1\). One then assumes ECDLP, ECCDH, and ECDDH at
the level Corollary~\ref{assumptions:cor:bitlength} predicts,
\emph{provided} the curve avoids the known structural attacks that
beat the generic bound. The screening criteria are as follows.

\begin{definition}[Secure-curve criteria]
\label{assumptions:def:secure-curve}
A curve \(E/\Fp\) of prime order \(q = \#E(\Fp)\) is admissible for
\(\lambda\)-bit DLP security if:
\begin{enumerate}
\item (\emph{Size / generic.}) The order satisfies
  \(\log_2 q \ge 2\lambda\), so Pollard rho costs at least
  \(2^{\lambda}\) operations
  (Corollary~\ref{assumptions:cor:bitlength}).
\item (\emph{Prime order.}) The order \(q\) is prime, so
  Pohlig--Hellman is vacuous and the cofactor is \(1\)
  (Corollary~\ref{assumptions:cor:large-prime}).
\item (\emph{Anti-MOV / large embedding degree.}) The multiplicative
  order \(k\) of \(p\) modulo \(q\)---the \emph{embedding
  degree}---is large. The MOV/Frey--R\"uck attack uses a pairing to
  embed \(E(\Fp)\) into \(\F_{p^{k}}^{\times}\) and runs the
  sub-exponential index calculus of
  Remark~\ref{assumptions:rem:index-calculus} there; the threat is
  real only for small \(k\), as on \emph{supersingular} curves---over
  a prime field \(p > 3\), those with trace \(t = 0\), so that
  \(\#E(\Fp) = p + 1\) divides \(p^{2} - 1\) and \(k \le 2\). A curve
  with \(t \ne 0\) is \emph{ordinary}.
\item (\emph{Anti-Smart / not anomalous.}) The order avoids
  \(q = p\): on an \emph{anomalous} curve, \(\#E(\Fp) = p\) with
  trace \(t = 1\), the Smart--Satoh--Araki--Semaev attack solves
  ECDLP in polynomial time.
\item (\emph{Twist security.}) The \emph{quadratic twist} \(E'\)---the
  curve \(d y^{2} = x^{3} + ax + b\) for a fixed nonsquare
  \(d \in \Fp\); every \(x \in \Fp\) is the \(x\)-coordinate of a
  point of \(E\) or of \(E'\)---should also have large prime order,
  so that computation cannot be pushed onto a weak twist by a fault
  or an invalid-point injection.
\end{enumerate}
\end{definition}

\begin{definition}[The Pasta curves]
\label{assumptions:def:pasta}
The \emph{Pasta} curves are a pair, \emph{Pallas} and \emph{Vesta},
in short Weierstrass form
\[
y^{2} = x^{3} + 5
\]
(so \(a = 0\), \(b = 5\), a \(j\)-invariant-\(0\) shape), defined over
two \(255\)-bit primes \(p\) and \(q\) arranged in a \emph{cycle}:
\[
\#\mathrm{Pallas} = q
  \ \ (\text{Pallas defined over } \Fp),
\qquad
\#\mathrm{Vesta} = p
  \ \ (\text{Vesta defined over } \Fq).
\]
The order of each curve equals the field of definition of the other.
Both primes are congruent to \(1 \bmod 2^{32}\), giving their
multiplicative groups large \(2\)-adic valuation for fast FFTs, and
each curve has cofactor \(1\): the whole group is of prime order. The
primes themselves are assembled digit by digit in the
\emph{Math Guide} (\S``Pallas and Vesta assembled''), and the
field/scalar bookkeeping of the cycle in \S``Base fields, scalar
fields, and the Pasta cycle'' there.
\end{definition}

\begin{remark}[Why a cycle, and what it buys]
\label{assumptions:rem:cycle}
The defining property---\(\#\mathrm{Pallas} = q\), the base field of
Vesta, and \(\#\mathrm{Vesta} = p\), the base field of Pallas---makes
(Pallas, Vesta) an \emph{amicable}, or cyclic, pair: one curve's
scalar field is the other's coordinate field. A proof system can
therefore use Pallas-arithmetic to verify statements about
Vesta-arithmetic and vice versa without non-native field emulation;
this is the mechanism behind recursive proof composition in Halo~2,
whose recursion alternates between the two curves. The cycle is an
\emph{efficiency} property, not a security one: it does not enter the
hardness assumptions, and each curve is screened independently
against Definition~\ref{assumptions:def:secure-curve}.
\end{remark}

\begin{proposition}[The Pasta curves against the criteria]
\label{assumptions:prop:pasta-secure}
The curves Pallas and Vesta satisfy criteria (1)--(4) of
Definition~\ref{assumptions:def:secure-curve} at
\(\lambda \approx 126\). Criterion (5) fails for both curves---neither
twist has prime order, and the Pallas twist falls below the
SafeCurves \(2^{100}\) twist-attack cost threshold---but the failure
is rendered moot by mandatory point validation.
\end{proposition}

\begin{proof}
\emph{Sketch, criterion by criterion.}

(1) Each order is a \(255\)-bit prime just above \(2^{254}\), so
\(\log_2 q \approx 254.0\) and the bare generic cost is
\(\sqrt q \approx 2^{127.0}\). Plain Pollard rho expects about
\(\sqrt{\pi q / 2} \approx 2^{127.3}\) operations; the \(\pm\)
automorphism speedup of Remark~\ref{assumptions:rem:rho-parallel}
lowers this to \(\sqrt{\pi q / 4} \approx 2^{126.8}\); and the
order-\(6\) automorphism group of these \(j\)-invariant-\(0\)
curves---the factor \(\sqrt m\) for an order-\(m\) automorphism group
of Remark~\ref{assumptions:rem:j0}, here an extra \(\sqrt 3\)---lowers
it to \(\sqrt{\pi q / 12} \approx 2^{126.0}\). Hence ``about \(126\)
bits'' of security, the same range as the \(128\)-bit curves of
Example~\ref{assumptions:exa:128}.

(2) Both orders are prime by construction: cofactor \(1\),
Pohlig--Hellman vacuous.

(3) Both curves are ordinary: the trace \(t = p + 1 - q\) is odd,
so \(t \ne 0\). The embedding degree \(k\) of each is enormous. For
Pallas, \(q - 1 = 2^{32} \cdot 3^{2} \cdot 1709 \cdot 24859 \cdot C\)
with \(C\) a \(194\)-bit composite, and
\(p^{(q-1)/32} \not\equiv 1\) and \(p^{(q-1)/\ell} \not\equiv 1
\pmod q\) for \(\ell \in \{3, 1709, 24859\}\), so the order \(k\) of
\(p\) modulo \(q\) is a multiple of
\(2^{28} \cdot 3^{2} \cdot 1709 \cdot 24859 > 2^{56}\); for Vesta,
\(p - 1 = 2^{32} \cdot 3 \cdot 463 \cdot C'\), and
\(q^{(p-1)/4} \not\equiv 1\), \(q^{(p-1)/\ell} \not\equiv 1 \pmod p\)
for \(\ell \in \{3, 463\}\) make \(k\) a multiple of
\(2^{31} \cdot 3 \cdot 463 > 2^{41}\).
Either extension field is astronomically beyond any feasible
index-calculus target; MOV/Frey--R\"uck does not apply.

(4) Neither curve is anomalous: \(\#\mathrm{Pallas} = q \ne p\) and
\(\#\mathrm{Vesta} = p \ne q\), so the trace \(t = p + 1 - q \ne 1\)
and the Smart attack does not apply.

(5) Twist security is the criterion the pair does not meet. The Pasta
search deliberately ignored it---the curves were generated with the
search flag \texttt{--ignoretwist} of the \texttt{amicable.sage}
search recorded in the README of the \texttt{zcash/pasta}
repository---and neither twist has prime order. Each \(x \in \Fp\)
supplies two affine points in total, on \(E\) or on \(E'\) (one on
each when \(x^{3} + b = 0\)), so \(\#E + \#E' = 2(p + 1)\) and
\(\#E' = 2(p + 1) - \#E = p + 1 + t\); hence
\(\#\mathrm{Pallas}' = 2(p + 1) - q\) and
\(\#\mathrm{Vesta}' = 2(q + 1) - p\), which factor as
\begin{align*}
\#\mathrm{Pallas}' &= 3^{2} \cdot 7 \cdot 8191 \cdot 85021
  \cdot 11540602760389 \cdot P_{176},\\
\#\mathrm{Vesta}' &= 3^{3} \cdot 7 \cdot 13 \cdot 2851 \cdot 30097
  \cdot P_{217},
\end{align*}
with \(P_{176}\) a \(176\)-bit prime and \(P_{217}\) a \(217\)-bit
prime. The best attack on the Pallas twist therefore costs about
\(\sqrt{P_{176}} \approx 2^{87.6}\), below the SafeCurves \(2^{100}\)
twist-security threshold; the Vesta twist clears the threshold at
about \(2^{108.2}\). The failure is moot in our setting: twist
attacks require an implementation that accepts an attacker-supplied
\(x\)-coordinate without validation (an \(x\)-only ladder) or that
skips its on-curve checks, and the deployed arithmetic does neither.
The \texttt{pasta\_curves} crate carries points in full coordinates
and has no \(x\)-only code path; its deserialisation
(\texttt{from\_bytes}) reconstructs \(y\)
by taking the square root of \(x^{3} + 5\), failing exactly when
\(x\) lies on the twist, where \(x^{3} + 5\) is a non-square, and
even \texttt{from\_bytes\_unchecked} delegates to the checked path
because a compressed encoding cannot skip the curve check, while
points built from raw coordinates are gated on
\texttt{is\_on\_curve}. This is the behaviour protocol specification
\href{https://zips.z.cash/protocol/protocol.pdf\#pallasandvesta}{\S~5.4.9.6} mandates: its decoding map returns \(\bot\) (failure) when
\(x^{3} + b\) has no square root, and an implementation must not
assume that the root exists or that the encoding lies on the curve.
Orchard deserialises its points through this
routine, and cofactor \(1\) leaves no
small subgroup on the curve itself. Invalid-curve and twist-point
injection are therefore excluded by \emph{validation}, not by twist
structure.

The standard generic analysis governs, and the security level is
about \(126\) bits.
\end{proof}

\begin{remark}[\(j\)-invariant \(0\) and endomorphisms]
\label{assumptions:rem:j0}
The shape \(y^{2} = x^{3} + b\) gives \(j\)-invariant \(0\), and
with it an extra automorphism of order \(3\) (\emph{Math Guide},
\S``The \(j\)-invariant and the GLV endomorphism''): for \(\omega\) a
primitive cube root of unity in \(\Fp\), the cheap endomorphism
\(\phi(x, y) = (\omega x, y)\) acts on the group as multiplication
by a fixed scalar \(\zeta\) with
\(\zeta^{2} + \zeta + 1 \equiv 0 \pmod q\) (the \(\lambda\) of the
\emph{Math Guide}'s proposition), enabling the GLV speedup for honest
scalar multiplication. The same endomorphism hands the attacker a
slightly larger automorphism group, and with it the constant-factor
Pollard-rho speedup used in
Proposition~\ref{assumptions:prop:pasta-secure}---a factor \(\sqrt m\)
for an order-\(m\) automorphism group---but that is a small constant
and does not change the \(\Theta(\sqrt q)\) scaling. Beyond that
constant, the extra automorphism is not known to weaken ECDLP; the
danger lies only in small embedding degree and in anomalousness,
both excluded above.
\end{remark}

\subsection{The quantum caveat: Shor's algorithm}
\label{assumptions:subsec:quantum}

Every assumption above holds the adversary to classical PPT
computation. Against a large fault-tolerant quantum computer the
picture changes qualitatively, and honesty requires stating exactly
what breaks and what survives.

\begin{theorem}[Shor, informal; stated without proof]
\label{assumptions:thm:shor}
There is a quantum algorithm that, given a cyclic group
\(G = \langle g\rangle\) of order \(q\) with efficient group
operations and a target \(h = g^{x}\), outputs \(x\) in time
polynomial in \(\log q\), succeeding with overwhelming probability.
Consequently DLP---and with it CDH and DDH---falls in quantum
polynomial time in \emph{any} group, elliptic-curve groups included.
\end{theorem}

We do not prove the theorem; the shape of the algorithm is
nevertheless worth recording. Discrete logarithm is an instance of
the \emph{hidden subgroup problem} over \((\Z/q\Z)^{2}\). The
function
\(f(\alpha, \beta) = g^{\alpha} h^{-\beta} = g^{\alpha - x\beta}\)
is constant exactly on the cosets of
\[
H = \bigl\{ (\alpha, \beta) : \alpha \equiv x\beta \pmod q \bigr\}
  = \langle (x, 1) \rangle \subseteq (\Z/q\Z)^{2},
\]
a line of slope \(x\). The algorithm prepares a uniform superposition
over \((\Z/q\Z)^{2}\), evaluates \(f\) into a second register, and
applies the quantum Fourier transform---the reversible change of
basis into the characters of \((\Z/q\Z)^{2}\), the homomorphisms
\(\chi_{u,v}(\alpha, \beta) = e^{2\pi i (u\alpha + v\beta)/q}\) into
the unit circle; measurement then samples characters trivial on
\(H\), i.e.\ pairs \((u, v)\) with \(u x + v \equiv 0 \pmod q\), and
a handful of samples reveal the slope \(x\) by classical linear
algebra modulo \(q\). The
whole procedure uses \(O(\poly(\log q))\) qubits and gates, and it
touches the group only through the black-box map
\((\alpha, \beta) \mapsto g^{\alpha} h^{-\beta}\), so it is agnostic
to representation: elliptic-curve DLP falls exactly as
\(\Fp^{\times}\) DLP does, and integer factoring falls to the
order-finding variant.

\begin{remark}[No contradiction with the square-root barrier]
\label{assumptions:rem:no-contradiction}
Shoup's lower bound (Theorem~\ref{assumptions:thm:shoup}) assumes the
generic model for \emph{classical} algorithms making sequential,
individual oracle queries. Shor's algorithm is non-generic in exactly
the relevant sense: it queries the group function in superposition
and exploits the periodicity the Fourier transform exposes, abilities
the classical generic model does not grant. The two results coexist:
the barrier bounds classical generic attackers; Shor is a quantum,
non-generic attacker.
\end{remark}

\begin{remark}[What Shor breaks and what survives]
\label{assumptions:rem:shor-scope}
Shor's algorithm breaks the entire discrete-log family and integer
factoring in polynomial time: RSA, finite-field Diffie--Hellman, and
elliptic-curve cryptography are all quantum-broken. What it does
\emph{not} break:
\begin{itemize}
\item \emph{Symmetric primitives and hash functions.} Grover's
  algorithm searches a space of size \(2^{n}\) in about \(2^{n/2}\)
  quantum queries---a quadratic, not exponential, speedup---so a
  cipher with a \(2\lambda\)-bit key, or an \(n = 2\lambda\)-bit hash
  against preimages, retains \(\lambda\) bits of \emph{quantum}
  security. Generic quantum collision finding instead costs
  \(\Theta(2^{n/3})\) queries, so in the query metric an output of
  about \(3\lambda\) bits is needed for \(\lambda\)-bit quantum
  collision security---though the algorithm needs comparably much
  quantum memory, and whether it beats classical parallel collision
  search in realistic cost models is debated. Doubling therefore
  suffices for key search and preimages, but for collisions only
  under that contested cost reading.
% Deviation from brief (y2-42): the brief's blanket "doubling
% key/output lengths suffices" fails for collisions in the query
% metric (Theta(2^{n/3}) generic quantum collision finding); narrowed
% to key search and preimages, with the collision bound's contested
% cost realism noted.
\item \emph{Post-quantum assumptions.} Hardness assumptions based on
  neither discrete logarithms nor factoring are not known to fall to
  Shor and underlie post-quantum cryptography.
\end{itemize}
The discrete-log assumptions of this volume are therefore adopted
under the standing premise that no cryptographically relevant quantum
computer exists. Within that premise, and against all classical
attacks, the Pasta curves deliver the approximately \(126\)-bit
security of Proposition~\ref{assumptions:prop:pasta-secure}---the
ground on which the Halo~2 constructions to come are built.
\end{remark}

\subsection{Standing assumptions}
\label{assumptions:subsec:standing}

\begin{remark}[Standing assumptions for the sequel]
\label{assumptions:rem:standing}
We collect the premises the rest of the volume relies on. For each
Pasta curve \(C \in \{\mathrm{Pallas}, \mathrm{Vesta}\}\) with prime
order \(q_C\) and generator \(g_C\):
\begin{enumerate}
\item ECDLP is hard in \(\langle g_C \rangle\): no classical PPT
  adversary finds \(x\) from \([x]g_C\) with non-negligible
  probability---concretely, every known attack costs at least
  \(2^{126}\) operations
  (Proposition~\ref{assumptions:prop:pasta-secure}).
\item ECCDH and, where needed, ECDDH hold in \(\langle g_C \rangle\)
  (Definitions~\ref{assumptions:def:cdh}
  and~\ref{assumptions:def:ddh}, transcribed additively).
\item The adversary is classical and probabilistic polynomial-time
  (Remark~\ref{assumptions:rem:shor-scope}).
\end{enumerate}
By the hierarchy of \S\ref{assumptions:subsec:hierarchy}, assuming
ECDDH is the strongest commitment and implies the rest; each later
construction invokes, by name, the weakest assumption it requires.
\end{remark}

\section{Hash functions and the random oracle model}\label{sec:hash}

The adversary of this section forges by \emph{coincidence}. She
prepares two contracts---one benign, one ruinous---whose digests
agree, has the benign one signed, and produces the ruinous one in
court; or she takes a digest minted for one purpose, say a
key-derivation output, and replays it in another role, say a
commitment, exploiting the fact that both are interchangeable bit
strings; or she simply walks a pseudo-random path through the digest
space until, square-root-of-the-space steps later, two inputs collide.
Denying her these coincidences is the business of the hash function: a
function compressing arbitrary data into a short fixed-length tag. It
underlies digital signatures, commitment schemes, message
authentication, key derivation, proof-of-work, Merkle trees, and the
Fiat--Shamir transform behind Halo~2's non-interactive arguments---so
many purposes for one primitive because an \emph{ideal} hash acts like
a deterministic source of fresh randomness keyed by its input. This
section builds the primitive from the ground up: its syntax; the three
classical security notions and the generic birthday attack that bounds
them; the random oracle model; domain separation; hashing into a
finite field and onto an elliptic curve; the arithmetisation-friendly
hash Poseidon that Orchard evaluates inside its circuit; and finally a
register of the transparent layer's SHA-256-derived hashes.

\subsection{Syntax}\label{hash:subsec:syntax}

Notation for strings is fixed once. The set of all finite binary
strings is \(\{0,1\}^{*}\); the strings of length exactly \(n\) form
\(\{0,1\}^{n}\); the length of a string \(x\) is \(\abs{x}\); and
concatenation is written \(x \,\|\, y\).

\begin{definition}[Hash function]\label{hash:def:hashfn}
A \emph{hash function} with output length \(n \in \N\) is a function
\[
H \colon \{0,1\}^{*} \longrightarrow \{0,1\}^{n}.
\]
Elements of the codomain are called \emph{digests}, \emph{hashes}, or
\emph{tags}. The function is \emph{compressing} on a domain
\(D \subseteq \{0,1\}^{*}\) if \(D\) contains strings longer than
\(n\); when the domain is all of \(\{0,1\}^{*}\), the function is
overwhelmingly many-to-one.
\end{definition}

A single fixed function such as SHA-256 is what practice deploys, but
a foundational wrinkle appears the moment one tries to say what its
security \emph{means}. For any fixed compressing \(H\) there
\emph{exist} inputs \(x \neq x'\) with \(H(x) = H(x')\)---the
pigeonhole principle guarantees it---and therefore a tiny program that
simply outputs one such hardcoded pair also exists. The asymptotic
statement ``no efficient adversary finds a collision'' is consequently
\emph{false} for every fixed compressing function: the hardcoding
program is efficient and always succeeds. (Writing that program down
is another matter---nobody knows a collision for SHA-256---but
existence suffices to falsify the universally quantified claim.) The
standard remedy indexes the function by a key.

\begin{definition}[Keyed hash family]\label{hash:def:keyed}
A \emph{keyed hash family} (or \emph{hash function ensemble}) is a
function
\[
H \colon \mathcal{K} \times \{0,1\}^{*} \longrightarrow \{0,1\}^{n},
\]
where \(\mathcal{K}\) is a finite key space. A key
\(s \in \mathcal{K}\) is sampled \emph{publicly} at random, fixed, and
we write \(H_{s}(\cdot) := H(s, \cdot)\) for the family member in
force.
\end{definition}

The key is \emph{not} secret. Its role is to index which member of the
family the adversary must attack: a hardcoded collision for one member
is useless against a freshly sampled other, so ``no efficient
adversary finds a collision for a random member'' is a meaningful
asymptotic statement.

\begin{remark}[Convention: keyed definitions, unkeyed deployments]
\label{hash:rem:keyed}
The convention of this volume follows the standard resolution of the
mismatch between theory and practice. Security notions are stated for
the keyed family, where the definitions are clean; every concrete
construction we deploy is unkeyed; and the key is mentally identified
with ``the choice of standard''---the sampling of \(s\) happened, once
and in public, when the function was designed. No formal theorem
bridges the two readings; the gap is the same one the random oracle
model of \S\ref{hash:subsec:rom} inhabits, and we flag it rather than
hide it.
\end{remark}

\begin{definition}[Extendable-output function]\label{hash:def:xof}
An \emph{extendable-output function} (XOF) is a primitive producing a
digest of caller-chosen length: syntactically a function
\[
H \colon \{0,1\}^{*} \times \N \longrightarrow \{0,1\}^{*},
\qquad \abs{H(x, \ell)} = \ell,
\]
whose outputs are \emph{prefix-consistent}: for
\(\ell \le \ell'\), the string \(H(x, \ell)\) is a prefix of
\(H(x, \ell')\), so a longer request extends a shorter one as a
stream.
\end{definition}

The canonical XOFs are SHAKE128 and SHAKE256 from the SHA-3 family.
An XOF is the natural tool wherever a \emph{controlled number} of
uniform bits is needed rather than a fixed digest---above all when
hashing into a finite field or onto a curve, where the output must be
reduced modulo a prime; \S\ref{hash:subsec:hashtofield} returns to
this.

\subsection{Security notions: preimage, second-preimage, and
collision resistance}\label{hash:subsec:games}

Each of the three classical security notions is a game between a
challenger and a PPT adversary \(\mathcal{A}\), in the template of
Definition~\ref{model:def:game}; the keyed family \(H\) is secure in a
given sense if every PPT adversary wins the corresponding game with
probability negligible in the security parameter \(\lambda\). All
sizes are functions of \(\lambda\): the output length
\(n = n(\lambda)\), the key length, and a polynomially bounded
challenge input length \(m = m(\lambda)\) from which the games sample
their random challenge \(x \xleftarrow{\$} \{0,1\}^{m}\). Recall that
a function \(\varepsilon(\lambda)\) is negligible,
\(\varepsilon \in \negl(\lambda)\), if it decays faster than the
reciprocal of every polynomial (\emph{Math Guide}, \S``Polynomial,
exponential, and negligible functions''; recalled in
\S\ref{model:subsec:indist}). The three notions form a strict
hierarchy, and no single definition of a ``good'' hash exists; which
notion an application needs is a per-application question that
Remark~\ref{hash:rem:which-notion} takes up.

\begin{definition}[Preimage resistance]\label{hash:def:pre}
The game \(\mathsf{PRE}^{\mathcal{A}}_{H}(\lambda)\) runs as follows.
\begin{enumerate}
\item The challenger samples \(s \xleftarrow{\$} \mathcal{K}\) and
  \(x \xleftarrow{\$} \{0,1\}^{m}\), computes \(y := H_{s}(x)\), and
  sends \((s, y)\) to the adversary.
\item The adversary outputs \(x' \in \{0,1\}^{*}\).
\item The game outputs \(1\) iff \(H_{s}(x') = y\).
\end{enumerate}
The family \(H\) is \emph{preimage resistant} (or \emph{one-way}) if
for every PPT \(\mathcal{A}\),
\[
\mathrm{Adv}^{\mathrm{pre}}_{H}(\mathcal{A}, \lambda)
:= \Pr\bigl[\mathsf{PRE}^{\mathcal{A}}_{H}(\lambda) = 1\bigr]
\;\in\; \negl(\lambda).
\]
\end{definition}

Note the winning condition: the adversary, given the key and the
digest of a random message, must output \emph{any} string hashing to
\(y\). Recovering the original \(x\) is not required; any preimage
wins.

\begin{definition}[Second-preimage resistance]\label{hash:def:sec}
The game \(\mathsf{SEC}^{\mathcal{A}}_{H}(\lambda)\) runs as follows.
\begin{enumerate}
\item The challenger samples \(s \xleftarrow{\$} \mathcal{K}\) and
  \(x \xleftarrow{\$} \{0,1\}^{m}\), and sends \((s, x)\) to the
  adversary.
\item The adversary outputs \(x' \in \{0,1\}^{*}\).
\item The game outputs \(1\) iff \(x' \neq x\) and
  \(H_{s}(x') = H_{s}(x)\).
\end{enumerate}
The family \(H\) is \emph{second-preimage resistant} if for every PPT
\(\mathcal{A}\),
\[
\mathrm{Adv}^{\mathrm{sec}}_{H}(\mathcal{A}, \lambda)
:= \Pr\bigl[\mathsf{SEC}^{\mathcal{A}}_{H}(\lambda) = 1\bigr]
\;\in\; \negl(\lambda).
\]
\end{definition}

Here the challenger \emph{gives} the adversary a specific random
message; the adversary must produce a different message colliding with
it.

\begin{definition}[Collision resistance]\label{hash:def:col}
The game \(\mathsf{COL}^{\mathcal{A}}_{H}(\lambda)\) runs as follows.
\begin{enumerate}
\item The challenger samples \(s \xleftarrow{\$} \mathcal{K}\) and
  sends \(s\) to the adversary.
\item The adversary outputs a pair \((x, x')\).
\item The game outputs \(1\) iff \(x \neq x'\) and
  \(H_{s}(x) = H_{s}(x')\).
\end{enumerate}
The family \(H\) is \emph{collision resistant} if for every PPT
\(\mathcal{A}\),
\[
\mathrm{Adv}^{\mathrm{col}}_{H}(\mathcal{A}, \lambda)
:= \Pr\bigl[\mathsf{COL}^{\mathcal{A}}_{H}(\lambda) = 1\bigr]
\;\in\; \negl(\lambda).
\]
\end{definition}

The three games differ in \emph{who chooses what}. In the preimage
game the challenger chooses the message and reveals only its digest;
in the second-preimage game the challenger chooses the message and
reveals it; in the collision game the adversary chooses \emph{both}
messages and need only make them collide. More adversarial choice
means an easier game and hence a stronger security demand on the
function; collision resistance, granting the most choice, is the
strongest of the three.

% TikZ figure placeholder --- fig:hash-hierarchy: implication diagram
% collision resistance => second-preimage resistance => preimage
% resistance (solid arrows down the chain, the second annotated with
% the additive 2^{n-m} loss and the compression hypothesis), plus a
% struck-through backward edge carrying the counterexample of
% Example~\ref{hash:ex:nonreversal} (drop-last-bit H'). Dedicated
% figures pass supplies the drawing.

\begin{proposition}[Hierarchy of notions]\label{hash:prop:hierarchy}
Let \(H\) be a keyed hash family that is compressing, with challenge
input length \(m > n\). Then:
\begin{enumerate}
\item collision resistance implies second-preimage resistance, with no
  loss; and
\item second-preimage resistance implies preimage resistance up to an
  additive loss of \(2^{n-m}\), which is negligible---so the
  implication holds outright---whenever \(H\) is sufficiently
  compressing, \(m - n = \omega(\log \lambda)\).
\end{enumerate}
Neither implication reverses in general.
\end{proposition}

\begin{proof}
(1) From an adversary \(\mathcal{A}\) against second-preimage
resistance, construct \(\mathcal{B}\) against collision resistance: on
input the key \(s\), the reduction \(\mathcal{B}\) samples
\(x \xleftarrow{\$} \{0,1\}^{m}\) \emph{itself}, runs
\(x' \gets \mathcal{A}(s, x)\)---the simulated view is exactly that of
\(\mathsf{SEC}\), since \(\mathcal{B}\) samples \(x\) from the same
distribution as the challenger---and outputs the pair \((x, x')\).
Whenever \(\mathcal{A}\) succeeds, \(x' \neq x\) and
\(H_{s}(x') = H_{s}(x)\), which is precisely a collision. Hence
\(\mathrm{Adv}^{\mathrm{col}}_{H}(\mathcal{B}, \lambda) \ge
\mathrm{Adv}^{\mathrm{sec}}_{H}(\mathcal{A}, \lambda)\): the reduction
is tight, and if the left side is negligible so is the right.

(2) Toward the contrapositive, let \(\mathcal{A}\) be an inverter; we
build a second-preimage finder \(\mathcal{B}\). On input \((s, x)\),
the reduction computes \(y := H_{s}(x)\), runs
\(x' \gets \mathcal{A}(s, y)\)---again a perfect simulation of
\(\mathsf{PRE}\), since \(x\) is uniform and \(y\) is its
digest---and outputs \(x'\). Whenever \(\mathcal{A}\) succeeds,
\(H_{s}(x') = y = H_{s}(x)\); what remains is to argue
\(x' \neq x\) with good probability. Fix \(\mathcal{A}\)'s coins, so
that its output on each digest is a single string
\(x' = \mathcal{A}(s, y)\). An input \(x\) can satisfy
\(x = \mathcal{A}(s, H_{s}(x))\) only if \(x\) is \emph{the} output
assigned to its own digest, so each of the \(2^{n}\) digests
contributes at most one such input, and at most \(2^{n}\) inputs in
all are ``fixed points'' of the composed map. Over uniform
\(x \xleftarrow{\$} \{0,1\}^{m}\), therefore,
\(\Pr[x' = x] \le 2^{n}/2^{m} = 2^{n-m}\); averaging over
\(\mathcal{A}\)'s coins preserves the bound. Hence
\[
\mathrm{Adv}^{\mathrm{sec}}_{H}(\mathcal{B}, \lambda)
\;\ge\;
\mathrm{Adv}^{\mathrm{pre}}_{H}(\mathcal{A}, \lambda) - 2^{n-m}.
\]
The loss is \emph{additive}, not multiplicative: nothing bounds
\(\Pr[x' \neq x]\) \emph{conditioned} on success, because
\(\mathcal{A}\)'s successes might concentrate on small fibres where
returning \(x\) itself is likely. Sufficient compression,
\(m - n = \omega(\log \lambda)\), makes \(2^{n-m}\) negligible and the
implication unconditional.

Non-reversal is witnessed by
Example~\ref{hash:ex:nonreversal} below.
\end{proof}

\begin{example}[Preimage-resistant but not collision-resistant]
\label{hash:ex:nonreversal}
Take any one-way family \(H\) and define \(H'_{s}(x) :=
H_{s}(\bar{x})\), where \(\bar{x}\) deletes the last bit of \(x\).
The family \(H'\) inherits one-wayness: a preimage challenge
\((s, y)\) for \(H\) is distributed exactly as one for \(H'\) at one
bit greater challenge length, and an \(H'\)-inverter's output
\(x''\) satisfies \(H_{s}(\overline{x''}) = y\), so deleting its last
bit hands the reduction a preimage under \(H\). Yet the pair
\(x \,\|\, 0\), \(x \,\|\, 1\) is
an immediate collision for every \(x\), found without any computation.
Hence preimage resistance does not imply collision resistance. The
other separations are constructed similarly; the implications of
Proposition~\ref{hash:prop:hierarchy} are genuinely one-directional.
\end{example}

\begin{remark}[Which application needs which notion]
\label{hash:rem:which-notion}
Password storage---keeping a digest of each password so that a stolen
table does not reveal them---is the natural home of \emph{preimage}
resistance: the thief holds digests, and any preimage lets her log
in. (The notion alone is not the whole story there: passwords are
guessable, so deployments hash each together with a unique random
string, its \emph{salt}, through a deliberately slow function, against
offline guessing.) A signature scheme that signs \(H(m)\), where an
adversary may later try to substitute a forged document for the
honestly chosen \(m\), needs \emph{second-preimage} resistance: the
message is fixed and honest, and the attacker must collide with it. A
scheme in which the \emph{signer herself} may cheat---preparing two
contracts with the same digest, signing the benign one, presenting the
malicious one---needs full \emph{collision} resistance, because there
the adversary controls both messages. Commitment schemes, Merkle
trees, and the Fiat--Shamir transform all fall into this last
category, which is why collision resistance is the headline property
of a cryptographic hash.
\end{remark}

\subsection{The birthday bound on generic collision finding}
\label{hash:subsec:birthday}

How hard can collision finding be? The generic attack---evaluate the
hash on many inputs and wait for a coincidence---succeeds far sooner
than intuition suggests, and its cost is governed by the birthday
calculation of the \emph{Math Guide}.

\begin{lemma}[Birthday bound for hashing]\label{hash:lem:birthday}
Let \(H \colon \{0,1\}^{*} \to \{0,1\}^{n}\) and \(N = 2^{n}\).
Evaluate \(H\) on \(q\) distinct inputs, in the idealised model where
the \(q\) digests behave as independent uniform samples from
\(\{0,1\}^{n}\), and let \(C\) be the event that two digests coincide.
Then
\[
\Pr[C] \;=\; 1 - \prod_{i=0}^{q-1}\Bigl(1 - \frac{i}{N}\Bigr),
\]
with the two-sided estimate
\[
1 - e^{-q(q-1)/(2N)} \;\le\; \Pr[C] \;\le\;
\frac{q(q-1)}{2N} \;\le\; \frac{q^{2}}{2N},
\]
and the lower bound gives \(\Pr[C] \ge 1/2\) once
\(q \ge 1.18\sqrt{N}\). A collision therefore appears with constant
probability after about \(\sqrt{N} = 2^{n/2}\) evaluations.
\end{lemma}

\begin{proof}
The exact product and both estimates are the birthday-bound
proposition of the \emph{Math Guide} (\S``The union bound and a
birthday calculation''), instantiated with set size \(N = 2^{n}\) and
sample count \(k = q\); we do not re-derive them. Only the
half-probability threshold is new. Setting the exponent of the lower
bound equal to \(-\ln 2\) gives \(q(q-1) = 2N \ln 2\), i.e.
\[
q \;\approx\; \sqrt{2 \ln 2}\,\sqrt{N} \;\approx\; 1.18\sqrt{N},
\]
the constant being \(\sqrt{2\ln 2} = 1.1774\ldots\); for
\(q \ge 1.18\sqrt{N}\) the exponential falls below \(1/2\) and
\(\Pr[C] \ge 1/2\).
\end{proof}

\begin{remark}[Reading the birthday bound]\label{hash:rem:birthday}
Three consequences deserve emphasis.

\emph{The square-root law.} An \(n\)-bit digest yields only \(n/2\)
bits of collision security in the sense of
Definition~\ref{model:def:bitsec}. For the conventional \(128\)-bit
collision security one must take \(n = 256\)---which is why SHA-256
(and not ``SHA-128'', which does not exist) is the workhorse digest
size, and why the transcript hashes of Fiat--Shamir-style protocols
must carry at least \(256\)-bit output.

\emph{The attack is practical.} The naive attack stores all
\(q \approx 2^{n/2}\) digests and sorts---\(2^{n/2}\) memory, the same
obstruction that made baby-step giant-step unrealistic in
Remark~\ref{assumptions:rem:bsgs-memory}. But collision finding, like
the discrete-logarithm walk of
Proposition~\ref{assumptions:prop:rho}, admits cycle finding: iterate
\(x \mapsto H(x)\) and detect the cycle with Floyd's or Brent's
algorithm, in the manner of Pollard's rho, or run the van
Oorschot--Wiener parallel collision search across many machines. These
achieve the same \(2^{n/2}\) time in essentially constant memory, so
the birthday attack is a genuine budget item, not a thought
experiment.

\emph{Preimages cost the square.} The best generic attack against
preimage and second-preimage resistance is brute force: about
\(2^{n}\) evaluations to hit a prescribed digest. Those notions
therefore enjoy full \(n\)-bit generic security---twice the collision
level---because the attacker cannot exploit birthday coincidences
among self-chosen pairs; the collision game is the only one that pays
the square-root discount.
\end{remark}

\subsection{The random oracle model}\label{hash:subsec:rom}

The model section introduced the random oracle model as one of the two
idealisations under which proofs proceed
(Definition~\ref{model:def:rom}), stated its two powers informally
(Remark~\ref{model:rem:rom-powers}), and fixed its status as a
heuristic with known limits (Remark~\ref{model:rem:rom-caveat}). With
hash syntax now in hand we can define the idealised object precisely
and \emph{prove} the powers.

\begin{definition}[Random oracle]\label{hash:def:ro}
A \emph{random oracle} with output length \(n\) is a function
\(\mathcal{O} \colon \{0,1\}^{*} \to \{0,1\}^{n}\) drawn uniformly at
random from the set of all such functions. Equivalently---and this is
the form every proof uses---the oracle is realised \emph{lazily}: a
table, initially empty, is maintained; on query \(x\), if \(x\) is in
the table the stored value is returned, and otherwise a fresh
\(\mathcal{O}(x) \xleftarrow{\$} \{0,1\}^{n}\) is sampled, recorded,
and returned.
\end{definition}

Three features characterise the object: outputs on distinct inputs are
independent and uniform; outputs on repeated inputs are consistent;
and access is \emph{only} by querying---the oracle has no description
shorter than its (exponentially long) truth table. A scheme lives in
the random oracle model (ROM) when every party, honest or adversarial,
interacts with a single shared \(\mathcal{O}\) in place of a concrete
hash, and security is proved relative to the random choice of
\(\mathcal{O}\); on deployment a concrete hash---SHA-256, say, or
BLAKE2b---\emph{instantiates} the oracle.

\begin{proposition}[The two powers of the ROM]\label{hash:prop:rom-powers}
In a ROM security proof, the reduction, which simulates
\(\mathcal{O}\) for the adversary, gains two capabilities that no
concrete hash affords.
\begin{enumerate}
\item \emph{Observability (extractability).} The reduction sees every
  input the adversary queries. Since querying is the only way to learn
  anything about \(\mathcal{O}(x)\), if the adversary's output depends
  on \(\mathcal{O}(x)\) then the reduction holds \(x\).
\item \emph{Programmability.} On a not-yet-queried input \(x\), the
  reduction may choose the response \(\mathcal{O}(x)\) to be any value
  of its liking, provided the value is distributed uniformly; the
  simulation remains perfect, and the reduction can thereby
  \emph{plant} a challenge---a value it must invert, say---inside an
  oracle answer.
\end{enumerate}
\end{proposition}

\begin{proof}
Both capabilities are immediate from the lazy realisation of
Definition~\ref{hash:def:ro}, because in the simulation the reduction
\emph{is} the table-keeper. Every query arrives at the reduction
before an answer exists, which is observability. And on a fresh input
the honest table-keeper would sample the answer uniformly and
independently of the adversary's view so far; a reduction that instead
inserts any value \emph{with that same distribution} produces a view
identical to the honest one, while the table keeps repeated queries
consistent. That is programmability.
\end{proof}

\begin{remark}[Why Fiat--Shamir needs the ROM]\label{hash:rem:fs-rom}
The paradigmatic consumer of both powers at once is the Fiat--Shamir
transform---the compiler that renders Halo~2's interactive protocol
non-interactive, developed in this volume under ``From identification
to signature via the Fiat--Shamir transform'' and in full generality
under ``The Fiat--Shamir transform: from interactive to
non-interactive''. The transform replaces the verifier's random
challenge \(e\) by \(e = \mathcal{O}(\text{transcript so far})\).
Soundness of the resulting non-interactive argument is proved by
\emph{rewinding}: the reduction observes (capability~1) the transcript
the prover commits to, then \emph{reprograms} (capability~2) the
oracle to return a fresh independent challenge on that same
transcript, extracting two accepting transcripts that share a prefix
and, from them, a witness---the forking technique previewed in
Remark~\ref{model:rem:forking}. Neither step is available against a
concrete hash, whose outputs the reduction can neither observe
selectively nor rewrite. This is the proof that lives \emph{only} in
the ROM, and the reason the model is indispensable to the
cryptography this monograph targets.
\end{remark}

\subsection{Domain separation and personalisation}
\label{hash:subsec:domainsep}

One physical hash serves many purposes in a deployed protocol: it
derives keys, commits to notes, tags spent notes, seeds Fiat--Shamir
challenges. Digests, however, are interchangeable bit strings, and an
output produced for one purpose can be \emph{replayed} as another
unless the purposes are cryptographically insulated. Domain separation
is the discipline that insulates them.

\begin{definition}[Domain separation]\label{hash:def:domainsep}
Let \(H\) be a hash modelled as a random oracle, and let
\(\{\tau_{i}\}_{i \in I}\) be a set of distinct, \emph{prefix-free}
domain tags (also called labels or personalisation strings), one per
intended use: no \(\tau_{i}\) is a prefix of another---as when all
tags have equal length, or when each is length-prefixed. The
\emph{use-\(i\) hash} is
\[
H^{(i)}(x) \;:=\; H(\tau_{i} \,\|\, x).
\]
\end{definition}

\begin{proposition}[Domain separation yields independent oracles]
\label{hash:prop:domainsep}
If \(H\) is a random oracle and the tags \(\{\tau_{i}\}\) are
prefix-free, then the family \(\{H^{(i)}\}_{i \in I}\) is distributed
exactly as a collection of \emph{independent} random oracles.
Consequently a digest obtained from \(H^{(i)}\) is---except with the
negligible probability of a generic collision---useless against any
use \(j \neq i\).
\end{proposition}

\begin{proof}
Prefix-freeness makes the domains \(\tau_{i} \,\|\, \{0,1\}^{*}\)
pairwise disjoint: a string with prefix \(\tau_{i}\) cannot also have
prefix \(\tau_{j}\) unless one tag is a prefix of the other. Inputs
queried through different \(H^{(i)}\) are therefore never equal. A
random oracle assigns independent uniform values to distinct inputs,
so restricting \(H\) to the disjoint slices yields, per slice, a
uniformly random function, independent across slices---precisely a
collection of independent random oracles. The replay conclusion
follows: the value \(H^{(i)}(x)\) reveals nothing about any
\(H^{(j)}(\cdot)\) with \(j \neq i\).
\end{proof}

Three idioms implement the discipline in deployed practice; Zcash uses
all three.

\begin{enumerate}
\item \emph{Context tagging.} Bind an ASCII string naming the use
  into every call---prepended, as in
  Definition~\ref{hash:def:domainsep}, or in any framing that keeps
  the tag unambiguously delimited from the data. The Orchard tag
  inventory is of this kind: the strings
  \texttt{z.cash:\allowbreak Orchard-cv},
  \texttt{z.cash:\allowbreak Orchard-NoteCommit}, and
  \texttt{z.cash:\allowbreak Orchard-CommitIvk} separate the value
  commitment, note commitment, and key commitment uses of one
  group-hash (protocol specification \href{https://zips.z.cash/protocol/protocol.pdf\#concretehomomorphiccommit}{\S\S~5.4.8.3}--\href{https://zips.z.cash/protocol/protocol.pdf\#concretesinsemillacommit}{5.4.8.4}), and
  \texttt{z.cash:\allowbreak Orchard-gd} separates diversified key
  generation (protocol specification \href{https://zips.z.cash/protocol/protocol.pdf\#concretediversifyhash}{\S~5.4.1.6}). The deployed
  framing mirrors the definition's rather
  than copying it: the curve hash of \S\ref{hash:subsec:hashtocurve}
  \emph{appends} its domain tag after the message, closed by a byte
  carrying the tag's length (protocol specification \href{https://zips.z.cash/protocol/protocol.pdf\#concretegrouphashpallasandvesta}{\S~5.4.9.8}); the
  length byte makes the tag self-delimiting, so distinct tags still
  induce disjoint effective inputs---all the proof of
  Proposition~\ref{hash:prop:domainsep} uses. The tags reach that
  hash by two routes, and the commitment tags by both: \texttt{-cv}
  and \texttt{-gd} serve directly as its domain tag, while each
  commitment tag, suffixed \texttt{-M}, is hashed as the
  \emph{message} under the fixed domain tag
  \texttt{z.cash:\allowbreak SinsemillaQ} to derive one of the
  commitment's two base points and, suffixed \texttt{-r}, serves
  directly as the domain tag deriving the other (protocol specification
  \href{https://zips.z.cash/protocol/protocol.pdf\#concretesinsemillahash}{\S\S~5.4.1.9} and \href{https://zips.z.cash/protocol/protocol.pdf\#concretesinsemillacommit}{5.4.8.4}).
\item \emph{Native personalisation.} BLAKE2b's parameter block carries
  a dedicated \(16\)-byte personalisation field that is mixed into the
  initialisation vector, so distinct personalisations separate the
  primitive's effective inputs before a single message byte is
  processed---the field for which Zcash favours BLAKE2 for its
  oracles. With BLAKE2b modelled as a random oracle of the pair
  (personalisation, message), the fixed-length field is a prefix-free
  tag and Proposition~\ref{hash:prop:domainsep} makes the resulting
  domains independent; a concrete BLAKE2b instantiation does not make
  that independence a fact ``by construction''. The deployed
  personalisations include \texttt{Zcash\_\allowbreak ExpandSeed} for
  key expansion (protocol specification \href{https://zips.z.cash/protocol/protocol.pdf\#concreteprfs}{\S~5.4.2}),
  \texttt{Zcash\_\allowbreak OrchardKDF} for the note-encryption
  key-derivation function (KDF, the object of
  \S\ref{symmetric:subsec:kdf}; protocol specification \href{https://zips.z.cash/protocol/protocol.pdf\#concreteorchardkdf}{\S~5.4.5.6}), and
  \texttt{Zcash\_\allowbreak Orchardock} for the outgoing cipher key
  (protocol specification \href{https://zips.z.cash/protocol/protocol.pdf\#concreteprfs}{\S~5.4.2})---all exactly \(16\) bytes.
\item \emph{Suffix and framing bits in the SHA-3 family.} FIPS~202
  reserves domain-separation suffix bits appended to the message
  before padding: \(01\) for SHA-3 hashing and \(1111\) for the SHAKE
  XOFs. SHA3-256 and SHAKE256 share the same underlying permutation,
  yet the framed inputs differ, so their effective inputs are
  disjoint---though their finite digests can still coincide by
  chance.
\end{enumerate}

The unifying requirement behind all three idioms is that distinct uses
induce \emph{disjoint effective inputs} to the underlying primitive.
An ad hoc tag that happens to be a prefix of another silently merges
two domains and reopens the replay attack---which is why
Definition~\ref{hash:def:domainsep} demands prefix-freeness rather
than mere distinctness. The Merkle trees of the final section apply
the same discipline \emph{within} a single structure, tagging tree
layers to foreclose a forgery that confuses leaves with internal
nodes; see ``Layer and domain separation'' there.

\subsection{Hashing to a field element}\label{hash:subsec:hashtofield}

The primitives ahead---commitments over the Pasta curves, Fiat--Shamir
challenges, key derivations---repeatedly need to hash a message
\emph{to a field element}: a near-uniform element of \(\Fp\) rather
than a bit string. The obstacle is arithmetic: a hash emits uniform
\emph{bits}, and reducing \(n\) uniform bits modulo a prime \(p\) is
not uniform on \(\Fp\), because \(2^{n}\) is not a multiple of
\(p\)---some residues receive one more preimage than others. The
remedy is to hash \emph{wide} and reduce.

\begin{construction}[Expand-then-reduce]\label{hash:con:expandreduce}
Let \(p\) be a prime of \(L = \ceil{\log_{2} p}\) bits, and fix a XOF,
or a domain-separated hash, \(\mathsf{expand}\), producing
arbitrary-length near-uniform output. To map a message \(m\) to
\(\Fp\), with a domain tag \(\tau\) and a security slack of \(k\)
bits (commonly \(k = 128\)):
\begin{enumerate}
\item compute \(b \gets \mathsf{expand}(\tau \,\|\, m,\; L + k)\), a
  string of \(L + k\) near-uniform bits;
\item interpret \(b\) as an integer
  \(B \in \{0, \dots, 2^{L+k} - 1\}\) and output
  \(B \bmod p \in \Fp\).
\end{enumerate}
\end{construction}

\begin{proposition}[Near-uniformity of expand-then-reduce]
\label{hash:prop:bias}
If \(\mathsf{expand}\) is a random oracle, the output of
Construction~\ref{hash:con:expandreduce} lies within statistical
distance (Definition~\ref{model:def:statdist}) at most
\[
\frac{p}{2^{\,L+k+1}} \;\le\; 2^{-k}
\]
of the uniform distribution on \(\Fp\). With \(k = 128\), no
algorithm whatsoever---efficient or not---distinguishes the output
from a uniform field element with advantage better than \(2^{-128}\).
\end{proposition}

\begin{proof}
The integer \(B\) is uniform on \(\{0, \dots, K - 1\}\) with
\(K = 2^{L+k}\), and the bias of reducing a uniform \(B\) modulo
\(p\) is exactly the situation of the modular-reduction proposition
of the \emph{Math Guide} (\S``Uniform sampling and the bias of
modular reduction''): each residue's probability differs from
\(1/p\) by at most \(1/K\), because the preimage counts
\(\floor{K/p}\) and \(\ceil{K/p}\) differ by one, and summing over
the \(p\) residues gives statistical distance at most
\(p/(2K) = p/2^{\,L+k+1}\)---the sharper form proved there. Since
\(p < 2^{L}\), this is less than \(2^{-(k+1)} \le 2^{-k}\). The final claim
is the variational characterisation of statistical distance
(Proposition~\ref{model:prop:statdist}): no decision procedure
exceeds the distance.
\end{proof}

The slack is not optional. With \(k = 0\)---a single \(L\)-bit block
reduced modulo \(p\)---the bias can be macroscopic: for \(p\) just
above a power of two, \(K = 2^{L}\) is barely twice \(p\), the
residues beyond the wrap-around point receive only one preimage while
the rest receive two, and each deprived residue carries as little as
half its uniform mass. The \(k\) extra bits are exactly what drives
this imbalance to negligibility.

Deployed practice packages the construction. The IETF hash-to-curve
specification (RFC~9380) defines
\(\mathsf{hash\_to\_field}(m, \mathit{count})\), producing
\(\mathit{count}\) field elements by calling an expansion
routine---\(\mathsf{expand\_message\_xmd}\) over a fixed-output hash,
or \(\mathsf{expand\_message\_xof}\) over a XOF---with a
\emph{mandatory} domain-separation tag, and reducing each chunk of
\[
L_{\mathrm{RFC}} \;=\; \Bigl\lceil
\bigl(\lceil \log_{2} p \rceil + k\bigr)/8 \Bigr\rceil
\ \text{bytes}
\]
modulo \(p\), where \(k\) is the suite's target security level in
bits. The RFC's own suites for \(255\)-bit fields---the curve25519
family---set \(k = 128\), hence \(48\)-byte chunks. The Zcash suite
for Pallas and Vesta is more generous: it takes \(k = 256\), reducing
each element from a full \(64\)-byte BLAKE2b-512 digest, so the slack
of Proposition~\ref{hash:prop:bias} is realised with room to spare
(the
\texttt{hash\_to\_field} routine expands with BLAKE2b in the XMD
pattern under the suite tag
\texttt{\{tag\}-\allowbreak\{curve\}\_\allowbreak
XMD:BLAKE2b\_\allowbreak SSWU\_RO\_} and
reduces two \(64\)-byte chunks each via
\texttt{from\_uniform\_bytes}; protocol specification \href{https://zips.z.cash/protocol/protocol.pdf\#concretegrouphashpallasandvesta}{\S~5.4.9.8}).
This is the concrete recipe for mapping identities to field elements,
and for producing the field elements that the next subsection feeds to
a curve. Halo~2's deployed Fiat--Shamir transcript uses the same
wide-output-then-reduce principle: each challenge is the \(64\)-byte
digest of a BLAKE2b state personalised with \texttt{Halo2-Transcript},
reduced by \texttt{from\_uniform\_bytes}; the order in which that state absorbs
messages and squeezes challenges is the \emph{Halo~2 Guide}'s,
\S``The complete Halo 2 protocol''.

\subsection{Hashing to a curve point}\label{hash:subsec:hashtocurve}

The volume ahead also needs to hash \emph{onto an elliptic curve}:
given a message, produce a point of \(E(\Fp)\) that behaves like a
uniform group element---the raw material for commitment generators
whose discrete logarithms nobody knows. Two requirements pull against
each other. The output should be near-uniform in the relevant
subgroup, so that the map can stand in for a random oracle into the
group in security arguments; and the map must run in \emph{constant
time}, because its input is often secret-adjacent.

The naive approach fails the second requirement instructively. Hash
the message to a candidate \(x \in \Fp\)
(Construction~\ref{hash:con:expandreduce}); test whether
\(g(x) = x^{3} + ax + b\) is a quadratic residue, so that a
\(y\)-coordinate exists; if not, increment \(x\) and retry. This
\emph{try-and-increment} loop is correct and roughly uniform, but the
number of trials depends on the message: about half of all \(x\) fail
the residue test, since the squares are half of \(\Fp^{\times}\)
(\emph{Math Guide}, \S``Quadratic residues and the Euler criterion'',
whose Legendre symbol \(\legendre{g(x)}{p}\) is exactly the test), so
the running time leaks information about \(m\)---an avenue for
side-channel attacks. Constant-time hashing to a curve requires a map
that succeeds on the \emph{first} try, for every input.

\begin{construction}[Simplified Shallue--van de
Woestijne--Ulas map]\label{hash:con:swu}
Let \(E \colon y^{2} = g(x) = x^{3} + ax + b\) over \(\Fp\) with
\(ab \neq 0\), and choose, as RFC~9380 requires, a constant
\(Z \in \Fp\) such that \(Z\) is a nonsquare, \(Z \neq -1\), the
polynomial \(g(X) - Z\) is irreducible over \(\Fp\), and
\(g(b/(Za))\) is a square. (The RFC selects \(Z\) mechanically, as
the candidate of smallest absolute value meeting these
conditions---a nothing-up-my-sleeve constant with no hidden design
freedom.) The map
\(\mathrm{SWU} \colon \Fp \to E(\Fp)\) sends \(u\) to a point as
follows. Compute the two candidate \(x\)-coordinates
\[
x_{1} \;=\; \frac{-b}{a}
  \Bigl(1 + \frac{1}{Z^{2}u^{4} + Zu^{2}}\Bigr),
\qquad
x_{2} \;=\; Z u^{2} x_{1},
\]
with the fallback \(x_{1} = b/(Za)\) when the denominator
\(Z^{2}u^{4} + Zu^{2}\) vanishes (as at \(u = 0\)). Set
\(x = x_{1}\) if \(g(x_{1})\) is a square and
\(x = x_{2}\) otherwise; let \(y\) be a square root of \(g(x)\)
(\emph{Math Guide}, \S``Square roots modulo a prime''), its sign fixed
by a convention---the deployed map matches the parity of \(y\) to
that of \(u\)---and output \((x, y)\).
\end{construction}

The construction contains no trial loop, hence runs in constant time;
what needs proof is that one of the two candidates always works.

\begin{proposition}[Correctness of the branch selection]
\label{hash:prop:swu}
Under the parameter conditions of Construction~\ref{hash:con:swu},
the map is well-defined for every \(u \in \Fp\): the selected value
\(g(x)\) is a square, so \(\mathrm{SWU}(u)\) is a genuine point of
\(E(\Fp)\).
\end{proposition}

\begin{proof}[Proof sketch]
Suppose first that \(D = Z^{2}u^{4} + Zu^{2} \neq 0\). The key
algebraic identity is
\[
g(x_{2}) \;=\; Z^{3} u^{6}\, g(x_{1}),
\]
a polynomial identity in \(u\): substitute \(x_{2} = Zu^{2}x_{1}\)
into \(g\) and simplify using the formula for \(x_{1}\) (we take the
computation as given). If \(g(x_{1}) = 0\), the first candidate
already works, with \(y = 0\). Otherwise apply multiplicativity of
the Legendre symbol (\emph{Math Guide}, \S``Quadratic residues and
the Euler criterion''):
\[
\legendre{g(x_{2})}{p}
= \legendre{Z^{3}u^{6}}{p} \legendre{g(x_{1})}{p}
= \legendre{Z}{p} \legendre{(Zu^{3})^{2}}{p} \legendre{g(x_{1})}{p}
= \legendre{Z}{p} \legendre{g(x_{1})}{p},
\]
since \(Z^{2}u^{6} = (Zu^{3})^{2}\) is a nonzero square, contributing
\(+1\) (note that \(D \neq 0\) forces \(u \neq 0\)). Because \(Z\) is
a nonsquare, \(\legendre{Z}{p} = -1\), so the two symbols are
opposite and the branch selection picks the square. If \(D = 0\), the
prescribed fallback has \(g(x_{1}) = g(b/(Za))\), a square by the
last condition on \(Z\). Thus every input produces a valid affine
point.
\end{proof}

A single SWU evaluation is not the end of the pipeline: the map is not
surjective onto \(E(\Fp)\), and its image is not uniformly
covered---it reaches only part of the curve, with a lumpy
distribution. The standardised pipeline repairs both defects.

\begin{construction}[The hash-to-curve pipeline]\label{hash:con:h2c}
The RFC~9380 map
\(\mathsf{hash\_to\_curve} \colon \{0,1\}^{*} \to E(\Fp)\) composes
four steps:
\begin{enumerate}
\item hash to the field: \(\mathsf{hash\_to\_field}(m, 2)\) produces
  two independent near-uniform field elements \(u_{0}, u_{1}\)
  (\S\ref{hash:subsec:hashtofield});
\item map each through SWU: \(Q_{0} = \mathrm{SWU}(u_{0})\),
  \(Q_{1} = \mathrm{SWU}(u_{1})\);
\item sum on the curve: \(R = Q_{0} + Q_{1}\). Summing two
  independent SWU images provably repairs the lumpiness: the sum
  lies within negligible statistical distance of uniform on
  \(E(\Fp)\), the theorem on which the pipeline's random-oracle
  claim rests;
\item clear the cofactor: multiply \(R\) by \(h = \#E(\Fp)/r\) to land
  in the prime-order subgroup of order \(r\) where the cryptography
  occurs.
\end{enumerate}
The result is a constant-time hash from messages to group elements
which, under the hash model and suite conditions of RFC~9380, can
instantiate a random oracle into the group. For the Pasta curves the
last step is vacuous---their
cofactor is \(1\) (Definition~\ref{assumptions:def:pasta}).
\end{construction}

% TikZ figure placeholder --- fig:hashtocurve: pipeline dataflow.
% Message m enters hash_to_field (tagged tau), which emits u0 and u1;
% each passes through the simplified SWU box onto the isogenous curve
% E' (drawn as a separate rail, labelled "ab != 0"); the two points
% are summed on E'; the 3-isogeny arrow carries the sum onto E
% (Pallas); a final faded box marks cofactor clearing, annotated
% "h = 1 for Pasta: skipped". Dedicated figures pass supplies the
% drawing.

\begin{remark}[The isogeny detour forced by \(j\)-invariant \(0\)]
\label{hash:rem:isogeny}
For the Pasta curves one step needs a detour. The SWU formula for
\(x_{1}\) divides by \(a\), and its numerator carries the factor
\(-b\): with \(a = 0\) the formula is undefined, while \(b = 0\)
makes \(x_{1}\) identically zero with \(g(x_{1}) = 0\), collapsing
the map to the constant point \((0, 0)\). The construction therefore
requires \(ab \neq 0\)---and Pallas and Vesta have \(a = 0\) (the
\(j\)-invariant-\(0\) shape \(y^{2} = x^{3} + 5\) of
Definition~\ref{assumptions:def:pasta}; curves with \(b = 0\),
\(j = 1728\), are excluded symmetrically). The obstruction is
intrinsic, not notational: the \(j\)-invariant is an isomorphism
invariant, and a Weierstrass change of variables rescales
\((a, b)\) to \((u^{4}a, u^{6}b)\) (\emph{Math Guide}, \S``The
\(j\)-invariant and the GLV endomorphism''), so \emph{no} model of
these curves attains \(ab \neq 0\). The standardised remedy evaluates
simplified SWU on an \emph{isogenous} auxiliary curve
\(E' \colon y^{2} = x^{3} + a'x + b'\) with \(a'b' \neq 0\), then
transports the result through an \emph{isogeny}
\(E' \to E\): a nonconstant map between curves given by polynomial
coordinate formulas and respecting the group law. The deployed
implementation applies a degree-\(3\) isogeny, and---because an
isogeny is a group homomorphism---sums \(Q_{0} + Q_{1}\) on \(E'\)
first and transports the single sum
(protocol specification
\href{https://zips.z.cash/protocol/protocol.pdf\#concretegrouphashpallasandvesta}{\S~5.4.9.8}).
\end{remark}

The culminating application ties together collision resistance, the
random-oracle idealisation, and hashing to a curve.

\begin{construction}[Nothing-up-my-sleeve generators]
\label{hash:con:nums}
To generate public generators \(G_{1}, \dots, G_{k}\) of a vector
commitment (the Pedersen vector commitment of
Construction~\ref{commitments:def:pedersen-vector} in the next
section) so that no one knows any discrete-logarithm relation among
them, derive each as
\[
G_{i} \;=\; \mathsf{hash\_to\_curve}(\tau \,\|\, i)
\]
from a fixed public domain tag \(\tau\)---often a human-readable
protocol name---and the index \(i\). Because the seed is a
transparent, auditable string with no hidden parameters, no party
could have engineered a backdoor relation in advance. This is exactly
how Halo~2 produces the generators of its Pedersen vector and
inner-product commitments (the latter
Construction~\ref{commitments:def:ipa}): every generator \(G_{i}\) is
\(\mathsf{hash\_to\_curve}(\texttt{"Halo2-Parameters"})\) applied to
a five-byte message---a zero byte, then the four-byte little-endian
encoding of \(i\)---the leading zero keeping the \(G_{i}\) disjoint
from the two auxiliary generators the same tag derives from the
one-byte messages \(1\) and \(2\), and Orchard's value
commitment derives its two bases from the tag
\texttt{z.cash:Orchard-cv} with the one-byte messages \texttt{"v"}
and \texttt{"r"}.
\end{construction}

\begin{remark}[Why NUMS generators are binding-safe]
\label{hash:rem:nums-binding}
Knowledge of a discrete-logarithm relation
\(\sum_{i} a_{i} G_{i} = 0\) with not all \(a_{i}\) zero would let a
committer open one commitment two ways, destroying the binding
property that the next section defines and relies on. The NUMS
derivation forecloses this: the generators are outputs of a function
modelled as a random oracle into the group, so they are distributed as
independent uniform group elements, and finding a discrete-logarithm
relation among independent uniform elements \emph{reduces} to
computing a discrete logarithm on \(E\), hard under the standing
assumptions of Remark~\ref{assumptions:rem:standing}; the reduction,
which embeds its challenge in every generator, is
Theorem~\ref{commitments:thm:vector-props} of the next section. Hence
no efficient party, the scheme designer included, can know such a
relation.
\end{remark}

\subsection{Arithmetisation-friendly hashing: Poseidon}
\label{hash:subsec:poseidon}

The hashes of this section so far are \emph{bit-oriented}: SHA-2 and
BLAKE2 shuffle and XOR \(32\)- and \(64\)-bit words. Inside an
arithmetic circuit over a large prime field \(\Fp\)---the setting of
proof systems such as Halo~2, whose cost model, developed in the
\emph{Halo~2 Guide}, \S``From computation to a grid of
constraints'', prices each field multiplication as a constraint---such
a hash is
prohibitively expensive: every Boolean operation must be re-expressed
through field arithmetic and range checks, and the tens of thousands
of resulting constraints dominate the circuit. Arithmetisation-\allowbreak
friendly hashes invert the design: the round function is a handful of
\emph{low-degree field operations}, native to the circuit, incurring
few constraints. Poseidon (Grassi, Khovratovich, Rechberger, Roy, and
Schofnegger, 2019) is the instance Orchard employs, and it arrives via
a mode of operation worth defining in general.

\begin{construction}[The sponge construction]\label{hash:con:sponge}
Fix a permutation \(f \colon \Fp^{t} \to \Fp^{t}\) on a state of
\(t\) field elements. Split the state into a \emph{rate} of \(r\)
elements---the only part input is added into and output is read
from---and a \emph{capacity} of \(c = t - r\) elements, never
accessed directly: the security reserve. To hash, \emph{absorb} the
message into the rate, \(r\) elements at a time, applying \(f\)
between blocks; then \emph{squeeze} output from the rate, applying
\(f\) between reads. Collision and preimage resistance hold up to
approximately \(\sqrt{p^{c}}\) work, as for a random sponge: the
birthday bound of Lemma~\ref{hash:lem:birthday} applied to the
\(p^{c}\) capacity states the adversary cannot see.
\end{construction}

% TikZ figure placeholder --- fig:sponge: Poseidon sponge dataflow.
% State drawn as a t-element column split by a horizontal rule into
% rate (top, r elements) and capacity (bottom, c elements); absorb
% phase shows message blocks added into the rate with an f-permutation
% box between consecutive blocks; squeeze phase reads output from the
% rate. The capacity lane runs uninterrupted left to right, visibly
% never touched by input or output arrows --- caption note: "the
% capacity is never exposed". Dedicated figures pass supplies the
% drawing.

\begin{construction}[The Poseidon permutation]\label{hash:con:poseidon}
The permutation \(f\) is a sequence of rounds, each comprising three
steps:
\begin{enumerate}
\item \emph{add round constants}: add a fixed vector in \(\Fp^{t}\)
  to the state;
\item \emph{S-box}: apply \(x \mapsto x^{\alpha}\) to state elements,
  where \(\alpha \ge 3\) is the smallest integer with
  \(\gcd(\alpha, p - 1) = 1\), so that the power map is a bijection
  on \(\Fp\); over the Pasta fields \(\alpha = 5\), because
  \(3 \mid p - 1\) there rules out \(\alpha = 3\);
\item \emph{mix}: multiply the state by a fixed
  maximum-distance-separable (MDS) matrix
  \(M \in \Fp^{t \times t}\)---one whose every square submatrix is
  invertible, giving the fastest possible diffusion of any element's
  change across the whole state.
\end{enumerate}
The cost-saving principle (the HADES strategy) distinguishes two
round types: a few \emph{full} rounds at the start and end apply the
S-box to all \(t\) elements, while in the many \emph{partial} rounds
between them the S-box touches a single element---the MDS mix still
diffusing it across the state. Round counts are chosen to stand
against the known algebraic (Gr\"obner-basis) attacks.
\end{construction}

\begin{remark}[Circuit cost]\label{hash:rem:poseidon-cost}
Inside a circuit, the only non-linear operation in all of Poseidon is
\(x \mapsto x^{5}\), verified by a couple of multiplication gates per
S-box; the round-constant additions and the MDS multiplication are
linear and almost free. A \(2\)-to-\(1\) Poseidon hash costs a few
hundred constraints, against tens of thousands for a bit-oriented
hash---the difference between a practical circuit and an impractical
one.
\end{remark}

\begin{example}[The Orchard Poseidon instance]\label{hash:ex:orchard}
Orchard fixes \(t = 3\) (rate \(2\), capacity \(1\)),
\(\alpha = 5\), and \(8\) full plus \(56\) partial rounds, over the
Pallas base field, at the \(128\)-bit security level---the
parametrisation named \texttt{P128Pow5T3} in the deployed source
(protocol specification \href{https://zips.z.cash/protocol/protocol.pdf\#poseidonhash}{\S~5.4.1.10}). It is used in
\emph{constant-input-length} mode as a two-input hash:
\[
\mathsf{PoseidonHash}(x, y) \;=\;
f\bigl([\,x,\; y,\; 2^{65}\,]\bigr)_{0},
\]
the first component of the state after a single permutation of the
input padded by the fixed capacity value
\(2^{65} = 2 \cdot 2^{64}\), which encodes the input length \(2\)
into the capacity so that no message padding is required
(the
\texttt{ConstantLength} domain sets the initial capacity element to
\(L \cdot 2^{64}\)). Keyed by its first argument, this hash serves as
the pseudorandom function
\[
F_{\mathsf{nk}}(\rho) \;=\; \mathsf{PoseidonHash}(\mathsf{nk}, \rho)
\]
---the PRF vocabulary is owned by ``Pseudorandom functions and
permutations'' later in this volume---deployed as
\(\mathsf{PRF}^{\mathsf{nfOrchard}}\) (protocol specification
\href{https://zips.z.cash/protocol/protocol.pdf\#concreteprfs}{\S~5.4.2}), from which Orchard constructs the
nullifier---a construction developed in the \emph{Ironwood Guide},
\S``The nullifier, and the loop back to minting''; this volume
supplies only the PRF. One contrast deserves note: unlike the Pedersen
and Sinsemilla hashes constructed in the next section, whose collision
resistance reduces to a clean assumption (discrete log), Poseidon's
security rests on cryptanalytic confidence in its round function---the
same epistemic footing as SHA-256, not the same as a reduction.
\end{example}

\subsection{The transparent layer's hashes}
\label{hash:subsec:transparent}

The Orchard shielded protocol hashes with BLAKE2b, Sinsemilla, and
Poseidon. The transparent layer, inherited from Bitcoin, hashes with
SHA-256 and two functions derived from it, none needing machinery
beyond \S\ref{hash:subsec:games} and the birthday bound of
\S\ref{hash:subsec:birthday}; this register fixes their names,
domains, and deployed uses for the volumes above.
\begin{itemize}
\item \emph{RIPEMD-160.} A \(160\)-bit iterated (Merkle--Damg\aa rd,
  the construction of \S\ref{prf:subsec:hmac}) hash,
  \(\{0,1\}^{*} \to \{0,1\}^{160}\), older than Bitcoin and
  inherited with its address format, whose collision resistance sits
  at the \(80\)-bit birthday bound of Lemma~\ref{hash:lem:birthday}.
  Zcash uses it only inside \(\mathsf{hash160}\) below and the script
  opcodes \texttt{OP\_RIPEMD160} and \texttt{OP\_HASH160}.
\item \emph{The address hash
  \(\mathsf{hash160} := \mathrm{RIPEMD160} \circ \mathrm{SHA256}\).}
  A pay-to-public-key-hash (P2PKH) address carries the \(20\)-byte
  \(\mathsf{hash160}\) of the \(33\)-byte compressed encoding of an
  ECDSA public key (the scheme is registered in
  \S\ref{signatures:subsec:transparent}), and a pay-to-script-hash
  (P2SH) address the \(\mathsf{hash160}\) of a \emph{redeem script},
  the spending condition the spender later reveals and satisfies
  (protocol specification \href{https://zips.z.cash/protocol/protocol.pdf\#transparentaddrencoding}{\S~5.6.1.1}). The P2PKH output script---the
  condition a spend of the output must satisfy---is
  \texttt{OP\_DUP OP\_HASH160 \(\langle h \rangle\) OP\_EQUALVERIFY
  OP\_CHECKSIG}, which recomputes the hash from the public key the
  spender reveals; wallets compute it through the
  \texttt{zcash\_transparent} crate, and a partially created
  transaction (PCZT), the multi-party construction format of the
  \emph{Wallet Guide}, \S``Partially created transactions (PCZT)'',
  carries on each transparent input preimage maps keyed by
  RIPEMD-160, SHA-256, \(\mathsf{hash160}\), and \(\mathsf{hash256}\)
  digests.
\item \emph{Double SHA-256 (SHA-256d, or \(\mathsf{hash256}\)).} The
  composition \(x \mapsto \mathrm{SHA256}(\mathrm{SHA256}(x))\)
  (protocol specification \href{https://zips.z.cash/protocol/protocol.pdf\#concretesha256}{\S~5.4.1.1}), whose outer application also
  closes the length-extension property of plain SHA-256
  (\S\ref{prf:subsec:hmac}). It hashes block headers; its
  leading four bytes are the checksum that Base58Check, the string
  encoding of transparent addresses, appends to the raw address bytes
  before converting to base~58 (protocol specification \href{https://zips.z.cash/protocol/protocol.pdf\#transparentaddrencoding}{\S~5.6.1.1}) and the payload checksum in
  every peer-to-peer message header; and it is the
  legacy transaction identifier: the txid of a version-4-or-earlier
  transaction is SHA-256d of its serialisation (protocol
  specification \href{https://zips.z.cash/protocol/protocol.pdf\#txnidentifiers}{\S~7.1.1}), whereas a version-5
  txid is the BLAKE2b digest tree of ZIP~244.
\end{itemize}

The toolbox now holds its first primitive in full: a hash with three
calibrated security notions, an idealisation that licenses proofs, a
discipline that insulates its many uses, and maps carrying its output
into fields, onto curves, and inside circuits. The next section
spends this capital immediately: commitments---the sealed envelopes
of the protocol---are built from exactly the hashes, generators, and
hardness assumptions now on the table.

\section{Commitment schemes}\label{sec:commitments}

Two adversaries stalk this section, one on each side of a sealed
envelope. Before the envelope is opened, the adversary is the
\emph{peeker}: the receiver, who holds the sealed commitment and tries
to distinguish which of two candidate messages it contains---to read
the bid before the auction closes, the vote before the tally. After
the envelope has been delivered, the adversary is the
\emph{equivocator}: the sender, who tries to open the same commitment
to a message other than the one sealed inside---to change the bid once
the rival's is public. A commitment scheme is the cryptographic
sealed, tamper-evident envelope that defeats both: once delivered, the
sender can no longer alter the message (\emph{binding}), yet the
recipient cannot read it until it is unsealed (\emph{hiding}).

The primitive is the natural next step after hashing: the two
constructions at the heart of this section, Pedersen and Sinsemilla,
need nothing beyond the discrete-logarithm hardness of
\S\ref{sec:assumptions} and the nothing-up-my-sleeve generators that
\S\ref{sec:hash} showed how to derive, and nearly everything built
later consumes them. Classical applications set the scene:
coin-flipping over a telephone (each party commits to a bit before
either reveals), zero-knowledge proofs (the prover commits to its
secrets and proves statements about the sealed values), and
verifiable secret sharing. Centrally for this volume, the
\emph{polynomial} commitments of \S\ref{commitments:subsec:poly} are the
cryptographic engine of modern succinct argument systems, including
the one deployed in Orchard: the value commitments checked by binding
signatures, the note commitments accumulated in Merkle trees, and the
witness polynomials of the Halo~2 proof system are all commitments in
the precise sense defined here.

The plan of the section follows the two adversaries. We first fix the
syntax of a commitment scheme and the three models of who may be
trusted to set it up. We then define hiding and binding as games in
the tradition of \S\ref{model:subsec:games}, each in three
quantitative flavours---perfect, statistical, computational---and
prove that the two perfections are jointly unattainable: every scheme
must choose which adversary to defeat unconditionally. The
constructions that follow are ordered by increasing algebraic
richness: hash-based commitments (opaque but transparent and
plausibly post-quantum), the Pedersen commitment and its vector
variant (perfectly hiding, computationally binding under discrete
log, and---decisively---linear), Sinsemilla (a commitment whose
internal operations are curve operations, built to be cheap inside a
proof), then polynomial commitment schemes, contrasting the
pairing-based KZG family with the inner-product family that Zcash
deploys, and the bridge from commitments to arguments.
A closing subsection returns to the classical application named
above---secret sharing---and builds its verifiable form from
coefficient commitments, for the threshold signatures of the
\emph{FROST Guide}.

\subsection{The sealed envelope}\label{commitments:subsec:envelope}

The envelope metaphor repays being taken slowly, because each of its
physical properties has an exact formal counterpart. The sender
writes a message on a card, seals it in the envelope, and hands the
envelope over. From that moment the receiver \emph{possesses} the
commitment: whatever the message was, it is now fixed---the sender
cannot swap the card, because the envelope is in the receiver's
hands and any tampering would show. Yet the receiver \emph{learns}
nothing: the envelope is opaque. Later, at a time of the protocol's
choosing, the sender opens the envelope---in the formalism, transmits
the message together with the randomness used to seal it---and the
receiver checks that the revealed message is consistent with the
envelope received earlier.

The two guarantees are directed at different parties and different
times. Hiding protects the \emph{sender} against a curious
\emph{receiver} during the window between commit and reveal; binding
protects the \emph{receiver} against a dishonest \emph{sender} at
reveal time. A scheme that is hiding but not binding is an envelope
with a false bottom: private, but worthless as evidence. A scheme
that is binding but not hiding is a glass envelope: tamper-evident,
but no secret survives it. The games of
\S\ref{commitments:subsec:hiding-binding} formalise exactly these two
failure modes, and the theorem of
\S\ref{commitments:subsec:impossibility} shows that against
computationally unbounded adversaries one of the two protections must
give way.

\subsection{Syntax}\label{commitments:subsec:syntax}

\begin{definition}[Commitment scheme]\label{commitments:def:scheme}
A \emph{commitment scheme} \(\Pi\) for a message space \(\mathcal M\)
and randomness space \(\mathcal R\) is a triple of algorithms
\((\mathrm{Setup}, \mathrm{Commit}, \mathrm{Open})\):
\begin{itemize}
\item The algorithm \(\mathrm{Setup}(1^\lambda)\) is PPT
  (\S\ref{model:subsec:adversaries}) and outputs \emph{public
  parameters} \(pp\), which implicitly determine the message space
  \(\mathcal M_{pp}\) and randomness space \(\mathcal R_{pp}\) and are
  an implicit input to the remaining algorithms.
\item The algorithm \(\mathrm{Commit}(pp, m; r)\), for
  \(m \in \mathcal M_{pp}\) and \(r \in \mathcal R_{pp}\), is
  deterministic given \(r\) and outputs a \emph{commitment} \(c\). We
  write \(\mathrm{Commit}(pp, m)\), without an explicit \(r\), for
  the probabilistic algorithm that samples
  \(r \xleftarrow{\$} \mathcal R_{pp}\) uniformly and returns \(r\)
  alongside \(c\).
\item The algorithm \(\mathrm{Open}(pp, c, m, r)\) is deterministic
  and outputs a bit: \(1\) (accept) or \(0\) (reject).
\end{itemize}
\emph{Correctness} requires that honestly produced commitments open:
for all \(pp\) in the support of \(\mathrm{Setup}(1^\lambda)\), all
\(m \in \mathcal M_{pp}\), and all \(r \in \mathcal R_{pp}\),
\[
\mathrm{Open}\bigl(pp,\, \mathrm{Commit}(pp, m; r),\, m,\, r\bigr)
  = 1 .
\]
\end{definition}

The pair \((m, r)\) is called the \emph{opening} (or
\emph{decommitment}) of \(c\). Every scheme in this volume is
\emph{canonically opening}: \(\mathrm{Open}(pp, c, m, r)\) merely
recomputes \(\mathrm{Commit}(pp, m; r)\) and checks equality with
\(c\). This is the adopted convention unless stated otherwise. For a
canonically opening scheme correctness holds automatically, and a
valid opening of \(c\) is precisely a pair \((m, r)\) with
\(\mathrm{Commit}(pp, m; r) = c\)---a reformulation the proofs below
use constantly.

\begin{remark}[The two phases, and a third concern]
\label{commitments:rem:phases}
A commitment scheme is used in two temporally separated phases. In
the \emph{commit phase} the sender transmits \(c\) and nothing else;
in the \emph{reveal phase} the sender transmits the opening
\((m, r)\) and the receiver runs \(\mathrm{Open}\). The two security
properties attach to the two phases: hiding constrains what the
receiver can learn \emph{after the commit phase}, while binding
constrains what the sender can do \emph{in the reveal phase}. A third
concern precedes both: the \emph{setup phase} that produces \(pp\).
Who runs \(\mathrm{Setup}\), and whether that party can be trusted,
proves decisive for some of the constructions ahead
(Remark~\ref{commitments:rem:setup-trust}).
\end{remark}

% FIGURE (planned, dedicated figures pass): fig:commit-phases ---
% commit/reveal timeline. A horizontal time axis with three marked
% intervals: setup (pp produced; trust question flagged), commit
% (c crosses from sender to receiver; hiding shields the interval
% after this point --- what the receiver holds reveals nothing), and
% reveal ((m, r) crosses; binding constrains this point --- the sender
% can produce no second valid opening). The peeker sits under the
% post-commit interval, the equivocator under the reveal point.

\begin{remark}[Three setup models]\label{commitments:rem:setup-trust}
Commitment schemes divide by how much trust their parameters demand.
\begin{enumerate}
\item \emph{Common reference string (CRS).} A trusted party runs
  \(\mathrm{Setup}\) and is assumed to \emph{erase} the randomness
  behind it. The parameters have internal structure, and knowledge of
  the discarded randomness---the \emph{trapdoor}---would break the
  scheme. Of the schemes treated in this section, only KZG
  (\S\ref{commitments:subsec:poly}) genuinely resides here.
\item \emph{Uniform random string (URS), or transparent setup.} The
  parameters consist only of public coins: uniformly random group
  elements with no known relations among them. There is no trapdoor
  and no trusted party---anyone can re-derive the parameters from a
  public seed. This is the regime of the Pedersen and inner-product
  commitments below. A Pedersen generator derived by
  hashing to the curve (Construction~\ref{hash:con:nums}) is
  transparent
  even though the textbook presentation phrases \(\mathrm{Setup}\) as
  sampling and discarding a secret.
\item \emph{Plain.} No parameters at all, or only a standardised hash
  function---the home of the hash-based commitment below.
\end{enumerate}
Transparent setup is highly desirable: a flawed or subverted trusted
setup can silently destroy binding, as the KZG trapdoor demonstrates
concretely in \S\ref{commitments:subsec:poly}.
\end{remark}

\subsection{Hiding and binding}
\label{commitments:subsec:hiding-binding}

Both properties are games in the sense of
Definition~\ref{model:def:game}. The hiding game arms the peeker: she
chooses the two candidate messages herself, receives a commitment to
one of them, and must say which.

\begin{definition}[Hiding game and hiding]
\label{commitments:def:hiding}
Let \(\Pi\) be a commitment scheme and
\(\mathcal A = (\mathcal A_1, \mathcal A_2)\) an adversary. For
\(b \in \{0,1\}\), the game
\(\mathsf{Hide}^{\mathcal A}_{\Pi}(\lambda, b)\) runs as follows.
\begin{enumerate}
\item The challenger runs \(pp \gets \mathrm{Setup}(1^\lambda)\).
\item The adversary \(\mathcal A_1(pp)\) outputs two messages
  \(m_0, m_1 \in \mathcal M_{pp}\) and state \(st\).
\item The challenger samples \(r \xleftarrow{\$} \mathcal R_{pp}\)
  and computes \(c = \mathrm{Commit}(pp, m_b; r)\).
\item The adversary \(\mathcal A_2(st, c)\) outputs a bit \(b'\),
  which is the output of the experiment.
\end{enumerate}
The \emph{hiding advantage} of \(\mathcal A\) is
\[
\mathrm{Adv}^{\mathrm{hide}}_{\Pi}(\mathcal A, \lambda)
:= \Bigl|
   \Pr\bigl[\mathsf{Hide}^{\mathcal A}_{\Pi}(\lambda, 0) = 1\bigr]
 - \Pr\bigl[\mathsf{Hide}^{\mathcal A}_{\Pi}(\lambda, 1) = 1\bigr]
   \Bigr| .
\]
The scheme is \emph{computationally hiding} if the advantage is
negligible (\S\ref{model:subsec:indist}) for every PPT
\(\mathcal A\); \emph{statistically hiding} if it is negligible even
for computationally unbounded \(\mathcal A\); and \emph{perfectly
hiding} if it is exactly \(0\) for every adversary and every
\(\lambda\).
\end{definition}

\begin{proposition}[Distributional characterisation of hiding]
\label{commitments:prop:hiding-dist}
A commitment scheme \(\Pi\) is perfectly hiding if and only
if for every \(pp\) in the support of \(\mathrm{Setup}\) and every
pair \(m_0, m_1 \in \mathcal M_{pp}\), the random variables
\[
\mathrm{Commit}(pp, m_0; r)
\quad\text{and}\quad
\mathrm{Commit}(pp, m_1; r),
\qquad r \xleftarrow{\$} \mathcal R_{pp},
\]
are identically distributed. It is statistically hiding if the
statistical distance (Definition~\ref{model:def:statdist}) between
them is negligible. (Only the perfect clause is an equivalence:
parameters of negligible \(\mathrm{Setup}\) probability could
exhibit large distance without granting any adversary noticeable
advantage.)
\end{proposition}

\begin{proof}
Fix \(pp\) and \(m_0, m_1\), and write \(D_{m}\) for the distribution
of \(\mathrm{Commit}(pp, m; r)\) over uniform \(r\). By the
variational characterisation of statistical distance
(Proposition~\ref{model:prop:statdist}),
\(\Delta(D_{m_0}, D_{m_1})\) equals \emph{exactly} the maximum
advantage, over all decision procedures---efficient or not---of
telling \(D_{m_0}\) from \(D_{m_1}\). An adversary who outputs the
pair \((m_0, m_1)\) on parameters \(pp\) and applies the optimal test
therefore achieves hiding advantage
\(\Delta(D_{m_0}, D_{m_1})\) whenever \(\mathrm{Setup}\) produces
\(pp\); conversely no adversary can exceed the expectation of this
quantity over \(pp\). Advantage exactly \(0\) for all adversaries is
thus equivalent to \(\Delta = 0\)---identical distributions---for
every attainable \(pp\) and every pair, and negligible statistical
distance yields negligible advantage for unbounded adversaries. This
coincidence of the game formulation with the distance formulation is
precisely why statistical hiding was \emph{defined} through the
unbounded-adversary game.
\end{proof}

The binding game arms the equivocator. He is not handed an honest
commitment; he manufactures his own, together with two openings.

\begin{definition}[Binding game and binding]
\label{commitments:def:binding}
For a commitment scheme \(\Pi\) and adversary \(\mathcal A\), the
game \(\mathsf{Bind}^{\mathcal A}_{\Pi}(\lambda)\) runs as follows.
\begin{enumerate}
\item The challenger runs \(pp \gets \mathrm{Setup}(1^\lambda)\).
\item The adversary \(\mathcal A(pp)\) outputs a commitment \(c\) and
  two openings \((m_0, r_0)\) and \((m_1, r_1)\).
\item The game outputs \(1\) iff \(m_0 \ne m_1\) and
  \(\mathrm{Open}(pp, c, m_0, r_0) =
   \mathrm{Open}(pp, c, m_1, r_1) = 1\).
\end{enumerate}
The \emph{binding advantage} of \(\mathcal A\) is
\[
\mathrm{Adv}^{\mathrm{bind}}_{\Pi}(\mathcal A, \lambda)
:= \Pr\bigl[\mathsf{Bind}^{\mathcal A}_{\Pi}(\lambda) = 1\bigr],
\]
a search game in the sense of Definition~\ref{model:def:game}. The
scheme is \emph{computationally}, \emph{statistically}, or
\emph{perfectly binding} according as the advantage is negligible for
every PPT adversary, negligible for every unbounded adversary, or
exactly \(0\) for every adversary and every \(\lambda\).
\end{definition}

For a canonically opening scheme the binding game has a useful
collision reformulation: for fixed \(pp\), two openings collide
precisely when
\(\mathrm{Commit}(pp, m_0; r_0) = \mathrm{Commit}(pp, m_1; r_1)\)
with \(m_0 \ne m_1\). Perfect binding then says that distinct
messages never share a commitment value at all---the message \(m\) is
a deterministic function of \(c\). Perfect binding is thus equivalent
to the map \(c \mapsto m\) being well defined on all attainable
commitments, a formulation the counting argument of
Remark~\ref{commitments:rem:info-impossibility} exploits.

\begin{remark}[The asymmetry of the quantifiers]
\label{commitments:rem:asymmetry}
The two games quantify over different objects. Hiding quantifies over
what can be \emph{extracted from} a single honestly produced
commitment: the challenger seals the envelope, and the adversary only
looks. Binding quantifies over what can be \emph{manufactured}: the
adversary builds an ambiguous commitment of its own choosing, under
no obligation to follow \(\mathrm{Commit}\) honestly. A scheme can be
strong in one respect and weak in the other, and the two properties
actively pull in opposite directions---the more commitment values a
message can take (good for hiding), the more opportunities for two
messages to share one (bad for binding). The next subsection makes
this tension a theorem.
\end{remark}

\subsection{Incompatibility of perfect hiding and perfect binding}
\label{commitments:subsec:impossibility}

\begin{theorem}[No scheme is perfectly hiding and perfectly binding]
\label{commitments:thm:impossibility}
Let \(\Pi\) be a canonically opening commitment scheme with
\(\abs{\mathcal M_{pp}} \ge 2\). If \(\Pi\) is perfectly hiding, then
it is not perfectly binding: indeed, a computationally unbounded
adversary wins the binding game with probability \(1\). Hence no
nontrivial commitment scheme is simultaneously perfectly hiding and
perfectly binding.
\end{theorem}

\begin{proof}
Fix any \(pp\) in the support of \(\mathrm{Setup}(1^\lambda)\) and
any two distinct messages \(m_0, m_1 \in \mathcal M_{pp}\). For a
message \(m\), let
\[
C_m := \bigl\{\, \mathrm{Commit}(pp, m; r) :
  r \in \mathcal R_{pp} \,\bigr\}
\]
be the set of attainable commitments to \(m\), and let \(D_m\) be the
distribution of \(\mathrm{Commit}(pp, m; r)\) for uniform \(r\).
Perfect hiding makes \(D_{m_0}\) and \(D_{m_1}\) identical
(Proposition~\ref{commitments:prop:hiding-dist}); identical
distributions have identical supports, so \(C_{m_0} = C_{m_1}\).
Pick any \(c\) in this common support. By definition of the supports
there exist \(r_0, r_1 \in \mathcal R_{pp}\) with
\[
\mathrm{Commit}(pp, m_0; r_0) = c = \mathrm{Commit}(pp, m_1; r_1),
\]
and then \(\bigl(c, (m_0, r_0), (m_1, r_1)\bigr)\) wins the binding
game: both openings are valid by canonical opening, and
\(m_0 \ne m_1\). The double opening \emph{exists} for every
attainable \(pp\); an unbounded adversary \emph{finds} it by
enumerating \(\mathcal R_{pp}\), so it wins
\(\mathsf{Bind}\) with probability \(1\), giving
\(\mathrm{Adv}^{\mathrm{bind}}_{\Pi} = 1 \ne 0\).
\end{proof}

\begin{remark}[The impossibility as a counting fact]
\label{commitments:rem:info-impossibility}
Information-theoretically the theorem is a statement about how many
bits the commitment carries. Perfect binding forces the map
\(c \mapsto m\) to be well defined, so \(c\) determines \(m\)
completely: the commitment carries at least
\(\log_2 \abs{\mathcal M}\) bits of information about the message.
Perfect hiding forces the distribution of \(c\) to be independent of
\(m\): the commitment carries \emph{zero} bits. The two demands are
compatible only when \(\log_2 \abs{\mathcal M} = 0\), that is,
\(\abs{\mathcal M} = 1\)---when there is nothing to commit to. Every
real scheme therefore chooses which property to make perfect (or
statistical) and settles for the other being computational. The
Pedersen commitment of \S\ref{commitments:subsec:pedersen} is
perfectly hiding and computationally binding; the hash-based
commitment of \S\ref{commitments:subsec:hash}, read in the random
oracle model, is computationally binding (from collision resistance)
and computationally hiding.
\end{remark}

\begin{remark}[Two regimes, two failure modes]
\label{commitments:rem:failure-modes}
The choice is not merely aesthetic; it decides \emph{how} the scheme
fails if its assumption falls. If binding is only computational, an
adversary with enough power---one able to compute discrete
logarithms, say---could \emph{retroactively change} a committed
value, and the soundness of any protocol built on top degrades to a
computational assumption. If hiding is only computational, such an
adversary could \emph{retroactively learn} a committed value:
privacy degrades the same way, and everlasting secrecy is lost the
day the assumption breaks, even for commitments published long
before. Which failure is worse is a judgement about the application:
whether integrity or long-term confidentiality is the greater
concern. A ledger that publishes its commitments forever has reason
to care about the second failure mode; a protocol whose soundness
guards live funds has reason to care about the first.
\end{remark}

\subsection{Hash-based commitments}\label{commitments:subsec:hash}

The simplest construction seals the envelope with a hash. Recall from
\S\ref{sec:hash} the two properties it needs: a hash function
\(H \colon \{0,1\}^* \to \{0,1\}^n\) is \emph{collision resistant}
(Definition~\ref{hash:def:col}) if no PPT adversary finds
\(x \ne x'\) with \(H(x) = H(x')\) except with negligible
probability, and \emph{preimage resistant} (one-way) if, given
\(H(x)\) for random \(x\) of appropriate length, no PPT adversary
finds any \(x'\) with \(H(x') = H(x)\) except with negligible
probability.

\begin{construction}[Hash commitment]\label{commitments:def:hash-commit}
Let \(H \colon \{0,1\}^* \to \{0,1\}^n\) be a hash function. The
\emph{hash commitment} has message space \(\mathcal M = \{0,1\}^*\)
and randomness space \(\mathcal R = \{0,1\}^{\ell}\) for a randomness
length \(\ell = \ell(\lambda)\) (for example \(\ell = 2n\)). Let
\(\mathsf{enc}(m, r)\) be an injective, self-delimiting encoding of
the pair---for example a canonical length encoding of \(m\), followed
by \(m\), followed by the fixed-length string \(r\). Define
\[
\mathrm{Commit}(m; r) := H\bigl(\mathsf{enc}(m, r)\bigr),
\]
opened canonically. There are no nontrivial public parameters beyond
the choice of \(H\)---the scheme lives in the plain model of
Remark~\ref{commitments:rem:setup-trust}.
\end{construction}

\begin{proposition}[Binding of the hash commitment]
\label{commitments:prop:hash-binding}
If \(H\) is collision resistant, the hash commitment is
computationally binding.
\end{proposition}

\begin{proof}
Suppose \(\mathcal A\) wins the binding game: it outputs \(c\) with
valid openings \((m_0, r_0)\), \((m_1, r_1)\) and \(m_0 \ne m_1\), so
\(H(\mathsf{enc}(m_0, r_0)) = c = H(\mathsf{enc}(m_1, r_1))\).
Injectivity of \(\mathsf{enc}\) makes the two hash inputs distinct
strings---this is exactly what the self-delimiting encoding buys;
without it, a boundary shift could present distinct pairs as the
same string. The adversary has therefore produced an
\(H\)-collision, and the reduction that runs \(\mathcal A\) and
outputs the pair
\(\bigl(\mathsf{enc}(m_0, r_0),\; \mathsf{enc}(m_1, r_1)\bigr)\)
wins the collision game with exactly the binding advantage of
\(\mathcal A\), in essentially the same time.
\end{proof}

\begin{proposition}[Hiding of the hash commitment, in the ROM]
\label{commitments:prop:hash-hiding}
Model \(H\) as a random oracle (Definition~\ref{hash:def:ro}). For
messages of length \(\poly(\lambda)\) and randomness length
\(\ell = \omega(\log \lambda)\), the hash commitment is
computationally hiding: an adversary making at most \(q\) oracle
queries has hiding advantage at most \(q \cdot 2^{-\ell}\).
\end{proposition}

\begin{proof}
The challenge commitment is \(c = H(\mathsf{enc}(m_b, r))\) for
uniform \(r \in \{0,1\}^{\ell}\). Because the oracle's outputs are
independent and uniform, \(c\) is a fresh uniform string---carrying
no information about \(b\)---\emph{unless} the adversary queries the
oracle at the exact point \(\mathsf{enc}(m_b, r)\). Before any such
query the randomness \(r\) is information-theoretically hidden: the
adversary's view is independent of it. Each of the \(q\) queries
therefore hits that point with probability at most \(2^{-\ell}\), and the union
bound (\emph{Math Guide}, \S``The union bound and a birthday
calculation'') caps the probability of any hit by
\(q \cdot 2^{-\ell}\). Conditioned on no hit, the adversary's view is
identically distributed in the two worlds \(b = 0\) and \(b = 1\), so
the advantage is at most \(q \cdot 2^{-\ell}\), negligible for
\(\ell = \omega(\log \lambda)\) and polynomial \(q\).
\end{proof}

\begin{remark}[Why the randomness is essential]
\label{commitments:rem:hash-salt}
Dropping \(r\) is fatal. The deterministic commitment \(c = H(m)\)
fails hiding for every low-entropy message space: to distinguish a
commitment to \(m_0\) from one to \(m_1\), the peeker simply hashes
her two candidates and compares. Nothing about collision resistance
or one-wayness prevents this---the attack never inverts anything. The
salt \(r\) contributes \(\ell\) uniformly random bits, so the hash
input is unpredictable even when the message is one of two known
values. The same phenomenon renders deterministic encryption insecure
for low-entropy plaintexts, and the same cure---randomisation---applies.
\end{remark}

\begin{remark}[Strength profile]\label{commitments:rem:hash-profile}
The hash commitment has two virtues. It is \emph{transparent}: no
trusted setup, indeed no parameters, and no algebraic structure for a
subverted setup to poison. And it is \emph{plausibly post-quantum}:
generic quantum collision finding takes \(\Theta(2^{n/3})\) queries
against an \(n\)-bit hash
(Remark~\ref{assumptions:rem:shor-scope}), so increasing the output
length compensates for the loss. Its weakness is the flip side of its
opacity: the commitment is an unstructured bit string. One cannot add
two hash commitments and obtain a commitment to the sum of the
messages; no algebraic relation on commitments reflects any relation
on the committed values. The Pedersen construction repairs exactly
this defect, at the price of resting on discrete log.
\end{remark}

\subsection{The Pedersen commitment}\label{commitments:subsec:pedersen}

We change algebraic scenery. Let \(\G\) be a cyclic group of prime
order \(q\), written \emph{additively}: for a point \(P \in \G\) and
scalar \(a \in \Fq\), the scalar multiple is \([a]P\), with
\([0]P = \mathcal O\) the identity, and
\(a \mapsto [a]P\) is a group homomorphism \(\Fq \to \G\).
Concretely, \(\G\) is the group of \(\Fp\)-points of a prime-order
elliptic curve (\emph{Math Guide}, \S``The group structure of
\(E(\Fp)\)''), and in this volume the Pallas curve of
\S\ref{assumptions:subsec:pasta}; but the only structural facts the
constructions use are that \(\G\) is isomorphic to \(\Fq\) as a group
and that discrete logarithms in \(\G\) are hard. Formally, a
\emph{group generator} \(\mathcal G\) on input \(1^\lambda\) outputs
\((\G, q, G)\) with \(q\) a \(\Theta(\lambda)\)-bit prime and \(G\) a
generator, and the DLog assumption for \(\mathcal G\) states that no
PPT adversary, given \((\G, q, G, [a]G)\) for uniform
\(a \xleftarrow{\$} \Fq\), outputs \(a\) except with negligible
probability---Definition~\ref{assumptions:def:dlp}, transcribed
additively as in \S\ref{assumptions:subsec:pasta}.

\begin{construction}[Pedersen commitment]\label{commitments:def:pedersen}
The algorithm \(\mathrm{Setup}(1^\lambda)\) runs
\((\G, q, G) \gets \mathcal G(1^\lambda)\), samples
\(s \xleftarrow{\$} \Fq^{\times}\), sets \(H := [s]G\), outputs
\(pp = (\G, q, G, H)\), and \emph{discards} \(s\). The message and
randomness spaces are \(\mathcal M = \mathcal R = \Fq\), and
\[
\mathrm{Commit}(pp, m; r) := [m]G + [r]H,
\]
opened canonically.
\end{construction}

The parameters are two generators \(G\) and \(H\) whose relative
discrete logarithm \(s = \log_G H\) is unknown to \emph{all}
parties, the committer included. Everything hinges on this one
unknown scalar: knowledge of \(s\) breaks binding
(Remark~\ref{commitments:rem:equivocation}), while perfect hiding,
as the proof below shows, does not even require ignorance of \(s\).
In deployment the sampled-and-discarded \(s\) of the textbook
presentation is replaced by deriving \(H\) through hashing to the
curve (Construction~\ref{hash:con:nums}): nobody ever \emph{holds}
\(s\),
and the setup is transparent
(Remark~\ref{commitments:rem:setup-trust}).

\begin{theorem}[Perfect hiding]\label{commitments:thm:pedersen-hiding}
The Pedersen commitment is perfectly hiding.
\end{theorem}

\begin{proof}
Fix any attainable \(pp = (\G, q, G, H)\); since
\(s \in \Fq^{\times}\), the point \(H\) is not \(\mathcal O\).
Because \(\G\) has prime order, the nonzero \(H\) generates \(\G\),
so \(r \mapsto [r]H\) is a bijection \(\Fq \to \G\). For uniform
\(r\), therefore, \([r]H\) is uniform on \(\G\)---and so is its
translate \([m]G + [r]H\), for \emph{every} message \(m\): adding
the fixed element \([m]G\) permutes the group and leaves the uniform
distribution invariant. The commitment distribution is uniform on
\(\G\) regardless of \(m\), so for any \(m_0, m_1\) the two
distributions coincide exactly, and every adversary's advantage is
\(0\) by Proposition~\ref{commitments:prop:hiding-dist}.
\end{proof}

\begin{theorem}[Computational binding under DLog]
\label{commitments:thm:pedersen-binding}
If the DLog assumption holds for \(\mathcal G\), the Pedersen
commitment is computationally binding. Concretely, from any adversary
\(\mathcal A\) that wins the binding game with probability
\(\varepsilon\) one constructs \(\mathcal B\) computing discrete
logarithms with probability
\((1 - 1/q)\,\varepsilon \ge \varepsilon - 1/q\), in essentially the
same running time.
\end{theorem}

\begin{proof}
First the algebra. A double opening gives
\([m_0]G + [r_0]H = [m_1]G + [r_1]H\) with \(m_0 \ne m_1\), so
\[
[m_0 - m_1]G = [r_1 - r_0]H .
\]
The case \(r_1 = r_0\) is impossible: it would force
\([m_0 - m_1]G = \mathcal O\), and since \(G\) has order \(q\) this
means \(m_0 \equiv m_1 \pmod q\), contradicting \(m_0 \ne m_1\) in
\(\Fq\). With \(r_1 - r_0\) invertible modulo the prime \(q\), the
relation solves for the hidden discrete logarithm:
\[
s = \log_G H = (m_0 - m_1)(r_1 - r_0)^{-1} \bmod q .
\]
A double opening \emph{is} the discrete log of \(H\).

Now the reduction. On a DLog challenge \((\G, q, G, Q)\), the
algorithm \(\mathcal B\) sets \(H := Q\), hands
\(pp = (\G, q, G, Q)\) to \(\mathcal A\), and, when \(\mathcal A\)
returns a double opening, outputs
\((m_0 - m_1)(r_1 - r_0)^{-1} \bmod q = \log_G Q\). One distribution
gap needs accounting: in the real scheme \(H = [s]G\) with
\(s \xleftarrow{\$} \Fq^{\times}\), so \(H\) is uniform on
\(\G \setminus \{\mathcal O\}\), whereas the challenge \(Q = [a]G\)
has \(a\) uniform on all of \(\Fq\), so \(Q\) is uniform on all of
\(\G\). The two distributions differ only on the event
\(Q = \mathcal O\), of probability \(1/q\)---their statistical
distance is exactly \(1/q\)---and on that event \(\mathcal A\) cannot
win anyway (with \(H = \mathcal O\) the relation above forces
\(m_0 = m_1\)). Conditioned on \(Q \ne \mathcal O\), the simulation
is perfect, so
\[
\mathrm{Adv}^{\mathrm{dlog}}(\mathcal B)
= \Bigl(1 - \tfrac{1}{q}\Bigr)\,
  \mathrm{Adv}^{\mathrm{bind}}_{\Pi}(\mathcal A)
\;\ge\; \varepsilon - \tfrac{1}{q},
\]
and \(\mathcal B\)'s overhead over \(\mathcal A\) is a constant
number of field operations.
\end{proof}

The same embed-and-extract pattern, scaled up to many generators,
proves the vector generalisation
(Theorem~\ref{commitments:thm:vector-props}).

\begin{remark}[Equivocation with the trapdoor]
\label{commitments:rem:equivocation}
Any party who knows \(s = \log_G H\) can open a Pedersen commitment
to \emph{any} message. Given \(c = [m]G + [r]H\) and a target
\(m'\), set
\[
r' := r + (m - m')\,s^{-1} \bmod q ;
\]
then
\([m']G + [r']H = [m']G + [r]H + [(m - m')]G = c\). This trapdoor
property is not a defect of the construction---it is precisely why
discarding \(s\) in \(\mathrm{Setup}\) is mandatory, and it is a
feature the theory exploits deliberately: zero-knowledge simulators
use exactly this equivocation to open commitments consistently with
a transcript produced out of order. The same phenomenon priced the
setup models of Remark~\ref{commitments:rem:setup-trust}: a
commitment scheme with a trapdoor is only as binding as the claim
that nobody holds it.
\end{remark}

\begin{theorem}[Additive homomorphism]
\label{commitments:thm:pedersen-homo}
For all \(m_0, m_1, r_0, r_1 \in \Fq\),
\[
\mathrm{Commit}(m_0; r_0) + \mathrm{Commit}(m_1; r_1)
= \mathrm{Commit}(m_0 + m_1;\, r_0 + r_1),
\]
and more generally, for all scalars \(\alpha, \beta \in \Fq\),
\[
[\alpha]\,\mathrm{Commit}(m_0; r_0)
 + [\beta]\,\mathrm{Commit}(m_1; r_1)
= \mathrm{Commit}(\alpha m_0 + \beta m_1;\,
                  \alpha r_0 + \beta r_1).
\]
\end{theorem}

\begin{proof}
Expand and regroup, using commutativity of \(+\) in \(\G\), the
identity \([a]P + [b]P = [a + b]P\), and the linearity
\([\alpha]([m]G + [r]H) = [\alpha m]G + [\alpha r]H\) of scalar
multiplication:
\begin{align*}
[\alpha]\bigl([m_0]G + [r_0]H\bigr)
 + [\beta]\bigl([m_1]G + [r_1]H\bigr)
&= [\alpha m_0]G + [\beta m_1]G + [\alpha r_0]H + [\beta r_1]H \\
&= [\alpha m_0 + \beta m_1]G + [\alpha r_0 + \beta r_1]H .
\qedhere
\end{align*}
\end{proof}

\begin{remark}[Value balancing: the homomorphism deployed]
\label{commitments:rem:homo-uses}
The homomorphism is what makes Pedersen commitments the right tool
for \emph{hidden arithmetic}, and Zcash's value balancing is the
canonical deployment. Each transaction value \(v_i\) is committed as
\(C_i = [v_i]G + [r_i]H\); a transaction balances iff the sum of its
input values equals the sum of its output values. By
Theorem~\ref{commitments:thm:pedersen-homo}, the verifier can check
this \emph{without learning any \(v_i\)}: the value components
cancel in
\(\sum C_{\mathrm{in}} - \sum C_{\mathrm{out}}\) precisely when the
value sums agree \emph{modulo the group order} \(q\), in which case
the difference equals \([r]H\) for the net randomness \(r\)---a
commitment to \(0\). Cancellation alone certifies balance only
modulo \(q\)---value sums differing by a multiple of \(q\) cancel
too---and the deployed protocol closes the gap with range checks
carried inside the proofs: the values of the notes (the shielded
value records) spent and created by each Action (Orchard's combined
spend-and-output unit) are proved to lie in
\(\{0, \dots, 2^{64}-1\}\) (the Action statement, protocol
specification \href{https://zips.z.cash/protocol/protocol.pdf\#actionstatement}{\S~4.18.4}; enforced by the note-commitment gadget's
\texttt{ValueCanonicity} gate), the committed net
value is the difference of two such values, and the number of
summands is bounded, so the sums cannot wrap (\href{https://zips.z.cash/protocol/protocol.pdf\#orchardbalance}{\S~4.14}). The prover
does not reveal \(r\); it proves \emph{knowledge} of \(r\) by using
\([r]H\) as the verification key of a signature over the
transaction---the \emph{binding signature}, treated with the
completed balance argument in \S\ref{signatures:subsec:binding}. In
deployed Orchard the instantiation is
\(\mathsf{ValueCommit}(v; \mathsf{rcv}) = [v]V + [\mathsf{rcv}]R\)
with hashed-to-curve generators \(V, R\)
(protocol specification \href{https://zips.z.cash/protocol/protocol.pdf\#concretehomomorphiccommit}{\S~5.4.8.3}); each Action publishes a
commitment to its \emph{net}
value change, and the verifier folds the published commitments with
a commitment to the transparent value balance to obtain the binding
verification key. The
homomorphism turns an arithmetic relation on hidden values into a
group equation on public commitments---a template the rest of the
volume reuses repeatedly.
\end{remark}

\subsection{Pedersen vector commitments}
\label{commitments:subsec:pedersen-vector}

A single group element can seal far more than a single scalar.

\begin{construction}[Pedersen vector commitment]
\label{commitments:def:pedersen-vector}
On input \(1^\lambda\) and a length \(n = \poly(\lambda)\),
\(\mathrm{Setup}\) outputs \((\G, q)\) together with generators
\(G_1, \dots, G_n, H \xleftarrow{\$} \G \setminus \{\mathcal O\}\),
sampled so that no nontrivial discrete-log relation among them is
known to any party---in practice each generator is derived from a
public seed by hashing into the curve, rejecting the identity if it
occurs, a nothing-up-my-sleeve procedure
(Construction~\ref{hash:con:nums}) that makes the setup transparent.
The
message space is \(\Fq^{n}\), the randomness space \(\Fq\), and for
\(\mathbf m = (m_1, \dots, m_n)\),
\[
\mathrm{Commit}(pp, \mathbf m; r)
:= \sum_{i=1}^{n} [m_i]G_i + [r]H
\;=\; \ip{\mathbf m}{\mathbf G} + [r]H,
\]
where \(\ip{\mathbf m}{\mathbf G} := \sum_i [m_i]G_i\) is the
multiscalar (mixed) inner product of the message vector with the
generator vector \(\mathbf G = (G_1, \dots, G_n)\). Opening is
canonical.
\end{construction}

\begin{theorem}[Properties of the vector commitment]
\label{commitments:thm:vector-props}
The Pedersen vector commitment is perfectly hiding, and
computationally binding under the DLog assumption in \(\G\).
\end{theorem}

\begin{proof}
\emph{Hiding.} Exactly as in
Theorem~\ref{commitments:thm:pedersen-hiding}: for uniform \(r\) the
blinding term \([r]H\) is uniform on \(\G\), so the commitment is
uniform on \(\G\) regardless of the message vector.

\emph{Binding.} A double opening with
\(\mathbf m \ne \mathbf m'\) gives
\(\ip{\mathbf m}{\mathbf G} + [r]H
 = \ip{\mathbf m'}{\mathbf G} + [r']H\), that is, a
\emph{nontrivial discrete-log relation}
\[
\sum_{i=1}^{n} [m_i - m_i']\,G_i + [r - r']\,H = \mathcal O
\]
with not all coefficients zero. The reduction \(\mathcal B\) must
extract a single logarithm from a relation among \(n + 1\)
generators, and it does so by embedding its challenge in \emph{every}
generator. On challenge \(Q = [s]G\), it samples
\(u_i, w_i \xleftarrow{\$} \Fq\) for each \(i\) and
\(u_H, w_H \xleftarrow{\$} \Fq\), sets
\[
G_i := [u_i]G + [w_i]Q, \qquad H := [u_H]G + [w_H]Q,
\]
and resamples any pair whose generator lands on \(\mathcal O\). Each
result is uniform on \(\G \setminus \{\mathcal O\}\)---the
\([u_i]G\) term alone makes the pre-rejection point uniform on
\(\G\)---so the simulated parameters follow the real
distribution. Crucially, the parameters reveal nothing about the
\(w\)-coordinates: for every candidate value of \(w_i\) there is
exactly one \(u_i\) consistent with the observed \(G_i\), so
conditioned on the adversary's entire view the \(w_i\) and \(w_H\)
remain independent and uniform. Write the recovered relation as
\(\sum_i [a_i]G_i + [a_H]H = \mathcal O\) with coefficient vector
\((a_1, \dots, a_n, a_H) \ne \mathbf 0\). Substituting the embeddings
and collecting the \(G\)- and \(Q\)-components gives
\[
c_0 + c_1 s \equiv 0 \pmod q,
\qquad
c_0 = \sum_i a_i u_i + a_H u_H,
\quad
c_1 = \sum_i a_i w_i + a_H w_H .
\]
Since the coefficient vector is nonzero and the \(w\)-coordinates
are uniform and independent of it, \(c_1\) is uniform on \(\Fq\),
hence nonzero with probability \(1 - 1/q\); in that case
\(s = -c_0\, c_1^{-1} \bmod q\), and \(\mathcal B\) outputs the
challenge logarithm. Altogether
\(\mathrm{Adv}^{\mathrm{dlog}}(\mathcal B)
 \ge (1 - 1/q)\,
 \mathrm{Adv}^{\mathrm{bind}}(\mathcal A)\).
\end{proof}

\begin{remark}[Compression and the cardinality trade-off]
\label{commitments:rem:compressing}
The vector commitment is \emph{compressing}: \(n\) field elements
(plus the randomness) map to a single group element, exponentially
shorter than the message for large \(n\). Compression is
incompatible with statistical binding for a counting reason: there
are \(q^{n}\) message vectors and only \(q\) commitment values, so
colliding openings exist in overwhelming abundance. This is exactly
why binding is only computational---an unbounded adversary finds a
colliding opening by solving the underlying discrete-log
relation---and it is the trade-off of
Theorem~\ref{commitments:thm:impossibility} made visible at the
level of cardinalities: a compressing commitment \emph{cannot}
determine its message, so it might as well hide it perfectly.
\end{remark}

\begin{remark}[Linearity survives]\label{commitments:rem:vector-homo}
The vector commitment remains additively homomorphic in both message
and randomness, componentwise:
\(\mathrm{Commit}(\mathbf m; r) + \mathrm{Commit}(\mathbf m'; r')
 = \mathrm{Commit}(\mathbf m + \mathbf m';\, r + r')\), by the same
regrouping as Theorem~\ref{commitments:thm:pedersen-homo}. This
linearity, applied to a \emph{committed} vector against \emph{public}
coefficient vectors, is the foundation of the inner-product arguments
of \S\ref{commitments:subsec:poly}, which prove statements about
\(\ip{\mathbf m}{\mathbf b}\) for public \(\mathbf b\) while
revealing only a single committed group element.
\end{remark}

\subsection{Sinsemilla: an algebraic hash-based commitment}
\label{commitments:subsec:sinsemilla}

Between the opaque hash commitment of
\S\ref{commitments:subsec:hash} and the fully linear Pedersen
commitment sits an intermediate species: \emph{algebraic hash
functions}, collision-resistant hashes whose internal operations are
elliptic-curve group operations. The point of the species is proof
cost. A zero-knowledge circuit that must verify a hash evaluation
pays for every internal operation; when those operations are curve
additions in the very field the proof system works over, the
verification is cheap, while collision resistance still reduces to
the discrete logarithm. Sinsemilla, deployed in Zcash Orchard, is
the archetype: a Pedersen-style hash over a sequence of message
chunks. We build up to it through its ancestor.

\begin{construction}[Pedersen hash]
\label{commitments:def:pedersen-hash}
Split the message \(M\) into \(k\) chunks \(m_1, \dots, m_k\), each
interpreted as a scalar in a bounded range, and fix \(k\) independent
generators \(P_1, \dots, P_k \in \G\) with unknown mutual discrete
logarithms. Define
\[
\mathsf{PedersenHash}(M) := \sum_{j=1}^{k} [m_j]P_j \;\in\; \G .
\]
This is the unblinded Pedersen vector commitment map used as a hash.
The output is a group element; composing with a fixed encoding of
group elements as bit strings yields an ordinary hash.
\end{construction}

\begin{proposition}[Collision resistance of the Pedersen hash]
\label{commitments:prop:pedersen-hash-cr}
Suppose no efficient algorithm finds a nontrivial relation
\(\sum_j [a_j]P_j = \mathcal O\) with not all \(a_j = 0\) and each
\(\abs{a_j}\) within the allowed chunk range---hardness that follows
from DLog when the \(P_j\) are independent uniform generators. Then
\(\mathsf{PedersenHash}\) is collision resistant.
\end{proposition}

\begin{proof}
A collision \(M \ne M'\), with chunk decompositions \((m_j)\) and
\((m_j')\), gives
\(\sum_j [m_j - m_j']\,P_j = \mathcal O\) with some coefficient
\(m_j - m_j' \ne 0\)---a nontrivial relation with coefficients in
the allowed range. The reduction from relation-finding to DLog is
the every-generator embedding of
Theorem~\ref{commitments:thm:vector-props}: set each
\(P_j = [u_j]G + [w_j]Q\) for the challenge \(Q\), and read the
unknown logarithm off the recovered relation.
\end{proof}

The Pedersen hash needs one independent generator per chunk, which
is costly for long inputs: the generator table grows linearly with
the maximum message length. Sinsemilla reorganises the computation
into an iterated two-operand form that reuses a small fixed
generator table---shared by all chunk positions---while preserving
the discrete-log collision reduction.

\begin{construction}[Sinsemilla hash]
\label{commitments:def:sinsemilla-hash}
Split the message \(M\) into pieces \(m_1, \dots, m_k\), each of
fixed width \(w \le 32\) bits, and identify a piece
\(v \in \{0,1\}^{w}\) with its integer value, bits read
little-endian (the specification's \(\mathsf{LEBS2IP}\)). The public
parameters are a base point \(Q \in \G\) and a table
\(S \colon \{0,1\}^{w} \to \G\) assigning to each possible piece
value its own generator, both obtained by hashing to the curve.
Write \(\mathsf{GroupHash}(D, M)\) for the hash-to-curve pipeline of
Construction~\ref{hash:con:h2c} applied to the byte string \(M\)
under the domain-separation tag built from the byte string
\(D\)---for Pallas the RFC~9380 suite tag
\(D \,\|\, \texttt{"-pallas\_XMD:BLAKE2b\_SSWU\_RO\_"}\) of
\S\ref{hash:subsec:hashtofield} (protocol specification
\href{https://zips.z.cash/protocol/protocol.pdf\#concretegrouphashpallasandvesta}{\S~5.4.9.8})---and write \(\mathrm{LE}_{32}(j)\) for the four-byte
little-endian encoding of an integer \(0 \le j < 2^{32}\). The table
is
\[
S(v) := \mathsf{GroupHash}\bigl(\texttt{"z.cash:SinsemillaS"},\,
  \mathrm{LE}_{32}(v)\bigr), \qquad v \in \{0,1\}^{w},
\]
with \(2^{w}\) entries shared by every instance, and the base point
is personalised (\S\ref{hash:subsec:domainsep}) to the enclosing
use---a byte string \(D\) naming it, such as a note-commitment or
Merkle-hash domain---by taking that string as the message:
\[
Q := \mathsf{GroupHash}\bigl(\texttt{"z.cash:SinsemillaQ"},\, D\bigr).
\]
Both derivations are nothing-up-my-sleeve
(Construction~\ref{hash:con:nums}), so no party knows any
discrete-log relation among \(\{S(v)\}_v \cup \{Q\}\)
(Remark~\ref{hash:rem:nums-binding}). The deployed instantiation
takes \(w = 10\) (the specification's \(k\)) over Pallas: the
\texttt{sinsemilla} crate fixes the two
personalisation strings and derives \(Q\) in
\texttt{HashDomain::\allowbreak new}; the table ships precomputed as
\texttt{SINSEMILLA\_S}, and the
crate's test \texttt{sinsemilla\_s} re-derives every entry from
\texttt{u32::\allowbreak to\_le\_bytes} of the index---exactly
\(\mathrm{LE}_{32}\); the circuit loads that table as its lookup;
protocol specification \href{https://zips.z.cash/protocol/protocol.pdf\#concretesinsemillahash}{\S~5.4.1.9}. The hash iterates over the pieces
with a running accumulator \(\mathrm{Acc} \in \G\):
\[
\mathrm{Acc}_0 := Q, \qquad
\mathrm{Acc}_j := \mathrm{Acc}_{j-1}
  + \bigl(\mathrm{Acc}_{j-1} + S(m_j)\bigr)
  = [2]\mathrm{Acc}_{j-1} + S(m_j),
\]
and outputs the point
\(\mathsf{SinsemillaHashToPoint}(M) := \mathrm{Acc}_k\). The
step costs two \emph{incomplete} curve additions---first
\(\mathrm{Acc}_{j-1} + S(m_j)\), then the result plus
\(\mathrm{Acc}_{j-1}\); no doubling is ever performed, and the
addition formula is undefined when its operands are equal or are
mutual negatives (\emph{Math Guide}, \S``Elliptic
curves'')---and this pair of additions is the operation the
circuit verifies in very few constraints. Unrolling the recurrence gives
the closed form
\[
\mathsf{SinsemillaHashToPoint}(M)
= \sum_{j=1}^{k} \bigl[2^{\,k-j}\bigr] S(m_j) + \bigl[2^{k}\bigr] Q :
\]
the integer coefficients are powers of two determined by position
alone, and the message enters only through \emph{which} generator
\(S(m_j)\) each piece selects.
\end{construction}

\begin{proposition}[Collision resistance of Sinsemilla]
\label{commitments:prop:sinsemilla-cr}
Model the table entries \(S(v)\) and the base point \(Q\) as
independent uniform generators (the nothing-up-my-sleeve
hash-to-curve derivation covers \(Q\) as well as the table). Then
any collision in \(\mathsf{SinsemillaHashToPoint}\), on inputs for
which its incomplete additions are defined, yields a nontrivial
discrete-log relation among \(\{S(v)\}_v \cup \{Q\}\), and any input
on which an incomplete addition is \emph{undefined} yields one
directly; either way a discrete logarithm follows via the
every-generator embedding of
Theorem~\ref{commitments:thm:vector-props} extended to \(Q\). The
point hash is thus collision resistant under the DLog assumption in
\(\G\), the exceptional cases of incomplete addition covered by the
same reduction rather than excluded.
\end{proposition}

\begin{proof}[Proof sketch]
Let \(M \ne M'\) collide. Suppose first that they have the same
piece count \(k\). Subtracting the unrolled forms cancels the
\([2^{k}]Q\) terms and leaves
\[
\sum_{j=1}^{k} \bigl[2^{\,k-j}\bigr]
  \bigl( S(m_j) - S(m_j') \bigr) = \mathcal O .
\]
Regroup over the distinct table generators: for each value
\(v \in \{0,1\}^{w}\), the net coefficient on \(S(v)\) is
\[
c_v \;=\; \sum_{j\,:\,m_j = v} 2^{\,k-j}
        \;-\; \sum_{j\,:\,m_j' = v} 2^{\,k-j},
\]
and the relation reads \(\sum_v [c_v]\,S(v) = \mathcal O\). Each
\(c_v\) is a difference of two subset-sums of the distinct powers
\(2^{k-1}, \dots, 2^{0}\), hence an integer with
\(\abs{c_v} < 2^{k}\). For \emph{admissible} message lengths---those
with \(2^{k}\) below the group order \(q\), a finite design-time
check---such a coefficient vanishes modulo \(q\) only if it vanishes
as an integer; and \(c_v = 0\) forces, by uniqueness of binary
expansions, the two subset-sums to be equal as sets of powers, that
is, the positions carrying value \(v\) in \(M\) and in \(M'\) to
coincide. Coincidence for every \(v\) would force \(M = M'\); since
\(M \ne M'\), some \(c_v \ne 0\) and the relation is nontrivial.

If instead the piece counts differ, say \(k > k'\), the
\(Q\)-constants do not cancel: the relation acquires the coefficient
\(2^{k} - 2^{k'}\) on \(Q\), nonzero and below the group order for
admissible lengths, so the relation is again nontrivial. In either
case, embedding a DLog challenge in every table generator \emph{and}
in \(Q\) recovers the challenge logarithm exactly as in
Proposition~\ref{commitments:prop:pedersen-hash-cr}.

Finally, the exceptional cases. An incomplete addition is undefined
only when its operands are equal or are mutual negatives, so an
input reaching one presents a coincidence
\([\alpha]\mathrm{Acc}_{j-1} + S(m_j) = \mathcal O\) with
\(\alpha \in \{-1, 1, 2\}\). Substituting the unrolled form of
\(\mathrm{Acc}_{j-1}\) makes this an explicit relation among \(Q\)
and the table generators, with coefficient
\(\alpha\,2^{\,j-1} \ne 0\) on \(Q\) and all magnitudes at most
\(2^{k}\)---below the group order for admissible lengths. An input
that merely triggers an exceptional case therefore hands over a
nontrivial relation itself, and the same embedding extracts the
challenge logarithm from it; the protocol specification proves
exactly this for the deployed instantiation (\href{https://zips.z.cash/protocol/protocol.pdf\#concretesinsemillahash}{\S~5.4.1.9}).
\end{proof}

The deployed hash differs from the point hash in two ways. First,
the deployed \(\mathsf{SinsemillaHash}\) applies the Pallas
coordinate extractor to \(\mathsf{SinsemillaHashToPoint}(M)\),
returning its \(x\)-coordinate rather than the point itself.
Collision resistance survives the extraction: equal
extracted \(x\)-coordinates mean the two points are equal or are
negatives of one another, and either signed equality yields the same
kind of nontrivial generator relation as in the proof above. Second,
the final piece is zero-padded to the piece width---the \texttt{Pad}
iterator of the \texttt{sinsemilla} crate pads
the message bits to a multiple of the piece width before
chunking---so two bit strings of different lengths within the same
final piece present identical piece sequences, and the deployed hash
claims collision resistance only \emph{between inputs of a fixed
length for a given personalisation}. This is exactly the property
the protocol specification requires and no more (protocol
specification \href{https://zips.z.cash/protocol/protocol.pdf\#concretesinsemillahash}{\S~5.4.1.9}: collision resistance is demanded
``between inputs of fixed length, for a given personalization
input''); each personalisation domain hashes inputs of a single
length, as the Merkle-tree hash \(\mathsf{MerkleCRH}\) of
\S\ref{sec:merkle} illustrates. The variable-length argument in the
proof above applies to the piece-aligned point presentation, where
the piece sequence is the message.

The step from hash to commitment is the Pedersen blinding term,
verbatim.

\begin{construction}[SinsemillaCommit]
\label{commitments:def:sinsemilla-commit}
Fix one further independent generator \(H'\), hashed to the curve
with its own domain string, so that no discrete-log relation between
\(H'\) and \(\{S(v)\}_v \cup \{Q\}\) is known. Define, for
\(r \xleftarrow{\$} \Fq\),
\[
\mathsf{SinsemillaCommit}(m; r)
:= \mathsf{SinsemillaHashToPoint}(m) + [r]H' .
\]
Hiding is inherited from Pedersen verbatim: by the argument of
Theorem~\ref{commitments:thm:pedersen-hiding}, \([r]H'\) is uniform
on \(\G\) for uniform \(r\), so the commitment distribution is
independent of \(m\), and the scheme is perfectly hiding. Binding
rests on discrete log: two valid openings to distinct messages yield
either a \(\mathsf{SinsemillaHashToPoint}\) collision---excluded by
Proposition~\ref{commitments:prop:sinsemilla-cr}---or a nontrivial
discrete-log relation between \(H'\) and the hash generators.
\end{construction}

Orchard's note commitment \(\mathsf{NoteCommit}\) instantiates
exactly this map: the \texttt{orchard} crate derives it through the
\texttt{sinsemilla} crate's \texttt{CommitDomain::\allowbreak commit}
(protocol specification
\href{https://zips.z.cash/protocol/protocol.pdf\#concretesinsemillacommit}{\S~5.4.8.4}). Its key commitment \(\mathsf{CommitIvk}\) composes the
same map with \(x\)-coordinate extraction---the specification's
\(\mathsf{SinsemillaShortCommit}\)---returning a base-field element
rather than a point.

\begin{remark}[Why Sinsemilla exists: cost inside a proof]
\label{commitments:rem:sinsemilla-circuit}
Sinsemilla's advantage is not speed outside a circuit---ordinary
hashes are faster there---but verification cost \emph{inside} one.
In a SNARK (\S\ref{sec:zk}) the prover re-expresses every
computational step as polynomial constraints; a bit-oriented hash such
as SHA-256, all Boolean choppings and modular additions foreign to the
proof field, costs tens of thousands of constraints per evaluation,
whereas Sinsemilla's per-piece cost is a handful of curve operations
chosen to lie in the very field the proof system already works over. The
result is a collision-resistant hash and commitment---used for note
commitments and Merkle trees in Orchard---simultaneously cheap to
prove and grounded in the same discrete-log hardness as the
surrounding protocol. The presentation here is the conceptual core;
the \emph{Ironwood Guide} takes up the concrete Orchard
instantiation---its per-use personalisation strings and the
in-circuit handling of incomplete addition.
\end{remark}

\subsection{Polynomial commitment schemes}
\label{commitments:subsec:poly}

The final refinement commits not to a value or a vector but to a
\emph{function}---a polynomial---and supports proving statements
about its values. This is the primitive on which the succinct
argument systems of \S\ref{sec:zk} stand.

\begin{definition}[Polynomial commitment scheme]
\label{commitments:def:pcs}
Fix a field \(\F\) and a degree bound \(d\). A \emph{polynomial
commitment scheme} (PCS) is a tuple
\((\mathrm{Setup}, \mathrm{Commit}, \mathrm{Open}, \mathrm{Eval},
\mathrm{Verify})\):
\begin{itemize}
\item The algorithm \(\mathrm{Setup}(1^\lambda, d)\) outputs public
  parameters \(pp\) supporting degree at most \(d\);
\item the algorithm \(\mathrm{Commit}(pp, f; r)\) outputs a
  commitment \(C\) to a polynomial \(f \in \F[X]_{\le d}\),
  optionally with hiding randomness \(r\);
\item the algorithm \(\mathrm{Open}(pp, C, f, r)\) outputs a bit,
  verifying that \(C\) commits to the \emph{whole} polynomial \(f\);
\item the algorithm \(\mathrm{Eval}(pp, f, r, z)\), for an
  evaluation point \(z \in \F\), outputs a claimed value
  \(y = f(z)\) together with an \emph{evaluation proof} \(\pi\);
\item the algorithm \(\mathrm{Verify}(pp, C, z, y, \pi)\) outputs a
  bit, checking the evaluation proof against the commitment.
\end{itemize}
Beyond the hiding and binding of the underlying commitment
(Definitions~\ref{commitments:def:hiding}
and~\ref{commitments:def:binding}), a PCS must satisfy
\emph{evaluation binding}: no PPT prover can produce a commitment
\(C\), a point \(z\), and two accepting proofs for distinct claimed
values \(y \ne y'\) at \(z\), except with negligible probability.
Stronger schemes additionally satisfy \emph{extractability} (proof
of knowledge): from any successful prover one can extract an actual
polynomial \(f\) of degree at most \(d\) consistent with all the
evaluations it gets accepted.
\end{definition}

\begin{remark}[Rigidity, and why evaluation binding is separate]
\label{commitments:rem:rigidity}
Why should one short commitment support probing at arbitrary points?
Because polynomials are \emph{rigid}: a nonzero polynomial of degree
at most \(d\) has at most \(d\) roots (\emph{Math Guide}, \S``The
Schwartz--Zippel lemma''), so two distinct polynomials of degree at
most \(d\) agree on at most \(d\) points. Answering a different
value at even a single point is the behaviour of a prover holding a
\emph{different polynomial}---there is no wiggle room between
functions. But a warning is in order: an accepting evaluation proof
is \emph{not} an opening of \(C\), so evaluation binding does
\emph{not} follow from the binding of \(\mathrm{Commit}\). It is a
separate cryptographic requirement, and each construction below
proves it from its own assumption---a strong Diffie--Hellman
assumption for KZG, discrete log via a folding extractor for the
inner-product scheme. The central design problem of a PCS is making
the evaluation proof \(\pi\) short and fast to verify; committing to
the coefficient list with any old commitment is easy, and useless.
\end{remark}

\subsubsection*{KZG: evaluation in the exponent}

The pairing-based scheme of Kate, Zaverucha, and Goldberg achieves
constant-size proofs by hiding a secret evaluation point inside the
parameters.

\begin{construction}[KZG commitment]\label{commitments:def:kzg}
Let \(\G_1, \G_2, \G_T\) be groups of prime order \(\abs{\F}\)
equipped with a nondegenerate bilinear pairing
\(e \colon \G_1 \times \G_2 \to \G_T\)
(Example~\ref{assumptions:exa:pairing}), with generators
\(G \in \G_1\) and \(\hat G \in \G_2\). The algorithm
\(\mathrm{Setup}\) samples a secret \(\tau \xleftarrow{\$} \F\),
publishes the \emph{structured reference string}
\[
pp = \bigl( [\tau^0]G,\, [\tau^1]G,\, \dots,\, [\tau^d]G,\;
            \hat G,\, [\tau]\hat G \bigr),
\]
and \emph{discards} \(\tau\). The commitment to
\(f = \sum_{i=0}^{d} a_i X^{i}\) is evaluation in the exponent:
\[
\mathrm{Commit}(pp, f)
:= \sum_{i=0}^{d} [a_i]\bigl([\tau^{i}]G\bigr)
 = [f(\tau)]G,
\]
computable from \(pp\) by anyone, without knowledge of \(\tau\).
\end{construction}

\begin{construction}[KZG evaluation proof]
\label{commitments:def:kzg-eval}
To prove \(f(z) = y\), the prover uses the factor theorem: \(z\) is
a root of \(f(X) - y\), so the quotient
\[
q(X) := \frac{f(X) - y}{X - z} \;\in\; \F[X]_{\le d-1}
\]
is a genuine polynomial \emph{precisely because} \(f(z) - y = 0\).
The proof is the commitment to the quotient,
\(\pi := [q(\tau)]G\), computed from the reference string. The
verifier checks the single pairing equation
\[
e\bigl(C - [y]G,\; \hat G\bigr)
\;\stackrel{?}{=}\;
e\bigl(\pi,\; [\tau]\hat G - [z]\hat G\bigr),
\]
which, by bilinearity, is exactly the in-the-exponent form of the
polynomial identity \(f(\tau) - y = q(\tau)\,(\tau - z)\).
\end{construction}

\begin{proposition}[KZG evaluation binding, sketched]
\label{commitments:prop:kzg-binding}
Evaluation binding of KZG rests on the \(d\)-strong Diffie--Hellman
assumption in the bilinear group: given the powers
\([\tau^{i}]G\), it is infeasible to produce
\([(\tau - z)^{-1}]G\) for any \(z\) of one's choosing.
\end{proposition}

\begin{proof}[Proof idea]
Two accepting proofs \(\pi, \pi'\) for distinct values
\(y \ne y'\) at the same point \(z\) satisfy the two pairing
equations; subtracting them (in the exponent) gives
\(\pi - \pi' = [(y' - y)(\tau - z)^{-1}]G\), so the forger's output
yields
\[
\bigl[(\tau - z)^{-1}\bigr]G
= \bigl[(y' - y)^{-1}\bigr](\pi - \pi'),
\]
exactly the element the assumption declares infeasible.
\end{proof}

\begin{remark}[The KZG trade]\label{commitments:rem:kzg-trade}
What KZG buys: commitments and evaluation proofs of \emph{constant}
size (one group element each) and \emph{constant} verification time
(two pairings), all independent of the degree \(d\). What it costs
is twofold. First, the setup is trusted in the strongest sense: the
secret \(\tau\) is a global trapdoor. Anyone who learns \(\tau\) can
forge accepting evaluation proofs for \emph{false} statements---for
any commitment \(C\), any point \(z\), and any desired value \(y\),
the element
\(\pi := [(\tau - z)^{-1}]\bigl(C - [y]G\bigr)\) passes the pairing
check, no knowledge of any polynomial required. This is the promised
demonstration that a subverted setup silently destroys binding
(Remark~\ref{commitments:rem:setup-trust}); evaluation binding
survives only if \(\tau\) is provably destroyed, in practice by a
multi-party \emph{powers-of-tau} ceremony whose trapdoor stays
secret as long as at least one participant is honest. Second, the
construction requires pairing-friendly curves, a more delicate and
less efficient setting than the ordinary curves of
\S\ref{assumptions:subsec:pasta}. Binding is computational either
way, resting on the bilinear assumption---consistent with
Theorem~\ref{commitments:thm:impossibility}.
\end{remark}

\subsubsection*{The inner-product alternative}

The transparent alternative commits to the coefficient vector with
the Pedersen vector commitment and replaces the pairing check by a
recursive argument.

\begin{construction}[IPA-style polynomial commitment]
\label{commitments:def:ipa}
Let \(\mathbf G = (G_0, \dots, G_d)\) and \(H\) be transparent
generators with unknown mutual discrete logs, derived from a public
seed (Construction~\ref{commitments:def:pedersen-vector}). Commit to
\(f = \sum_{i=0}^{d} a_i X^{i}\) by the Pedersen vector commitment
of its coefficient vector \(\mathbf a = (a_0, \dots, a_d)\):
\[
\mathrm{Commit}(pp, f; r)
:= \sum_{i=0}^{d} [a_i]G_i + [r]H
 = \ip{\mathbf a}{\mathbf G} + [r]H .
\]
No pairings, no trusted setup.
\end{construction}

Evaluation is an inner product: writing
\(\mathbf z = (1, z, z^{2}, \dots, z^{d})\) for the public vector of
powers of the query point,
\[
f(z) = \sum_{i=0}^{d} a_i z^{i} = \ip{\mathbf a}{\mathbf z},
\]
so proving an evaluation reduces to proving an inner-product
relation about a committed vector against a public vector. The
\emph{inner-product argument} (IPA) proves it recursively: each
round folds the two halves of the coefficient vector---and of the
generator vector---into vectors of half the dimension by taking a
random linear combination dictated by a verifier challenge, sending
two group elements per round (the cross-term commitments \(L_j\) and
\(R_j\)) to account for the mixed terms; after
\(\log_2 (d + 1)\) rounds a single scalar remains, and the verifier
checks one final equation tying it to the folded commitment and the
folded generators.

\begin{proposition}[Properties of the IPA-based PCS]
\label{commitments:prop:ipa-props}
The inner-product polynomial commitment has: transparent setup
(public-coin parameters, no trapdoor); evaluation proofs of
\(O(\log d)\) group elements; and \(O(d)\)-group-operation
verification, batchable across proofs. Its binding and evaluation
binding reduce to the discrete logarithm assumption in \(\G\).
\end{proposition}

\begin{proof}[Proof idea]
Binding of the commitment itself is
Theorem~\ref{commitments:thm:vector-props}. For evaluation binding
and extractability, one shows the folding protocol is \emph{tree
special-sound}: from a suitable tree of accepting transcripts,
branching over the verifier's challenge at each round, an extractor
works back up the recursion and computes a coefficient vector
\(\mathbf a\) consistent with the commitment and the claimed inner
product---a notion this sketch uses only informally and
\S\ref{sec:zk} defines precisely. Two distinct extracted
openings of the same commitment would yield a nontrivial
discrete-log relation among \(\mathbf G, H\), contradicting DLog by
the every-generator embedding. The \(O(\log d)\) proof size is the
halving recursion; the \(O(d)\) verification is the final
multiscalar multiplication over the unfolded generators.
\end{proof}

\begin{remark}[KZG versus IPA, and the Halo lineage]
\label{commitments:rem:kzg-vs-ipa}
The two families mark out the design space. KZG gives constant
proof size and constant verification, at the price of a structured
trusted setup and pairing-friendly curves; the IPA gives transparent
setup over a plain curve and rests on nothing beyond discrete log,
at the price of logarithmic proofs and linear-time verification.
The Halo~2 proof system deployed in Zcash Orchard builds on the
inner-product family precisely to avoid trusted setup, and recovers
practical verification cost through \emph{batching}: the deployed
verifier folds the final multiscalar multiplications of many proofs
into one shared multiexponentiation---each proof's deferred check is
scaled by a fresh random factor and accumulated, and a single
combined check is evaluated once. Batching amortises the per-proof group
operations across a batch, though per-proof verifier work remains
linear in \(d\). The Halo lineage's accumulation technique extends
the same deferral across \emph{recursively composed} proofs, one
proof attesting to another's deferred check---the capability the
Pasta cycle of Remark~\ref{assumptions:rem:cycle} was designed to
serve---but deployed Zcash does not use recursion: each Orchard
proof is verified (in batch) directly.
\end{remark}

\begin{remark}[Hiding variants]\label{commitments:rem:pcs-hiding}
As stated, the KZG commitment is binding but not hiding: it is a
deterministic function of \(f\), so committing twice
to the same polynomial is visible, and a low-entropy polynomial can
be found by search---the failure of
Remark~\ref{commitments:rem:hash-salt} in new clothing. That is
acceptable when the polynomial encodes only public data. Whenever it
encodes secrets---as the witness polynomials of a zero-knowledge
proof do---randomness is required: KZG is augmented with a blinding
summand \([r][\gamma]G\) using an extra setup element
\([\gamma]G\), while the IPA commitment of
Construction~\ref{commitments:def:ipa} already carries its \([r]H\)
and is perfectly hiding by
Theorem~\ref{commitments:thm:vector-props}. The resulting regime is
the perfect-hiding/computational-binding profile of Pedersen
(Theorems~\ref{commitments:thm:pedersen-hiding}
and~\ref{commitments:thm:pedersen-binding}), lifted from field
elements to polynomials.
\end{remark}

\subsubsection*{From commitments to arguments}

The polynomial commitment is where the section's main line hands
over to the proof systems of \S\ref{sec:zk}, and the hand-over can
be stated in one paragraph. A polynomial commitment scheme turns an
information-theoretic proof object---a polynomial identity holding
over a large field---into a succinct cryptographic argument. The
prover commits to its witness polynomials; the verifier challenges
at random points; evaluation proofs then convince the verifier that
the committed polynomials satisfy the required identities, without
the polynomials ever being revealed. Soundness is supplied by
Schwartz--Zippel (\emph{Math Guide}, \S``The Schwartz--Zippel
lemma''): a nonzero polynomial of low degree is almost surely
nonzero at a random point, so an identity that holds at the
challenge almost surely holds identically. Binding and succinctness
are supplied by the commitment: the prover is locked to its
polynomials \emph{before} seeing the challenge, and the verifier
handles only short commitments and proofs in place of the
polynomials themselves. How this compilation is carried out in
full---and what ``argument of knowledge'' precisely means---is the
business of \S\ref{sec:zk}, whose SNARK compilation subsection
develops the evaluation-argument role sketched here.

\subsection{Secret sharing}\label{commitments:subsec:secret-sharing}

The section opened by naming verifiable secret sharing among the
classical applications of commitments; it closes by constructing it,
because the threshold signatures of the \emph{FROST Guide} stand on
it and nothing below that volume otherwise provides it. The setting
is a prime field \(\Fq\)---in the sequel the scalar field of the
group \(\G\) of \S\ref{commitments:subsec:pedersen}---and the tool is
Lagrange interpolation (\emph{Math Guide}, \S``Lagrange
interpolation''): a polynomial of degree at most \(t - 1\) over
\(\Fq\) is determined by its values at any \(t\) distinct points, and
any \(t\) prescribed values at distinct points are attained by
exactly one such polynomial.

\begin{construction}[Shamir \(t\)-of-\(n\) secret sharing]
\label{commitments:def:shamir}
Fix a threshold \(t\) and a party count \(n\) with
\(1 \le t \le n < q\). To share a secret \(s \in \Fq\), the dealer
samples coefficients \(a_1, \dots, a_{t-1} \xleftarrow{\$} \Fq\)
independently and uniformly and forms
\[
f(X) := s + a_1 X + a_2 X^{2} + \dots + a_{t-1} X^{t-1}
\;\in\; \Fq[X], \qquad f(0) = s ,
\]
a polynomial of degree at most \(t - 1\) with constant term the
secret. Party \(i \in \{1, \dots, n\}\) receives the \emph{share}
\(s_i := f(i)\), sent privately; the identifiers \(1, \dots, n\) are
distinct nonzero elements of \(\Fq\) because \(n < q\).
\emph{Reconstruction} from the shares of any set \(I\) of \(t\)
parties is interpolation at \(0\):
\[
s \;=\; f(0) \;=\; \sum_{i \in I} \lambda_{I,i}\, s_i ,
\qquad
\lambda_{I,i} := \prod_{j \in I,\; j \ne i} \frac{j}{j - i} ,
\]
where \(\lambda_{I,i} = \ell_i(0)\) is the Lagrange basis polynomial
for the nodes \(I\) evaluated at \(0\), computable from the
identifiers alone.
\end{construction}

Correctness of reconstruction is the interpolation theorem: the \(t\)
points \((i, s_i)_{i \in I}\) determine a unique polynomial of degree
at most \(t - 1\), which is \(f\), and the Lagrange formula
\(f(X) = \sum_{i \in I} s_i\, \ell_i(X)\) evaluated at \(X = 0\)
gives the displayed sum. The coalition may be chosen after the
shares are dealt, since the weights depend only on \(I\).

\begin{theorem}[Perfect secrecy below the threshold]
\label{commitments:thm:shamir-secrecy}
Let \(I\) be any set of at most \(t - 1\) identifiers. For every
secret \(s \in \Fq\), the shares \((s_i)_{i \in I}\) of
Construction~\ref{commitments:def:shamir} are jointly uniform on
\(\Fq^{\abs{I}}\). In particular their distribution does not depend
on \(s\): fewer than \(t\) shares are independent of the secret, and
an adversary of unbounded power holding them has advantage exactly
\(0\) at telling any two candidate secrets apart.
\end{theorem}

\begin{proof}
It suffices to treat \(\abs{I} = t - 1\), since fewer shares are a
projection of these. Fix \(s\) and \(I\), and consider the map
\[
\Sigma_s \colon \Fq^{t-1} \to \Fq^{t-1}, \qquad
(a_1, \dots, a_{t-1}) \mapsto \bigl(f(i)\bigr)_{i \in I} ,
\]
from the dealer's coefficient vector to the coalition's shares. It
is injective: two coefficient vectors with the same shares give two
polynomials of degree at most \(t - 1\) that agree at the \(t\)
distinct points \(0\) (where both equal \(s\)) and \(I\), hence are
equal by uniqueness of interpolation, so the coefficient vectors
coincide. It is surjective: for any target \((y_i)_{i \in I}\),
Lagrange interpolation through the \(t\) points \((0, s)\) and
\((i, y_i)_{i \in I}\) produces a polynomial of degree at most
\(t - 1\) with constant term \(s\), whose remaining coefficients are
a preimage. So \(\Sigma_s\) is a bijection---for \emph{every}
\(s\)---and it carries the uniform distribution on coefficient
vectors to the uniform distribution on \(\Fq^{t-1}\). The share
vector is therefore uniform whatever the secret, and for any
\(s_0 \ne s_1\) the two share distributions coincide exactly; by the
variational characterisation of statistical distance
(Proposition~\ref{model:prop:statdist}) no procedure, efficient or
not, distinguishes them.
\end{proof}

The theorem is the perfect-hiding half of a commitment
(Proposition~\ref{commitments:prop:hiding-dist} in different
clothing): the coalition's view is a uniformly random point of
\(\Fq^{t-1}\), consistent with every candidate secret through
exactly one polynomial each. What Shamir sharing lacks is the
binding half. Nothing stops a dishonest dealer from handing out
values that lie on no common polynomial of degree at most \(t - 1\),
so that different coalitions of \(t\) parties reconstruct different
``secrets''---and a party receiving a share cannot tell. Verifiable
secret sharing repairs this by having the dealer commit to the
polynomial.

\begin{construction}[Feldman verifiable secret sharing]
\label{commitments:def:feldman}
Let \((\G, q, G)\) be as in \S\ref{commitments:subsec:pedersen}, so
that \(\Fq\) is the scalar field of \(\G\). The dealer shares \(s\)
as in Construction~\ref{commitments:def:shamir}, with
\(f(X) = \sum_{k=0}^{t-1} a_k X^{k}\) and \(a_0 = s\), and
additionally \emph{broadcasts} the coefficient commitments
\[
\Phi_k := [a_k]\,G \;\in\; \G, \qquad k = 0, \dots, t - 1 ,
\]
the coefficients in the exponent. Party \(i\), on receiving its
private share \(s_i\), accepts it iff
\[
[s_i]\,G \;\stackrel{?}{=}\; \sum_{k=0}^{t-1} \bigl[i^{k}\bigr]\Phi_k .
\]
An honest share passes: since \(a \mapsto [a]G\) is a homomorphism
\(\Fq \to \G\),
\(\sum_k [i^{k}]\Phi_k = \sum_k [a_k i^{k}]G = [f(i)]G\).
\end{construction}

\begin{proposition}[What the Feldman check certifies]
\label{commitments:prop:feldman-check}
Because \(G\) generates the prime-order group \(\G\), the map
\(a \mapsto [a]G\) is a bijection \(\Fq \to \G\), so the broadcast
vector \((\Phi_0, \dots, \Phi_{t-1})\)---however the dealer chose
it---determines a unique polynomial
\(f^{*}(X) := \sum_k a_k^{*} X^{k}\) with
\(a_k^{*} = \log_G \Phi_k\), of degree at most \(t - 1\). Party
\(i\)'s check passes iff \(s_i = f^{*}(i)\). Consequently every
coalition of \(t\) parties holding accepted shares reconstructs the
same polynomial \(f^{*}\) and the same secret
\(s^{*} = f^{*}(0) = \log_G \Phi_0\), and a party handed an
inconsistent share detects it alone and with certainty. The check
certifies nothing about the value of \(s^{*}\), which the dealer
chooses freely.
\end{proposition}

\begin{proof}
Both claims are injectivity of \(a \mapsto [a]G\). The right-hand
side of the check equals \([f^{*}(i)]G\) by the homomorphism, so the
equation \([s_i]G = [f^{*}(i)]G\) holds iff \(s_i = f^{*}(i)\).
Accepted shares are thus evaluations of the one polynomial
\(f^{*}\), and \(t\) of them interpolate to it by uniqueness.
\end{proof}

In the language of this section, the vector \((\Phi_k)_k\) is a
\emph{perfectly binding} commitment to the polynomial \(f\)---the
coefficient list committed entry by entry under the bijection
\(a \mapsto [a]G\)---and the share check is an evaluation check
against it, verifiable by anyone because it is linear in the
coefficients. Perfect binding costs hiding, as
Theorem~\ref{commitments:thm:impossibility} demands.

\begin{remark}[The leak, and where it is harmless]
\label{commitments:rem:feldman-leak}
Feldman sharing publishes \(\Phi_0 = [s]G\). Secrecy below the
threshold is no longer perfect: from the broadcast alone an
unbounded adversary computes \(s = \log_G \Phi_0\). For an efficient
coalition of \(t - 1\) parties the leak is exactly \([s]G\) and
nothing more: its entire view---shares and broadcast---is computable
from \([s]G\) and \(t - 1\) uniformly random field elements, because
the shares are uniform by
Theorem~\ref{commitments:thm:shamir-secrecy} and the \(\Phi_k\) are
then determined by interpolating in the exponent through
\((0, [s]G)\) and the points \((i, [s_i]G)\). Recovering \(s\) from
that view is the discrete-logarithm problem
(Definition~\ref{assumptions:def:dlp}). The leak is acceptable
exactly when \(s\) is itself a discrete-logarithm secret key whose
public key \([s]G\) is published anyway; then \(\Phi_0\) \emph{is}
the public key and nothing new escapes. When the secret must stay
perfectly hidden, Pedersen verifiable secret sharing replaces each
\(\Phi_k\) by the Pedersen commitment \([a_k]G + [b_k]H\) of
Construction~\ref{commitments:def:pedersen}, restoring perfect
hiding at the price of computational binding---off this volume's
path and not constructed here. Threshold signing is the consumer:
in the \emph{FROST Guide} each signer holds a Shamir share of a
signing key, the coefficient commitments broadcast during key
generation are Feldman's, and the Lagrange weights
\(\lambda_{I,i}\) recombine per-signer contributions into a single
signature.
\end{remark}

\section{Pseudorandom functions, block ciphers, and format-preserving
encryption}\label{sec:prf}

The adversary of this section sits before a black box. She may submit
any input she likes and read off the answer; she may let each query
depend on everything the box has said so far; and when she has asked
all she cares to ask, she must announce a single bit: does a short
secret key drive the box, or blind chance? If no efficient strategy
tells the difference, then \(\lambda\) bits of key have bought a
perfect counterfeit of an object too large to exist in the world---a
uniformly random function---and everything downstream that needs
random-looking values on demand can draw them from that key. If she
\emph{can} tell, forgery follows at once. Key a Merkle--Damg\aa rd
hash by naive concatenation and she extends an authentication tag
she was shown onto a message nobody authenticated, without ever
learning the key. Derive
the diversifiers behind Orchard's shielded addresses by any map she
can predict or invert, and she links every address of one wallet to
its siblings. One adversary---the \emph{oracle distinguisher}---thus
carries the whole section, and every definition in it is the same
distinguishing game at a different strength: pseudorandom functions
against random functions, pseudorandom permutations against random
permutations, strong pseudorandom permutations against a two-sided
oracle. The switching lemma prices the substitution of one ideal
object for the other; AES supplies the deployed permutation; keyed
BLAKE2 supplies the PRFs this section inventories; and
format-preserving encryption closes the section by manufacturing, from
a PRF that is not invertible, the keyed bijection on an 11-byte format
that Orchard's address derivation demands.

Notation is that of the provable-security model. Strings of length
\(\ell\) form \(\{0,1\}^{\ell}\), all finite strings form
\(\{0,1\}^{*}\), and \(\abs{x}\) is the length of \(x\); for a finite
set \(S\), the notation \(x \xleftarrow{\$} S\) draws \(x\) uniformly
from \(S\), independently of everything else. Adversaries are PPT
machines in the security parameter \(\lambda\)
(\S\ref{model:subsec:adversaries}), a negligible function is
\(\negl(\lambda)\) (recalled in \S\ref{model:subsec:indist}), and
every security notion below is a distinguishing game in the sense of
Definition~\ref{model:def:game}, scored by the difference
\[
\mathrm{Adv}(\mathcal{A})
\;=\;
\bigl|\Pr[\mathcal{A}\text{ outputs }1\text{ in }\mathsf{Exp}_{1}]
     -\Pr[\mathcal{A}\text{ outputs }1\text{ in }\mathsf{Exp}_{0}]\bigr|
\]
between the two experiments the game distinguishes. One device remains
to be formalised: the black box itself.

\begin{definition}[Oracle access]\label{prf:def:oracle}
For a function \(f\), the notation \(\mathcal{A}^{f}\) denotes an
adversary granted \emph{oracle access} to \(f\): a black box to which
\(\mathcal{A}\) may submit inputs \(x\) of its choice, receiving
\(f(x)\) in a single computation step, and from which it learns
nothing about \(f\) beyond the answers to the queries it makes. The
adversary's running time always bounds the number of its oracle
queries, which is hence polynomial in \(\lambda\).
\end{definition}

\subsection{Pseudorandom functions and permutations}
\label{prf:subsec:prf}

Pseudorandomness is measured against an ideal object, which must come
first: the ``truly random function'' has to be made precise before
anything can be said to imitate it.

\begin{definition}[Random function]\label{prf:def:randfun}
Let \(\mathsf{Func}(D, R)\) denote the set of all functions from a
finite set \(D\) to a finite set \(R\). A \emph{random function} from
\(D\) to \(R\) is an element \(F \xleftarrow{\$} \mathsf{Func}(D, R)\),
chosen uniformly at random. Equivalently, the values
\(\{F(x)\}_{x \in D}\) are independent and uniformly distributed over
\(R\).
\end{definition}

There are \(\abs{R}^{\abs{D}}\) functions in \(\mathsf{Func}(D, R)\),
an astronomically large set for the domains of interest: for
\(D = R = \{0,1\}^{128}\) the count is \(2^{128 \cdot 2^{128}}\).
No such function can be stored or transmitted---its
truth table is the only description it has, and the table does not fit
in the universe. The essential property of a pseudorandom function is
that it replaces this object by a short key while preserving its
observable behaviour. The clean mental model of the ideal object is
\emph{lazy sampling}, exactly as for the random oracle of
Definition~\ref{hash:def:ro}: the oracle keeps a table, initially
empty; on a fresh query \(x\) it samples
\(F(x) \xleftarrow{\$} R\), records the pair, and returns it; on a
repeated query it returns the recorded value. This produces exactly
the distribution of Definition~\ref{prf:def:randfun} while only ever
materialising the points actually queried.

\begin{definition}[Keyed function]\label{prf:def:keyed}
A \emph{keyed function} is a deterministic, efficiently computable map
\[
F \colon \{0,1\}^{\lambda} \times \{0,1\}^{\ell(\lambda)}
  \longrightarrow \{0,1\}^{m(\lambda)},
\qquad
(k, x) \longmapsto F_{k}(x),
\]
where the input length \(\ell\) and output length \(m\) are
polynomially bounded functions of the key length \(\lambda\). Fixing
the key \(k\) yields the function
\(F_{k} \colon \{0,1\}^{\ell(\lambda)} \to \{0,1\}^{m(\lambda)}\).
\end{definition}

\begin{definition}[Pseudorandom function]\label{prf:def:prf}
Let \(F\) be a keyed function. The distinguishing game
\(\mathsf{PRF}^{\mathcal{A}}_{F}(\lambda)\) runs as follows.
\begin{enumerate}
\item The challenger samples a bit \(b \xleftarrow{\$} \{0,1\}\). If
  \(b = 1\) it samples \(k \xleftarrow{\$} \{0,1\}^{\lambda}\) and
  sets \(\mathcal{O} := F_{k}\); if \(b = 0\) it sets
  \(\mathcal{O} := R\) for
  \(R \xleftarrow{\$} \mathsf{Func}(\ell(\lambda), m(\lambda))\),
  realised by lazy sampling, where
  \(\mathsf{Func}(\ell, m) :=
  \mathsf{Func}(\{0,1\}^{\ell}, \{0,1\}^{m})\).
\item The adversary \(\mathcal{A}(1^{\lambda})\), receiving the
  security parameter in unary to convey the parameter size, queries
  \(\mathcal{O}\) adaptively.
\item The adversary outputs a bit \(b'\); the game outputs \(1\) iff
  \(b' = b\).
\end{enumerate}
The family \(F\) is a \emph{pseudorandom function} (PRF) if for every
PPT adversary \(\mathcal{A}\),
\[
\mathrm{Adv}^{\mathrm{prf}}_{F}(\mathcal{A}, \lambda)
\;:=\;
\Bigl|
\Pr_{k \xleftarrow{\$} \{0,1\}^{\lambda}}
  \bigl[\mathcal{A}^{F_{k}}(1^{\lambda}) = 1\bigr]
-
\Pr_{R \xleftarrow{\$} \mathsf{Func}(\ell(\lambda), m(\lambda))}
  \bigl[\mathcal{A}^{R}(1^{\lambda}) = 1\bigr]
\Bigr|
\;\in\; \negl(\lambda).
\]
The two probabilities are the adversary's output distributions in the
worlds \(b = 1\) and \(b = 0\), so this difference equals twice the
bit-guessing advantage
\(\bigl|\Pr[\mathsf{PRF}^{\mathcal{A}}_{F}(\lambda) = 1] -
\tfrac{1}{2}\bigr|\) of Definition~\ref{model:def:game}; the factor
of two is immaterial to negligibility, and the difference form is the
one every bound below is stated in.
\end{definition}

Stated informally: an adversary handed a black box that is either
\(F_{k}\) for a secret random key or a freshly sampled truly random
function, and allowed to query it adaptively, cannot tell which it
holds. Two features of the definition carry the weight. First, the
algorithm defining \(F\) is public; only the sampled key is hidden.
The adversary obtains oracle access to the resulting keyed instance
\(F_{k}\) through Definition~\ref{prf:def:oracle}---never the key,
only its input--output behaviour. Second, the comparison object is the
gigantic, unstorable random function; the PRF compresses it into
\(\lambda\) bits while remaining indistinguishable to every efficient
observer.

That last qualifier is not removable: pseudorandomness is an
inherently computational notion. Information theory cannot afford the
trade. As \(k\) ranges over its \(2^{\lambda}\) values, the keyed
function \(F_{k}\) takes at most \(2^{\lambda}\) distinct values in
\(\mathsf{Func}(D, R)\)---a vanishing fraction of all
\(\abs{R}^{\abs{D}}\) functions---so the two distributions of
Definition~\ref{prf:def:prf} are statistically almost as far apart as
distributions can be, and an \emph{unbounded} adversary could in
principle distinguish them. Security holds only against
computationally bounded distinguishers, in precisely the manner of
the PRG separation in the proof of
Proposition~\ref{model:prop:hierarchy}.

\begin{remark}[Necessity of adaptivity]\label{prf:rem:adaptive}
The adversary chooses its queries; it need not commit to them in
advance, and later queries may depend on earlier answers. This
\emph{adaptive} access is the strongest reasonable model and the one
downstream constructions rely on. A weaker non-adaptive definition, in
which the adversary fixes all queries before seeing any answer, is
genuinely weaker---families exist that pass every non-adaptive test
yet fall to an adaptive distinguisher---and is insufficient for many
applications.
\end{remark}

A block cipher is not merely a keyed function but a keyed
\emph{permutation}: each key must yield an invertible map, or
decryption would be impossible. The relevant ideal object is therefore
a random permutation.

\begin{definition}[Keyed permutation]\label{prf:def:keyedperm}
A keyed function
\(F \colon \{0,1\}^{\lambda} \times \{0,1\}^{\ell} \to
\{0,1\}^{\ell}\) with equal input and output length is a \emph{keyed
permutation} if for every key \(k\) the map \(F_{k}\) is a bijection
of \(\{0,1\}^{\ell}\), and both \(F_{k}\) and \(F_{k}^{-1}\) are
efficiently computable given \(k\). The set of all permutations of
\(\{0,1\}^{\ell}\) is \(\mathsf{Perm}(\ell)\), of cardinality
\((2^{\ell})!\).
\end{definition}

\begin{definition}[Pseudorandom permutation]\label{prf:def:prp}
Let \(F\) be a keyed permutation. The game
\(\mathsf{PRP}^{\mathcal{A}}_{F}(\lambda)\) is the game of
Definition~\ref{prf:def:prf} with the ideal world replaced: for
\(b = 0\) the challenger sets \(\mathcal{O} := P\) for
\(P \xleftarrow{\$} \mathsf{Perm}(\ell)\), a uniformly random
permutation. The family \(F\) is a \emph{pseudorandom permutation}
(PRP) if for every PPT adversary \(\mathcal{A}\),
\[
\mathrm{Adv}^{\mathrm{prp}}_{F}(\mathcal{A}, \lambda)
\;:=\;
\Bigl|
\Pr_{k}\bigl[\mathcal{A}^{F_{k}}(1^{\lambda}) = 1\bigr]
-
\Pr_{P \xleftarrow{\$} \mathsf{Perm}(\ell)}
  \bigl[\mathcal{A}^{P}(1^{\lambda}) = 1\bigr]
\Bigr|
\;\in\; \negl(\lambda).
\]
\end{definition}

\begin{definition}[Strong pseudorandom permutation]\label{prf:def:sprp}
A keyed permutation \(F\) is a \emph{strong pseudorandom permutation}
(SPRP) if no PPT adversary distinguishes even when granted oracle
access to \emph{both} the permutation and its inverse: for every PPT
\(\mathcal{A}\),
\[
\mathrm{Adv}^{\mathrm{sprp}}_{F}(\mathcal{A}, \lambda)
\;:=\;
\Bigl|
\Pr_{k}\bigl[\mathcal{A}^{F_{k},\, F_{k}^{-1}}(1^{\lambda}) = 1\bigr]
-
\Pr_{P \xleftarrow{\$} \mathsf{Perm}(\ell)}
  \bigl[\mathcal{A}^{P,\, P^{-1}}(1^{\lambda}) = 1\bigr]
\Bigr|
\;\in\; \negl(\lambda).
\]
\end{definition}

The distinction between a PRP and an SPRP is exactly the availability
of a decryption oracle. A block cipher deployed where the attacker can
obtain decryptions of chosen ciphertexts must be an SPRP; this is the
natural security target for a block cipher, and the one the next
subsection adopts.

From the outside, the two ideal objects are nearly the same thing.
Query a length-preserving oracle on \(q\) distinct inputs. If the
oracle is a random function \(R \in \mathsf{Func}(\ell, \ell)\), the
answers are \(q\) independent uniform strings, and output collisions
occur. If the oracle is a random permutation
\(P \in \mathsf{Perm}(\ell)\), the answers are \(q\) \emph{distinct}
uniform strings, and no collision can occur. This is the \emph{only}
statistical difference visible to any observer, and it is the entire
content of the switching lemma.

\begin{lemma}[PRP/PRF switching lemma]\label{prf:lem:switching}
Let \(\ell \ge 1\) and let \(\mathcal{A}\) be any adversary---even a
computationally unbounded one---making at most \(q\) queries to its
oracle. Then
\[
\Bigl|
\Pr_{P \xleftarrow{\$} \mathsf{Perm}(\ell)}
  \bigl[\mathcal{A}^{P} = 1\bigr]
-
\Pr_{R \xleftarrow{\$} \mathsf{Func}(\ell, \ell)}
  \bigl[\mathcal{A}^{R} = 1\bigr]
\Bigr|
\;\le\; \frac{q(q-1)}{2^{\ell+1}}.
\]
\end{lemma}

We use the lemma as a statement only. Its proof is a standard
game-hopping argument (Remark~\ref{model:rem:game-hopping}) bounding
the probability of the single distinguishing event just
identified---an output collision among the \(q\) answers---and the
bound is the birthday-type count \(\binom{q}{2}/2^{\ell}\) familiar
from Lemma~\ref{hash:lem:birthday}, whose underlying calculation is
the \emph{Math Guide}'s (\S``The union bound and a birthday
calculation'').

The consequence is a licence: a secure PRP on a large domain is also a
secure PRF whenever the number of queries stays well below the
birthday bound \(2^{\ell/2}\), because the adversary's total advantage
against the function world exceeds its advantage against the
permutation world by at most \(q(q-1)/2^{\ell+1}\). For AES, with
\(\ell = 128\), the switching term \(q^{2}/2^{129}\) is negligible for
any realistic query budget. This is the licence by which the FF1
construction of \S\ref{prf:subsec:fpe} treats AES---a
permutation---as the PRF its round function requires.

\subsection{Block ciphers and AES}\label{prf:subsec:blockcipher}

\begin{definition}[Block cipher]\label{prf:def:blockcipher}
A \emph{block cipher} with block length \(b\) and key length
\(\kappa\) is a keyed permutation
\[
E \colon \{0,1\}^{\kappa} \times \{0,1\}^{b}
  \longrightarrow \{0,1\}^{b},
\]
together with its inverse \(D = E^{-1}\) satisfying
\(D_{k}(E_{k}(x)) = x\) for all \(k, x\). Here \(E_{k}\) is
\emph{encryption} and \(D_{k}\) \emph{decryption} under key \(k\). The
security goal is that \(E\) be a strong pseudorandom permutation
(Definition~\ref{prf:def:sprp}).
\end{definition}

A block cipher is a concrete, fast, fixed-parameter instantiation of
the SPRP abstraction, and the fixing of parameters changes the
language of its security claims. The asymptotic PRP definition speaks
of a family indexed by \(\lambda\); a real block cipher fixes \(b\)
and \(\kappa\) once and for all---\(b = 128\) and
\(\kappa \in \{128, 192, 256\}\) for AES---and offers \emph{concrete}
security in the resource-explicit style of
Definition~\ref{model:def:concrete}: an explicit function of the
adversary's resources bounds the advantage of the best known attacks.

\begin{definition}[Concrete \((t, q, \varepsilon)\)-security]
\label{prf:def:concrete}
A block cipher \(E\) is a \emph{\((t, q, \varepsilon)\)-secure} PRP if
every adversary running in time at most \(t\) and making at most \(q\)
oracle queries has
\(\mathrm{Adv}^{\mathrm{prp}}_{E}(\mathcal{A}) \le \varepsilon\);
analogously for SPRP and PRF security.
\end{definition}

\begin{construction}[AES]\label{prf:con:aes}
The Rijndael cipher of Daemen and Rijmen, standardised as AES in
FIPS~197, has block length \(b = 128\) bits and key lengths
\(\kappa \in \{128, 192, 256\}\), with \(N_{r} \in \{10, 12, 14\}\)
rounds respectively. Structurally it is a
\emph{substitution--permutation network} (SPN): the \(128\)-bit
state, viewed as a \(4 \times 4\) array of bytes---elements of
\(\F_{2^{8}}\)---is first whitened by an initial round-key XOR, then
passes through \(N_{r}\) rounds built from four layers:
\begin{enumerate}
\item \textbf{SubBytes}, a nonlinear byte substitution---the
  \emph{S-box}---built on field inversion \(a \mapsto a^{-1}\) in
  \(\F_{2^{8}}\) (composed with a fixed affine map), chosen so that
  every input difference maps to any given output difference with
  only small probability, and every linear approximation of it holds
  with probability close to one half;
\item \textbf{ShiftRows}, a transposition of bytes within the array;
\item \textbf{MixColumns}, an invertible column-wise linear mix; and
\item \textbf{AddRoundKey}, the XOR of a \(128\)-bit round key that a
  key schedule expands from \(k\).
\end{enumerate}
The first \(N_{r} - 1\) rounds compose all four layers in this order;
the final round omits MixColumns. The two permutation layers together
guarantee that after two rounds every output byte depends on every
input byte. Each layer---the initial whitening included---is
individually invertible, so \(E_{k}\) is a permutation as
Definition~\ref{prf:def:blockcipher} requires, and decryption applies
the inverses in reverse order.
\end{construction}

\begin{remark}[Confusion, diffusion, and the status of AES security]
\label{prf:rem:confusiondiffusion}
The SPN design is the modern embodiment of Shannon's two principles.
The nonlinear S-box supplies \emph{confusion}: making the relationship
between key and ciphertext as complex as possible. The ShiftRows and
MixColumns layers supply \emph{diffusion}: spreading the influence of each
input bit over the whole output. Iterating the two for ten or more
rounds is what gives AES its empirical resistance to all known
attacks. No proof that AES is a PRP exists; like every practical block
cipher, its security is a \emph{conjecture}, corroborated by decades
of failed cryptanalysis but resting on no reduction to a cleaner
problem. Every theorem downstream therefore takes ``\(E\) is a
PRP/SPRP'' as an assumption and reasons from it---the same epistemic
footing on which Remark~\ref{model:rem:rom-caveat} placed the concrete
hash functions.
\end{remark}

\subsection{PRFs from hash functions: length extension and keyed
BLAKE2}\label{prf:subsec:hmac}

Block ciphers are one source of PRFs; cryptographic hash functions are
another, often more convenient one, being unkeyed and ubiquitous. Here
the oracle distinguisher does not merely distinguish---she forges. The
naive keying of a hash \(H\) by prepending the key,
\[
F_{k}(m) \;=\; H(k \,\|\, m),
\]
is unsafe for the iterated Merkle--Damg\aa rd hashes (MD5, SHA-1,
SHA-256)---which pad the message into blocks, fold the blocks one at
a time into a fixed-width chaining state through a compression
function, and emit the final state---because of the
\emph{length-extension attack}: from
\(H(k \,\|\, m)\) and \(\abs{m}\) alone, the adversary computes
\(H(k \,\|\, m \,\|\, \mathrm{pad} \,\|\, m')\) for a suffix \(m'\) of
her choosing, \emph{without knowing \(k\)}. The reason is structural:
a Merkle--Damg\aa rd digest is exactly the internal chaining state
after absorbing \(k \,\|\, m\), so the published output is a resumable
snapshot of the iteration, and the attacker simply continues absorbing
blocks from where the honest computation stopped. Used as a message
authentication code (MAC)---a keyed function whose output tag certifies a
message, defined under ``Message authentication codes'' in the next
section---the naive construction thus surrenders a valid tag on
\(m \,\|\, \mathrm{pad} \,\|\, m'\), a message the key holder never
authenticated.

The classical repair is HMAC, the nested double call
\[
\mathrm{HMAC}_{k}(m)
\;=\;
H\bigl((k \oplus \mathrm{opad}) \,\|\,
       H((k \oplus \mathrm{ipad}) \,\|\, m)\bigr)
\]
for two fixed pad constants: the inner hash absorbs the message under
one key-derived value, and the outer hash re-keys the fixed-length
result, so the emitted digest is no longer a resumable chaining state.
Bellare's analysis shows that HMAC is a PRF whenever the underlying
compression function is, up to a birthday term \(q^{2}/2^{\ell}\) in
the chaining-state length \(\ell\); as a PRF it is in particular a
secure message authentication code, and it is the building block of
the standardised key-derivation function HKDF, noted in passing under
``Key-derivation functions'' in the next section. This summary is all
we need, because the shielded protocol uses HMAC nowhere: its hash of
choice admits direct keying.

Modern hash functions of the sponge family
(Construction~\ref{hash:con:sponge}) and of the HAIFA family---the
Merkle--Damg\aa rd variant whose compression function also takes the
count of bits hashed so far, the lineage of BLAKE2---do not suffer
length extension and so admit a far simpler keying. The BLAKE2 design
is built to be keyed \emph{natively}: the construction loads the key
length into its parameter block and prepends the key, padded to a
full block, as the first message block, and its compression function
mixes in a block counter and finalisation flags. (Its successor
BLAKE3 is natively keyed too, by a different mechanism: the key words
replace the initial state, and a dedicated flag marks the keyed
mode.) One simply writes
\[
F_{k}(m) \;=\; \mathrm{BLAKE2}\bigl(m;\ \mathrm{key} = k\bigr)
\]
and obtains a PRF directly---no nested double call. No
length-extension attack remains to repair: the finalisation flag,
XORed into the state for the last block, means the emitted digest is
\emph{not} the raw chaining value and cannot be used to resume the
computation.

\begin{remark}[Zcash's use of keyed and personalised BLAKE2]
\label{prf:rem:blake2}
The Zcash protocol uses BLAKE2b and BLAKE2s pervasively, exploiting
two parameters of the BLAKE2 design: the \emph{key}, for PRF and MAC
use as above, and the \emph{personalisation} string, a
domain-separation tag baked into the initial state---\(16\) bytes for
BLAKE2b, \(8\) bytes for BLAKE2s. Personalisation makes the same input
hashed for two purposes---nullifier derivation versus note-commitment
derivation, say---yield independent-looking outputs, so each use site
is effectively an independent PRF: the practical realisation of the
independent-oracle-per-context idiom that
Proposition~\ref{hash:prop:domainsep} proved for tagged random
oracles, obtained here before a single message byte is processed and
cheaper and cleaner than HMAC, because BLAKE2 provides keying and
domain separation as first-class features. The deployed inventory
spans both widths. On the BLAKE2b side, the key-expansion PRF
\(\mathsf{PRF}^{\mathsf{expand}}\) is BLAKE2b-512 personalised with
\texttt{Zcash\_\allowbreak ExpandSeed} (protocol specification
\href{https://zips.z.cash/protocol/protocol.pdf\#concreteprfs}{\S``Pseudo Random
Functions''}), evaluated on the spending key to derive Orchard's key
tree; the
note-encryption KDF is personalised with
\texttt{Zcash\_\allowbreak OrchardKDF}
(\href{https://zips.z.cash/protocol/protocol.pdf\#concreteorchardkdf}{\S``Orchard Key Derivation''})
and the outgoing
cipher key with \texttt{Zcash\_\allowbreak Orchardock}
(\href{https://zips.z.cash/protocol/protocol.pdf\#concreteprfs}{\S``Pseudo Random Functions''})---all \(16\)-byte tags. On the BLAKE2s
side, the nullifier PRF and fixed-base generators of Sapling---the
shielded protocol preceding Orchard---use \(8\)-byte personalisations
such as \texttt{Zcash\_\allowbreak nf} and
\texttt{Zcash\_\allowbreak PH}
(\href{https://zips.z.cash/protocol/protocol.pdf\#concreteprfs}{\S``Pseudo Random Functions''} and
\href{https://zips.z.cash/protocol/protocol.pdf\#concretepedersenhash}{\S``Pedersen Hash Function''}).
\end{remark}

\subsection{Format-preserving encryption and FF1}\label{prf:subsec:fpe}

All constructions so far operate on bit strings of a length the
cipher's block size dictates. The deployed need that closes this
section does not. Orchard derives the \emph{diversifiers} behind its
shielded addresses---one \(11\)-byte value per address, many
addresses per spending key---by encrypting a sequential address index
with
a keyed bijection on the set of \(11\)-byte strings, so that the
published diversifiers of one wallet look like unrelated random values
rather than the consecutive integers they conceal. An adversary who
could invert the map, or merely distinguish it from a random
permutation of the format, would link a wallet's addresses to one
another. What is needed is precisely a block cipher whose ``block'' is
an exotic finite set; the classical commercial need---encrypting a
\(16\)-digit card number into a \(16\)-digit card number so that
legacy databases, field-length constraints, and checksum formats
continue to accept it---is the same problem. This is the problem
\emph{format-preserving encryption} (FPE) solves.

\begin{definition}[Format-preserving encryption]\label{prf:def:fpe}
Let \(\mathcal{X}\) be a finite set, the \emph{format}---for example
\(\mathcal{X} = \{0, 1, \dots, 9\}^{16}\), the \(16\)-digit strings. A
\emph{format-preserving encryption scheme} on \(\mathcal{X}\) is a
keyed permutation
\(\mathcal{E} \colon \{0,1\}^{\kappa} \times \mathcal{X} \to
\mathcal{X}\) such that for every key \(k\) the map
\(\mathcal{E}_{k} \colon \mathcal{X} \to \mathcal{X}\) is a bijection
with efficiently computable inverse \(\mathcal{D}_{k}\). Security
requires \(\mathcal{E}\) to be a (strong) pseudorandom permutation on
\(\mathcal{X}\)---Definition~\ref{prf:def:prp} with
\(\mathsf{Perm}(\mathcal{X})\) in place of
\(\mathsf{Perm}(\{0,1\}^{\ell})\): no efficient adversary
distinguishes \(\mathcal{E}_{k}\) from a uniformly random permutation
of \(\mathcal{X}\), even with access to the inverse.
\end{definition}

The defining feature is that the domain \(\mathcal{X}\) is
\emph{arbitrary}---in particular \(\abs{\mathcal{X}}\) need not be a
power of two, so a plain block cipher does not apply. What is required
is a keyed pseudorandom \emph{bijection} on a set whose size is, say,
\(10^{16}\). Determinism is forced, and it separates FPE from the
randomised encryption modes of the next section: the format
\(\mathcal{X}\) has no room to store a per-message random value, so
FPE is the deterministic, format-preserving analogue of a block cipher
on the alphabet of \(\mathcal{X}\). It necessarily leaks equality of
plaintexts, as any deterministic cipher must; this is the unavoidable
cost of preserving format with no expansion, and it is acceptable for
the tokenisation use cases FPE targets---and immaterial for
diversifier derivation, where each index is encrypted once.

Formats reduce to integers. Any format of the shape ``strings of
length \(\nu\) over an alphabet of radix \(\rho\)'' is in bijection
with \(\{0, 1, \dots, \rho^{\nu} - 1\}\) by reading the string as a
base-\(\rho\) numeral, so it suffices to build a keyed pseudorandom
permutation of \(\Z/M\Z\) for an arbitrary modulus
\(M = \rho^{\nu}\). The FF1 construction below splits
\(M = a \cdot b\) into two factors of nearly equal size and runs a
\emph{Feistel network} whose two halves live in \(\Z/a\Z\) and
\(\Z/b\Z\). The Feistel network is the structural device that turns a
round function---not necessarily invertible---into a permutation; we
present the classical bitwise case first, then the modular
generalisation FF1 uses.

\begin{definition}[Balanced Feistel network]\label{prf:def:feistel}
Let \(F_{1}, \dots, F_{r} \colon \{0,1\}^{w} \to \{0,1\}^{w}\) be
arbitrary \emph{round functions}. The \(r\)-round Feistel network
\(\Psi[F_{1}, \dots, F_{r}]\) is the map on \(\{0,1\}^{2w}\) defined
as follows: write the input as \((L_{0}, R_{0})\) with
\(L_{0}, R_{0} \in \{0,1\}^{w}\), iterate
\[
L_{i} = R_{i-1},
\qquad
R_{i} = L_{i-1} \oplus F_{i}(R_{i-1}),
\qquad i = 1, \dots, r,
\]
and output \((L_{r}, R_{r})\).
\end{definition}

% TikZ figure placeholder --- fig:feistel: one balanced Feistel round
% and its inversion. Left panel: the forward round --- inputs L_{i-1}
% (left rail) and R_{i-1} (right rail); R_{i-1} branches into the
% round-function box F_i, whose output meets L_{i-1} at an XOR node;
% the rails cross so that L_i = R_{i-1} and R_i = L_{i-1} xor
% F_i(R_{i-1}). Right panel: the same wiring read bottom-to-top,
% recovering R_{i-1} = L_i and L_{i-1} = R_i xor F_i(L_i); annotate
% that F_i is evaluated FORWARD in both directions --- read backwards
% the diagram is the decryption direction. Dedicated figures pass
% supplies the drawing.

\begin{proposition}[Feistel is a bijection regardless of the round
functions]\label{prf:prop:feistelbijection}
For any round functions \(F_{1}, \dots, F_{r}\)---injective or
not---the map \(\Psi[F_{1}, \dots, F_{r}]\) is a bijection of
\(\{0,1\}^{2w}\), inverted by running the rounds backwards:
\[
R_{i-1} = L_{i},
\qquad
L_{i-1} = R_{i} \oplus F_{i}(L_{i}),
\qquad i = r, \dots, 1.
\]
\end{proposition}

\begin{proof}
It suffices to show each round is invertible, the composition of
bijections being a bijection. The \(i\)-th round sends
\((L_{i-1}, R_{i-1}) \mapsto
(R_{i-1},\, L_{i-1} \oplus F_{i}(R_{i-1}))\). Given the output
\((L_{i}, R_{i})\), one recovers \(R_{i-1} = L_{i}\) directly, and
then
\(L_{i-1} = R_{i} \oplus F_{i}(R_{i-1}) = R_{i} \oplus F_{i}(L_{i})\),
using that XOR is its own inverse. The recovery uses only
\emph{forward} evaluations of \(F_{i}\)---at no point is an inverse of
the round function required, which is the crucial point permitting a
permutation to be built out of a one-way-ish PRF. The displayed
backward recursion is exactly this inversion applied round by round.
\end{proof}

This proposition is the structural heart of format-preserving
encryption: it manufactures a \emph{keyed bijection} from a keyed
function that is not itself a bijection. In FF1 the round functions
derive from a PRF---AES iterated in cipher-block-chaining fashion, as
the construction below specifies---so the network is invertible
without ever inverting the PRF: AES is only ever evaluated in the
forward direction, in encryption and decryption alike.

Bijectivity is free; pseudorandomness is not. The Luby--Rackoff
theorem is the classical statement that enough Feistel rounds over a
PRF yield a pseudorandom permutation.

\begin{theorem}[Luby--Rackoff]\label{prf:thm:lubyrackoff}
Let the round functions \(F_{1}, \dots, F_{r}\) be independent PRFs on
\(\{0,1\}^{w}\). Against an adversary making \(q\) queries, the
Feistel network \(\Psi\) is:
\begin{enumerate}
\item with \(r = 3\) rounds, a pseudorandom permutation (forward
  queries only); and
\item with \(r = 4\) rounds, a \emph{strong} pseudorandom permutation
  (forward and inverse queries);
\end{enumerate}
in each case with distinguishing advantage \(O(q^{2}/2^{w})\) above
the PRF advantage of the round functions.
\end{theorem}

We use the theorem as a statement, without proof, but the shape of the
argument is worth recording, because it explains both the birthday
term and the round counts. One first replaces each PRF round function
by a truly random function, at the cost of \(r\) PRF advantage terms.
For three rounds one then shows that, unless two of the at most \(q\)
queries collide in an intermediate Feistel value---a birthday event of
probability \(O(q^{2}/2^{w})\)---the middle round's random function is
evaluated at distinct inputs, so it produces independent uniform
values that perfectly mask the relationship between input and output
halves, and the network's answers are distributed exactly as a random
\emph{function}'s---independent and uniform on \(\{0,1\}^{2w}\). A
final application of the switching lemma
(Lemma~\ref{prf:lem:switching}), at a further cost of
\(q(q-1)/2^{2w+1}\), carries the comparison from the random function
to the random permutation; the term is absorbed by the same
\(O(q^{2}/2^{w})\) bound. The fourth round adds the symmetry needed to resist
inverse queries as well, upgrading PRP to SPRP. The full argument is a
transcript analysis bounding the probability of any internal
collision; the leftover term is the stated birthday bound.

\begin{remark}[Why FF1 uses ten rounds rather than four]
\label{prf:rem:rounds}
Theorem~\ref{prf:thm:lubyrackoff} gives security only up to
\(q \approx 2^{w/2}\) queries, and in the FPE setting \(w\) is
tiny---one half of a \(16\)-digit number---so the birthday bound
\(q^{2}/2^{w}\) is weak. The standardised FF1 therefore prescribes
\emph{ten} rounds rather than the minimal four, as a security margin
against the more powerful attacks known for small-domain Feistel
networks, which can sometimes beat the generic birthday bound. The
classification of this design claim is worth making explicit: the
round count itself is \emph{specified}---NIST SP~800-38G fixes ten
rounds---while a security proof matching the margin is an \emph{open
problem}; the analysis that exists (Luby--Rackoff and its small-domain
refinements) does not certify the full strength the margin is hoped to
buy. What the round count does \emph{not} affect is invertibility,
which Proposition~\ref{prf:prop:feistelbijection} guarantees for any
round count and any round function.
\end{remark}

\begin{construction}[The FF1 mode of NIST SP~800-38G]
\label{prf:con:ff1}
The FF1 mode instantiates an \emph{unbalanced, modular, tweaked} Feistel
network of ten rounds to build a PRP on \(\Z/M\Z\) with
\(M = \rho^{\nu}\), the radix-\(\rho\) numeral reading of a
length-\(\nu\) string. It supports a \emph{tweak} \(T\): a non-secret
string that selects, in effect, one permutation from a family indexed
by \(T\), so that the same key encrypts the same plaintext differently
under different tweaks---useful for binding a token to its context.
Its operation:
\begin{enumerate}
\item \textbf{Split.} Set \(u = \floor{\nu/2}\) and \(v = \nu - u\).
  Split the plaintext numeral string \(X\) into a left part \(A\) of
  \(u\) symbols and a right part \(B\) of \(v\) symbols, with integer
  values in \(\{0, \dots, \rho^{u} - 1\}\) and
  \(\{0, \dots, \rho^{v} - 1\}\) respectively. The two halves live in
  the moduli \(a = \rho^{u}\) and \(b = \rho^{v}\); the network is
  unbalanced when \(\nu\) is odd.
\item \textbf{Round function from a PRF.} For each round
  \(i = 0, 1, \dots, 9\), assemble a byte string \(Q\) encoding the
  tweak \(T\), the round index \(i\), and the current right half
  \(B\); prepend a fixed parameter block \(P\) encoding the radix, the
  lengths, and the tweak length; and apply
  \(\mathrm{PRF}(P \,\|\, Q)\), where \(\mathrm{PRF}\) is AES in
  CBC-MAC mode: set \(Y_{0} = 0^{128}\), iterate
  \(Y_{j} = \mathrm{AES}_{k}(Y_{j-1} \oplus X_{j})\) over the
  \(16\)-byte blocks \(X_{j}\) of \(P \,\|\, Q\), and output the final
  block \(Y_{m}\) (SP~800-38G, Algorithm~6). The iteration chains AES
  through an evolving state exactly as the Merkle--Damg\aa rd
  iteration of \S\ref{prf:subsec:hmac} chains a compression function;
  emitting the final state is safe \emph{here} because the PRF's
  inputs form a \emph{prefix-free} set: the leading block \(P\)
  encodes the lengths of everything that follows, so two distinct
  inputs either differ already in their first block or agree on all
  lengths, and in neither case is one a proper prefix of the other.
  Prefix-free inputs are precisely the setting in which plain CBC-MAC
  is proved to be a PRF (a result of Petrank and Rackoff), and no
  length-extension query can arise.
  Further AES calls then expand the result to a long enough
  digest, which is read as an integer \(y\).
\item \textbf{Modular Feistel step.} Update the halves by the modular
  analogue of Definition~\ref{prf:def:feistel}:
  \[
  C \;=\; (A + y) \bmod m,
  \qquad
  A \leftarrow B,
  \qquad
  B \leftarrow C,
  \]
  where the modulus \(m\) alternates between \(a = \rho^{u}\) and
  \(b = \rho^{v}\) with the parity of the round, so that the sum lands
  in the correct half's domain; modular addition replaces XOR.
\item \textbf{Output.} After ten rounds, concatenate the final \(A\)
  and \(B\) and render the result back as a length-\(\nu\)
  radix-\(\rho\) string.
\end{enumerate}
Decryption runs the rounds in reverse, replacing modular addition by
modular subtraction, \(A \leftarrow (C - y) \bmod m\), exactly
mirroring the backward recursion of
Proposition~\ref{prf:prop:feistelbijection}. One quantitative detail
of step~2 deserves note: the standard sizes the expanded digest at
least four bytes beyond what the larger half requires, so the integer
\(y\) ranges over at least \(2^{32}\) multiples of the modulus, and
the bias of the reduction \(y \bmod m\) away from uniform is bounded
by the modular-reduction bias of the \emph{Math Guide} (\S``Uniform
sampling and the bias of modular reduction'') at below \(2^{-33}\)
per round.
\end{construction}

\begin{proposition}[FF1 is a keyed bijection on \(\mathcal{X}\)]
\label{prf:prop:ff1bijection}
For every key \(k\) and tweak \(T\), the FF1 map
\(\mathcal{E}_{k,T} \colon \mathcal{X} \to \mathcal{X}\) is a
bijection of the format set \(\mathcal{X}\), and
\(\mathcal{D}_{k,T}\) is its two-sided inverse.
\end{proposition}

\begin{proof}
Each round is the map \((A, B) \mapsto (B, (A + y) \bmod m)\), where
\(y = y(i, T, B)\) depends only on the round index, the tweak, and the
right half \(B\)---\emph{not} on \(A\). Fix \(B\), \(i\), and \(T\);
then \(y\) is a constant and \(A \mapsto (A + y) \bmod m\) is a
translation of \(\Z/m\Z\), hence a bijection, inverted by subtracting
the same constant. The full round is therefore invertible: from the
output \((B, C)\) recover \(B\) directly, recompute \(y = y(i, T, B)\)
by a \emph{forward} PRF evaluation, and set
\(A = (C - y) \bmod m\). Each even-indexed round is a bijection of
\(\Z/a\Z \times \Z/b\Z\) onto \(\Z/b\Z \times \Z/a\Z\) and each
odd-indexed round a bijection back, the two products coinciding when
\(\nu\) is even; the composition of the ten rounds, an even number, is
a bijection of \(\Z/a\Z \times \Z/b\Z\). The numeral map
\((A, B) \mapsto A \cdot b + B\) identifies
\(\Z/a\Z \times \Z/b\Z\) with \(\Z/M\Z\) as a bijection \emph{of sets
only}, not of groups: the factors satisfy
\(\gcd(a, b) = \rho^{\min(u, v)} > 1\), so the Chinese remainder
theorem does not apply, and only bijectivity is used. Transporting
through this map and the numeral bijection
\(\mathcal{X} \cong \Z/M\Z\), the composition \(\mathcal{E}_{k,T}\)
is a bijection of \(\mathcal{X}\). Running the rounds backward with
subtraction inverts it on both sides, so
\(\mathcal{D}_{k,T} = \mathcal{E}_{k,T}^{-1}\).
\end{proof}

\begin{remark}[Bijectivity from algebra, security from the PRF]
\label{prf:rem:fpe-separation}
The construction separates its two obligations cleanly, and the
separation is why Feistel is the standard route to FPE. Invertibility
comes for free from the structure:
Propositions~\ref{prf:prop:feistelbijection}
and~\ref{prf:prop:ff1bijection} hold for \emph{any} round function
whatsoever, including a broken or constant one. The PRF---AES---is
responsible not for invertibility but for \emph{pseudorandomness}:
making \(\mathcal{E}_{k,T}\) look like a uniformly random permutation
of \(\mathcal{X}\) in the sense of Definition~\ref{prf:def:fpe}, per
the Luby--Rackoff \emph{heuristic}---Theorem~\ref{prf:thm:lubyrackoff}
adapted, beyond what it literally proves, to the modular, multi-round,
tweaked setting.
\end{remark}

% Register deviation (y4-13): the FF1 parameters stay here, cited to the
% spec section, because Ironwood's ``Diversified addresses'' omits them.
\begin{remark}[The domain-size caveat, and the single Orchard use]
\label{prf:rem:ff1-deployed}
Two facts complete the picture, one cautionary and one deployed.

The domain-size caveat is real. The published SP~800-38G requires
\(\rho^{\nu} \ge 100\) and recommends \(\rho^{\nu} \ge 10^{6}\)---its
draft Revision~1 would make \(10^{6}\) mandatory---because for very
small domains an attacker can recover or distinguish the keyed
permutation by exhaustively cataloguing input--output pairs, and known
message-recovery attacks on small-domain FPE (Bellare--Hoang--Tessaro
and successors) erode the security margin. For the moderately large
domains FPE is meant for, and with its ten rounds, FF1 is the
NIST-standardised, AES-backed realisation of a keyed pseudorandom
bijection on an arbitrary finite format.

The single use of FF1 in Orchard is the diversifier derivation with
which this subsection opened: FF1-AES256 maps an address index to the
\(11\)-byte diversifier embedded in a shielded address, so that one
spending key yields many mutually unlinkable addresses. The deployed
instance (protocol specification
\href{https://zips.z.cash/protocol/protocol.pdf\#concreteprps}{\S``Pseudo Random
Permutations''}) fixes radix \(\rho = 2\), length \(\nu = 88\), and the
empty tweak, and reads the \(11\)-byte index as a little-endian bit
string, so the domain is the \(2^{88}\) strings of \(88\)
bits---comfortably beyond the caveat above. The address built on the
diversifier, and the case for a keyed permutation rather than a hash
in this role, are the business of the \emph{Ironwood Guide}
(\S``Diversified addresses'').
\end{remark}

The toolbox has gained its keyed compartment: one distinguishing game
at three strengths, a standard permutation conjectured to win it, a
natively keyed hash that turns domain-separated digests into
independent PRFs, and a Feistel construction that bends a PRF into a
bijection on any finite format. The next section spends these tools on
the channel itself: symmetric encryption, authentication, and the
authenticated-encryption composition that carries every shielded note,
each built by running a PRF in a mode that converts its
indistinguishability into secrecy or integrity.

\section{Symmetric encryption, AEAD, and key derivation}\label{sec:symmetric}

The adversary of this section owns the channel. Every ciphertext
passes through her hands, and she attacks on two fronts at once. As an
eavesdropper she tries to \emph{distinguish}: to tell an encryption of
one note plaintext from an encryption of another, or to recognise
that two ciphertexts hide the same message. As a forger she tries to
\emph{modify}: a stream cipher's ciphertext is XOR-malleable, so
flipping a ciphertext bit flips the corresponding plaintext bit
without any key material at all; she can also shift bytes between a
packet's public header and its encrypted body, replay a ciphertext
under a different context, or manufacture a wholly new ciphertext that
the receiver accepts. She is patient enough to exploit a single reused
pad---two ciphertexts under one pad hand her the XOR of the
plaintexts---and mathematically literate enough to exploit a key that
is merely high-entropy rather than uniform. This section arms the
honest parties against all of it, one guarantee at a time:
confidentiality from the one-time pad through the nonce-based stream
cipher ChaCha20 and the chosen-plaintext game; integrity from message
authentication codes through universal hashing and Carter--Wegman;
their composition as authenticated encryption with associated data,
culminating in the ChaCha20-Poly1305 scheme that Zcash's
note-encryption layer deploys; and finally key derivation, the
manufacture of the uniform key that every theorem in between takes as
its hypothesis.

Throughout, an \emph{adversary} is a probabilistic algorithm
\(\mathcal{A}\); the notation \(y \gets \mathcal{A}(x)\) denotes
running \(\mathcal{A}\) on input \(x\) and assigning its output to
\(y\), and \(x \xleftarrow{\$} S\) denotes sampling \(x\) uniformly
from the finite set \(S\). Every adversary in the computational
definitions is probabilistic polynomial time (PPT) in the security
parameter \(\lambda\), and \(\negl(\lambda)\) denotes a negligible
function (\emph{Math Guide}, \S``Polynomial, exponential, and
negligible functions''; recalled in \S\ref{model:subsec:indist}).

\paragraph{The one-time pad and EAV security.}
The information-theoretic baseline is the \emph{one-time pad}: to
encrypt \(m \in \{0,1\}^{L}\) under a uniformly random key
\(k \in \{0,1\}^{L}\) used exactly once, output \(c = m \oplus k\).
Since \(m \oplus k\) is uniform whatever \(m\) is, the ciphertext
distribution is independent of the message, and even an unbounded
eavesdropper learns nothing---Shannon's \emph{perfect secrecy}. The
price is severe on both counts the name advertises: the key must be
as long as the message, and it must never be reused, for two
ciphertexts under one pad reveal
\(c \oplus c' = m \oplus m'\)---the \emph{two-time-pad} failure, which
the adversary of the opening paragraph collects for free. The
computational relaxation is \emph{EAV security} (security against an
eavesdropper): an adversary who chooses two equal-length messages
\(m_{0}, m_{1}\) and sees a single ciphertext
\(\mathrm{Enc}_{k}(m_{b})\) for a hidden uniform bit \(b\) guesses
\(b\) with advantage at most \(\negl(\lambda)\)---the one-ciphertext,
no-oracle restriction of the chosen-plaintext game formalised in
\S\ref{symmetric:subsec:cpa}. Replacing the truly random pad by the
output of a pseudorandom generator stretched from a short seed---the
\emph{pseudo-one-time pad}---achieves EAV security with a
fixed-length key, at the cost of resting on a pseudorandomness
assumption.

\subsection{Stream ciphers and ChaCha20}\label{symmetric:subsec:stream}

A \emph{stream cipher} is the practical instantiation of the
pseudo-one-time pad: a keyed, seekable generator producing a
\emph{keystream} that XORs into the plaintext. To encrypt many
messages under one key without violating the one-time-pad no-reuse
rule, modern stream ciphers take an additional public input called a
\emph{nonce} (number used once) and behave like a pseudorandom
function of the nonce, in the sense of Section~\ref{sec:prf}.

\begin{definition}[Nonce-based stream cipher]\label{symmetric:def-stream}
A \emph{nonce-based stream cipher} is a deterministic function
\[
\mathrm{KS} : \mathcal{K} \times \mathcal{N} \times \N
\to \{0,1\}^{\ast}
\]
producing, for key \(k\), nonce \(\nu\), and length \(L\), a keystream
\(\mathrm{KS}(k, \nu, L) \in \{0,1\}^{L}\). Encryption of
\(m \in \{0,1\}^{L}\) is
\(\mathrm{Enc}_{k}(\nu, m) = m \oplus \mathrm{KS}(k, \nu, |m|)\) with
ciphertext \((\nu, c)\); decryption recomputes the keystream and XORs
it off. The security goal is that for distinct nonces the keystreams
are jointly indistinguishable from independent uniform strings,
provided no nonce repeats under a fixed key.
\end{definition}

The stream cipher Zcash deploys---inside the AEAD of
\S\ref{symmetric:subsec:chachapoly}, in the note-encryption layer---is
\textsf{ChaCha20}, designed by Bernstein and standardised in
RFC~8439. Its core is a reversible mixing of four \(32\)-bit words
built only from modular addition, XOR, and rotation (an ``ARX''
design), which avoids data-dependent table lookups and hence the
timing side channels that table-based ciphers must engineer around.

\begin{construction}[ChaCha20]\label{symmetric:con-chacha}
The \textsf{ChaCha20} state is a \(4 \times 4\) matrix of \(32\)-bit
words, viewed as a vector
\((x_{0}, \dots, x_{15}) \in (\Z/2^{32}\Z)^{16}\); the symbol \(+\)
denotes addition modulo \(2^{32}\), \(\oplus\) XOR, and
\(x \lll r\) left rotation by \(r\) bits. The \emph{quarter-round}
\(\mathrm{QR}(a, b, c, d)\) updates four words by
\[
  \begin{aligned}
    a &\mathrel{+}= b; & d &\mathrel{\oplus}= a; & d &\lll= 16; \\
    c &\mathrel{+}= d; & b &\mathrel{\oplus}= c; & b &\lll= 12; \\
    a &\mathrel{+}= b; & d &\mathrel{\oplus}= a; & d &\lll= 8;  \\
    c &\mathrel{+}= d; & b &\mathrel{\oplus}= c; & b &\lll= 7.
  \end{aligned}
\]
The state is initialised from the \(128\)-bit constant
``\texttt{expand 32-byte k}'' (words \(x_{0}\)--\(x_{3}\)), a
\(256\)-bit key (\(x_{4}\)--\(x_{11}\)), a \(32\)-bit block counter
(\(x_{12}\)), and a \(96\)-bit nonce (\(x_{13}\)--\(x_{15}\)). One
\emph{double round} applies \(\mathrm{QR}\) to the four columns and
then to the four diagonals; the \(20\) rounds of \textsf{ChaCha20}
are ten double rounds. Writing \(X\) for the initial state and
\(R(X)\) for the state after the rounds, the \(64\)-byte keystream
block is the \emph{feed-forward} sum \(R(X) + X\) (word-wise, modulo
\(2^{32}\)), serialised in little-endian order.
\end{construction}

\begin{remark}[Feed-forward and counter mode]
\label{symmetric:rem-chacha-design}
The rounds \(R\) form a bijection---each quarter-round is invertible,
being a composition of invertible word operations---so without the
final addition the block function would be an invertible permutation,
and an adversary holding a keystream block could run the rounds
backwards. The feed-forward sum \(R(X) + X\) destroys invertibility;
it is the standard Davies--Meyer-style move for turning a permutation
into a one-way mixing function. Incrementing the block counter
\(x_{12}\) produces successive \(64\)-byte keystream blocks, so
\(\mathrm{KS}(k, \nu, L)\) is the concatenation of the blocks for
counters \(0, 1, 2, \dots\); this \emph{counter mode} is what makes
the cipher seekable and parallelisable. The security status is a
conjecture, in the same heuristic position as every concrete cipher
of Section~\ref{sec:prf}: \textsf{ChaCha20} is \emph{conjectured} to
be a secure PRF of the pair (nonce, counter), and no attack better
than brute force over the \(256\)-bit key is known. Reusing a
\((k, \nu)\) pair reproduces the same keystream and re-creates the
two-time-pad failure; nonce uniqueness is a hard requirement, not a
hygiene recommendation.
\end{remark}

\subsection{Chosen-plaintext security}\label{symmetric:subsec:cpa}

EAV security models an adversary who sees a single ciphertext. The
adversary of this section is stronger: she can influence which
messages get encrypted---by triggering payments, injecting plaintext
into a session, or simply guessing---and she observes many
ciphertexts. The corresponding notion is security against
\emph{chosen-plaintext attack} (CPA), modelled by granting her oracle
access to the encryption algorithm, in the game template of
Definition~\ref{model:def:game}.

\begin{definition}[CPA security]\label{symmetric:def-cpa}
For a symmetric encryption scheme \(\Pi\) with key generator
\(\mathrm{Gen}\), adversary \(\mathcal{A}\), and bit
\(b \in \{0,1\}\), the game
\(\mathsf{CPA}^{\mathcal{A}}_{\Pi}(\lambda, b)\) runs as follows.
\begin{enumerate}
\item The challenger samples \(k \gets \mathrm{Gen}(1^{\lambda})\)
  and gives \(\mathcal{A}\) oracle access to
  \(\mathrm{Enc}_{k}(\cdot)\).
\item At some point \(\mathcal{A}\) submits a challenge pair
  \((m_{0}, m_{1})\) of equal-length messages and receives
  \(c^{\star} \gets \mathrm{Enc}_{k}(m_{b})\).
\item The adversary may continue querying the oracle and finally
  outputs a bit \(b'\), which is the output of the game.
\end{enumerate}
The scheme \(\Pi\) is \emph{CPA-secure} if for every PPT
\(\mathcal{A}\),
\[
\mathrm{Adv}^{\mathrm{cpa}}_{\Pi}(\mathcal{A}, \lambda)
:= \bigl|
   \Pr\bigl[\mathsf{CPA}^{\mathcal{A}}_{\Pi}(\lambda, 0) = 1\bigr]
 - \Pr\bigl[\mathsf{CPA}^{\mathcal{A}}_{\Pi}(\lambda, 1) = 1\bigr]
   \bigr|
\;\in\; \negl(\lambda).
\]
\end{definition}

The first consequence of giving the adversary an encryption oracle is
a structural one: encryption cannot be a function of the message
alone.

\begin{proposition}[Determinism precludes CPA security]
\label{symmetric:prop-deterministic}
No scheme whose \(\mathrm{Enc}\) is deterministic and stateless is
CPA-secure.
\end{proposition}

\begin{proof}
The adversary queries the oracle on two distinct messages
\(m_{0} \ne m_{1}\), recording \(c_{0} = \mathrm{Enc}_{k}(m_{0})\)
and \(c_{1} = \mathrm{Enc}_{k}(m_{1})\). Decryption correctness
forces \(c_{0} \ne c_{1}\): a single ciphertext cannot decrypt to
both messages. She submits the challenge \((m_{0}, m_{1})\), receives
\(c^{\star} = \mathrm{Enc}_{k}(m_{b})\), and outputs \(0\) if
\(c^{\star} = c_{0}\) and \(1\) otherwise. Since \(\mathrm{Enc}\) is
deterministic, \(c^{\star} = c_{b}\), and since \(c_{0} \ne c_{1}\)
the comparison identifies \(b\): the guess is always correct and the
advantage is \(1\).
\end{proof}

Consequently a CPA-secure scheme must be randomised or nonce-based. A
nonce-based stream cipher used with a fresh nonce per
encryption---random, or drawn from a counter---attains CPA security
under the PRF conjecture for its keystream generator, because each
encryption is then effectively a fresh one-time pad. We record the
notion here and defer the routine PRF-based proof to the AEAD
analysis (Theorem~\ref{symmetric:thm-chachapoly}), where
confidentiality and integrity are treated together.

\begin{remark}[Semantic security]\label{symmetric:rem-semantic}
Goldwasser and Micali's original notion, \emph{semantic security},
demands that whatever an efficient adversary can compute about the
plaintext from the ciphertext, she could have computed without the
ciphertext, from the message length alone. Semantic security is
provably equivalent to the indistinguishability notions above; we use
indistinguishability throughout because it is far easier to
manipulate in reductions. The content is the same: the ciphertext is
computationally useless for learning anything about the message
beyond its length.
\end{remark}

\subsection{Message authentication codes}\label{symmetric:subsec:mac}

CPA security protects confidentiality and says nothing about
integrity, and for the XOR-based stream ciphers above the gap is not
hypothetical: flipping a bit of the ciphertext flips the
corresponding plaintext bit, so the channel-owning adversary can
convert a payment of one amount into a payment of another without
ever touching the key. Detecting tampering requires a second
primitive, keyed like the cipher but aimed at the forger rather than
the eavesdropper: the \emph{message authentication code}. Here we
build this symmetric, shared-key form of authentication, the form
that the AEAD needs.

\begin{definition}[Message authentication code]\label{symmetric:def-mac}
A \emph{message authentication code} (MAC) is a pair
\((\mathrm{Mac}, \mathrm{Vrfy})\) with key space \(\mathcal{K}\):
algorithm \(\mathrm{Mac}_{k}(m)\) outputs a tag \(t\), and
\(\mathrm{Vrfy}_{k}(m, t) \in \{0,1\}\) accepts (\(1\)) or rejects
(\(0\)). Correctness requires
\(\mathrm{Vrfy}_{k}(m, \mathrm{Mac}_{k}(m)) = 1\) for all \(k, m\).
When \(\mathrm{Mac}\) is deterministic, canonical verification
recomputes the tag and checks equality.
\end{definition}

\begin{definition}[Existential and strong unforgeability]
\label{symmetric:def-eufcma}
For a MAC \(\Pi\) and adversary \(\mathcal{A}\), the game
\(\mathsf{Forge}^{\mathcal{A}}_{\Pi}(\lambda)\) runs as follows.
\begin{enumerate}
\item The challenger samples \(k \xleftarrow{\$} \mathcal{K}\)
  and gives \(\mathcal{A}\) oracle access to
  \(\mathrm{Mac}_{k}(\cdot)\); let \(Q\) be the set of messages
  queried.
\item The adversary outputs a pair \((m^{\star}, t^{\star})\); the
  game outputs \(1\) iff
  \(\mathrm{Vrfy}_{k}(m^{\star}, t^{\star}) = 1\) and
  \(m^{\star} \notin Q\).
\end{enumerate}
The MAC is \emph{existentially unforgeable under chosen-message
attack} (EUF-CMA) if for every PPT \(\mathcal{A}\),
\[
\mathrm{Adv}^{\mathrm{forge}}_{\Pi}(\mathcal{A}, \lambda)
:= \Pr\bigl[\mathsf{Forge}^{\mathcal{A}}_{\Pi}(\lambda) = 1\bigr]
\;\in\; \negl(\lambda).
\]
A \emph{strongly} unforgeable (sUF-CMA) MAC additionally forbids
producing a new valid tag on an already-queried message: the
adversary wins on \((m^{\star}, t^{\star})\) as long as the oracle
never returned that exact \emph{pair}. We write
\(\mathrm{Adv}^{\mathrm{s\text{-}forge}}_{\Pi}(\mathcal{A}, \lambda)\)
for the advantage in this variant.
\end{definition}

The MAC Zcash uses, \textsf{Poly1305}, is not a
pseudorandom-function MAC but a \emph{one-time},
\emph{universal-hash} MAC: it provides information-theoretic security
per nonce and is extremely fast, at the cost of a key that
authenticates only a single message. The enabling abstraction is a
hash family whose \emph{differences} are unpredictable.

\begin{definition}[\(\varepsilon\)-almost-\(\Delta\)-universal hash
family]\label{symmetric:def-axu}
Let \((\mathcal{G}, +)\) be a finite abelian group. A family of
functions \(\{h_{r} : \mathcal{X} \to \mathcal{G}\}_{r \in
\mathcal{R}}\) is \emph{\(\varepsilon\)-almost-\(\Delta\)-universal}
(\(\varepsilon\)-A\(\Delta\)U) if for all distinct
\(x, x' \in \mathcal{X}\) and every \(\delta \in \mathcal{G}\),
\[
\Pr_{r \xleftarrow{\$} \mathcal{R}}
  \bigl[\, h_{r}(x) - h_{r}(x') = \delta \,\bigr]
\;\le\; \varepsilon.
\]
Taking \(\delta = 0\) recovers ordinary
\(\varepsilon\)-universality---low collision probability---while the
\(\Delta\) condition controls \emph{additive differences}, which is
what makes the one-time MAC below secure. With
\(\mathcal{G} = \{0,1\}^{n}\) under XOR this is the classical
almost-XOR-universality; \textsf{Poly1305} takes
\(\mathcal{G} = \Z/2^{128}\Z\) under addition modulo \(2^{128}\).
\end{definition}

The Carter--Wegman paradigm turns any A\(\Delta\)U family into a
one-time MAC by masking the hash value with a fresh uniform pad.

\begin{theorem}[Carter--Wegman one-time MAC]\label{symmetric:thm-cw}
Let \(\{h_{r}\}\) be \(\varepsilon\)-A\(\Delta\)U into
\((\mathcal{G}, +)\). Define a MAC with key \((r, s)\), where
\(r \xleftarrow{\$} \mathcal{R}\) and the pad
\(s \xleftarrow{\$} \mathcal{G}\), by
\(\mathrm{Mac}_{r,s}(m) = h_{r}(m) + s\). If \((r, s)\)
authenticates at most one message, then any adversary---even a
computationally unbounded one---who sees one message--tag pair
\((m, t)\) forges a valid pair \((m^{\star}, t^{\star})\) with
\(m^{\star} \ne m\) with probability at most \(\varepsilon\).
\end{theorem}

\begin{proof}
The adversary observes \(t = h_{r}(m) + s\) and must output
\((m^{\star}, t^{\star})\) with \(m^{\star} \ne m\) and
\(t^{\star} = h_{r}(m^{\star}) + s\). Subtracting the two tag
equations eliminates the pad: the forgery succeeds iff
\[
h_{r}(m^{\star}) - h_{r}(m) \;=\; t^{\star} - t.
\]
Condition on the adversary's view. The pad \(s\) is uniform and
independent of \(r\), so given the single observed tag \(t\), every
value of \(h_{r}(m)\) is equally consistent---each candidate value
determines a unique \(s\)---and the observation therefore leaks
nothing about \(r\): the posterior on \(r\) remains uniform on
\(\mathcal{R}\). The adversary's difference
\(\delta := t^{\star} - t\) and distinct pair \((m, m^{\star})\) are
thus fixed independently of \(r\), and the A\(\Delta\)U property
gives
\(\Pr[\, h_{r}(m^{\star}) - h_{r}(m) = \delta \,] \le \varepsilon\).
\end{proof}

The MAC \textsf{Poly1305} instantiates the paradigm with a
polynomial-evaluation hash over the prime field \(\F_{p}\) for
\(p = 2^{130} - 5\) (prime-field arithmetic as in the \emph{Math
Guide}, \S``The prime field and two inversion routes'').

\begin{construction}[Poly1305]\label{symmetric:def-poly1305}
Fix the prime \(p = 2^{130} - 5\) and work in \(\F_{p}\). The key is
a pair \((r, s)\): \(r\) is a \(128\)-bit value subjected to a fixed
bit-\emph{clamping} (the top four bits of bytes \(3, 7, 11, 15\) and
the bottom two bits of bytes \(4, 8, 12\) are forced to zero), and
\(s\) is a \(128\)-bit pad. To authenticate a message, split it into
\(\ell\) blocks \(c_{1}, \dots, c_{\ell}\) of at most \(16\) bytes;
encode each block as an integer in \([0, 2^{128})\) and add
\(2^{8b}\), where \(b\) is the block's byte length---this sets a
\(1\) bit just above the block's bytes, making every coefficient
nonzero and the block encoding injective in both content and byte
length. Interpreting
\(r\) and each \(c_{i}\) as elements of \(\F_{p}\), the hash is the
polynomial evaluation
\[
h_{r}(m) \;=\; \sum_{i=1}^{\ell} c_{i}\, r^{\,\ell - i + 1}
\;\in\; \F_{p},
\]
computed by Horner's rule as
\(\bigl(\cdots((c_{1} r + c_{2}) r + c_{3}) r \cdots\bigr) r\). The
tag is
\[
\mathrm{Mac}_{r,s}(m)
\;=\; \bigl( (h_{r}(m) \bmod p) + s \bigr) \bmod 2^{128},
\]
the field hash reduced and added to the pad modulo \(2^{128}\).
\end{construction}

\begin{proposition}[Poly1305 is almost-\(\Delta\)-universal]
\label{symmetric:prop-poly1305-axu}
Restricted to messages of at most \(L\) bytes---hence at most
\(\ell_{\max} = \ceil{L/16}\) blocks---the reduced \textsf{Poly1305}
hash family \(\{m \mapsto h_{r}(m) \bmod p\}\), read into
\(\Z/2^{128}\Z\), is \(\varepsilon\)-almost-\(\Delta\)-universal with
respect to addition modulo \(2^{128}\), with
\[
\varepsilon \;\le\; \frac{8\, \ell_{\max}}{2^{106}}.
\]
The loss over the ideal \(\ell_{\max} \cdot p^{-1}\) accounts for the
clamping of \(r\)---which restricts \(r\) to a subset of size
\(2^{106}\)---and for the final reduction modulo \(2^{128}\).
\end{proposition}

\begin{proof}
\emph{Sketch.} Fix two distinct messages \(m, m'\) with block
encodings \((c_{i})\) and \((c'_{j})\), padded to a common length
\(\ell \le \ell_{\max}\) by leading zero coefficients without
changing their values. The block encoding guarantees that
\(m \ne m'\) yields distinct coefficient vectors: genuine
coefficients are nonzero, so differing block counts set a nonzero
coefficient against a zero padding one, while equal block counts
leave two blocks differing in content or in final-block byte
length, hence in encoding. The difference of
hashes is
\[
h_{r}(m) - h_{r}(m')
\;=\; \sum_{i=1}^{\ell} (c_{i} - c'_{i})\, r^{\,\ell - i + 1}
\;=:\; P(r) \;\in\; \F_{p},
\]
a nonzero polynomial in \(r\) of degree at most \(\ell\). Now count
the ways a target difference \(\delta \in \Z/2^{128}\Z\) can be hit.
The reduced hash values lie in \([0, p)\), so their integer
difference lies in the interval \((-p,\, p)\), which has length
\(2p < 8 \cdot 2^{128}\); each residue class modulo \(2^{128}\)
therefore contains at most \(\ceil{2p/2^{128}} = 8\) candidate
integer differences \(u\). For each such \(u\), the equation
\(P(r) = u\) in \(\F_{p}\) has at most \(\ell\) roots---a polynomial
of degree \(\le \ell\) over a field has at most \(\ell\) roots
(\emph{Math Guide}, \S``Roots and the factor theorem'')---so at most
\(8\ell\) values of \(r\) produce the difference \(\delta\). The
clamped \(r\) is uniform over a set of \(2^{106}\) values (clamping
zeroes \(22\) of the \(128\) bits), so it lands on one of those
roots with probability at most \(8\ell/2^{106} \le
8\ell_{\max}/2^{106}\).
\end{proof}

The bound is comfortably small in practice: for messages of
\(16\,384\) bytes (\(\ell_{\max} = 1024\) blocks) it gives
\(\varepsilon \le 8 \cdot 1024/2^{106} = 2^{-93}\), against the ideal
\(\ell_{\max}/p \approx 2^{-120}\)---the clamping and reduction
losses cost \(27\) bits of the margin and leave it enormous.
Combining Proposition~\ref{symmetric:prop-poly1305-axu} with the
Carter--Wegman Theorem~\ref{symmetric:thm-cw}: \textsf{Poly1305} with
a fresh, uniformly random pad \(s\) per message---uniformity, not
mere unpredictability, is what the theorem's pad-elimination step
uses---is a one-time MAC with forgery probability at most
\(8\ell_{\max}/2^{106}\) per verification attempt. In the AEAD below,
\textsf{ChaCha20} itself generates the one-time key \((r, s)\)
pseudorandomly from the nonce---uniform once the cipher is idealised
as a random function---so the per-nonce one-time guarantee is
exactly what is needed.

\subsection{Authenticated encryption with associated data}
\label{symmetric:subsec:aead}

We now combine the two guarantees. The composite goal is
\emph{authenticated encryption}: the adversary should be unable to
learn anything about plaintexts (CPA security) \emph{and} unable to
produce any new ciphertext that decrypts to anything other than the
distinguished rejection symbol \(\bot\), a value distinct from every
message (ciphertext integrity). Real protocols additionally require
ciphertexts to bind to public context---packet headers, version
bytes, an ephemeral public key---that travels in the clear but must
be authenticated; this is the \emph{associated data}.

\begin{definition}[AEAD scheme]\label{symmetric:def-aead}
An \emph{authenticated encryption with associated data} (AEAD)
scheme is a pair \((\mathrm{Enc}, \mathrm{Dec})\) with key space
\(\mathcal{K}\), nonce space \(\mathcal{N}\), and associated-data
space \(\mathcal{D}\): algorithm \(\mathrm{Enc}_{k}(\nu, a, m)\)
returns a ciphertext \(c\), and \(\mathrm{Dec}_{k}(\nu, a, c)\)
returns a message or \(\bot\). Correctness:
\(\mathrm{Dec}_{k}(\nu, a, \mathrm{Enc}_{k}(\nu, a, m)) = m\) for all
\(k, \nu, a, m\). The scheme authenticates but does not encrypt the
associated data \(a\).
\end{definition}

\begin{definition}[Ciphertext integrity, INT-CTXT]
\label{symmetric:def-intctxt}
For an AEAD scheme \(\Pi\) and adversary \(\mathcal{A}\), the game
\(\mathsf{Ctxt}^{\mathcal{A}}_{\Pi}(\lambda)\) runs as follows.
\begin{enumerate}
\item The challenger samples \(k \xleftarrow{\$} \mathcal{K}\)
  and gives \(\mathcal{A}\) two oracles: an encryption oracle
  \(\mathrm{Enc}_{k}(\cdot, \cdot, \cdot)\), recording the set \(Q\)
  of returned triples \((\nu, a, c)\), and a verification oracle
  that on \((\nu, a, c)\) returns \(1\) if
  \(\mathrm{Dec}_{k}(\nu, a, c) \ne \bot\) and \(0\) otherwise. The
  adversary must never repeat a nonce to the \emph{encryption}
  oracle (the \emph{nonce-respecting} model); verification queries
  may reuse nonces freely.
\item The game outputs \(1\) iff some verification query
  \((\nu^{\star}, a^{\star}, c^{\star}) \notin Q\) is answered
  \(1\).
\end{enumerate}
The scheme \(\Pi\) has \emph{ciphertext integrity} if for every PPT
\(\mathcal{A}\),
\(\Pr\bigl[\mathsf{Ctxt}^{\mathcal{A}}_{\Pi}(\lambda) = 1\bigr]
\in \negl(\lambda)\).
\end{definition}

\begin{definition}[Authenticated encryption]\label{symmetric:def-ae}
An AEAD scheme is \emph{secure} (is \emph{an AEAD}) if it is both
CPA-secure and has ciphertext integrity
(Definition~\ref{symmetric:def-intctxt}), where the CPA game of
Definition~\ref{symmetric:def-cpa} is run with
\(k \xleftarrow{\$} \mathcal{K}\) and extended so that the adversary
supplies the nonce and associated data of every oracle query and of
the challenge, subject to the \emph{nonce-respecting}
restriction---a restriction on the adversary, not a property of the
scheme---that no nonce is used twice across oracle queries and
challenge. The restriction is necessary: an AEAD may encrypt
deterministically given \((\nu, a, m)\)---ChaCha20-Poly1305
does---and then querying the oracle at \((\nu, a, m_{0})\) and
submitting the challenge \((\nu, a, m_{0}, m_{1})\) would identify
the bit by comparing ciphertexts.
\end{definition}

The standard route from a CPA-secure cipher and a strongly
unforgeable MAC to an AEAD is \emph{encrypt-then-MAC}: encrypt the
message, compute a MAC over the ciphertext together with the
associated data and nonce, and append the tag; decryption verifies
the tag \emph{before} decrypting, and any tag failure returns
\(\bot\). The contrasting orders---MAC-then-encrypt and
encrypt-and-MAC---do \emph{not} generically achieve authenticated
encryption; encrypt-then-MAC does.

\begin{theorem}[Encrypt-then-MAC yields authenticated encryption]
\label{symmetric:thm-etm}
Let \(\Pi_{E} = (\mathrm{Enc}, \mathrm{Dec})\) be a CPA-secure
encryption scheme and \(\Pi_{M} = (\mathrm{Mac}, \mathrm{Vrfy})\) a
strongly unforgeable (sUF-CMA) MAC, with \emph{independent} keys
\(k_{E}, k_{M}\). Define the composite scheme \(\Pi'\) by
\[
\mathrm{Enc}'_{k_{E}, k_{M}}(\nu, a, m) = (c, t),
\qquad
c \gets \mathrm{Enc}_{k_{E}}(\nu, m),
\quad
t \gets \mathrm{Mac}_{k_{M}}(\nu \,\|\, a \,\|\, c),
\]
with \(\mathrm{Dec}'\) returning \(\bot\) unless
\(\mathrm{Vrfy}_{k_{M}}(\nu \|a\| c,\, t) = 1\), in which case it
returns \(\mathrm{Dec}_{k_{E}}(\nu, c)\). Then \(\Pi'\) is a secure
AEAD: it is CPA-secure and has ciphertext integrity.
\end{theorem}

\begin{proof}
\emph{Ciphertext integrity.} Suppose a PPT \(\mathcal{A}\) making at
most \(q_{v}\) verification queries wins the game: some fresh query
\((\nu^{\star}, a^{\star}, (c^{\star}, t^{\star}))
\notin Q\) decrypts successfully, i.e.
\(\mathrm{Vrfy}_{k_{M}}(\nu^{\star} \|a^{\star}\| c^{\star},
t^{\star}) = 1\). Each encryption-oracle call on \((\nu, a, m)\)
returned \((c, t)\) and corresponds to one MAC query on the string
\(w = \nu \|a\| c\) returning tag \(t\). Because the framing
\(\nu \|a\| c\) parses uniquely (the component lengths are fixed or
prepended), the freshness of \((\nu^{\star}, a^{\star}, c^{\star})\)
means the MAC oracle never produced the string--tag pair
\((w^{\star}, t^{\star})\) with
\(w^{\star} := \nu^{\star} \|a^{\star}\| c^{\star}\): either
\(w^{\star}\) was never queried, or it was queried and returned a
different tag. Either way \((w^{\star}, t^{\star})\) is a fresh valid
pair---a \emph{strong} forgery against \(\Pi_{M}\). The reduction
\(\mathcal{B}\) runs \(\mathcal{A}\), answering encryption queries
with its own MAC oracle and a self-chosen \(k_{E}\). Verification
queries \(\mathcal{B}\) cannot answer, lacking \(k_{M}\), so it
guesses: it samples a uniform index \(i \le q_{v}\) in advance,
answers each fresh verification query before the \(i\)-th with
\(0\) (queries on triples in \(Q\) are answered \(1\), correctly),
halts at the \(i\)-th fresh query, and outputs that query's
\((w^{\star}, t^{\star})\) as its forgery. On the event that
\(\mathcal{A}\) wins and \(i\) hits her \emph{first} successful
fresh query---the guess is right with probability
\(1/q_{v}\)---every earlier \(0\) answer was correct, the
simulation is perfect, and \(\mathcal{B}\) wins, so
\[
\Pr\bigl[\mathsf{Ctxt}^{\mathcal{A}}_{\Pi'}(\lambda) = 1\bigr]
\;\le\; q_{v} \cdot
\mathrm{Adv}^{\mathrm{s\text{-}forge}}_{\Pi_{M}}(\mathcal{B},
\lambda)
\;\in\; \negl(\lambda),
\]
the factor \(q_{v}\) being polynomial.

\emph{CPA security.} The tag is computed from \(c\) (together with
\(k_{M}, \nu, a\), and the MAC's own coins) and so carries no
information about the challenge bit beyond what \(c\) already
carries. Formally, the
reduction \(\mathcal{B}'\) samples \(k_{M}\) itself, forwards
\(\mathcal{A}\)'s oracle and challenge queries to its own
\(\Pi_{E}\) oracle to obtain the \(c\)-components, and appends tags
it computes with \(k_{M}\). This simulates \(\Pi'\) perfectly, so
\(\mathrm{Adv}^{\mathrm{cpa}}_{\Pi'}(\mathcal{A}, \lambda) =
\mathrm{Adv}^{\mathrm{cpa}}_{\Pi_{E}}(\mathcal{B}', \lambda) \in
\negl(\lambda)\).
\end{proof}

\begin{remark}[Why the ordering matters]\label{symmetric:rem-etm}
Encrypt-then-MAC lets the receiver reject forged ciphertexts
\emph{without decrypting}, and this is what blocks chosen-ciphertext
and padding-oracle attacks: a rejected ciphertext never reaches the
decryption logic, so the decryption logic cannot leak. By contrast,
MAC-then-encrypt verifies only after decrypting, exposing the
decryption routine to attacker-chosen ciphertexts---the historical
source of a long line of padding-oracle breaks. Ciphertext integrity
is also exactly the property that upgrades CPA security to
chosen-ciphertext (CCA) security, the CPA game of
Definition~\ref{symmetric:def-cpa} with a decryption oracle added
that refuses only the challenge ciphertext: an
authenticated-encryption scheme is automatically CCA-secure, because
that oracle is useless to the adversary---every ciphertext she did
not legitimately obtain decrypts to \(\bot\) except with negligible
probability.
\end{remark}

\subsection{The ChaCha20-Poly1305 AEAD}
\label{symmetric:subsec:chachapoly}

We now assemble the concrete scheme of RFC~8439, the AEAD by which
Zcash encrypts every shielded note ciphertext and outgoing ciphertext
(protocol specification,
\href{https://zips.z.cash/protocol/protocol.pdf\#concretesym}{\S``Symmetric Encryption''}). It is an
encrypt-then-MAC composition of the \textsf{ChaCha20} stream cipher
with the \textsf{Poly1305} one-time MAC, with one economical
twist: the cipher derives \emph{both} the keystream and the one-time
MAC key from the single \((k, \nu)\) input.

% FIGURE (planned, dedicated figures pass): fig:aead ---
% ChaCha20-Poly1305 dataflow. Left: (key k, nonce nu) feeding the
% ChaCha20 block function; a column of keystream blocks for counters
% 0, 1, 2, ...; block 0 peeled off and split into (r, s) [clamp on r
% marked], blocks >= 1 XORed into the plaintext to give c. Right: the
% Poly1305 lane hashing pad16(a) || pad16(c) || len(a) || len(c)
% under r, then adding s to give the tag t. Output (c, t). Caption
% notes the counter-0 reservation as domain separation.

\begin{construction}[ChaCha20-Poly1305]\label{symmetric:def-chachapoly}
Fix a \(256\)-bit key \(k\) and a \(96\)-bit nonce \(\nu\).
Encryption of a message \(m\) with associated data \(a\) proceeds as
follows.
\begin{enumerate}
\item \emph{One-time MAC key.} Run \textsf{ChaCha20} with key \(k\),
  nonce \(\nu\), and block counter \(0\); take the first \(32\) bytes
  of the keystream block as the \textsf{Poly1305} one-time key
  \((r, s)\)---clamp the first \(16\) bytes to form \(r\); the next
  \(16\) are \(s\). Discard the rest of this block.
\item \emph{Encrypt.} Run \textsf{ChaCha20} with key \(k\) and nonce
  \(\nu\) starting at block counter \(1\), producing the keystream
  \(\mathrm{KS}_{1}\) (the keystream of
  Remark~\ref{symmetric:rem-chacha-design} taken from counter \(1\)
  onward), and set \(c = m \oplus \mathrm{KS}_{1}(k, \nu, |m|)\).
\item \emph{Authenticate.} Form the \textsf{Poly1305} input by
  concatenating: \(a\) padded with zeros to a \(16\)-byte boundary,
  \(c\) padded with zeros to a \(16\)-byte boundary, the \(8\)-byte
  little-endian length of \(a\), and the \(8\)-byte little-endian
  length of \(c\). Compute the tag
  \(t = \mathrm{Poly1305}_{r,s}(\text{that input})\).
\end{enumerate}
The ciphertext is \((c, t)\). Decryption recomputes \((r, s)\) and
the tag from \((k, \nu, a, c)\), compares it to \(t\) in
\emph{constant time}, returns \(\bot\) on mismatch, and otherwise
outputs \(m = c \oplus \mathrm{KS}_{1}(k, \nu, |c|)\).
\end{construction}

\begin{remark}[Length encoding prevents splicing]
\label{symmetric:rem-length-encoding}
The trailing \(8\)-byte lengths of \(a\) and \(c\) are essential.
Without them, the zero-padding that aligns \(a\) and \(c\) to
\(16\)-byte boundaries would let an adversary shift bytes between the
associated data and the ciphertext, or pad-extend one of them, while
leaving the \textsf{Poly1305} input string unchanged---breaking the
binding between \(a\) and \(c\) even though every authenticated byte
is intact. Appending the explicit lengths makes the MAC input an
injective encoding of the pair \((a, c)\), the same injectivity
that the per-block length encoding supplies within one message in
the A\(\Delta\)U analysis
(Proposition~\ref{symmetric:prop-poly1305-axu}).
\end{remark}

\begin{theorem}[Security of ChaCha20-Poly1305]
\label{symmetric:thm-chachapoly}
Model \textsf{ChaCha20} as a pseudorandom function \(F_{k}(\nu, j)\)
from (nonce, block counter) pairs to \(64\)-byte blocks. In the
nonce-respecting model, \textsf{ChaCha20-Poly1305} is a secure AEAD.
Concretely, consider a PPT adversary making \(q\) encryption queries
and \(q_{v}\) verification queries, each of whose \textsf{Poly1305}
inputs---padded associated data, padded ciphertext, and the length
block---comprises at most \(\ell_{\max}\) sixteen-byte blocks,
\(\ell_{\max} = \ceil{|a|/16} + \ceil{|c|/16} + 1\). Her CPA
advantage is at most twice the best PRF-distinguishing advantage
against \(F\) at comparable cost, and she forges---wins the
INT-CTXT game---with probability at most
\[
\mathrm{Adv}^{\mathrm{prf}}_{F}(\mathcal{B}, \lambda)
\;+\; \frac{8\, q_{v}\, \ell_{\max}}{2^{106}},
\]
where \(\mathcal{B}\) is a PRF distinguisher of comparable cost.
\end{theorem}

\begin{proof}
\emph{Sketch.} Replace \(F_{k}\) by a truly random function
\(\Phi\) from (nonce, counter) pairs to \(64\)-byte blocks. Any
noticeable change in the adversary's success probability yields a
PRF distinguisher of comparable cost---it simulates the whole game
using its function oracle---so each replacement costs one
\(\mathrm{Adv}^{\mathrm{prf}}_{F}\) term; the CPA game is played
twice (once per challenge bit), whence the factor two there. We
analyse the idealised scheme.

\emph{Confidentiality.} Since the adversary is nonce-respecting,
every encryption query uses a fresh \(\nu\), and the keystream blocks
\(\Phi(\nu, 1), \Phi(\nu, 2), \dots\) are uniform and independent of
everything else in her view; the ciphertext
\(c = m \oplus \mathrm{KS}\) is a genuine one-time pad and reveals
nothing about \(m\) beyond its length. The tag is a deterministic
function of \(c, a, \nu\) and the (independent) key \((r, s)\), and
adds no information. The idealised CPA advantage is therefore zero.

\emph{Integrity.} For each nonce \(\nu\), the block
\(\Phi(\nu, 0)[0{:}32]\) is uniform and independent across distinct
nonces; clamping its first half samples \(r\) from
\textsf{Poly1305}'s prescribed restricted set, while the second half
\(s\) stays uniform, so the resulting one-time keys are independent
across nonces. Because counter \(0\) is reserved for this key while
encryption consumes counters \(\ge 1\), each one-time key is also
independent of the keystream that produced its ciphertext. Per
nonce, the setting is
exactly the Carter--Wegman one-time MAC of
Theorem~\ref{symmetric:thm-cw} instantiated with the
\(\varepsilon\)-A\(\Delta\)U \textsf{Poly1305} hash of
Proposition~\ref{symmetric:prop-poly1305-axu},
\(\varepsilon \le 8\ell_{\max}/2^{106}\). Consider any single fresh
verification query \((\nu^{\star}, a^{\star}, c^{\star},
t^{\star})\). If \(\nu^{\star}\) was never used for encryption, the
pad \(s\) for \(\nu^{\star}\) is uniform and entirely unknown, so the
tag is correct with probability \(2^{-128}\). If \(\nu^{\star}\)
equals the nonce of some encryption query---nonce reuse is allowed in
\emph{verification} queries, only encryption is nonce-respecting---
then the adversary has seen exactly one tag under that \((r, s)\),
and freshness of \((a^{\star}, c^{\star})\) makes this a forgery on a
new message under a one-time key, succeeding with probability at most
\(\varepsilon\) by Theorem~\ref{symmetric:thm-cw}. The union bound
over the \(q_{v}\) verification queries (\emph{Math Guide}, \S``The
union bound and a birthday calculation'') gives forgery probability
at most \(8 q_{v} \ell_{\max}/2^{106}\) in the idealised scheme;
adding back the PRF-switching cost yields the stated bound.
\end{proof}

\begin{remark}[Single key, two roles]\label{symmetric:rem-single-key}
The scheme reuses one \(256\)-bit key for both encryption and
authentication, seemingly violating the independent-keys hypothesis
of Theorem~\ref{symmetric:thm-etm}. Reserving block counter \(0\) for
the \textsf{Poly1305} key and counters \(\ge 1\) for the keystream is
exactly what restores independence: once the cipher is idealised as
a truly random function---the switch the proof of
Theorem~\ref{symmetric:thm-chachapoly} pays for with a
PRF-distinguishing term---outputs at distinct counters are
independent, so \((r, s)\) is independent of the keystream and the
two-key analysis applies. For the PRF itself the independence is
computational, not literal. This
is a recurring design pattern---domain separation by a counter or
label, the same discipline as \S\ref{hash:subsec:domainsep}, lets one
master key safely play several roles.
\end{remark}

\begin{remark}[Wrong-key rejection and key commitment]
\label{symmetric:rem-wrong-key}
The upper volumes consume this AEAD in a \emph{trial-decryption}
pattern: a receiver detects the notes addressed to it by attempting
decryption of every ciphertext under each of its keys, relying on
decryption under a wrong key to return \(\bot\)
(protocol specification,
\href{https://zips.z.cash/protocol/protocol.pdf\#concretesym}{\S``Symmetric Encryption''}). For
\emph{honestly produced}
ciphertexts the scheme delivers this. Fix a ciphertext
\((c, t) = \mathrm{Enc}_{k}(\nu, a, m)\) and a key \(k'\)
independent of \(k\), and idealise the cipher under \(k\) and under
\(k'\) as independent random functions, as in the proof of
Theorem~\ref{symmetric:thm-chachapoly}. The one-time key
\((r', s')\) that decryption under \(k'\) derives from block
\((\nu, 0)\) is then uniform and independent of \((c, t)\): the pad
\(s'\) alone makes the recomputed tag match \(t\) with probability
exactly \(2^{-128}\), and even conditioned on one exposed tag under
\((r', s')\), Theorem~\ref{symmetric:thm-cw} bounds acceptance by
the A\(\Delta\)U slack \(8 \ell_{\max}/2^{106}\) of
Proposition~\ref{symmetric:prop-poly1305-axu}, with \(\ell_{\max}\)
as in Theorem~\ref{symmetric:thm-chachapoly}. Decryption under an
independent key therefore rejects except with probability at most
\(2^{-128} + 8 \ell_{\max}/2^{106}\) per attempt, plus two
PRF-distinguishing terms for the two idealisations.

The guarantee stops at honest senders, and the boundary deserves
care. CPA security and INT-CTXT are \emph{single-key} notions:
nothing in Definition~\ref{symmetric:def-ae} constrains a ciphertext
crafted by a \emph{malicious sender} who knows several keys. The
scheme is in fact not \emph{key-committing}---the separate notion
demanding that no ciphertext verify under two distinct keys---and a
sender who chooses keys \(k_{1}, k_{2}\) herself can craft a single
\((c, t)\) that decrypts validly, to different plaintexts, under
both: with \((r_{i}, s_{i})\) the one-time key each \(k_{i}\)
derives, the hash \(h_{r_{i}}\) of
Construction~\ref{symmetric:def-poly1305} is linear in the blocks of
\(c\), so fixing \(t\), leaving two blocks of \(c\) free, and
solving over \(\F_{p}\) the two linear equations that pin the hash
under \(r_{i}\) to \(t - s_{i}\) (retrying until both solutions
encode valid blocks) makes both keys accept \((c, t)\)---the
multi-key collision of Len, Grubbs, and Ristenpart's
partitioning-oracle attack. The upper volumes' trial-decryption
arguments therefore never rest on the tag alone against adversarial
senders: their
note-plaintext consistency checks---recomputing the note commitment
from the decrypted plaintext and comparing it against the
commitment the transaction carries---are what close this gap.
\end{remark}

\subsection{Key-derivation functions}\label{symmetric:subsec:kdf}

Theorem~\ref{symmetric:thm-chachapoly} assumes a uniformly random
\(256\)-bit key. In practice keys may come from sources that are
\emph{high-entropy but not uniform}: a Diffie--Hellman shared secret
is a structured group encoding, not a uniform bit string, and a
hardware noise source may spread its unpredictability unevenly
across its bits. (A passphrase is a different, usually
\emph{low}-entropy source, and needs a deliberately costly password
KDF rather than the machinery of this subsection.) The adversary is
owed nothing here---if the honest parties feed a non-uniform key
into a scheme proved secure for uniform keys, the proof simply does
not apply, and structure in the key is structure she may exploit. A
\emph{key-derivation function} (KDF) converts an appropriate source
secret into one or more pseudorandom keys. The right entropy measure
is worst-case, not average-case.

\begin{definition}[Min-entropy]\label{symmetric:def-minentropy}
The \emph{min-entropy} of a random variable \(X\) is
\[
H_{\infty}(X) \;=\; -\log_{2} \max_{x} \Pr[X = x];
\]
equivalently, the best chance of guessing \(X\) in one try is
\(2^{-H_{\infty}(X)}\). Given side information \(Z\), the
\emph{conditional min-entropy} is
\[
\tilde H_{\infty}(X \mid Z) \;=\;
-\log_{2} \E_{z \gets Z} \max_{x} \Pr[X = x \mid Z = z],
\]
so that \(2^{-\tilde H_{\infty}(X \mid Z)}\) is the best chance of
guessing \(X\) from \(Z\). Recall the \emph{statistical distance}
\(\Delta(X, Y) = \tfrac12 \sum_{u} |\Pr[X = u] - \Pr[Y = u]|\) of
Definition~\ref{model:def:statdist} (\emph{Math Guide},
\S``Statistical distance''); we say \(X\) is
\emph{\(\epsilon\)-close to uniform} on a set \(S\) if
\(\Delta(X, U_{S}) \le \epsilon\).
\end{definition}

\begin{definition}[Key-derivation function and its security]
\label{symmetric:def-kdf}
A \emph{key-derivation function} is a function
\[
\mathrm{KDF} : \mathcal{X} \times \mathcal{I} \times \{0,1\}^{\ast}
\times \N \to \{0,1\}^{\ast}
\]
taking a source secret (a sample of a random variable \(\Sigma\) over
\(\mathcal{X}\)), a non-secret \emph{salt}
\(\mathrm{slt} \in \mathcal{I}\), an \emph{info} or context string
\(\mathrm{ctx}\), and a desired output length \(L\). The security
goal must specify the source model. For an arbitrary source of
conditional min-entropy at least \(\gamma\) given the adversary's
side information, an extractor-style guarantee uses
a salt sampled uniformly and independently of the source and
requires the pair
\(\bigl(\mathrm{slt},\,
\mathrm{KDF}(\Sigma, \mathrm{slt}, \mathrm{ctx}, L)\bigr)\) to be
computationally indistinguishable from \((\mathrm{slt},\, U_{L})\),
even given \(\mathrm{ctx}\) and that side information. No fixed
deterministic unsalted function can provide this guarantee for
\emph{every} high-min-entropy source; an unsalted KDF therefore
rests on a narrower premise, such as input keying material that is
already computationally pseudorandom, or that is hard to query in a
random-oracle model.
\end{definition}

The notion a KDF targets is that of a seeded randomness extractor
whose output need only \emph{look} uniform.

\begin{definition}[Strong computational extractor]
\label{symmetric:def-extractor}
A keyed function
\(\mathrm{Ext} : \mathcal{I} \times \mathcal{X} \to \{0,1\}^{n}\) is
a \emph{strong \((\gamma, \epsilon)\)-computational extractor} if for
every source \(\Sigma\) over \(\mathcal{X}\) of conditional
min-entropy at least \(\gamma\) given auxiliary information \(Z\),
and a uniformly random salt
\(\mathrm{slt} \xleftarrow{\$} \mathcal{I}\), the pair
\(\bigl(\mathrm{slt},\, \mathrm{Ext}_{\mathrm{slt}}(\Sigma)\bigr)\)
is computationally indistinguishable from
\(\bigl(\mathrm{slt},\, U_{n}\bigr)\), even given \(Z\), with
distinguishing advantage at most \(\epsilon\). ``Strong'' means the
salt is revealed to the distinguisher.
\end{definition}

Information theory can realise the notion unconditionally; we state
the classical result \emph{without proof}, since the deployed
mechanism below rests on random-oracle modelling instead.

\begin{theorem}[Leftover hash lemma; stated without proof]
\label{symmetric:thm-lhl}
Let \(\{g_{r} : \mathcal{X} \to \{0,1\}^{n}\}\) be a
\(2^{-n}\)-universal hash family with a uniform public seed \(R\).
For any \(X\) with \(H_{\infty}(X) \ge \gamma\),
\[
\Delta\bigl( (R,\, g_{R}(X)),\; (R,\, U_{n}) \bigr)
\;\le\; \tfrac12 \sqrt{2^{\,n - \gamma}}.
\]
Thus investing \(\gamma \ge n + 2 \log_{2}(1/\epsilon)\) bits of
min-entropy buys \(n\) output bits that are \(\epsilon\)-close to
uniform---a strong \((\gamma, \epsilon)\)-extractor in the
statistical, not merely computational, sense.
\end{theorem}

\paragraph{HKDF, in passing.}
The standardised general-purpose KDF is Krawczyk's \textsf{HKDF}
(RFC~5869), built from HMAC (\S\ref{prf:subsec:hmac}) in an
\emph{extract-then-expand} architecture: a salted HMAC call first
distils suitable non-uniform input keying material into a short
pseudorandom key (the extractor stage, justified along leftover-hash
lines), and a counter-chained sequence of HMAC calls then expands
that key, under a context string, up to RFC~5869's limit of \(255\)
hash-output-length blocks (the expander stage, justified by HMAC's
PRF security). The resulting KDF is a standard choice across many
protocol source models and output lengths, though its precise
guarantee still depends on assumptions about the source and about
HMAC. Zcash does not use it---indeed the shielded protocol uses HMAC
nowhere: neither the \texttt{orchard} nor the
\texttt{zcash\_note\_encryption} crate carries an HMAC dependency
(HMAC enters the wallet stack through the transparent BIP-32
derivation and BIP-39 mnemonic layers and, transitively, through the
optional Tor client, never through a shielded crate)---because its
requirements are narrower and a single hash call meets them in the
random-oracle analysis.

\paragraph{The Zcash KDF.}
Orchard derives the note-encryption key by one personalised BLAKE2b
call:
\[
K_{\mathrm{enc}}
\;=\; \mathrm{KDF}(\mathsf{sharedSecret}, \mathsf{epk})
\;=\; \mathrm{BLAKE2b}\text{-}256\bigl(
  \texttt{"Zcash\_OrchardKDF"};\;
  \mathsf{sharedSecret} \,\|\, \mathsf{epk}
\bigr),
\]
where \(\mathsf{sharedSecret}\) is the Diffie--Hellman shared point
of the note-encryption key agreement,
\(\mathsf{epk}\) is the sender's ephemeral public key, and the quoted
string occupies BLAKE2b's \(16\)-byte personalisation field
(the function
\texttt{kdf\_orchard} hashes the serialised shared point
followed by the ephemeral-key bytes with a \(32\)-byte output;
protocol specification,
\href{https://zips.z.cash/protocol/protocol.pdf\#concreteorchardkdf}{\S``Orchard Key Derivation''}). The
analysis models this personalised hash as a random oracle: by the
domain-separation principle of \S\ref{hash:subsec:domainsep}
(Proposition~\ref{hash:prop:domainsep}, via the native-personalisation
idiom catalogued there), the personalisation makes it an oracle
independent of every other hash use in the protocol, and the derived
key is uniform in the adversary's view \emph{unless} she queries
that oracle at the hidden \(\mathsf{sharedSecret}\)---a query whose
argument she must first compute from the two public keys, an
instance of the computational Diffie--Hellman problem of
\S\ref{assumptions:subsec:dh}. Folding \(\mathsf{epk}\) into the
input binds the derived key to this particular ephemeral, the role
the context string plays in \textsf{HKDF}. No salt is used because
the claim is this narrower
Diffie--Hellman-plus-random-oracle one---the
unsalted premise anticipated in
Definition~\ref{symmetric:def-kdf}---not extraction from every
high-min-entropy source.

\begin{remark}[The role of the KDF in the larger protocol]
\label{symmetric:rem-kdf-role}
In the Zcash note-encryption layer, a Diffie--Hellman exchange on the
Pallas curve produces a shared group element, and the single BLAKE2b
call above condenses that element into the symmetric key fed to
\textsf{ChaCha20-Poly1305}. The two halves of this section meet
exactly here: key derivation manufactures the uniform key that the
AEAD security theorem (Theorem~\ref{symmetric:thm-chachapoly}) takes
as its hypothesis, and the chain from a non-uniform shared secret to
an authenticated ciphertext is complete.
Section~\ref{sec:keyagreement} assembles that chain end to end as
hybrid public-key encryption.
\end{remark}

\section{Key agreement}\label{sec:keyagreement}

Two parties who have never met share no secret, and the adversary of
this section intends to keep it that way. She owns the channel
between them. As an eavesdropper she records the two group elements
they exchange and tries to \emph{distinguish} the key they derive
from a random string---a key she can distinguish is a key whose every
later use leaks. As an active attacker she does better than listen:
she intercepts each party's contribution, substitutes her own, and
runs one honest-looking session with each victim, so that both
``secure channels'' terminate at her desk and she reads and rewrites
every message in transit---forging, as we shall prove, without
violating any hardness assumption at all. And on a public ledger her
seat is permanent: every note ciphertext is published for ever, and
she may scan them all at leisure, holding each against every address
she can enumerate. This section builds the machinery that denies her
all three postures: the Diffie--Hellman protocol atop the
exponentiation asymmetry of Section~\ref{sec:assumptions}; the
static/ephemeral taxonomy that turns an interactive exchange into a
one-shot flow; the key-agreement scheme abstraction that the deployed
code programs against; passive security, which we prove equal to the
DDH advantage \emph{exactly}; the man-in-the-middle attack and the
three external mechanisms of authentication; the elliptic-curve
instantiation with its validation duties; the hybrid composition with
a key-derivation function and an AEAD; key privacy, which keeps a
ciphertext silent about its recipient; and the deployed endpoint,
shielded-note delivery as hybrid public-key encryption carried in
band on the ledger.

\subsection{The Diffie--Hellman protocol}
\label{keyagreement:subsec:dh}

The raw material is the asymmetry established in
Section~\ref{sec:assumptions}. In a cyclic group
\(G = \langle g\rangle\) of order \(q\), the map
\[
\exp_g \colon \Z/q\Z \to G, \qquad x \mapsto g^{x},
\]
is a group isomorphism---it carries addition of exponents to
multiplication of elements---that is cheap to evaluate, at
\(O(\log q)\) group operations by repeated squaring, but believed
infeasible to invert in the groups this volume uses: inverting it is
precisely the discrete logarithm problem
(Definition~\ref{assumptions:def:dlp}). Everything in this section
depends only on the group structure and the hardness of certain
problems in it, never on the presentation of the group; the reference
instances are a large prime-order subgroup of \(\Fp^{\times}\) and
the elliptic-curve point groups of
\S\ref{keyagreement:subsec:ecdh}. Diffie and Hellman's
observation is that the asymmetry alone already lets two strangers
manufacture a shared secret in public.

\begin{construction}[Diffie--Hellman key exchange]
\label{keyagreement:def:dhke}
Fix public parameters \((G, g, q)\) with \(G = \langle g\rangle\) of
order \(q\), known to all parties (and to the adversary).
\begin{enumerate}
\item Alice samples \(a \xleftarrow{\$} \Z/q\Z\) uniformly, computes
  \(A = g^{a}\), and sends \(A\) to Bob.
\item Bob samples \(b \xleftarrow{\$} \Z/q\Z\) uniformly, computes
  \(B = g^{b}\), and sends \(B\) to Alice.
\item Alice computes \(k_{A} = B^{a}\); Bob computes
  \(k_{B} = A^{b}\).
\end{enumerate}
The value \(k_{A} = k_{B}\) is the \emph{shared secret}; the
transmitted elements \(A\) and \(B\) are the \emph{public shares}
(public keys), and the exponents \(a\) and \(b\) are the
\emph{private keys}, never sent.
\end{construction}

\begin{proposition}[Correctness of Diffie--Hellman]
\label{keyagreement:prop:dh-correct}
In Construction~\ref{keyagreement:def:dhke} both parties compute the
same group element,
\(k_{A} = k_{B} = g^{ab}\).
\end{proposition}

\begin{proof}
Unwinding the definitions,
\[
k_{A} = B^{a} = (g^{b})^{a} = g^{ba} = g^{ab} = (g^{a})^{b}
      = A^{b} = k_{B},
\]
where \(ba = ab\) because integer multiplication is commutative, and
exponents may be read modulo \(q\) because \(g\) has order \(q\).
\end{proof}

The shared secret \(g^{ab}\) is thus a well-defined function of the
transcript \((A, B)\): it can be computed by anyone who knows
\emph{one} of the two private keys, and---conjecturally---by no one
who knows only the public shares. Making that conjecture precise, and
seeing exactly how much it buys, is the business of
\S\ref{keyagreement:subsec:passive}. First, a structural precaution.

\begin{remark}[Small subgroups and the prime-order rule]
\label{keyagreement:rem:primeorder}
If \(G\) has small subgroups, an adversary can exploit them: by
substituting for a public share an element of small order she
confines the victim's computed secret to a small subgroup and learns
the victim's exponent modulo a small factor of \(q\)---the
Pohlig--Hellman leakage of
Corollary~\ref{assumptions:cor:large-prime} in protocol form. The
remedy is the standing rule of Section~\ref{sec:assumptions}: prefer
\(G\) of \emph{prime} order \(q\). Then the only subgroups are
\(\{1\}\) and \(G\) itself, every non-identity element is a
generator, and no nontrivial confinement is possible. When the
natural ambient group does not have prime order---\(\Fp^{\times}\)
has composite order \(p - 1\)---one works inside a prime-order
subgroup and has each recipient \emph{validate} incoming elements as
members: check \(B^{q} = 1\) and \(B \ne 1\) before exponentiating.
The elliptic-curve analogue of this validation appears in
Remark~\ref{keyagreement:rem:ec-validation}.
\end{remark}

\subsection{Static and ephemeral keys}
\label{keyagreement:subsec:staticephemeral}

Nothing in Construction~\ref{keyagreement:def:dhke} says how long a
key pair lives, and the protocol changes character with the answer.

\begin{definition}[Static and ephemeral keys]
\label{keyagreement:def:static}
A Diffie--Hellman key pair \((x, g^{x})\) is \emph{static}
(long-term) if it is generated once and reused across many sessions,
typically published and bound to an identity; it is \emph{ephemeral}
if it is generated freshly for a single session and discarded
immediately afterwards.
\end{definition}

Three protocol shapes result, and all three matter to what follows.

\begin{itemize}
\item \emph{Ephemeral--ephemeral (DHE).} Both parties use fresh
  keys. Once \(a\) and \(b\) are erased, no one---the parties
  included---can recover \(g^{ab}\): this is \emph{forward secrecy},
  the guarantee that a later compromise of all stored long-term state
  cannot decrypt recorded past sessions, because the only secrets
  that ever existed were the deleted ephemerals.
\item \emph{Static--ephemeral.} The recipient Bob holds a static pair
  \((b, B = g^{b})\), with \(B\) known in advance; Alice contributes
  a fresh ephemeral \((a, A)\). The shared secret \(g^{ab}\) then
  requires \emph{no interaction from Bob at agreement time}: Alice
  computes \(g^{ab} = B^{a}\) from the published \(B\), attaches
  \(A\), and sends everything in a single flow; Bob recovers
  \(g^{ab} = A^{b}\) whenever he next comes online. This is exactly
  the shape of hybrid encryption
  (\S\ref{keyagreement:subsec:hpke}) and of shielded-note delivery
  (\S\ref{keyagreement:subsec:notes}): sender online once, recipient
  possibly offline, no round trip.
\item \emph{Static--static.} Both keys are long-term, so \(g^{ab}\)
  is a fixed function of the two identities---the same in every
  session. There is no forward secrecy, and the never-varying secret
  must be combined with fresh per-session randomness (a nonce) before
  it may key any symmetric scheme; using it directly would reuse one
  key across all time.
\end{itemize}

The static--ephemeral case is the conceptual bridge from key
exchange, an interactive protocol, to public-key encryption, a
one-shot algorithm. A static public key \(B = g^{b}\) functions
exactly as an encryption public key: anyone can derive a secret
shared with its owner by sending a single ephemeral share. The rest
of the section develops that bridge until, by
\S\ref{keyagreement:subsec:notes}, it carries deployed traffic.

\subsection{The key-agreement scheme abstraction}
\label{keyagreement:subsec:abstraction}

Deployed code does not manipulate exponents; it programs against an
interface. We therefore package the non-interactive
(static--ephemeral) pattern as a scheme with named algorithms, so
that later constructions---and the deployed note-encryption
machinery---can invoke key agreement without caring which group sits
underneath.

\begin{definition}[Key-agreement scheme]
\label{keyagreement:def:kascheme}
A (non-interactive) \emph{key-agreement scheme} over a public
parameter set \(\mathsf{pp}\) consists of a private-key space
\(\mathcal{S}\), a public-key space \(\mathcal{P}\), a shared-secret
space \(\mathcal{K}\), and three algorithms:
\begin{itemize}
\item \(\mathrm{KeyGen}(\mathsf{pp}) \to s\), randomised, outputting
  a private key \(s \in \mathcal{S}\);
\item \(\mathrm{DerivePublic}(\mathsf{pp}, s) \to p\), deterministic,
  mapping a private key to its public key \(p \in \mathcal{P}\);
\item \(\mathrm{Agree}(\mathsf{pp}, s, p') \to k\), deterministic,
  mapping one's own private key and a counterparty's public key to a
  shared secret \(k \in \mathcal{K}\).
\end{itemize}
Correctness demands that for all \(\mathsf{pp}\) and all \(s, t\) in
the support of \(\mathrm{KeyGen}(\mathsf{pp})\),
\begin{equation}
\label{keyagreement:eq:correctness}
\mathrm{Agree}\bigl(\mathsf{pp}, s,
  \mathrm{DerivePublic}(\mathsf{pp}, t)\bigr)
=
\mathrm{Agree}\bigl(\mathsf{pp}, t,
  \mathrm{DerivePublic}(\mathsf{pp}, s)\bigr).
\end{equation}
\end{definition}

\begin{construction}[Diffie--Hellman as a key-agreement scheme]
\label{keyagreement:def:dhscheme}
Fix \(\mathsf{pp} = (G, g, q)\) with \(G = \langle g\rangle\) of
prime order \(q\). The scheme \(\mathrm{DH}_{G,g}\) takes
\(\mathcal{S} = \Z/q\Z\) (or \((\Z/q\Z) \setminus \{0\}\), to avoid
the degenerate case whose public key is the identity),
\(\mathcal{P} = \mathcal{K} = G\), and
\[
\mathrm{KeyGen} \colon a \xleftarrow{\$} \mathcal{S},
\qquad
\mathrm{DerivePublic}(\mathsf{pp}, a) = g^{a},
\qquad
\mathrm{Agree}(\mathsf{pp}, a, B) = B^{a}.
\]
Correctness is exactly
Proposition~\ref{keyagreement:prop:dh-correct}:
\(\mathrm{Agree}(\mathsf{pp}, s, g^{t}) = (g^{t})^{s} = g^{ts}
= g^{st} = (g^{s})^{t} = \mathrm{Agree}(\mathsf{pp}, t, g^{s})\).
\end{construction}

\begin{remark}[The deployed interface generalises the base point]
\label{keyagreement:rem:deployed-abstraction}
The Zcash note-encryption machinery programs against precisely this
abstraction, with one generalisation: the deployed
\(\mathrm{DerivePublic}\) takes the base point as an \emph{explicit
argument} rather than fixing a protocol-wide generator, so that
Sapling and Orchard derive keys over per-address \emph{diversified
bases} (protocol specification
\href{https://zips.z.cash/protocol/protocol.pdf\#abstractkeyagreement}{\S``Key Agreement''}; the
\texttt{zcash\_note\_encryption} crate's \texttt{Domain} trait's
\texttt{ka\_derive\_public}, \texttt{ka\_agree\_enc}, and
\texttt{ka\_agree\_dec} operate on abstract associated types, and
Orchard's implementation of \texttt{derive\_public} in the
\texttt{orchard} crate receives the base
\(g_{d}\) as a parameter). Those bases are put to work in
\S\ref{keyagreement:subsec:notes}, and the address machinery that
produces them is detailed in the \emph{Ironwood Guide}
(\S``Diversified addresses''). What matters here is that the
interface still invokes an abstract key agreement and never
inspects the underlying group---which is why the same construction
has transported across three groups, from Sprout's Curve25519 to
Sapling's Jubjub to Orchard's Pallas.
\end{remark}

\begin{remark}[The abstraction boundary, and what falls outside it]
\label{keyagreement:rem:boundary}
Any concrete scheme satisfying
equation~\eqref{keyagreement:eq:correctness} can be swapped in
without changing the surrounding logic: a prime-field group, an
elliptic curve, or a post-quantum non-interactive key agreement such
as the CSIDH group action. A \emph{key-encapsulation mechanism}
(KEM), by contrast, cannot satisfy the correctness equation at all:
encapsulation requires the recipient's public key before it can
produce its ciphertext, so no counterparty-independent
\(\mathrm{DerivePublic}\) exists. The static--ephemeral one-flow
\emph{pattern} does generalise to any KEM---the encapsulation
ciphertext plays the role of the ephemeral public key---but at the
cost of changing the interface the surrounding logic programs
against.
\end{remark}

\subsection{Passive security and the Diffie--Hellman assumptions}
\label{keyagreement:subsec:passive}

The weakest adversary is the eavesdropper: she reads the channel but
does not alter it, and her entire view of an honest session is the
transcript \((g^{a}, g^{b})\). Both formalisations of her task were
laid down in \S\ref{assumptions:subsec:dh}. The \emph{computational}
problem is to produce the shared secret: writing
\(\mathrm{Adv}_{\mathrm{cdh}}(\mathcal{A})\) for the success
probability of \(\mathcal{A}\) in the CDH game of
Definition~\ref{assumptions:def:cdh},
\[
\mathrm{Adv}_{\mathrm{cdh}}(\mathcal{A})
= \Pr_{a,b}\bigl[\mathcal{A}(g^{a}, g^{b}) = g^{ab}\bigr].
\]
The \emph{decisional} problem is to distinguish
\((g^{a}, g^{b}, g^{ab})\) from \((g^{a}, g^{b}, g^{c})\) with
advantage \(\mathrm{Adv}_{\mathrm{ddh}}(\mathcal{A})\)
(Definition~\ref{assumptions:def:ddh}). The hierarchy
\(\mathrm{DDH} \le \mathrm{CDH} \le \mathrm{DLP}\) was established in
\S\ref{assumptions:subsec:hierarchy}
(Theorems~\ref{assumptions:thm:dlp-le-cdh}
and~\ref{assumptions:thm:cdh-le-ddh}), together with the
pairing-based separation of DDH from CDH
(Example~\ref{assumptions:exa:pairing}). The distinction is not
pedantry. Hardness of CDH asserts only that the \emph{whole} secret
cannot be computed; keying a symmetric cipher requires that the
secret be indistinguishable from a uniform group element, so that no
individual bit or predicate of it leaks. The following classical
break shows the gap is real in the most familiar group.

\begin{remark}[A concrete DDH break: the Legendre symbol]
\label{keyagreement:rem:legendre}
Take \(G = \Fp^{\times}\), the \emph{full} multiplicative group, with
generator \(g\). The Legendre symbol \(\legendre{h}{p}\) is
efficiently computable by Euler's criterion
\(h^{(p-1)/2} \bmod p\) and is multiplicative (\emph{Math Guide},
\S``Quadratic residues and the Euler criterion''). A generator of
\(\Fp^{\times}\) is a nonsquare, so
\(\legendre{g^{a}}{p} = (-1)^{a}\): the symbol of a public share
reveals the \emph{parity} of its exponent. The parities of \(a\) and
\(b\) determine the parity of \(ab\), hence
\(\legendre{g^{ab}}{p} = (-1)^{ab}\) is determined by the two
observed symbols. The distinguisher computes the predicted symbol of
the shared secret from \((g^{a}, g^{b})\), compares it with the
symbol of the third component, and outputs ``Diffie--Hellman'' on a
match. On a genuine triple the match is certain; on a random triple
the third component's symbol is an independent uniform sign, matching
with probability \(1/2\). The distinguishing advantage is the
constant \(1/2\)---even though CDH in \(\Fp^{\times}\) is believed
hard. The CDH/DDH gap therefore needs no pairings; the oldest group
in the subject exhibits it. The cure is the prime-order rule of
Remark~\ref{keyagreement:rem:primeorder}: in a subgroup of odd prime
order \(q\), every element \(h\) is a square
(\(h = (h^{(q+1)/2})^{2}\)), the symbol is identically \(+1\), and
the leak vanishes. This is one of the two reasons---the other being
Pohlig--Hellman---that DDH is only ever assumed in prime-order
groups.
\end{remark}

With the assumption in place, passive security is not merely implied
by DDH; it \emph{is} DDH, advantage for advantage. We state the game
in the volume's tradition and prove the identification.

\begin{definition}[Passive security of a key-agreement scheme]
\label{keyagreement:def:passive}
For a key-agreement scheme \(\Sigma\) over \(\mathsf{pp}\) and an
adversary \(\mathcal{A}\), the game
\(\mathsf{KA\text{-}IND}^{\mathcal{A}}_{\Sigma}(\lambda)\) runs as
follows.
\begin{enumerate}
\item The challenger samples
  \(s, t \gets \mathrm{KeyGen}(\mathsf{pp})\) independently and
  computes the transcript
  \(\bigl(\mathrm{DerivePublic}(\mathsf{pp}, s),\,
  \mathrm{DerivePublic}(\mathsf{pp}, t)\bigr)\).
\item It sets
  \(k_{0} = \mathrm{Agree}(\mathsf{pp}, s,
  \mathrm{DerivePublic}(\mathsf{pp}, t))\), samples
  \(k_{1} \xleftarrow{\$} \mathcal{K}\) uniformly, draws a hidden bit
  \(\beta \xleftarrow{\$} \{0,1\}\), and sends the transcript
  together with \(k_{\beta}\).
\item The adversary outputs a bit \(\beta'\).
\end{enumerate}
The advantage is
\(\mathrm{Adv}^{\mathrm{ka\text{-}ind}}_{\Sigma}(\mathcal{A})
= \abs{\Pr[\beta' = 1 \mid \beta = 0]
- \Pr[\beta' = 1 \mid \beta = 1]}\); the scheme is \emph{passively
secure} if every PPT adversary's advantage is negligible in
\(\lambda\).
\end{definition}

\begin{theorem}[Passive security of Diffie--Hellman]
\label{keyagreement:thm:passive}
For the scheme \(\mathrm{DH}_{G,g}\) of
Construction~\ref{keyagreement:def:dhscheme} with
\(\mathrm{KeyGen}\) uniform on \(\Z/q\Z\), every adversary
\(\mathcal{A}\) satisfies
\[
\mathrm{Adv}^{\mathrm{ka\text{-}ind}}_{\mathrm{DH}_{G,g}}(\mathcal{A})
\;=\;
\mathrm{Adv}_{\mathrm{ddh}}(\mathcal{A}),
\]
identically and with running time preserved exactly. Under the DDH
assumption the shared secret of an honest ephemeral--ephemeral
session is therefore pseudorandom given the transcript: a passive
adversary learns nothing about it that she could not have guessed
about a random group element.
\end{theorem}

\begin{proof}
The two games are the same game. In
\(\mathsf{KA\text{-}IND}\) for \(\mathrm{DH}_{G,g}\), the transcript
is \((g^{a}, g^{b})\) with \(a, b\) uniform and independent; the real
key is \(k_{0} = g^{ab}\); and the uniform alternative
\(k_{1} \xleftarrow{\$} G\) is exactly \(g^{c}\) for uniform
independent \(c\), because \(\exp_g\) is a bijection. The challenge
\((g^{a}, g^{b}, k_{\beta})\) is thus verbatim a sample of
\(\mathcal{D}_{\mathrm{dh}}\) or \(\mathcal{D}_{\mathrm{rand}}\)
according to \(\beta\)---the DDH distinguishing game of
Definition~\ref{assumptions:def:ddh}. The identity map converts each
adversary into the other, so the advantages coincide and no running
time is added.
\end{proof}

\begin{remark}[Why pseudorandomness, not mere unpredictability]
\label{keyagreement:rem:pseudorandom}
Every use of a Diffie--Hellman secret in this volume feeds it into a
key-derivation function to produce a symmetric key. What should the
group deliver to that function? Hardness of CDH
guarantees the secret is hard to produce \emph{in full}, but not that
it lacks efficiently predictable structure: the Legendre leak of
Remark~\ref{keyagreement:rem:legendre} is exactly such structure.
Hardness of DDH guarantees the secret is as good as a random group
element (Theorem~\ref{keyagreement:thm:passive})---no predicate of it
is efficiently computable from the transcript. In practice one often
assumes only CDH and models the hash-based KDF as a random oracle:
the derived key is then uniform unless the adversary queries the
oracle at the hidden shared secret, a route made precise in
Theorem~\ref{keyagreement:thm:hpke-secure}.
Either way the design goal is the same: a derived symmetric key
indistinguishable from uniform, because that is the hypothesis every
theorem of Section~\ref{sec:symmetric} consumes.
\end{remark}

\subsection{Active attacks and the role of authentication}
\label{keyagreement:subsec:active}

Passive security is the most the bare protocol can offer, and the
gap between it and what a channel needs is best exhibited as a
construction. The adversary now controls the wire.

\begin{definition}[Man-in-the-middle attack on unauthenticated
Diffie--Hellman]
\label{keyagreement:def:mitm}
An active adversary \(\mathcal{M}\) controlling the channel between
Alice and Bob runs two independent Diffie--Hellman sessions, one with
each victim. She samples \(m \xleftarrow{\$} \Z/q\Z\) and computes
\(g^{m}\). When Alice sends \(A = g^{a}\) intended for Bob,
\(\mathcal{M}\) intercepts it and forwards \(g^{m}\) to Bob in its
place; when Bob replies \(B = g^{b}\), she intercepts it and forwards
\(g^{m}\) to Alice. Alice computes
\(k_{A} = (g^{m})^{a} = g^{am}\), which \(\mathcal{M}\) also computes
as \(A^{m}\); symmetrically Bob computes
\(k_{B} = (g^{m})^{b} = g^{bm}\), which \(\mathcal{M}\) computes as
\(B^{m}\).
\end{definition}

% TikZ figure placeholder --- fig:mitm: man in the middle relaying two
% independent DH sessions. Three columns (Alice, M, Bob); Alice's A
% and Bob's B each terminate at M, who substitutes g^m in both
% directions; the two half-session keys g^{am} and g^{bm} annotated at
% M, with "believes she shares with Bob" / "believes he shares with
% Alice" callouts. Dedicated figures pass supplies the drawing.

\begin{proposition}[The attack defeats unauthenticated
Diffie--Hellman]
\label{keyagreement:prop:mitm}
After the attack of Definition~\ref{keyagreement:def:mitm}, Alice and
Bob each hold a secret they believe is shared with the other but is
in fact shared with \(\mathcal{M}\); and \(\mathcal{M}\), knowing
both \(g^{am}\) and \(g^{bm}\), can decrypt, read, modify, and
re-encrypt every subsequent message in each direction, transparently
to both parties. No assumption on \(G\)---not even ideal hardness of
the discrete logarithm problem---prevents this.
\end{proposition}

\begin{proof}
Correctness of Diffie--Hellman
(Proposition~\ref{keyagreement:prop:dh-correct}), applied once per
session, gives agreement between Alice and \(\mathcal{M}\) on
\(g^{am}\) and between Bob and \(\mathcal{M}\) on \(g^{bm}\).
The adversary knows her own exponent \(m\) and received both \(A\)
and \(B\), so she computes both secrets directly---no problem
instance is ever inverted, which is why no hardness assumption is
relevant. Neither victim verifies the \emph{origin} of the group
element they received; each transcript is a syntactically valid run
of Construction~\ref{keyagreement:def:dhke}, so nothing in either
party's view distinguishes the attacked execution from an honest one.
With \(k_{A}\) and \(k_{B}\) in hand, \(\mathcal{M}\) decrypts
traffic from Alice under \(k_{A}\), re-encrypts it under \(k_{B}\)
for Bob, and symmetrically in reverse, editing at will in between.
\end{proof}

The lesson deserves stating structurally: \emph{passive security
concerns the secrecy of the shared secret; active security
additionally requires the authenticity of the public shares}.
Confidentiality of a channel is worthless if the channel is
established with the wrong party. The bare protocol
(Construction~\ref{keyagreement:def:dhke}) provides no
authentication---nothing in \(g^{a}\) or \(g^{b}\) testifies to who
sent it---so the remedy must bind each public share to its sender's
identity, letting the recipient reject a substituted share. Every
such binding is \emph{external} to the key exchange itself.

\begin{remark}[Three mechanisms supply authentication]
\label{keyagreement:rem:auth-supplied}
All deployed remedies fall into three families, each external to the
bare exchange.
\begin{enumerate}
\item \emph{Signed Diffie--Hellman.} Each party signs its ephemeral
  share with a long-term signing key whose verification key is
  certified---by a certificate authority (a third party whose
  signature on the key the parties already trust) or by
  trust-on-first-use (accepting the first key seen and rejecting
  later changes).
  The adversary of Proposition~\ref{keyagreement:prop:mitm} cannot
  forge Alice's signature on \(g^{m}\), so her substitution is
  rejected. Digital signatures are the subject of
  Section~\ref{sec:signatures}; this composition is the structure of
  authenticated TLS.
\item \emph{Static keys that are themselves authenticated.} If Bob's
  static public key \(B = g^{b}\) reaches Alice authentically---a
  trusted directory, a pinned key---then a static--ephemeral exchange
  toward \(B\) already authenticates Bob to Alice: only the holder of
  \(b\) can complete it. Note the direction: this authenticates the
  \emph{recipient}, not the sender. It is the model relevant to
  public-key encryption, and to shielded-note delivery, where the
  recipient's address \emph{is} an authenticated public key
  (\S\ref{keyagreement:subsec:notes}).
\item \emph{Out-of-band verification.} The parties compare a short
  fingerprint of the exchanged public keys over a channel that is
  authentic even if low-bandwidth---digits read aloud on a call, a QR
  code scanned in person.
\end{enumerate}
In every case the key-agreement security goal remains the passive
guarantee of Theorem~\ref{keyagreement:thm:passive}; authentication
is a separable layer composed around the exchange, not a property of
the group operation.
\end{remark}

\subsection{Elliptic-curve Diffie--Hellman}
\label{keyagreement:subsec:ecdh}

The development so far is deliberately group-agnostic; we now
instantiate it on the groups the deployed protocol actually uses.
Recall from the \emph{Math Guide} (\S``Elliptic curves'') that an
elliptic curve \(E\) over a field \(F\) of characteristic neither
\(2\) nor \(3\) is given in short Weierstrass form
\(y^{2} = x^{3} + ax + b\) with nonzero discriminant, and that its
\(F\)-rational points \(E(F)\) form a finite abelian group under
chord-and-tangent addition with the point at infinity
\(\mathcal{O}\) as identity (\emph{Math Guide}, \S``The group
structure of \(E(\Fp)\)''). The group is written additively, so
exponentiation \(g \mapsto g^{x}\) becomes \emph{scalar
multiplication} \(P \mapsto [x]P\), computable in \(O(\log x)\) point
additions by double-and-add. Elliptic-curve Diffie--Hellman is the
transliteration.

\begin{construction}[ECDH]
\label{keyagreement:def:ecdh}
Fix an elliptic curve \(E/F\), a base point \(P \in E(F)\) of large
prime order \(q\) generating \(G = \langle P\rangle\), and cofactor
\(h = \abs{E(F)}/q\). In the abstraction of
Definition~\ref{keyagreement:def:kascheme}:
\[
\mathrm{KeyGen} \colon a \xleftarrow{\$} \Z/q\Z,
\qquad
\mathrm{DerivePublic}(a) = [a]P,
\qquad
\mathrm{Agree}(a, Q) = [a]Q.
\]
The shared secret of parties with private keys \(a\) and \(b\) is
\([a]([b]P) = [ab]P = [b]([a]P)\). The associated hard problems are
the discrete-logarithm, CDH, and DDH definitions of
Section~\ref{sec:assumptions}
(Definitions~\ref{assumptions:def:dlp}, \ref{assumptions:def:cdh},
and~\ref{assumptions:def:ddh}) read with \([x]P\) in place of
\(g^{x}\)---the problems ECDLP (\emph{Math Guide}, \S``The
elliptic-curve discrete logarithm problem''), ECCDH, and ECDDH of
\S\ref{assumptions:subsec:pasta}.
\end{construction}

\begin{proposition}[Correctness of ECDH]
\label{keyagreement:prop:ecdh-correct}
The scheme of Construction~\ref{keyagreement:def:ecdh} satisfies the
key-agreement correctness
condition~\eqref{keyagreement:eq:correctness}.
\end{proposition}

\begin{proof}
Scalar multiplication is cyclic-group exponentiation in additive
dress: the group is abelian, and \(\Z\) acts on
\(\langle P\rangle\) through \(\Z/q\Z\) because \([q]P =
\mathcal{O}\), so \([x]([y]P) = [xy]P\) for all integers \(x, y\).
Hence
\(\mathrm{Agree}(a, [b]P) = [ab]P = [ba]P = \mathrm{Agree}(b, [a]P)\),
which is condition~\eqref{keyagreement:eq:correctness}.
\end{proof}

\begin{remark}[Why elliptic curves]
\label{keyagreement:rem:why-ec}
Well-chosen elliptic-curve groups admit only the generic
\(O(\sqrt q)\) attacks of \S\ref{assumptions:subsec:generic}, whereas
index calculus solves the discrete logarithm in \(\Fp^{\times}\) in
subexponential time
(Remark~\ref{assumptions:rem:index-calculus}); a \(256\)-bit curve
therefore delivers the security level that a multi-thousand-bit
prime field requires (Example~\ref{assumptions:exa:128}). That
efficiency, together with the availability of curves engineered for
cheap in-circuit arithmetic, is why the Pallas curve of
\S\ref{assumptions:subsec:pasta} underlies Orchard.
\end{remark}

\begin{remark}[Subgroup and curve validation]
\label{keyagreement:rem:ec-validation}
A recipient computing \([a]Q\) on an attacker-supplied point \(Q\)
must guard against two pitfalls, the elliptic-curve forms of the
validation duty in Remark~\ref{keyagreement:rem:primeorder}.
\begin{enumerate}
\item \emph{Invalid-curve attack.} The received coordinates may
  satisfy not \(E\)'s equation but that of a different, weaker curve
  which shares the same addition formulas (the chord and tangent
  slopes of the \emph{Math Guide}, \S``Explicit affine formulas'',
  involve \(a\) but never \(b\)); scalar multiplication then runs
  happily in
  the wrong group, whose order may be smooth. The implementation must
  check the curve equation on every incoming point.
\item \emph{Small-subgroup attack.} When the cofactor \(h > 1\), an
  adversary submits a point \(Q\) of small order dividing \(h\),
  confining \([a]Q\) to a tiny subgroup and learning \(a\) modulo a
  small factor. The defences are \emph{cofactor clearing}---multiply
  incoming points by \(h\)---or choosing a curve with \(h = 1\).
\end{enumerate}
Prime-order curves such as Pallas sidestep the cofactor issue
entirely: with \(h = 1\) there is no small subgroup to hit, and only
the on-curve check remains. The deployed arithmetic enforces exactly
that check at deserialisation, as documented in
Proposition~\ref{assumptions:prop:pasta-secure}, whose compressed
encoding cannot even represent a point off the curve.
\end{remark}

\subsection{Hybrid public-key encryption}
\label{keyagreement:subsec:hpke}

Three obstacles separate the Diffie--Hellman shared group element
from a usable encryption scheme. First, the shared secret is a
\emph{structured group element}---a curve point with coordinates
satisfying an equation---not a uniform string of key bits. Second,
the exchange by itself transports no message: it manufactures a
secret, and stops. Third, we want the static--ephemeral
non-interactive flow of \S\ref{keyagreement:subsec:staticephemeral},
so that the recipient need not be online. All three are resolved by
one composition: key agreement for the secret, a key-derivation
function for uniformity, an authenticated cipher for the
message---\emph{hybrid public-key encryption} (HPKE). The two symmetric
ingredients were built in Section~\ref{sec:symmetric}; we recall them
in the specialised shape the composition needs.

\begin{definition}[KDF, recalled and specialised]
\label{keyagreement:def:kdf}
A \emph{key-derivation function} here is a deterministic function
\[
\mathrm{KDF} \colon \mathcal{K} \times \{0,1\}^{*}
\to \{0,1\}^{\ell},
\]
mapping input keying material \(k \in \mathcal{K}\) together with a
context string to an \(\ell\)-bit key. This specialises
Definition~\ref{symmetric:def-kdf} (\S\ref{symmetric:subsec:kdf})
by omitting the salt and fixing the output length \(\ell\). For
Theorem~\ref{keyagreement:thm:hpke-secure} we model the function as
a random oracle (\S\ref{hash:subsec:rom}): distinct inputs then have
independent uniform outputs, and security follows provided the
adversary never queries the oracle at the hidden Diffie--Hellman
secret and context. This is a source-specific random-oracle
guarantee, not a generic claim that a deterministic function
extracts uniform bits from every high-min-entropy source.
\end{definition}

\begin{definition}[AEAD, recalled]
\label{keyagreement:def:aead}
An \emph{authenticated encryption scheme with associated data} with
key space \(\{0,1\}^{\ell}\)---matching the KDF output---is a pair
\((\mathrm{Enc}, \mathrm{Dec})\) with
\(\mathrm{Enc}_{K}(N, \mathit{ad}, m) = c\) and
\(\mathrm{Dec}_{K}(N, \mathit{ad}, c) \in \{m\} \cup \{\bot\}\),
where \(N\) is a nonce, \(\mathit{ad}\) is associated data
authenticated but not encrypted, and \(\bot\) denotes rejection
(Definition~\ref{symmetric:def-aead}). Correctness:
\(\mathrm{Dec}_{K}(N, \mathit{ad},
\mathrm{Enc}_{K}(N, \mathit{ad}, m)) = m\). Security is
authenticated-encryption security
(Definition~\ref{symmetric:def-ae}): encryptions of any two
equal-length messages are CPA-indistinguishable
(Definition~\ref{symmetric:def-cpa}), \emph{and} no adversary can
produce a fresh ciphertext decrypting to anything but \(\bot\)
(ciphertext integrity, Definition~\ref{symmetric:def-intctxt})---
confidentiality and integrity simultaneously.
\end{definition}

\begin{construction}[DH-based hybrid public-key encryption]
\label{keyagreement:def:hpke}
Fix ECDH over \(G = \langle P\rangle\) of prime order \(q\)
(Construction~\ref{keyagreement:def:ecdh}, in the interface of
Definition~\ref{keyagreement:def:kascheme}), a secure KDF, and a
secure AEAD as above. The recipient holds a static key pair
\((b, B = [b]P)\) with \(B\) published. To \emph{encrypt} a message
\(m\) to \(B\):
\begin{enumerate}
\item sample an ephemeral \(e \xleftarrow{\$} \Z/q\Z\) and set
  \(E_{\mathrm{pk}} = [e]P\);
\item compute the shared secret
  \(S = \mathrm{Agree}(e, B) = [eb]P\);
\item derive the symmetric key
  \(K = \mathrm{KDF}(S, \mathit{context})\), where
  \(\mathit{context}\) includes \(E_{\mathrm{pk}}\) and the
  protocol's labels;
\item output \((E_{\mathrm{pk}}, c)\) with
  \(c = \mathrm{Enc}_{K}(N, \mathit{ad}, m)\) for a fixed or derived
  nonce \(N\)---fixed is safe because each \(K\) is used once.
\end{enumerate}
To \emph{decrypt} \((E_{\mathrm{pk}}, c)\) with \(b\): compute
\(S' = [b]E_{\mathrm{pk}}\), derive
\(K' = \mathrm{KDF}(S', \mathit{context})\) with
\(E_{\mathrm{pk}}\) read from the ciphertext, and output
\(\mathrm{Dec}_{K'}(N, \mathit{ad}, c)\), which is \(m\) on an
honest ciphertext; on a tampered one, ciphertext integrity makes
the output \(\bot\) except with negligible probability. Anyone holding
\(B\) can encrypt; only the holder of \(b\) can decrypt.
\end{construction}

% TikZ figure placeholder --- fig:hpke: hybrid encryption dataflow.
% Sender lane: e sampled, E_pk = [e]P published; B (recipient's
% static key) enters Agree to give S = [eb]P; S and context (E_pk,
% labels) enter KDF to give K; K, nonce N, ad, and m enter AEAD Enc
% to give c; wire carries (E_pk, c). Recipient lane mirrors it:
% [b]E_pk -> S' -> KDF -> K' -> Dec -> m or bot. Dedicated figures
% pass supplies the drawing.

\begin{theorem}[Correctness of hybrid encryption]
\label{keyagreement:thm:hpke-correct}
For every message \(m\) and every honestly generated ciphertext
\((E_{\mathrm{pk}}, c)\) of
Construction~\ref{keyagreement:def:hpke}, decryption returns \(m\).
\end{theorem}

\begin{proof}
Correctness of ECDH
(Proposition~\ref{keyagreement:prop:ecdh-correct}) gives
\(S' = [b]([e]P) = [be]P = [eb]P = [e]B = S\). The KDF is
deterministic and both sides supply the same context---the recipient
recovers \(E_{\mathrm{pk}}\) from the ciphertext itself---so
\(K' = K\). Correctness of the AEAD
(Definition~\ref{keyagreement:def:aead}) finishes:
\(\mathrm{Dec}_{K}(N, \mathit{ad},
\mathrm{Enc}_{K}(N, \mathit{ad}, m)) = m\).
\end{proof}

\begin{theorem}[Passive security of hybrid encryption]
\label{keyagreement:thm:hpke-secure}
Model the KDF as a random oracle
(Definition~\ref{hash:def:ro}) and let the AEAD be secure in the
sense of Definition~\ref{symmetric:def-ae}. Consider a passive
adversary \(\mathcal{A}\) who sees the public key \(B\) and a single
challenge ciphertext: she names two equal-length messages
\(m_{0}, m_{1}\), receives the encryption of \(m_{\beta}\) for a
hidden uniform bit \(\beta\), and guesses \(\beta\), with CPA
advantage \(\mathrm{Adv}_{\mathrm{cpa}}(\mathcal{A})\) defined as in
Definition~\ref{symmetric:def-cpa}. Under CDH in \(G\) no PPT
adversary has non-negligible advantage; concretely, if
\(\mathcal{A}\) makes at most \(q_{H}\) random-oracle queries, then
\[
\mathrm{Adv}_{\mathrm{cpa}}(\mathcal{A})
\;\le\;
2\,q_{H} \cdot \mathrm{Adv}_{\mathrm{cdh}}(\mathcal{B})
\;+\;
\mathrm{Adv}_{\mathrm{ae}}(\mathcal{C}),
\]
for explicit reductions \(\mathcal{B}\) against CDH and
\(\mathcal{C}\) against the AEAD's authenticated-encryption security
(of which the proof consumes only the confidentiality half, so
\(\mathrm{Adv}_{\mathrm{ae}}(\mathcal{C})\) is the CPA advantage of
Definition~\ref{symmetric:def-cpa} read against the AEAD), each
running in essentially the same time as \(\mathcal{A}\).
\end{theorem}

\begin{proof}
The proof is a game hop.

\emph{Game 0} is the real CPA game: the challenge ciphertext is
\((E_{\mathrm{pk}}, c)\) with \(E_{\mathrm{pk}} = [e]P\),
\(S = [eb]P\), \(K = \mathrm{KDF}(S, \mathit{context})\), and
\(c = \mathrm{Enc}_{K}(N, \mathit{ad}, m_{\beta})\). Say
\(\mathcal{A}\) \emph{wins} a game when her guess \(\beta'\) equals
\(\beta\).

\emph{Game 1} replaces \(K\) by an independent uniform
\(\ell\)-bit string, leaving everything else unchanged.

\emph{The hop.} In the random oracle model the two games proceed
identically unless the adversary queries the oracle at exactly the
point \((S, \mathit{context}) = ([eb]P, \mathit{context})\): until
that query, the oracle's value there is a uniform string she has
never seen, which is precisely what Game~1 hands out. But producing
\(S = [eb]P\) from the view \((B, E_{\mathrm{pk}}) = ([b]P, [e]P)\)
is exactly solving CDH. The reduction \(\mathcal{B}\) embeds its CDH
challenge as \((B, E_{\mathrm{pk}})\) and simulates Game~1 around
it: it draws \(\beta\) and a uniform \(\ell\)-bit key \(K\) itself,
encrypts \(m_{\beta}\) under \(K\) to form the challenge
ciphertext, and implements the oracle by lazy sampling---a fresh
uniform string for each new query, the stored answer on a repeated
one, so that the oracle remains a consistent function. The
simulation is a perfect run of Game~1, and the critical query
occurs in it with the same probability as in Game~0, the two games
being identical until it happens. One obstruction remains: in a
pairing-free group \(\mathcal{B}\) cannot
\emph{recognise} which of the queries carries \([eb]P\)
(Example~\ref{assumptions:exa:pairing} shows recognising it is the
DDH-test a gap group would supply). So \(\mathcal{B}\) picks one of
the at most \(q_{H}\) queries uniformly at random and outputs its
first component; whenever the critical query occurs,
\(\mathcal{B}\)'s choice is correct with probability at least
\(1/q_{H}\). Hence
\[
\bigl|\Pr[\mathcal{A} \text{ wins Game } 0]
- \Pr[\mathcal{A} \text{ wins Game } 1]\bigr|
\;\le\;
q_{H} \cdot \mathrm{Adv}_{\mathrm{cdh}}(\mathcal{B}).
\]
The factor \(q_{H}\) is a genuine guessing-type security loss, in
the tightness taxonomy of Definition~\ref{model:def:tightness}.
Assuming \emph{gap}-CDH---CDH hardness even given a DDH-decision
oracle, which would let \(\mathcal{B}\) test each query---restores
tightness; we state the loose bound because it is what plain CDH
buys, and honesty about the loss is part of the bound.

\emph{Game 1 analysed.} The AEAD key is now uniform and independent
of everything else in the adversary's view, which is exactly the
hypothesis of the AEAD's confidentiality game. The reduction
\(\mathcal{C}\) forwards \((m_{0}, m_{1})\) to its AEAD challenger,
embeds the returned ciphertext as \(c\) alongside a self-generated
\(E_{\mathrm{pk}}\), and relays \(\mathcal{A}\)'s guess; the AEAD
challenger's hidden bit plays \(\beta\), so
\(\abs{\Pr[\mathcal{A} \text{ wins Game } 1] - \tfrac12}
= \tfrac12\,\mathrm{Adv}_{\mathrm{ae}}(\mathcal{C})\).

\emph{Accounting.} The difference-form advantage of
Definition~\ref{symmetric:def-cpa} is twice the bit-guessing
advantage \(\abs{\Pr[\mathcal{A} \text{ wins}] - \tfrac12}\) of
Definition~\ref{model:def:game}, so the triangle inequality across
the hop gives
\[
\mathrm{Adv}_{\mathrm{cpa}}(\mathcal{A})
\;\le\;
2\,\bigl(q_{H} \cdot \mathrm{Adv}_{\mathrm{cdh}}(\mathcal{B})
+ \tfrac12\,\mathrm{Adv}_{\mathrm{ae}}(\mathcal{C})\bigr)
\;=\;
2\,q_{H} \cdot \mathrm{Adv}_{\mathrm{cdh}}(\mathcal{B})
+ \mathrm{Adv}_{\mathrm{ae}}(\mathcal{C}),
\]
the stated bound.

\emph{Alternative hypothesis.} The random oracle can be traded
away, but not for Definition~\ref{symmetric:def-kdf} alone: that
guarantee is stated for a uniform independent salt, which the
specialised KDF of Definition~\ref{keyagreement:def:kdf} omits, and
the definition denies any generic unsalted guarantee. The trade
needs an explicit hypothesis on the deployed function---the
narrower premise Definition~\ref{symmetric:def-kdf} allows for
unsalted use: for a \emph{uniform} group element \(S\), the output
\(\mathrm{KDF}(S, \mathit{context})\) is computationally
indistinguishable from a uniform \(\ell\)-bit string given the
context. Under that hypothesis and DDH in place of CDH,
Theorem~\ref{keyagreement:thm:passive} first replaces the shared
secret by a uniform group element at cost
\(\mathrm{Adv}_{\mathrm{ddh}}\), and the stated
indistinguishability then replaces \(K\) by a uniform string---no
oracle and no \(q_{H}\) factor, at the price of the stronger
decisional assumption on the group and an explicit pseudorandomness
assumption on the KDF.
\end{proof}

\begin{remark}[Every ingredient is necessary]
\label{keyagreement:rem:necessary}
Dropping any piece of Construction~\ref{keyagreement:def:hpke}
breaks the composition.
\begin{itemize}
\item \emph{Without the KDF}, one would key the cipher with the raw
  group element \(S\). Its bit-representation is non-uniform---curve
  points do not fill the encoding space---and in non-DDH groups it is
  partially predictable, the Legendre-symbol leakage of
  Remark~\ref{keyagreement:rem:legendre} being the classical
  instance. The KDF standardises the format, extracts uniform bits,
  and, by folding \(E_{\mathrm{pk}}\) into its context, ties the
  derived key to this particular ephemeral share.
\item \emph{Without AEAD integrity}, a passively confidential but
  malleable cipher would let an active adversary tamper with \(c\)
  undetectably---the XOR-malleability of
  Section~\ref{sec:symmetric}'s stream ciphers is total. Ciphertext
  integrity (Definition~\ref{symmetric:def-intctxt}) makes any
  modification decrypt to \(\bot\) except with negligible
  probability, giving the symmetric payload its chosen-ciphertext
  robustness.
\item \emph{Without the ephemeral key}---a static--static
  variant---every message to \(B\) would begin from the same shared
  secret. The fresh per-message ephemeral, with \(E_{\mathrm{pk}}\)
  bound into the KDF context, gives each ciphertext separate key
  material, so nothing relies on nonce uniqueness under a reused
  key. What it does \emph{not} supply is forward secrecy against
  later compromise of the recipient's static key \(b\): from a
  recorded \(E_{\mathrm{pk}}\), an attacker who learns \(b\)
  recomputes \([b]E_{\mathrm{pk}}\) and with it the ciphertext's
  key. Erasing \(e\) buys only the sender-side guarantee that the
  key derives from no stored secret of the sender.
\end{itemize}
\end{remark}

\subsection{Key privacy}
\label{keyagreement:subsec:keyprivacy}

Theorem~\ref{keyagreement:thm:hpke-secure} answers the eavesdropper
and leaves the scanner untouched. Her posture is different: she holds
a published ciphertext against every public key she can enumerate
and asks not what it says but \emph{to whom}. Confidentiality is
silent on the question---a scheme may hide the message perfectly and
still stamp the recipient's key on every ciphertext. The property
that denies her is \emph{key privacy}, due to Bellare, Boldyreva,
Desai, and Pointcheval. We state it for a generic public-key scheme,
then prove it for the bare static--ephemeral flow with the message
folded in---ElGamal encryption, the shape whose key privacy the upper
volumes consume.

\begin{definition}[Key privacy under chosen-plaintext attack, IK-CPA]
\label{keyagreement:def:ikcpa}
A \emph{public-key encryption scheme} \(\Pi\) over \(\mathsf{pp}\)
consists of a randomised
\(\mathrm{KeyGen}(\mathsf{pp}) \to (\mathsf{sk}, \mathsf{pk})\), a
randomised \(\mathrm{Enc}_{\mathsf{pk}}(m) \to c\), and a
deterministic \(\mathrm{Dec}_{\mathsf{sk}}(c)\) with
\(\mathrm{Dec}_{\mathsf{sk}}(\mathrm{Enc}_{\mathsf{pk}}(m)) = m\) for
every key pair in the support of \(\mathrm{KeyGen}\) and every
message \(m\). For an adversary \(\mathcal{A}\) the game
\(\mathsf{IK\text{-}CPA}^{\mathcal{A}}_{\Pi}(\lambda)\) runs as
follows.
\begin{enumerate}
\item The challenger samples
  \((\mathsf{sk}_{0}, \mathsf{pk}_{0}), (\mathsf{sk}_{1},
  \mathsf{pk}_{1}) \gets \mathrm{KeyGen}(\mathsf{pp})\)
  independently and sends \((\mathsf{pk}_{0}, \mathsf{pk}_{1})\).
\item The adversary outputs a message \(m\).
\item The challenger draws a hidden bit
  \(\beta \xleftarrow{\$} \{0,1\}\) and sends
  \(c \gets \mathrm{Enc}_{\mathsf{pk}_{\beta}}(m)\).
\item The adversary outputs a bit \(\beta'\).
\end{enumerate}
The advantage is
\(\mathrm{Adv}^{\mathrm{ik\text{-}cpa}}_{\Pi}(\mathcal{A})
= \abs{\Pr[\beta' = 1 \mid \beta = 0]
- \Pr[\beta' = 1 \mid \beta = 1]}\); the scheme is
\emph{key-private} (IK-CPA) if every PPT adversary's advantage is
negligible in \(\lambda\). The message is the adversary's own, so
nothing about it is hidden; what the game hides is the key.
\end{definition}

\begin{construction}[ElGamal encryption]
\label{keyagreement:def:elgamal}
Fix \(\mathsf{pp} = (G, P, q)\) as in
Construction~\ref{keyagreement:def:ecdh}, written additively, and
take group elements as messages; write \(\mathrm{EG}\) for the
scheme. A key pair carries its own base \(H\): in the
\emph{fixed-base} variant \(H = P\); in the
\emph{per-key-base} variant \(\mathrm{KeyGen}\) samples
\(H \xleftarrow{\$} G \setminus \{\mathcal{O}\}\), the generalised
interface of Remark~\ref{keyagreement:rem:deployed-abstraction}.
\begin{itemize}
\item \(\mathrm{KeyGen}\): choose \(H\) as above and
  \(x \xleftarrow{\$} \Z/q\Z\); output \(\mathsf{sk} = x\) and
  \(\mathsf{pk} = (H, Y)\) with \(Y = [x]H\).
\item \(\mathrm{Enc}_{(H, Y)}(M)\): sample
  \(r \xleftarrow{\$} \Z/q\Z\); output
  \((R, C) = \bigl([r]H,\; [r]Y + M\bigr)\).
\item \(\mathrm{Dec}_{x}(R, C)\): output \(C - [x]R\).
\end{itemize}
Correctness is the scalar identity in the proof of
Proposition~\ref{keyagreement:prop:ecdh-correct}, valid for any base
in \(G\): \([x]R = [x]([r]H) = [r]([x]H) = [r]Y\), so
\(C - [x]R = M\). The scheme is the static--ephemeral flow of
\S\ref{keyagreement:subsec:staticephemeral} with the shared secret
\([rx]H\) spent as a one-time pad on a group element; the fixed-base
variant is ElGamal's original scheme transposed from
\(\Fp^{\times}\) to a prime-order group.
\end{construction}

\begin{theorem}[ElGamal is key-private under DDH]
\label{keyagreement:thm:ikcpa}
For ElGamal over \(G\) (Construction~\ref{keyagreement:def:elgamal},
either base variant), every adversary \(\mathcal{A}\) yields two
distinguishers \(\mathcal{B}_{0}, \mathcal{B}_{1}\) against DDH in
\(G\) (Definition~\ref{assumptions:def:ddh}, read with \([x]P\) in
place of \(g^{x}\) as in Construction~\ref{keyagreement:def:ecdh}),
each running in the time of \(\mathcal{A}\) plus one key generation
and four scalar multiplications, with
\[
\mathrm{Adv}^{\mathrm{ik\text{-}cpa}}_{\mathrm{EG}}(\mathcal{A})
\;\le\;
\mathrm{Adv}_{\mathrm{ddh}}(\mathcal{B}_{0})
+ \mathrm{Adv}_{\mathrm{ddh}}(\mathcal{B}_{1}).
\]
Under the DDH assumption ElGamal is therefore IK-CPA: a ciphertext is
computationally independent of the public key it was made for.
\end{theorem}

\begin{proof}
The proof is a hybrid through a game that uses no key at all. Write
\(\mathsf{H}_{0}\) and \(\mathsf{H}_{1}\) for the IK-CPA game with
\(\beta\) fixed to \(0\) and to \(1\), and \(p_{0}, p_{1}\) for the
probability that \(\mathcal{A}\) outputs \(1\) in each, so that
\(\mathrm{Adv}^{\mathrm{ik\text{-}cpa}}_{\mathrm{EG}}(\mathcal{A})
= \abs{p_{0} - p_{1}}\). Let \(\mathsf{H}_{*}\) be the game whose
challenger, after the same key generation and the same message
\(M\) from \(\mathcal{A}\), answers with
\((R, C) = (U_{1}, U_{2} + M)\) for
\(U_{1}, U_{2} \xleftarrow{\$} G\) independent and uniform, and let
\(p_{*}\) be the corresponding probability. The challenge of
\(\mathsf{H}_{*}\) is computed without reference to either key, and
the triangle inequality
\(\abs{p_{0} - p_{1}} \le \abs{p_{0} - p_{*}} + \abs{p_{*} - p_{1}}\)
reduces the claim to bounding each hop by a DDH advantage.

\emph{The hop \(\mathsf{H}_{0} \to \mathsf{H}_{*}\).} The
distinguisher \(\mathcal{B}_{0}\) receives a triple
\((X, Y, Z) = ([a]P, [b]P, [c]P)\) with \(c = ab\) or \(c\) uniform.
It samples \(s \xleftarrow{\$} (\Z/q\Z)^{\times}\) in the
per-key-base variant and sets \(s = 1\) in the fixed-base variant,
plants the challenge in key~\(0\) as
\[
H_{0} = [s]P, \qquad Y_{0} = [s]X = [a]H_{0},
\]
generates \((x_{1}, (H_{1}, Y_{1}))\) honestly, and hands
\(\mathcal{A}\) the pair \(((H_{0}, Y_{0}), (H_{1}, Y_{1}))\). For
invertible \(s\) the base \([s]P\) is uniform on
\(G \setminus \{\mathcal{O}\}\), and \(Y_{0} = [a]H_{0}\) with
\(a\) uniform is an honest public key whose private key \(a\) the
distinguisher never learns---nor needs, the game offering no
decryption. On receiving \(M\) it replies with
\[
(R, C) = \bigl([s]Y,\; [s]Z + M\bigr)
\]
and outputs \(\mathcal{A}\)'s bit. If \(c = ab\), then
\(R = [sb]P = [b]H_{0}\) and
\([s]Z = [sab]P = [b]([a]H_{0}) = [b]Y_{0}\), so \((R, C)\) is an
honest encryption under key~\(0\) with ephemeral \(r = b\) uniform:
the view of \(\mathcal{A}\) is exactly \(\mathsf{H}_{0}\). If \(c\)
is uniform, then \([s]Z = [sc]P\) is uniform in \(G\) and
independent of everything else because \(s\) is invertible, and
\(R = [b]H_{0}\) is uniform in \(G\) and independent of everything
else because \(H_{0}\) generates \(G\) and \(b\) is independent of
\(a, c, s\), of key~\(1\), and of \(M\), which \(\mathcal{A}\) chose
before seeing \(R\): the view is exactly \(\mathsf{H}_{*}\). Hence
\(\abs{p_{0} - p_{*}} = \mathrm{Adv}_{\mathrm{ddh}}(\mathcal{B}_{0})\).

\emph{The hop \(\mathsf{H}_{*} \to \mathsf{H}_{1}\)} is the same
with the challenge planted in key~\(1\) and key~\(0\) generated
honestly, giving \(\mathcal{B}_{1}\) with
\(\abs{p_{*} - p_{1}} = \mathrm{Adv}_{\mathrm{ddh}}(\mathcal{B}_{1})\).
Each distinguisher performs one honest key generation and the four
scalar multiplications \([s]P, [s]X, [s]Y, [s]Z\) beyond running
\(\mathcal{A}\). Summing the two hops gives the bound.
\end{proof}

\begin{remark}[What key privacy buys upstairs]
\label{keyagreement:rem:keyprivacy-deployed}
The hybrid transfers verbatim to
Construction~\ref{keyagreement:def:hpke}: planting the challenge in
the recipient's key and the ephemeral share exactly as above makes
the shared secret \(S\) uniform, and once it is, the pair
\((E_{\mathrm{pk}}, c)\) with
\(c = \mathrm{Enc}_{\mathrm{KDF}(S, \mathit{context})}
(N, \mathit{ad}, m)\) is computed without reference to the
recipient's key---\emph{provided} the KDF context names no recipient
key, which it must not. The deployed KDF takes
\((\mathsf{sharedSecret}, \mathsf{epk})\)
(\S\ref{symmetric:subsec:kdf}): the shared secret is the keying
material, the context is \(\mathsf{epk}\) under a personalisation
label, and no key of the recipient enters. The protocol
specification requires the composed note-encryption scheme to be
key-private (\href{https://zips.z.cash/protocol/protocol.pdf\#abstractkdf}{\S``Key Derivation''}): a chain scanner holding a note
ciphertext against every address she can enumerate learns nothing about which
one it was sent to. The second use is the group-element share of
diversified-address unlinkability. A diversified address
\((G_{d}, \mathsf{pk_d}) = (G_{d}, [\mathsf{ivk}]\,G_{d})\)
(Remark~\ref{keyagreement:rem:notes}) is a per-key-base ElGamal
public key, and a fresh address of the same \(\mathsf{ivk}\) over a
new base \(G_{d'}\) has, when that base is modelled as a uniform
group element, the distribution of an encryption of \(\mathcal{O}\)
under that key; telling which of two authorities a third address
belongs to is then the IK-CPA game, and
Theorem~\ref{keyagreement:thm:ikcpa} answers it. The protocol
specification sets up that correspondence, modelling the group hash
behind the base as a random oracle
(\href{https://zips.z.cash/protocol/protocol.pdf\#concretediversifyhash}{\S``\(\mathsf{DiversifyHash}^{\mathsf{Sapling}}\) and
\(\mathsf{DiversifyHash}^{\mathsf{Orchard}}\) Hash Functions''}); the
\emph{Wallet Guide} states the resulting unlinkability requirement
(\S``Diversified addresses'').
\end{remark}

\subsection{Shielded-note delivery as deployed hybrid encryption}
\label{keyagreement:subsec:notes}

The chain assembled in this section is not an analogy for what the
deployed protocol does; it is the deployed protocol, run in band on
the public ledger.

\begin{remark}[Note delivery is HPKE on the ledger]
\label{keyagreement:rem:notes}
A shielded address embeds a static public key
\(\mathsf{pk_d} = [\mathsf{ivk}]\,G_{d}\) on the protocol's curve
(Pallas, in Orchard), where the private scalar \(\mathsf{ivk}\) is
the recipient's incoming viewing key. The base \(G_{d}\) is not a
protocol-wide generator but a per-address \emph{diversified base},
the hash-to-curve image (\S\ref{hash:subsec:hashtocurve}) of the
address's diversifier---the generalised interface of
Remark~\ref{keyagreement:rem:deployed-abstraction} exists precisely
to accommodate it (\texttt{DiversifiedTransmissionKey} is derived as
\(\mathsf{pk_d} = [\mathsf{ivk}]\,g_{d}\)); the \emph{Ironwood
Guide} develops the key hierarchy (\S``Viewing keys'') and address
encoding (\S``Diversified addresses'') behind these bases. The
sender runs Construction~\ref{keyagreement:def:hpke}
over that base, with one deployed refinement: the ephemeral scalar
\(\mathsf{esk}\) is not sampled directly but derived by
\(\mathsf{PRF}^{\mathsf{expand}}\) (Remark~\ref{prf:rem:blake2})
from the note's random seed (ZIP~212; \texttt{RandomSeed::esk\_inner} applies
\texttt{PrfExpand::ORCHARD\_ESK}), so that the uniform-ephemeral
hypothesis of Theorem~\ref{keyagreement:thm:hpke-secure} is met
through the PRF's pseudorandomness and the recipient can later
re-derive \(\mathsf{esk}\) from the delivered seed. The sender
publishes \(\mathsf{epk} = [\mathsf{esk}]\,G_{d}\) in the
transaction, derives
\[
K = \mathrm{KDF}\bigl([\mathsf{esk} \cdot \mathsf{ivk}]\,G_{d},\;
\mathsf{epk}\bigr)
\]
by the personalised BLAKE2b call of \S\ref{symmetric:subsec:kdf},
and AEAD-encrypts the note plaintext---a version byte, the
recipient's diversifier, the value, the note's random seed, and the
memo---under \(K\) with \textsf{ChaCha20-Poly1305}
(Construction~\ref{symmetric:def-chachapoly}), attaching the
ciphertext to the transaction. The nonce is fixed to zero and the
associated data is empty, safe because each \(K\) keys exactly one
encryption (protocol specification
\href{https://zips.z.cash/protocol/protocol.pdf\#saplingandorchardencrypt}{\S``Encryption (Sapling and
Orchard)''}; the \texttt{zcash\_note\_encryption} crate's
\texttt{encrypt\_note\_plaintext}, which chains
\texttt{ka\_agree\_enc}, \texttt{kdf}, and
\texttt{ChaCha20Poly1305} with the zero nonce; the Orchard
instantiation of the \texttt{Domain} trait is in the
\texttt{orchard} crate).

Chain scanners see \((\mathsf{epk}, c)\) in every shielded output,
but lacking \(\mathsf{ivk}\) they learn nothing
(Theorem~\ref{keyagreement:thm:hpke-secure}); the recipient
\emph{trial-decrypts}, deriving
\(K' = \mathrm{KDF}([\mathsf{ivk}]\,\mathsf{epk}, \mathsf{epk})\)
for each incoming transaction and keeping those outputs whose AEAD
tag verifies \emph{and} whose plaintext passes two consistency
checks: the note commitment recomputed from it must match the
commitment on chain, and the \(\mathsf{esk}\) re-derived from its
random seed must reproduce the published \(\mathsf{epk}\)---the
notes addressed to them (protocol specification
\href{https://zips.z.cash/protocol/protocol.pdf\#decryptivk}{\S``Decryption using
an Incoming Viewing Key (Sapling and Orchard)''}).
Non-interactivity is essential, and
it is exactly the static--ephemeral flow of
\S\ref{keyagreement:subsec:staticephemeral}: the recipient need
never have been online when the note was sent. Authenticity of the
static key---the concern of \S\ref{keyagreement:subsec:active}---is
supplied by mechanism~(2) of
Remark~\ref{keyagreement:rem:auth-supplied}: the address is part of
the recipient's verified payment credential, so the
man-in-the-middle substitution has nowhere to stand; and, as noted
there, this authenticates the recipient rather than the sender. The
complementary guarantees on the note's \emph{contents}---that the
value is well-formed and conserved, that the commitment on chain
matches the plaintext delivered---come not from the encryption but
from the accompanying zero-knowledge proof.
\end{remark}

The primitive toolbox now covers the delivery of a secret to a
party who was never online to negotiate it. What travels alongside
that ciphertext---the authorising signatures over the transaction
that carries it---is the business of the next section.

\section{Digital signatures}\label{sec:signatures}

The adversary of this section is the \emph{forger}. It watches the
network accept spend after spend, each authorised by a signature; it
may coax the legitimate signer into signing any messages of its
choosing; and it wins if it can then exhibit one signature the signer
never made---a spend of coins it does not control, an authorisation
conjured from public data alone. Every verifier in a consensus
network checks every signature, so the forger's target is public and
its victory is total: a single forged spend authorisation steals, and
a single forged balance certificate mints value from nothing. The
primitive that defeats it, the digital signature, is the last
cryptographic object of this volume built directly on the discrete
logarithm; it will also turn out, in its Schnorr form, to be the
first \emph{proof of knowledge} we construct---the bridge to the
proof systems of Section~\ref{sec:zk}.

\subsection{Signatures versus MACs}\label{signatures:subsec:vsmac}

The message authentication codes of
\S\ref{symmetric:subsec:mac} already defeat a forger, but only
between two parties who share a secret key: verification runs
\(\mathrm{Vrfy}_{k}\), so whoever can \emph{check} a tag can also
\emph{mint} one. Authentication by MAC is therefore private and
repudiable---a verifier cannot exhibit the tag to a third party as
evidence, since the verifier could have produced it. A blockchain
inverts every one of these constraints. A transaction's
authorisation must be checked by every full node, none of whom the
signer has met and none of whom may be capable of forging it; and
the check must be transferable, so that a block, once validated,
convinces everyone else.

A \emph{digital signature} meets these demands by splitting the key.
Signing uses a secret key \(sk\); verification uses a matching
\emph{public} key \(pk\) that may be published to the world.
Anybody can verify; only the holder of \(sk\) can sign; and because
verification needs no secret, a valid signature is
\emph{non-repudiable} evidence that the holder of \(sk\) signed.
The price is computational: where the Carter--Wegman MAC of
\S\ref{symmetric:subsec:mac} was information-theoretically secure,
every signature scheme in this section rests on the hardness of the
discrete logarithm.

Zcash's Orchard protocol uses two closely related signature schemes,
both Schnorr-type schemes over an elliptic curve, and both derived
from the one identification protocol this section develops
(protocol specification \href{https://zips.z.cash/protocol/protocol.pdf\#concretereddsa}{\S~5.4.7}):
\begin{itemize}
\item the \emph{spend authorisation signature}, whose key pair can be
  \emph{re-randomised} so that successive uses are unlinkable
  (protocol specification \href{https://zips.z.cash/protocol/protocol.pdf\#spendauthsig}{\S~4.15}), and
\item the \emph{binding signature}, which simultaneously certifies
  the consistency of a transaction's value commitments and proves
  knowledge of a discrete logarithm (protocol specification \href{https://zips.z.cash/protocol/protocol.pdf\#orchardbalance}{\S~4.14}).
\end{itemize}

\paragraph{Conventions for this section.}
Throughout, \(\G\) is a cyclic group of prime order \(q\), written
additively, with a fixed generator \(G\); scalar multiplication is
\([x]P\) for \(x \in \Fq\) and \(P \in \G\). Concretely \(\G\) is the
Pallas group of \S\ref{assumptions:subsec:pasta}, whose order is the
\(255\)-bit prime \(q\), so scalars live in the field \(\Fq\).
Because \(\G\) has prime order and \(G\) generates it, the map
\(x \mapsto [x]G\) is a bijection \(\Fq \to \G\): the discrete
logarithm of any \(P \in \G\) to base \(G\) always \emph{exists} and
is \emph{unique}. Hardness of computing it is the DLog assumption of
\S\ref{assumptions:subsec:dlp}, transcribed additively.

\subsection{Syntax and security goal}\label{signatures:subsec:syntax}

\begin{definition}[Digital signature scheme]\label{signatures:def:scheme}
A \emph{digital signature scheme} is a triple of PPT algorithms
\(\Pi = (\mathrm{KeyGen}, \mathrm{Sign}, \mathrm{Verify})\) with, for
each security parameter \(\lambda\), a message space
\(\mathcal M_\lambda\):
\begin{itemize}
\item \(\mathrm{KeyGen}(1^{\lambda})\) outputs a key pair
  \((pk, sk)\);
\item \(\mathrm{Sign}(sk, m)\), for \(m \in \mathcal M_\lambda\),
  outputs a signature \(\sigma\);
\item \(\mathrm{Verify}(pk, m, \sigma)\) is deterministic and
  outputs a bit.
\end{itemize}
\emph{Correctness} requires that for every \(\lambda\), every
\((pk, sk)\) in the support of \(\mathrm{KeyGen}(1^{\lambda})\), and
every \(m \in \mathcal M_\lambda\),
\[
\Pr\bigl[\mathrm{Verify}(pk, m, \mathrm{Sign}(sk, m)) = 1\bigr] = 1,
\]
the probability over the coins of \(\mathrm{Sign}\).
Probability-\(1\) correctness is called \emph{perfect}; a scheme
achieving only overwhelming-probability correctness tolerates a
\(\negl(\lambda)\) failure rate. Every scheme in this section is
perfectly correct.
\end{definition}

The security goal transcribes the MAC forgery game of
Definition~\ref{symmetric:def-eufcma} into the public-key setting:
the adversary now \emph{receives} the verification key, and its
oracle access models a signer who can be induced to sign arbitrary
messages.

\begin{definition}[Existential unforgeability under chosen-message
attack]\label{signatures:def:eufcma}
For a signature scheme \(\Pi\) and adversary \(\mathcal A\), the game
\(\mathsf{Forge}^{\mathcal A}_{\Pi}(\lambda)\) runs as follows.
\begin{enumerate}
\item The challenger runs
  \((pk, sk) \gets \mathrm{KeyGen}(1^{\lambda})\).
\item The adversary \(\mathcal A\) runs on input \(pk\) with oracle
  access to \(\mathrm{Sign}(sk, \cdot)\); let \(Q\) be the set of
  messages it submits to the oracle, and let
  \((m^{\star}, \sigma^{\star})\) be its output.
\item The game outputs \(1\) iff
  \(\mathrm{Verify}(pk, m^{\star}, \sigma^{\star}) = 1\) and
  \(m^{\star} \notin Q\).
\end{enumerate}
Scheme \(\Pi\) is \emph{existentially unforgeable under
chosen-message attack} (EUF-CMA) if for every PPT \(\mathcal A\),
\[
\mathrm{Adv}^{\mathrm{forge}}_{\Pi}(\mathcal A, \lambda)
:= \Pr\bigl[\mathsf{Forge}^{\mathcal A}_{\Pi}(\lambda) = 1\bigr]
\;\in\; \negl(\lambda).
\]
\end{definition}

Two features of the definition deserve emphasis. The oracle access
is \emph{adaptive} and bounded only by the adversary's polynomial
running time: it may choose each query after seeing all previous
signatures. And the forgery is \emph{existential}: the adversary
wins with a valid signature on \emph{any} message never submitted to
the oracle, however meaningless---the definition does not ask the
forged message to be useful, so security under it protects even
messages the signer would consider absurd.

\begin{definition}[Strong unforgeability]\label{signatures:def:sufcma}
The \emph{strong} variant, sUF-CMA, replaces the winning condition
\(m^{\star} \notin Q\) by: the pair \((m^{\star}, \sigma^{\star})\)
was not among the message--signature pairs returned by the oracle.
\end{definition}

Strong unforgeability forbids producing even a \emph{new} signature
on an \emph{already-signed} message. Plain EUF-CMA permits a scheme
in which, given one valid signature, anyone can cheaply compute a
second, different valid signature on the same message; such
\emph{malleability} is harmless when a signature only certifies
origin, but fatal whenever a signature doubles as a unique
token---certain anti-replay mechanisms, for instance, identify a
message by its signature and break if signatures are malleable.

\subsection{The hash-and-sign paradigm}
\label{signatures:subsec:hashandsign}

The textbook mechanisms---RSA, Schnorr, ECDSA---sign only
fixed-size inputs: an element of \(\Z/N\Z\) for RSA, a scalar in
\(\Fq\) for Schnorr and ECDSA. Arbitrary messages are accommodated
by hashing first, and the composition is provably safe.

\begin{theorem}[Hash-and-sign]\label{signatures:thm:hashandsign}
Let \(\Pi'\) be a signature scheme with message space
\(\mathcal Y_\lambda\), and let
\(H \colon \{0,1\}^{*} \to \mathcal Y_\lambda\) be drawn from a
collision-resistant family (Definition~\ref{hash:def:col}). Define
\(\Pi\) with message space \(\{0,1\}^{*}\) by
\[
\mathrm{Sign}(sk, m) := \mathrm{Sign}'(sk, H(m)),
\qquad
\mathrm{Verify}(pk, m, \sigma) := \mathrm{Verify}'(pk, H(m), \sigma).
\]
If \(\Pi'\) is EUF-CMA and \(H\) is collision resistant, then
\(\Pi\) is EUF-CMA; concretely, for every adversary \(\mathcal A\)
against \(\Pi\) there are comparably efficient \(\mathcal B_1\)
against \(H\) and \(\mathcal B_2\) against \(\Pi'\) with
\[
\Pr[\mathcal A \text{ wins}]
\;\le\; \mathrm{Adv}^{\mathrm{col}}_{H}(\mathcal B_1)
      + \mathrm{Adv}^{\mathrm{forge}}_{\Pi'}(\mathcal B_2).
\]
\end{theorem}

\begin{proof}
Let \((m^{\star}, \sigma^{\star})\) be a valid forgery of
\(\mathcal A\) on a fresh message \(m^{\star} \notin Q\), and split
on the digest.
\emph{Case 1: the digest is reused}, i.e.\ \(H(m^{\star}) = H(m)\)
for some queried \(m \in Q\). Then \(m^{\star} \ne m\) but their
digests agree, so the pair \((m^{\star}, m)\) is a collision in
\(H\); the reduction \(\mathcal B_1\), which generates a key pair
with \(\mathrm{KeyGen}'\) itself and answers \(\mathcal A\)'s
signing queries by signing digests under that secret key---the
collision game grants no signing oracle---outputs it.
\emph{Case 2: the digest is fresh}, i.e.\ \(H(m^{\star})\) differs
from the digest of every queried message. Then
\(\bigl(H(m^{\star}), \sigma^{\star}\bigr)\) is a valid forgery
against \(\Pi'\) on a message never submitted to \emph{its} signing
oracle, and \(\mathcal B_2\) outputs it. The two cases are
exhaustive, so the union bound on the winning event gives the stated
inequality.
\end{proof}

\begin{remark}[What ``hash-and-sign'' hides]
\label{signatures:rem:hash-more}
In all the Schnorr-type schemes below, the hash absorbs more than
the message: the public key and the per-signature commitment value
enter the hash alongside \(m\). Hashing the public key prevents
key-substitution (duplicate-signature) attacks, in which a signature
valid under one key is exhibited as valid under another; hashing the
commitment is exactly what the Fiat--Shamir transform of
\S\ref{signatures:subsec:fiatshamir} requires. ``Hash-and-sign'' in
practice therefore means hashing more than the message;
Theorem~\ref{signatures:thm:hashandsign} captures only the
message-compression part of the discipline.
\end{remark}

\subsection{The Schnorr identification protocol}
\label{signatures:subsec:schnorrid}

Every scheme in the rest of this section is a costume worn by a
single interactive protocol, in which a prover convinces a verifier
that it \emph{knows} the discrete logarithm of a public point---
without revealing anything about it. We develop the protocol
self-contained, with concrete statements and proofs; the general
vocabulary it exemplifies (\emph{Sigma-protocols}, of which it is
the canonical instance: three-move, public-coin, honest-verifier
zero-knowledge (HVZK) proofs of knowledge) is deliberately deferred to
Section~\ref{sec:zk}, where it is developed for arbitrary relations.

The object of the protocol is the discrete-log relation
\[
R_{\mathrm{dl}}
= \bigl\{\, (X, x) \in \G \times \Fq \;:\; X = [x]G \,\bigr\}:
\]
the \emph{statement} is a public point \(X\), the \emph{witness} its
discrete logarithm \(x\).

\begin{construction}[Schnorr identification]
\label{signatures:con:schnorrid}
Fix \(\G = \langle G \rangle\) of prime order \(q\). The prover
holds a secret key \(x \in \Fq\); both parties hold the public key
\(X = [x]G\). The protocol has three moves.
\begin{enumerate}
\item \emph{Commitment.} The prover samples
  \(r \xleftarrow{\$} \Fq\) uniformly and sends \(R = [r]G\). We
  call \(R\) the \emph{commitment} (or \emph{nonce point}) and \(r\)
  the \emph{nonce}.
\item \emph{Challenge.} The verifier samples
  \(c \xleftarrow{\$} \Fq\) uniformly and sends \(c\).
\item \emph{Response.} The prover sends \(s = r + cx\) in \(\Fq\).
\end{enumerate}
The verifier accepts iff
\begin{equation}\label{eq:signatures-verify}
[s]G \;=\; R + [c]X .
\end{equation}
A triple \((R, c, s)\) satisfying \eqref{eq:signatures-verify} is an
\emph{accepting transcript} for \(X\). The verifier's only message
is a uniformly random value chosen independently of \(R\), so the
protocol is \emph{public-coin}: the verifier keeps no secrets and
exercises no discretion.
\end{construction}

Three properties make this protocol what it is: honest runs always
accept; answering two distinct challenges on one commitment is
possible only for a party who knows \(x\); and a transcript reveals
nothing about \(x\). We prove each in turn.

\begin{proposition}[Perfect completeness]
\label{signatures:prop:completeness}
If prover and verifier are honest, the verifier accepts with
probability \(1\).
\end{proposition}

\begin{proof}
With \(R = [r]G\), \(X = [x]G\) and \(s = r + cx\), linearity of
scalar multiplication gives
\[
[s]G = [r + cx]G = [r]G + [cx]G = R + [c]\bigl([x]G\bigr)
     = R + [c]X,
\]
which is exactly \eqref{eq:signatures-verify}. The computation holds
identically for every choice of \(r\) and \(c\).
\end{proof}

\begin{theorem}[2-special soundness]
\label{signatures:thm:specialsound}
There is an efficient \emph{extractor} that, given two accepting
transcripts \((R, c, s)\) and \((R, c', s')\) for the same statement
\(X\) with the same commitment \(R\) but distinct challenges
\(c \ne c'\), outputs the discrete logarithm \(x\) of \(X\).
\end{theorem}

\begin{proof}
Both transcripts satisfy \eqref{eq:signatures-verify}; subtracting
the two equations in \(\G\) cancels \(R\):
\[
[s - s']G \;=\; [c - c']X \;=\; [(c - c')x]G .
\]
Since \(G\) generates a group of prime order \(q\), equality of the
two multiples forces the equality of scalars
\(s - s' = (c - c')x\) in \(\Fq\). The hypothesis \(c \ne c'\) makes
\(c - c'\) a nonzero element of the field \(\Fq\), hence invertible,
so the extractor outputs
\[
x \;=\; (s - s')\,(c - c')^{-1} \in \Fq .
\]
Verification that \(X = [x]G\) is direct, and the whole computation
costs a single field inversion.
\end{proof}

Theorem~\ref{signatures:thm:specialsound} is the central mechanism
of this section. It makes the protocol a proof \emph{of knowledge}:
below, it converts a convincing prover into an extractable witness;
after the Fiat--Shamir transform it is the source of unforgeability;
and in \S\ref{signatures:subsec:romforking} it reappears with its
sign reversed, as the reason a signer must never reuse a nonce.

\begin{remark}[From special soundness to knowledge, by rewinding]
\label{signatures:rem:rewinding}
Suppose a (possibly dishonest) prover makes the verifier accept with
probability \(\varepsilon\) noticeably larger than \(1/q\)---the
probability of simply guessing the single challenge for which a
prepared response works. Run the prover once: record its commitment
\(R\), feed it a uniform challenge \(c\), and note whether its
response accepts. Now \emph{rewind} it to the moment just after it
sent \(R\)---restore its internal state, which is possible because
we run the prover as a subroutine---and feed a fresh independent
challenge \(c'\). A standard averaging argument shows that with
probability roughly \(\varepsilon^{2}\), both runs accept and
\(c \ne c'\); the argument is the \emph{heavy-row lemma}, an
instance of Markov's inequality applied to the matrix whose rows are
the prover's random coins and whose columns are challenges
(\emph{Math Guide}, \S``Random variables and expectation''). The
two accepting transcripts share \(R\), so the extractor of
Theorem~\ref{signatures:thm:specialsound} recovers \(x\). Making the
verifier accept with non-trivial probability is therefore, up to
rewinding, the same as \emph{knowing} \(x\). The same rewinding
argument, sharpened into the forking lemma, carries the signature
analysis of \S\ref{signatures:subsec:romforking}.
\end{remark}

\begin{theorem}[Special honest-verifier zero knowledge]
\label{signatures:thm:hvzk}
There is an efficient simulator \(\mathrm{Sim}\) that, on input a
statement \(X \in \G\) and a challenge \(c \in \Fq\), outputs an
accepting transcript \((R, c, s)\) distributed \emph{identically} to
an honest execution conditioned on the challenge being \(c\).
Consequently, for a uniform challenge, the distribution of simulated
transcripts coincides exactly with the distribution of real ones.
\end{theorem}

\begin{proof}
The simulator works \emph{backwards}: it samples
\(s \xleftarrow{\$} \Fq\) uniformly, sets
\[
R := [s]G - [c]X,
\]
and outputs \((R, c, s)\)---accepting by construction, since
\eqref{eq:signatures-verify} rearranges to exactly this definition
of \(R\).

Compare distributions with the challenge fixed to \(c\). In a real
honest transcript, \(r\) is uniform on \(\Fq\), \(R = [r]G\), and
\(s = r + cx\); the map \(r \mapsto r + cx\) is a translation of
\(\Fq\), hence a bijection, so \(s\) is uniform on \(\Fq\), and
\(R = [s - cx]G = [s]G - [c]X\) is the same deterministic function
of \(s\) that the simulator uses. The real transcript \emph{is}:
sample \(s\) uniform, set \(R = [s]G - [c]X\)---precisely the
simulated distribution. The two distributions are identical, not
merely statistically close, and averaging over a uniform \(c\)
preserves the identity.
\end{proof}

\begin{remark}[Why honest-verifier zero knowledge suffices]
\label{signatures:rem:hvzk-scope}
The simulator of Theorem~\ref{signatures:thm:hvzk} matches the real
distribution only when the challenge is sampled independently of
\(R\). A cheating verifier could instead choose \(c\) as a function
of \(R\), and against such a verifier the simulator---which commits
to \(R\) only \emph{after} seeing \(c\)---no longer obviously works.
For our purposes honest-verifier zero knowledge is exactly enough:
after the Fiat--Shamir transform the challenge is a \emph{hash} of
\(R\), and in the random oracle model that hash behaves like an
honest, \(R\)-independent coin. The deeper point of the theorem is
what the simulator's existence \emph{proves}: since anyone, without
\(x\), can manufacture transcripts with the true distribution, a
transcript carries no information about \(x\) beyond what the
statement \(X\) already reveals.
\end{remark}

\subsection{From identification to signature via the Fiat--Shamir
transform}\label{signatures:subsec:fiatshamir}

An identification protocol convinces one verifier, once,
interactively. A signature must convince everyone, forever, on its
own. The \emph{Fiat--Shamir transform} removes the interaction by
making the prover compute the challenge itself---as a hash of
everything the verifier would have seen before issuing it, with the
message folded in so that the resulting object certifies \(m\). The
transform is generic: it converts \emph{any} Sigma-protocol into a
non-interactive proof, and, with a message in the hash, into a
signature. This subsection owns only the Schnorr instance; the
general transform, its subtleties and its failure modes are the
business of Section~\ref{sec:zk} (\S``The Fiat--Shamir transform:
from interactive to non-interactive'').

\begin{construction}[Schnorr signature]
\label{signatures:con:schnorrsig}
Fix \(\G = \langle G \rangle\) of prime order \(q\) and a hash
function \(H \colon \{0,1\}^{*} \to \Fq\). Messages are arbitrary
bit strings.
\begin{itemize}
\item \(\mathrm{KeyGen}(1^{\lambda})\): sample
  \(x \xleftarrow{\$} \Fq\), set \(X = [x]G\), output
  \((pk, sk) = (X, x)\).
\item \(\mathrm{Sign}(x, m)\): sample \(r \xleftarrow{\$} \Fq\), set
  \(R = [r]G\), compute the challenge
  \[
  c = H(R \,\|\, X \,\|\, m),
  \]
  set \(s = r + cx\) in \(\Fq\), and output \(\sigma = (R, s)\).
\item \(\mathrm{Verify}(X, m, (R, s))\): recompute
  \(c = H(R \,\|\, X \,\|\, m)\) and accept iff
  \([s]G = R + [c]X\).
\end{itemize}
\end{construction}

\begin{proposition}[Perfect correctness]
\label{signatures:prop:sig-correct}
The Schnorr signature scheme is perfectly correct.
\end{proposition}

\begin{proof}
For an honest signature \((R, s)\) with \(R = [r]G\),
\(c = H(R \,\|\, X \,\|\, m)\) and \(s = r + cx\), the verifier
recomputes the same \(c\)---the challenge is a deterministic
function of \((R, X, m)\), all of which the verifier holds---and the
check \([s]G = R + [c]X\) holds by the completeness computation of
Proposition~\ref{signatures:prop:completeness}.
\end{proof}

\begin{remark}[Two conventions]\label{signatures:rem:sig-conventions}
Two presentational choices in
Construction~\ref{signatures:con:schnorrsig} deserve note.
\begin{enumerate}
\item The signature transmits \((R, s)\) and omits \(c\), which the
  verifier recomputes. The space-saving variant transmits
  \((c, s)\) instead: the verifier recovers \(R = [s]G - [c]X\)
  (the simulator's equation, read as reconstruction) and checks
  \(H(R \,\|\, X \,\|\, m) = c\). The two are interchangeable; the
  deployed schemes below use the \((R, s)\) form.
\item The hash takes the public key \(X\) alongside \(R\) and
  \(m\)---\emph{key-prefixing}. The single-key unforgeability proof
  of \S\ref{signatures:subsec:romforking} does not require it, but
  it forecloses attacks that relate signatures under different keys
  (Remark~\ref{signatures:rem:hash-more}), it is standard in modern
  instantiations such as EdDSA, and it earns its keep in the
  re-randomised setting of \S\ref{signatures:subsec:rerand}.
\end{enumerate}
\end{remark}

\subsection{Security in the random oracle model and the forking
lemma}\label{signatures:subsec:romforking}

What must a forger of Schnorr signatures accomplish? A valid pair
\((R^{\star}, s^{\star})\) on a fresh message is an accepting
transcript of the identification protocol whose challenge is
\(H(R^{\star} \| X \| m^{\star})\)---a value the forger cannot
choose freely if the hash is honest. The security argument makes
``honest'' precise by idealising \(H\), then turns the rewinding of
Remark~\ref{signatures:rem:rewinding} into a quantitative tool.

\begin{definition}[Random oracle, general range]
\label{signatures:def:rom-range}
The random oracle of Definition~\ref{hash:def:ro} generalises
verbatim to an arbitrary finite range: a random oracle with range
\(Y\) is a function \(\mathcal O \colon \{0,1\}^{*} \to Y\) drawn
uniformly from the set of all such functions, realised lazily by
sampling each new answer \(\mathcal O(z) \xleftarrow{\$} Y\) on
first query. For Schnorr signatures the range is the challenge
space \(Y = \Fq\), and the ROM for signatures models
\(H(\cdot) = \mathcal O(\cdot)\), with every party---signer,
verifier, forger, and reduction---sharing oracle access
(Definition~\ref{model:def:rom}).
\end{definition}

The model remains a heuristic, with the uninstantiability caveat and
the working stance of Remark~\ref{model:rem:rom-caveat}: a ROM proof
rules out all attacks
that treat the hash as a black box---the overwhelming majority---and
it is the standard yardstick for Fiat--Shamir signatures. The
reductions below exploit both powers of
Proposition~\ref{hash:prop:rom-powers}: \emph{programmability} (the
reduction chooses oracle answers on the fly, provided they appear
uniform) and \emph{observability} (the reduction sees every query
the adversary makes).

The first tool disposes of the signing oracle: in the ROM, the
zero-knowledge simulator of Theorem~\ref{signatures:thm:hvzk}
answers signing queries without the secret key.

\begin{lemma}[Signature simulation]\label{signatures:lem:simulation}
In the ROM there is an efficient simulator that, given only the
public key \(X\) and the ability to program \(H\), answers signing
queries with signatures distributed identically to genuine ones,
except that it aborts with probability at most
\(q_s(q_s + q_h)/q\) over an attack making \(q_s\) signing and
\(q_h\) hash queries.
\end{lemma}

\begin{proof}
To sign \(m\) without \(x\): sample \(c \xleftarrow{\$} \Fq\) and
\(s \xleftarrow{\$} \Fq\) uniformly, set \(R := [s]G - [c]X\)---the
HVZK simulator of Theorem~\ref{signatures:thm:hvzk}---then
\emph{program}
\[
H(R \,\|\, X \,\|\, m) := c,
\]
and output \((R, s)\). The signature verifies by construction, and
by Theorem~\ref{signatures:thm:hvzk} its distribution matches a real
signature's exactly: in both cases \(s\) is uniform, and \(c\) is
uniform because a fresh oracle answer is.

Programming fails only if the entry \(H(R \,\|\, X \,\|\, m)\) is
already defined---by an earlier adversary hash query or an earlier
simulated signature. The point \(R = [s]G - [c]X\) is a fixed
translate of \([s]G\) for uniform \(s\), hence uniform over the
\(q\) elements of \(\G\); at most \(q_h + q_s\) oracle entries exist
at any point of the attack, so a given signing query collides with
probability at most \((q_h + q_s)/q\), and the union bound over
\(q_s\) signing queries gives the stated abort probability.
Conditioned on no abort, the adversary's view is identical to the
real game.
\end{proof}

Lemma~\ref{signatures:lem:simulation} reduces forgery to a pure
discrete-log problem: the reduction can play the whole EUF-CMA game
knowing only \(X\). What remains is to convert one forgery into
\emph{two} accepting transcripts on the same commitment---the input
that Theorem~\ref{signatures:thm:specialsound} demands. The
quantitative engine is the forking lemma, in the game-based form due
to Bellare and Neven.

% FIGURE (planned, dedicated figures pass): fig:forking --- rewinding
% a forger to two accepting transcripts sharing a prefix. Two
% horizontal runs of the same algorithm B with the same coins rho,
% drawn as parallel timelines; identical prefix c_1..c_{I-1} shaded as
% a shared band; the fork at index I where run 1 receives c_I and run
% 2 receives a fresh c'_I; both runs ending at the same critical query
% R*||X||m* with different challenge answers, yielding transcripts
% (R*, c*, s*) and (R*, c'*, s'*) feeding the special-soundness
% extractor box that emits x.

\begin{theorem}[General forking lemma]\label{signatures:thm:forking}
Fix an integer \(Q \ge 1\), a finite set \(C\) with
\(\abs{C} = h \ge 2\), and a randomised input generator
\(\mathrm{Gen}\). Let \(B\) be a randomised algorithm that on input
\(y\) and challenges \(c_1, \dots, c_Q \in C\) returns a pair
\((I, \mathit{out})\) with \(I \in \{0, 1, \dots, Q\}\), where
\(I = 0\) means failure, and let
\[
\mathit{acc}
:= \Pr\bigl[I \ge 1 :
   y \gets \mathrm{Gen};\;
   c_1, \dots, c_Q \xleftarrow{\$} C;\;
   (I, \mathit{out}) \gets B(y, c_1, \dots, c_Q)\bigr].
\]
Define the forking algorithm \(F_B\) on input \(y\): pick coins
\(\rho\) for \(B\) and \(c_1, \dots, c_Q \xleftarrow{\$} C\); run
\((I, \mathit{out}) \gets B(y, c_1, \dots, c_Q; \rho)\); if
\(I = 0\), return \(\bot\); otherwise resample
\(c'_I, \dots, c'_Q \xleftarrow{\$} C\) afresh and run
\((I', \mathit{out}') \gets
  B(y, c_1, \dots, c_{I-1}, c'_I, \dots, c'_Q; \rho)\)
with the \emph{same} coins \(\rho\); if \(I' = I\) and
\(c_I \ne c'_I\), return \((I, \mathit{out}, \mathit{out}')\), else
\(\bot\). Then
\[
\mathit{frk}
:= \Pr\bigl[F_B(y) \ne \bot : y \gets \mathrm{Gen}\bigr]
\;\ge\; \mathit{acc}\left(\frac{\mathit{acc}}{Q}
        - \frac{1}{h}\right),
\]
the probability also over the coins of \(F_B\).
\end{theorem}

\begin{proof}[Proof idea]
Condition on the coins \(\rho\) and the prefix
\(c_1, \dots, c_{I-1}\) up to the forking point. For each fixed
prefix there is (averaging over the suffix) a set \(S\) of
\emph{good} values of \(c_I\) on which \(B\) succeeds with index
\(I\); write \(p = \abs{S}/h\). The fork succeeds at this prefix iff
the first draw \(c_I\) and the independent re-draw \(c'_I\) both
land in \(S\) and differ: two independent uniform draws land in
\(S\) together with probability \(p^{2}\), and subtracting the
diagonal of equal draws costs at most \(p \cdot 1/h\), so the
per-prefix fork probability is at least \(p(p - 1/h)\). Taking
expectations over prefixes,
\(\mathbb E[p^{2}] \ge (\mathbb E[p])^{2}\), since
\(\mathbb E[p^{2}] - (\mathbb E[p])^{2}
= \mathbb E[(p - \mathbb E[p])^{2}] \ge 0\) by linearity of
expectation (\emph{Math Guide}, \S``Random variables and
expectation''); this squares the average success probability. The
index \(I\) ranging over \(Q\) possible positions introduces the
factor \(1/Q\), and collecting terms yields
\(\mathit{frk} \ge \mathit{acc}^{2}/Q - \mathit{acc}/h\). The
squaring step is where the \emph{squared} probability---and hence
the tightness loss examined below---enters.
\end{proof}

\begin{theorem}[EUF-CMA security of Schnorr signatures in the ROM]
\label{signatures:thm:schnorr-euf}
Model \(H\) as a random oracle with range \(\Fq\). If an adversary
\(\mathcal A\) runs in time \(t\), makes at most \(q_h\) hash
queries and \(q_s\) signing queries, and forges with probability
\(\varepsilon\), then there is an algorithm \(\mathcal D\) running
in time approximately \(2t\) that computes discrete logarithms in
\(\G\) with probability
\[
\varepsilon' \;\ge\;
\frac{\varepsilon^{2}}{q_h + q_s + 1}
\;-\; \frac{q_h + q_s + 1}{q}
\;-\; \frac{q_s (q_h + q_s)}{q}.
\]
In particular, if DLog is hard in \(\G\), the Schnorr signature
scheme is EUF-CMA in the ROM.
\end{theorem}

\begin{proof}[Proof structure]
Wrap \(\mathcal A\) as the algorithm \(B\) of
Theorem~\ref{signatures:thm:forking}, with input \(y = X\)---the
DLog challenge, so \(\mathrm{Gen}\) outputs \(X = [x]G\) for uniform
\(x\)---challenge set \(C = \Fq\) (so \(h = q\)), and
\(Q = q_h + q_s + 1\) prepared challenges. Algorithm \(B\) runs
\(\mathcal A\)
on \(pk = X\), answering the \(j\)-th \emph{distinct} hash query
with \(c_j\) from its challenge list and answering signing queries
by the simulator of Lemma~\ref{signatures:lem:simulation}, whose
programming aborts cost the \(q_s(q_h + q_s)/q\) term.

Suppose \(\mathcal A\) outputs a valid fresh forgery
\((m^{\star}, (R^{\star}, s^{\star}))\), and let
\(c^{\star} = H(R^{\star} \| X \| m^{\star})\). The critical hash
input \(R^{\star} \| X \| m^{\star}\) must actually have been
queried during the run: otherwise \(c^{\star}\) is a uniform value
\(\mathcal A\) has never seen, and the verification equation
\([s^{\star}]G = R^{\star} + [c^{\star}]X\) then holds only with
probability \(1/q\) over the oracle's choice---the
\((q_h + q_s + 1)/q\) term collects these guessing events (the
\(+1\) covering the query \(B\) itself makes to verify the forgery).
So \(B\) outputs the index \(I\) of the critical query together with
\(\mathit{out} = (R^{\star}, c^{\star}, s^{\star}, m^{\star})\), and
\(\mathit{acc} \ge \varepsilon\) minus the loss terms above.

Now fork. Both runs of \(F_B\) share the coins \(\rho\) and the
challenge prefix \(c_1, \dots, c_{I-1}\), which together determine
\(\mathcal A\)'s behaviour up to the moment it produces the critical
query; both runs therefore query the \emph{same} critical input
\(R^{\star} \| X \| m^{\star}\) at index \(I\), but receive
different answers \(c^{\star} \ne c'^{\star}\). Success of the fork
yields two accepting transcripts
\((R^{\star}, c^{\star}, s^{\star})\) and
\((R^{\star}, c'^{\star}, s'^{\star})\) with the same commitment and
distinct challenges, whereupon the extractor of
Theorem~\ref{signatures:thm:specialsound} computes
\[
x = (s^{\star} - s'^{\star})(c^{\star} - c'^{\star})^{-1},
\]
the discrete logarithm of \(X\). Substituting \(\mathit{acc}\),
\(h = q\) and \(Q = q_h + q_s + 1\) into
\(\mathit{frk} \ge \mathit{acc}^{2}/Q - \mathit{acc}/h\) gives the
stated bound; the running time is two runs of \(\mathcal A\) plus
bookkeeping.
\end{proof}

\begin{example}[The tightness loss, in numbers]
\label{signatures:exa:tightness}
The dominant term of Theorem~\ref{signatures:thm:schnorr-euf} is
\(\varepsilon' \gtrsim \varepsilon^{2}/q_h\): the reduction runs the
forger twice and \emph{squares} its success probability, divided by
the number of hash queries. The reduction is therefore far from
tight (Definition~\ref{model:def:tightness}). Concretely, a forger
succeeding with probability \(\varepsilon = 2^{-40}\) after
\(q_h = 2^{60}\) hash queries yields a guaranteed DLog-solver
advantage of only about
\[
\varepsilon^{2}/q_h \;=\; 2^{-80}/2^{60} \;=\; 2^{-140},
\]
a factor \(2^{100}\) weaker than \(\varepsilon\) itself. The
guarantee runs in the useful direction---any forger yields
\emph{some} DLog solver, so DLog hardness still implies
security---but the quantitative content is feeble. Two lessons
follow, and they must be kept apart. The conservative lesson
(Remark~\ref{model:rem:loss-matters}): a designer who wants
\emph{this theorem} to underwrite a target signature-security level
must provision \(\G\) with roughly twice as many bits of DLog
security as the target, since the reduction hands back only the
square of the forger's advantage; and the squaring is intrinsic to
the rewinding route---it enters at the squaring step of
Theorem~\ref{signatures:thm:forking}, so any analysis that runs the
forger twice and extracts by special soundness pays it. What the
theorem does \emph{not} establish is a matching attack: the loss
belongs to the reduction, not provably to the scheme, so doubled
parameters are the price of this proof, not a demonstrated
necessity, and parameter selection in practice weighs concrete
attack models and query bounds as well. Tighter reductions exist
once the rewinding route is abandoned---under the \emph{one-more}
discrete logarithm assumption (that no efficient adversary, given
\(\ell + 1\) random elements of \(\G\) and \(\ell\) queries to a
discrete-logarithm oracle, outputs all \(\ell + 1\) logarithms),
under DDH (Definition~\ref{assumptions:def:ddh}) together with
key-prefixing, or in the idealised algebraic models of
Remark~\ref{model:rem:other-models}.
\end{example}

\begin{proposition}[Nonce reuse is fatal]
\label{signatures:prop:noncereuse}
Suppose a signer produces two Schnorr signatures \((R, s)\) and
\((R, s')\) on distinct messages \(m \ne m'\) with the \emph{same}
commitment \(R\)---a nonce \(r\) reused outright, or drawn twice
from a faulty source. Then, except with probability \(1/q\) over
the challenge hash, anyone can compute the secret key from the two
signatures.
\end{proposition}

\begin{proof}
The challenges \(c = H(R \,\|\, X \,\|\, m)\) and
\(c' = H(R \,\|\, X \,\|\, m')\) differ except with probability
\(1/q\), the chance that the random oracle answers the two distinct
inputs identically. Given \(c \ne c'\), the pairs \((R, c, s)\) and
\((R, c', s')\) are two accepting transcripts sharing a commitment
with distinct challenges, and the extractor of
Theorem~\ref{signatures:thm:specialsound}---run now by the
\emph{adversary}---outputs
\(x = (s - s')(c - c')^{-1}\).
\end{proof}

Special soundness, the source of unforgeability, is thus also its
dark mirror: the security of the scheme rests entirely on the nonce
being fresh and unpredictable for every signature. This is not a
theoretical worry---it is the classic operational failure of
deployed Schnorr-family schemes---and it motivates deriving nonces
\emph{deterministically} from the secret key and message, the
defining idea of EdDSA, to which we now turn.

\subsection{RedDSA: re-randomisable Schnorr for Orchard}
\label{signatures:subsec:reddsa}

\begin{construction}[EdDSA, schematically]
\label{signatures:con:eddsa}
EdDSA retains the Schnorr signature algebra of
Construction~\ref{signatures:con:schnorrsig} unchanged and hardens
two implementation surfaces.
\begin{enumerate}
\item \emph{Deterministic nonces.} The nonce is derived as a hash of
  a secret prefix (part of the expanded secret key) and the message,
  rather than sampled fresh: no faulty random-number generator can
  repeat a nonce across distinct messages and trigger
  Proposition~\ref{signatures:prop:noncereuse}.
\item \emph{Key-prefixing.} The challenge hash absorbs the public
  key alongside \(R\) and \(m\), as in
  Construction~\ref{signatures:con:schnorrsig} and for the reasons
  of Remark~\ref{signatures:rem:sig-conventions}.
\end{enumerate}
In the ROM the deterministic nonce is, to anyone ignorant of the
secret prefix, a fresh uniform value per message, so the
unforgeability proof of Theorem~\ref{signatures:thm:schnorr-euf}
carries over essentially unchanged.
\end{construction}

Zcash's signature scheme generalises EdDSA in two directions, and
the generalisations are the whole point: the base point becomes a
parameter, and the key becomes shiftable.

\begin{construction}[RedDSA]\label{signatures:con:reddsa}
RedDSA is the Schnorr/EdDSA scheme with the base point
\(P \in \G\) a \emph{parameter} of the scheme: different uses sign
with respect to different generators.
\begin{itemize}
\item \(\mathrm{KeyGen}\): sample \(x \xleftarrow{\$} \Fq\), output
  secret key \(x\) and public key \(X = [x]P\).
\item \(\mathrm{Sign}(x, m)\): sample a fresh random byte string
  \(T\) and derive the nonce \(r = H(T \,\|\, X \,\|\, m)\). This
  is \emph{randomised}, hash-based nonce generation: unlike EdDSA's
  derivation, it is not deterministic from the signing key and
  message, so its safety rests on \(T\) being fresh and
  unpredictable, the hashing ensuring only that a repeated \(T\)
  cannot repeat a nonce across distinct messages. Set
  \(R = [r]P\), \(c = H(R \,\|\, X \,\|\, m) \in \Fq\),
  \(s = r + cx\), and output \(\sigma = (R, s)\).
\item \(\mathrm{Verify}(X, m, (R, s))\): accept iff
  \([s]P = R + [c]X\) with \(c = H(R \,\|\, X \,\|\, m)\).
\end{itemize}
For a fixed base point, RedDSA \emph{is} Schnorr/EdDSA, and its
EUF-CMA security in the ROM is
Theorem~\ref{signatures:thm:schnorr-euf} with \(G\) replaced by
\(P\).
\end{construction}

\begin{remark}[RedDSA as deployed]\label{signatures:rem:redpallas}
The deployed scheme follows Construction~\ref{signatures:con:reddsa}
line for line (protocol specification \href{https://zips.z.cash/protocol/protocol.pdf\#concretereddsa}{\S~5.4.7}). The randomiser
\(T\) is \(80\)
uniform bytes, and both hashes are the function
\(\mathsf{H}^{\circledast}\): BLAKE2b-512 personalised with
\texttt{Zcash\_RedPallasH}, its \(64\)-byte output reduced to
\(\Fq\) by wide reduction---
the expand-then-reduce discipline of
Construction~\ref{hash:con:expandreduce}, with the \(512\)-bit width
making the modular bias negligible. The nonce hash absorbs
\(T \,\|\, X \,\|\, m\) and the challenge hash
\(R \,\|\, X \,\|\, m\), key-prefixed in both cases.

Orchard instantiates RedDSA \emph{twice}, both times as
\emph{RedPallas}---RedDSA over the Pallas curve---and the two
instantiations differ \emph{only} in the base point:
\begin{itemize}
\item the \emph{spend authorisation} base point, derived by hashing
  to the curve with domain \texttt{z.cash:\allowbreak Orchard} and input
  \texttt{G}, for the re-randomisable signature of
  \S\ref{signatures:subsec:rerand}; and
\item the \emph{value-commitment randomness} base point, derived
  with domain \texttt{z.cash:Orchard-cv} and input \texttt{r}, for
  the binding signature of \S\ref{signatures:subsec:binding}. This
  is byte-for-byte the randomness base of Orchard's value commitment---
  an identity that \S\ref{signatures:subsec:binding} shows is the
  entire design.
\end{itemize}
The curve varies only across protocol generations: Sapling's
\emph{RedJubjub} is the same scheme over the Jubjub curve;
the curve, its two base
points and the \(\mathsf{H}^{\circledast}\) personalisation
(\texttt{Zcash\_RedJubjubH} against \texttt{Zcash\_RedPallasH}) are
all that distinguish the two.
\end{remark}

\subsection{Key re-randomisation and unlinkability}
\label{signatures:subsec:rerand}

A shielded protocol faces a linkage problem that no ordinary
signature scheme addresses. Spending a note requires proving spend
authority, and the obvious design---publish the spender's public key
and a signature under it---would stamp every spend by the same
wallet with the same key, linking them all. Orchard's answer
exploits the algebraic structure of the RedDSA public key: because
the key is a homomorphic image of the secret, it can be
\emph{shifted} into a fresh disguise for every spend.

\begin{definition}[Key re-randomisation]\label{signatures:def:rerand}
Fix a base point \(P\). Given a secret key \(x \in \Fq\) with public
key \(X = [x]P\), and a \emph{randomiser} \(\alpha \in \Fq\), the
re-randomised signing and verification keys are
\[
rsk := x + \alpha \in \Fq,
\qquad
rk := X + [\alpha]P \in \G .
\]
\end{definition}

Orchard publishes a fresh \(rk\) per spend, signs with the shifted
secret key \(rsk\), and proves in zero knowledge that \(rk\) is a
valid re-randomisation of an authorised public key.

\begin{proposition}[Re-randomisation consistency]
\label{signatures:prop:rerand-consistent}
The pair \((rsk, rk)\) is a valid RedDSA key pair: \(rk = [rsk]P\),
so a RedDSA signature under \(rsk\) verifies under \(rk\).
\end{proposition}

\begin{proof}
Linearity of scalar multiplication gives
\[
[rsk]P = [x + \alpha]P = [x]P + [\alpha]P = X + [\alpha]P = rk .
\]
The pair is thus exactly \(\mathrm{KeyGen}\)-shaped, and the
correctness of signing and verification
(Proposition~\ref{signatures:prop:sig-correct}) applies verbatim.
\end{proof}

The enabling structure deserves a sentence of its own. The public
key is a homomorphic image \([\cdot]P\) of the scalar, so shifting
the secret by \(\alpha\) shifts the public key by
\([\alpha]P\)---a quantity computable from \(\alpha\) and \(P\)
alone, \emph{without} the secret key. A verifier, or a circuit, can
therefore check the relationship between \(X\) and \(rk\) given
\(\alpha\), while only the spender who knows \(x\) can sign under
the shifted key.

\begin{theorem}[Perfect unlinkability of re-randomised keys]
\label{signatures:thm:unlink}
Fix a base point \(P\). For any public key \(X \in \G\), if
\(\alpha \xleftarrow{\$} \Fq\) is uniform then
\(rk = X + [\alpha]P\) is uniformly distributed over \(\G\).
Consequently, in the game where an adversary---of unbounded
computational power, and permitted to choose two public keys
\(X_0, X_1\) itself---receives \(rk = X_b + [\alpha]P\) for a
challenger's uniform bit \(b\) and fresh uniform \(\alpha\), and
must guess \(b\): the distributions of \(rk\) for \(b = 0\) and
\(b = 1\) are \emph{identical}, and no adversary guesses \(b\) with
probability better than \(1/2\).
\end{theorem}

\begin{proof}
The point \(P\) generates \(\G\), which has prime order \(q\), so
\(\alpha \mapsto [\alpha]P\) is a bijection \(\Fq \to \G\); for
uniform \(\alpha\), the point \([\alpha]P\) is uniform over \(\G\).
Translation by the fixed element \(X\) is a bijection of \(\G\) and
carries the uniform distribution to itself, so
\(rk = X + [\alpha]P\) is uniform over \(\G\)---for \emph{every}
\(X\), hence independently of \(X\). The two conditional
distributions of \(rk\) coincide exactly, so the adversary's view is
statistically independent of \(b\); no test distinguishes identical
distributions, and the best any distinguisher can do is guess,
succeeding with probability exactly \(1/2\).
\end{proof}

The guarantee is information-theoretic---perfect, resting on no
computational assumption---but its \emph{precondition} matters:
\(\alpha\) must be sampled uniformly and freshly per use, and must
be kept secret. In Orchard the transaction builder samples
\(\alpha \xleftarrow{\$} \Fq\) afresh for each spend
and reveals it only to the
proving circuit and the signer.

\begin{remark}[Scope of unlinkability]
\label{signatures:rem:unlink-scope}
Theorem~\ref{signatures:thm:unlink} covers only the re-randomised
\emph{key}. Full unlinkability of a spend further requires that the
accompanying signature and zero-knowledge proof leak nothing beyond
their statements, that the randomiser \(\alpha\) stays secret, and
that no other transaction field correlates two spends. In Orchard
the protocol publishes \(rk\) with each Action (the unit that
spends one note and creates one) while proving, inside the Halo~2
circuit, knowledge of a witness for \(rk = X + [\alpha]P\) with
\(X\) a public key whose holder is authorised to spend the note
(Orchard's \emph{spend validating key})---the circuit constrains
\(rk\) as a public input---and the network
verifies the spend authorisation signature against the public
\(rk\). Signing under \(rk\) requires \(rsk = x + \alpha\), that is,
knowledge of \emph{both} \(x\) and \(\alpha\); only the legitimate
key-holder in possession of this spend's \(\alpha\) can sign.
\end{remark}

\begin{remark}[Unforgeability under re-randomised keys, and the
key-prefixing subtlety]\label{signatures:rem:rerand-euf}
Does unforgeability survive re-randomisation? A forgery \((R, s)\)
under \(rk\) satisfies \([s]P = R + [c]\,rk\) with
\(c = H(R \,\|\, rk \,\|\, m)\). One is tempted to subtract
\([\alpha]P\) from the verification equation:
\([s - c\alpha]P = R + [c]X\), apparently a forgery under the
original key \(X\). But the identity holds with
\(c = H(R \,\|\, rk \,\|\, m)\), while the base verifier recomputes
\(H(R \,\|\, X \,\|\, m)\)---the plain subtraction converts
forgeries only for the \emph{non-key-prefixed} variant of Schnorr.
For key-prefixed RedDSA one argues in the ROM by either of two
routes. (i) The pair
\((rsk, rk) = (x + \alpha,\, X + [\alpha]P)\) is itself a
base-scheme key pair (Proposition~\ref{signatures:prop:rerand-consistent})
whose public key is uniform (Theorem~\ref{signatures:thm:unlink}),
so a forger against \(rk\) is \emph{literally} a forger against the
base scheme at the key \(rk\); the reduction, which knows
\(\alpha\), recovers \(x = rsk - \alpha\) from the extracted
\(rsk\). (ii) When the adversary sees signatures under several
randomisers \(\alpha_i\), one re-runs the forking argument of
Theorem~\ref{signatures:thm:schnorr-euf} on the critical query
\(H(R^{\star} \,\|\, rk \,\|\, m^{\star})\), simulating signing
under the remaining keys by oracle programming
(Lemma~\ref{signatures:lem:simulation}). A pitfall marks the
boundary: a naive translation that programs
\(H(R \,\|\, rk \,\|\, m) := H_{\mathrm{base}}(R \,\|\, X \,\|\, m)\)
must be a carefully injective domain swap, and it fails outright
with multiple randomisers---distinct \(rk_i\)-prefixed inputs would
collide on a single \(X\)-prefixed point, a detectable deviation
from a random oracle.
\end{remark}

\subsection{Binding signatures as a proof of knowledge of a discrete
logarithm}\label{signatures:subsec:binding}

The second RedPallas instantiation guards a different treasure: not
who may spend, but whether value is conserved. A shielded
transaction hides its amounts inside commitments; the forger of this
subsection is a counterfeiter who would arrange commitments whose
hidden values do not balance---minting money---while still
producing whatever certificate the protocol demands. The defence
consumes the homomorphism established in
\S\ref{commitments:subsec:pedersen} and turns this section's central
lesson---a Schnorr signature is a proof of knowledge of a discrete
logarithm---from analysis into design.

\paragraph{The value-commitment setup.}
Shielded protocols commit to value with the homomorphic Pedersen
commitment
\[
\mathrm{cm}(v, \rho) \;=\; [v]V + [\rho]P,
\]
where \(V\) and \(P\) are \emph{independent} generators of
\(\G\)---independent meaning that no party knows the discrete
logarithm of one to the base of the other, arranged by deriving both
by hashing to the curve (Construction~\ref{hash:con:nums})---and
\(\rho\) is a random blinding scalar. (In the notation of the
\texttt{orchard} crate and of
Remark~\ref{commitments:rem:homo-uses}, the randomness base here
called \(P\) is the generator \(R\); this section reserves \(R\) for
the signature nonce point. The base \(P\) is exactly the binding
base point of Remark~\ref{signatures:rem:redpallas}.) By
Theorem~\ref{commitments:thm:pedersen-homo}, the group sum of
commitments commits to the sum of values with the sum of blindings.

Sapling commits \emph{per note}: a transaction carries input
commitments \(\mathrm{cm}(v^{\mathrm{in}}_i, \rho^{\mathrm{in}}_i)\)
and output commitments
\(\mathrm{cm}(v^{\mathrm{out}}_j, \rho^{\mathrm{out}}_j)\), and the
net commitment is
\[
B \;=\; \sum_i \mathrm{cm}(v^{\mathrm{in}}_i, \rho^{\mathrm{in}}_i)
   \;-\; \sum_j \mathrm{cm}(v^{\mathrm{out}}_j,
                            \rho^{\mathrm{out}}_j).
\]
Orchard has \emph{no} per-note value commitments: each Action---the
unit that simultaneously spends a note of value \(v^{\mathrm{old}}\)
and creates a note of value \(v^{\mathrm{new}}\)---commits once to
its \emph{net} value change, publishing
\[
\mathrm{cv}^{\mathrm{net}}
= \mathrm{cm}(v^{\mathrm{old}} - v^{\mathrm{new}}, \rho)
= [v^{\mathrm{old}} - v^{\mathrm{new}}]V + [\rho]P,
\]
and \(B = \sum_i \mathrm{cv}^{\mathrm{net},i}\) over the
transaction's Actions (protocol specification \href{https://zips.z.cash/protocol/protocol.pdf\#concretehomomorphiccommit}{\S~5.4.8.3}). The
Action structure itself---how spends and outputs
are fused into these units and assembled into a transaction---is
the business of the \emph{Ironwood Guide}. Either way, the
homomorphism collapses the pile:
\[
B \;=\; [\Delta_v]V + [\bar\rho]P,
\]
where \(\Delta_v\) is the net value---
\(\sum_i v^{\mathrm{in}}_i - \sum_j v^{\mathrm{out}}_j\) in Sapling,
\(\sum_i (v^{\mathrm{old}}_i - v^{\mathrm{new}}_i)\) in Orchard---
and \(\bar\rho\) is the net blinding. The transaction balances
precisely when \(\Delta_v\) equals the public balancing value
\(v^{\mathrm{bal}}\); for a transaction with internal balance
\(\Delta_v = 0\), the \(V\)-component vanishes and
\(B = [\bar\rho]P\). Three observations set up the construction.
Every element of \(\G\) is some multiple of \(P\), since \(P\)
generates the prime-order group, so what distinguishes a balanced
transaction is not the \emph{existence} of a representation
\(B = [\beta]P\) but the ability to compute one. When the values
balance, the signer knows the multiplier: it is the net blinding
\(\bar\rho\), assembled from the individual blindings. When they do
not, computing any \(P\)-multiple representation of \(B\) from the
published commitments would require knowing \(\log_P V\), assumed
hard.

\begin{definition}[Binding signature]\label{signatures:def:bindingsig}
With \(B\) and \(\bar\rho\) as above and the transaction balanced,
so that \(B = [\bar\rho]P\), set the \emph{binding verification key}
\(bvk := B\) and the \emph{binding signing key}
\(bsk := \bar\rho\). The binding signature is a RedDSA signature
over the base point \(P\), with signing key \(\bar\rho\) and
verification key \(B\), on a message that is the hash of the
transaction (the \emph{sighash}). The verifier \emph{recomputes}
\(B\) from the publicly listed value commitments and the public
balancing value, then checks
\([s]P = R + [c]B\) with \(c = H(R \,\|\, B \,\|\, \text{sighash})\).
\end{definition}

\begin{theorem}[The binding signature is a proof of knowledge of
\(\log_P B\)]\label{signatures:thm:bindingpok}
In the ROM, a valid binding signature on a transaction is a proof of
knowledge of the discrete logarithm of the net commitment \(B\) to
the base \(P\); combined with the soundness of the accompanying
zero-knowledge proofs (per Action in Orchard, per note in Sapling),
which fix \emph{and range-check} the committed values, it forces
the transaction to balance.
\end{theorem}

\begin{proof}
\emph{Part 1: proof of knowledge.} The binding signature is exactly
a Fiat--Shamir Schnorr proof for the relation \(R_{\mathrm{dl}}\)
with base \(P\) and statement \(B\). By special soundness
(Theorem~\ref{signatures:thm:specialsound}) together with the
forking argument of Theorem~\ref{signatures:thm:schnorr-euf}---
specialised to a single statement and \emph{no} signing oracle---any
algorithm producing a valid signature under \(B\) with
non-negligible probability can be rewound to two accepting
transcripts \((R, c, s)\), \((R, c', s')\) with \(c \ne c'\), from
which the extractor recovers
\(\bar\rho = (s - s')(c - c')^{-1}\) with \(B = [\bar\rho]P\).
Producing a valid binding signature is therefore equivalent, up to
the standard rewinding loss, to \emph{knowing} \(\log_P B\).

\emph{Part 2: knowledge forces balance.} Write
\(B = [\Delta_v]V + [\bar\rho']P\) for the \emph{true} net value and
net blinding that the commitments imply; these are well defined once
the accompanying proofs fix the committed openings---in Orchard the
per-Action circuit fixes \(v^{\mathrm{old}}\), \(v^{\mathrm{new}}\)
and the opening of \(\mathrm{cv}^{\mathrm{net}}\), in Sapling the
per-note proofs fix each \(v_i, \rho_i\). If the signer knows
\(\bar\rho\) with \(B = [\bar\rho]P\), subtracting the two
expressions for \(B\) gives
\[
[\Delta_v]V \;=\; [\bar\rho - \bar\rho']P,
\]
where \(\Delta_v\) enters as a scalar, that is, as its residue in
\(\Fq\). Whenever \(\Delta_v \not\equiv 0 \pmod q\), that residue is
a nonzero element of \(\Fq\), hence invertible, so
\[
V \;=\; \bigl[(\bar\rho - \bar\rho')\,\Delta_v^{-1}\bigr]P
\]
---exhibiting \(\log_P V\) and contradicting the independence of
\(V\) and \(P\). Therefore \(\Delta_v \equiv 0 \pmod q\).

\emph{Part 3: from congruence to equality.} Balance in \(\Fq\)
alone is not yet economic balance: an integer net value equal to a
nonzero multiple of \(q\) would pass the test above while minting
\(q\) units at a stroke. The gap is closed by the \emph{range
checks} in the same accompanying proofs: in Orchard the per-Action
circuit constrains \(v^{\mathrm{old}}\) and \(v^{\mathrm{new}}\) to
\(\{0, \dots, 2^{64}-1\}\) (the Action statement, protocol
specification \href{https://zips.z.cash/protocol/protocol.pdf\#actionstatement}{\S~4.18.4}), and in Sapling the per-note proofs
constrain each \(v_i\) likewise. A transaction with \(n\) Actions
(count notes instead for Sapling) therefore has
\(\abs{\Delta_v} \le n\,(2^{64}-1)\) as an integer,
vastly below the \(255\)-bit \(q\) for any transaction that can be
serialised, and the only multiple of \(q\) in that range is \(0\).
Hence \(\Delta_v = 0\) over \(\Z\), not merely in \(\Fq\).
\end{proof}

In deployment the transaction is not internally balanced: it
declares a public balancing value \(v^{\mathrm{bal}}\), and the
verifier signs off against \(B - [v^{\mathrm{bal}}]V\), folding the
public value into the verification key---the \texttt{orchard}
crate's \texttt{Bundle::\allowbreak binding\_validating\_key}
subtracts a commitment to \(v^{\mathrm{bal}}\) with zero blinding
from the sum of the Actions' \(\mathrm{cv}^{\mathrm{net}}\)
(protocol specification \href{https://zips.z.cash/protocol/protocol.pdf\#orchardbalance}{\S~4.14}). The
argument of Theorem~\ref{signatures:thm:bindingpok} applies verbatim
to the folded key and forces
\(\Delta_v \equiv v^{\mathrm{bal}} \pmod q\); the balancing value is
encoded as a signed \(64\)-bit integer, so the range checks promote
the congruence to the integer equality
\(\Delta_v = v^{\mathrm{bal}}\) exactly as in Part~3.

\begin{remark}[Why a signature, and not just a proof]
\label{signatures:rem:why-signature}
The single pair \((R, s)\) accomplishes three things at once.
First, it proves knowledge of \(\bar\rho\), certifying balance, by
Theorem~\ref{signatures:thm:bindingpok}. Second, because the
Fiat--Shamir challenge hashes the transaction sighash, the same
signature \emph{authenticates} the transaction: altering any signed
field changes \(c\) and invalidates \((R, s)\), so the binding
signature doubles as an integrity check binding the shielded
components together---the origin of the name. Third, the
construction needs no trusted setup and no separate circuit for the
balance check: it is pure elliptic-curve arithmetic that anyone can
verify. The homomorphism of the commitment turns the arithmetic
statement ``the values balance'' into the algebraic statement
``\(B\) is a known multiple of \(P\)'', and the Schnorr signature is
the off-the-shelf proof of knowledge for precisely that
statement---the purest illustration in the protocol that a digital
signature \emph{is} a proof of knowledge of a discrete logarithm.
\end{remark}

\begin{remark}[One scheme, two meanings]
\label{signatures:rem:two-uses}
Set the two RedPallas uses side by side. The spend authorisation
signature verifies under a re-randomised key
\(rk = X + [\alpha]P\), chosen to \emph{hide} which long-term key
authorised the spend: the signing key \(rsk = x + \alpha\) encodes
authority, and \(\alpha\) buys unlinkability
(Theorem~\ref{signatures:thm:unlink}). The binding signature
verifies under a key \(B = [\bar\rho]P\) that the spender does
\emph{not} choose: the transaction's own commitments force it, the
``secret key'' \(\bar\rho\) is the net blinding factor, and signing
proves balance (Theorem~\ref{signatures:thm:bindingpok}). Both are
the same three lines of Schnorr arithmetic; what differs is the
\emph{meaning} assigned to the discrete logarithm---authority in one
case, value-conservation in the other. The unifying lesson of this
section: once a public key is the image of a secret under
\(x \mapsto [x]P\), a Schnorr signature is a compact,
non-interactive proof of knowing that secret, and the designer
chooses what knowing it \emph{means}.
\end{remark}

\subsection{The transparent layer's schemes: ECDSA and BIP-340 over
secp256k1}\label{signatures:subsec:transparent}

Two signature schemes in the series live outside the Pallas group:
the transparent layer authorises spends with ECDSA, and the issuance
bundle of Zcash Shielded Assets is authorised with a BIP-340 Schnorr
signature. Both run over secp256k1, and both are registered here at
the level a consumer needs---syntax, deployed conventions, and where
their security stands---rather than constructed from scratch.

The curve secp256k1 is the short Weierstrass curve
\(y^{2} = x^{3} + 7\) over the prime field \(\F_{p}\) with
\(p = 2^{256} - 2^{32} - 977\), with a fixed base point \(G\) of
prime order \(n\) (a \(256\)-bit prime) and cofactor \(1\); every
constant is a named parameter of SEC~2. In SEC~2's terminology it is
a \emph{Koblitz curve}, one with \(a = 0\), a choice that buys
implementation speed and plays no role in what follows. The group
law, point encodings, and the discrete-logarithm problem on such a
curve are the general machinery of the \emph{Math Guide},
\S``Elliptic curves'' and \S``The elliptic-curve discrete logarithm
problem''; the hardness assumption is
Definition~\ref{assumptions:def:dlog} in the group
\(\langle G \rangle\).

\begin{construction}[ECDSA over secp256k1, schematically]
\label{signatures:con:ecdsa}
Let \(e \in \Z/n\Z\) be the integer read from the \(32\)-byte digest
of the message, and write \(x(R)\) for the \(x\)-coordinate of a
point, read as an integer; the secret scalar is written \(d\), since
\(x\) is taken.
\begin{itemize}
\item \(\mathrm{KeyGen}\): sample \(d \xleftarrow{\$} [1, n-1]\), set
  \(X = [d]G\), output \((X, d)\).
\item \(\mathrm{Sign}(d, e)\): sample \(k \xleftarrow{\$} [1, n-1]\),
  set \(R = [k]G\) and \(r = x(R) \bmod n\), then
  \(s = k^{-1}(e + r d) \bmod n\); restart if \(r = 0\) or \(s = 0\);
  output \((r, s)\).
\item \(\mathrm{Verify}(X, e, (r, s))\): require
  \(r, s \in [1, n-1]\), compute
  \(R' = [e s^{-1}]G + [r s^{-1}]X\), and accept iff \(R'\) is not
  the identity and \(x(R') \equiv r \pmod n\).
\end{itemize}
Correctness: \(R' = [s^{-1}(e + rd)]G = [k]G = R\). The nonce \(k\)
must be fresh and secret, for the reason of
Proposition~\ref{signatures:prop:noncereuse} in different algebra:
two signatures \((r, s)\) and \((r, s')\) under one \(k\) on distinct
digests \(e \ne e'\) give \(s - s' = k^{-1}(e - e')\), hence
\(k = (e - e')(s - s')^{-1}\) and then \(d = (sk - e)\,r^{-1}\).
\end{construction}

Where Schnorr's challenge is a hash of \(R\), ECDSA's is the
coordinate \(r\) itself, so the scheme is not a Fiat--Shamir
transform and the forking-lemma route to
Theorem~\ref{signatures:thm:schnorr-euf} does not apply.
Unforgeability analyses exist---in the generic group model, and in
the ROM under additional assumptions on the map
\(R \mapsto x(R) \bmod n\)---but they lie outside this volume's
path: the transparent layer inherits ECDSA's security from Bitcoin's
track record rather than from a reduction.

\begin{remark}[Deployment conventions: the sighash, DER, and low
\(s\)]\label{signatures:rem:ecdsa-conventions}
Three conventions fix how the transparent layer uses
Construction~\ref{signatures:con:ecdsa}, which protocol
specification \href{https://zips.z.cash/protocol/protocol.pdf\#abstractsig}{\S~4.1.7} names as the script signature scheme.
\begin{enumerate}
\item \emph{The message is the sighash.} The digest \(e\) is the
  \(32\)-byte SIGHASH transaction hash of protocol specification
  \href{https://zips.z.cash/protocol/protocol.pdf\#sighash}{\S~4.10} (ZIP~244 for version-5 transactions), so the sighash
  algorithm plays the role of the hash in
  Theorem~\ref{signatures:thm:hashandsign}. Signer and verifier both
  feed it in as a digest.
\item \emph{DER encoding and the sighash-type byte.} On the wire a
  signature is the DER encoding of \((r, s)\)---an ASN.1
  \texttt{SEQUENCE} of two minimally encoded
  \texttt{INTEGER}s---followed by one trailing byte, the
  \emph{sighash type}, selecting which parts of the transaction the
  sighash commits to (\(\mathsf{SIGHASH\_ALL}\) for every signature
  the transaction builder emits; a signer working from a partially
  created transaction
  honours whichever valid type the input records). The signer
  appends the byte; the verifier splits it off and
  enforces strict DER unconditionally, the BIP-66 rule, so a
  malformed encoding fails script evaluation. The public key travels
  as its \(33\)-byte compressed encoding, the bytes that
  \(\mathsf{hash160}\) commits to (\S\ref{hash:subsec:transparent}).
\item \emph{Low-\(s\) normalisation.} If \((r, s)\) verifies then so
  does \((r, n - s)\), so a third party can re-encode a valid
  signature and change a pre-ZIP-244 transaction id; the
  normalisation takes the representative \(s \le (n-1)/2\) (BIP-62),
  which the deployed signer always emits (\texttt{secp256k1} crate,
  \texttt{sign\_ecdsa}, whose libsecp256k1 backend documents the
  created signature as always in lower-\(s\) form) and which script
  evaluation checks only under the \texttt{LowS} flag, one the
  consensus check
  does not set (\texttt{zcash\_script} crate,
  \texttt{src/external/\allowbreak pubkey.rs},
  \texttt{check\_low\_s}; \texttt{zebra-script} crate,
  \texttt{src/lib.rs}).
\end{enumerate}
\end{remark}

\begin{construction}[BIP-340 Schnorr over secp256k1, schematically]
\label{signatures:con:bip340}
Write
\(H_{\tau}(m) = \mathrm{SHA256}(\mathrm{SHA256}(\tau) \,\|\,
\mathrm{SHA256}(\tau) \,\|\, m)\) for the \emph{tagged hash} with
ASCII tag \(\tau\)---domain separation in the sense of
Definition~\ref{hash:def:domainsep}, the doubled tag digest making
the prefix a whole \(64\)-byte SHA-256 block---and read digests as
integers modulo \(n\).
\begin{itemize}
\item \emph{Keys.} Sample \(d \xleftarrow{\$} [1, n-1]\) and set
  \(X = [d]G\). The public key is the \(32\)-byte coordinate
  \(x(X)\) alone (\emph{x-only}): the verifier reconstructs \(X\) as
  the point with that abscissa and \emph{even} \(y\), and the signer
  replaces \(d\) by \(n - d\) when \(y(X)\) is odd, keeping the pair
  consistent.
\item \(\mathrm{Sign}(d, m)\): derive
  \(k = H_{\mathtt{BIP0340/nonce}}(t \,\|\, x(X) \,\|\, m) \bmod n\)
  with \(t = d \oplus H_{\mathtt{BIP0340/aux}}(a)\) for auxiliary
  randomness \(a\); set \(R = [k]G\), negating \(k\) when \(y(R)\) is
  odd; compute
  \(e = H_{\mathtt{BIP0340/challenge}}(x(R) \,\|\, x(X) \,\|\, m)
  \bmod n\) and \(s = k + e d \bmod n\); output the \(64\) bytes
  \((x(R), s)\).
\item \(\mathrm{Verify}(x(X), m, (r, s))\): lift \(x(X)\) to the
  even-\(y\) point \(X\), recompute \(e\), set
  \(R' = [s]G - [e]X\), and accept iff \(R'\) is not the identity,
  \(y(R')\) is even, and \(x(R') = r\).
\end{itemize}
\end{construction}

\begin{remark}[BIP-340 against
Construction~\ref{signatures:con:schnorrsig}; the ZSA consumer]
\label{signatures:rem:bip340}
The algebra is the Schnorr signature of
Construction~\ref{signatures:con:schnorrsig} in its \((R, s)\) form
with key-prefixing (Remark~\ref{signatures:rem:sig-conventions}),
under \(\G = \langle G \rangle\) of order \(n\) and
\(H = H_{\mathtt{BIP0340/challenge}}\); the unforgeability proof of
\S\ref{signatures:subsec:romforking} carries over with that tagged
hash modelled as the random oracle. What differs is convention:
points travel by abscissa alone with parity fixed by the even-\(y\)
rule, saving a byte per key and the parity of \(R\) per signature
and letting the same \(32\) bytes serve as key encoding and as hash
input; the challenge hash is domain-separated by tag; and the nonce
is deterministic with optional auxiliary randomness, as in
Construction~\ref{signatures:con:eddsa}. The in-series consumer is
the issuance authorisation signature of Zcash Shielded Assets:
ZIP~227 (\emph{specified}, in a draft ZIP) instantiates
\(\mathsf{IssueAuthSig}\) as BIP-340 over secp256k1, prefixes both
the validating-key encoding and the signature with a scheme byte
\(\mathtt{0x00}\), and signs the transaction's sighash. The
reference implementation, not yet deployed, is the QED-it fork of
the \texttt{orchard} crate (commit \texttt{cf801a5}),
\texttt{src/issuance/\allowbreak auth.rs}, whose
\texttt{ZSASchnorr} signs through the \texttt{secp256k1} crate's
BIP-340 routine with zeroed auxiliary randomness. The issuance
protocol itself is the business of the \emph{ZSA Guide}.
\end{remark}

\section{Interactive proofs, zero knowledge, and SNARKs}\label{sec:zk}

The adversary of this section wears two faces, one on each side of a
conversation. Facing the verifier stands the \emph{cheating prover}: a
party with no witness---no spendable note, no authorisation, no
satisfying assignment---who would nevertheless talk the verifier into
accepting, and thereby mint value, spend another's coins, or certify a
false computation to a network in which every validator takes the
proof's word for the hidden data. Facing the prover stands the
\emph{curious verifier}: a party who follows the protocol, or deviates
from it arbitrarily, in order to squeeze from the interaction more
than the single bit it is owed---to distinguish which witness the
prover holds, and so unmask the spender, link the transactions, or
read the amounts a shielded protocol exists to hide. The two threats
pull in opposite directions: defeating the first means the transcript
must pin the prover down; defeating the second means the transcript
must reveal nothing. That both can be defeated at once---and, more
remarkably still, by a proof \emph{shorter than the statement it
certifies}---is the subject of this section, and the reason a private
transaction can be checked by millions of verifiers who learn nothing
from it but its validity.

\subsection{Three questions}\label{zk:subsec:questions}

Three questions organise everything that follows.

\emph{Can a prover convince a sceptical verifier that a statement is
true?} Not by assertion---the verifier trusts nothing---but by a
protocol at whose end the verifier is justified in accepting, because
no cheating strategy could have survived it. The answer is the theory
of \emph{interactive proofs} and \emph{arguments}, refined into
\emph{proofs of knowledge}: protocols after which the verifier is
entitled to conclude not merely that a witness exists but that the
prover \emph{possesses} one.

\emph{Can the prover convince without revealing why the statement is
true?} The witness---the factorisation, the colouring, the discrete
logarithm, the transcript of the computation---must stay hidden; the
verifier must finish the protocol knowing the statement holds and
nothing else. The answer is \emph{zero knowledge}, made precise by
the simulation paradigm.

\emph{Can the proof be small---dwarfed by the statement it
certifies---and cheap to check?} A blockchain verifier cannot afford
to re-execute every computation it accepts, nor even to read a proof
as long as the computation's description. The answer is the
\emph{SNARK}: the succinct non-interactive argument of knowledge.

The development answers the questions in order and culminates in the
architecture underlying the Halo~2 proving system deployed in Orchard:
a \emph{polynomial interactive oracle proof}, compiled into an
argument by a \emph{polynomial commitment scheme}
(\S\ref{commitments:subsec:poly}) and made non-interactive by the
\emph{Fiat--Shamir transform}. Every earlier section of this volume
feeds in: the games and reductions of \S\ref{sec:model}, the discrete
logarithm of \S\ref{sec:assumptions}, the random oracle of
\S\ref{sec:hash}, the commitments of \S\ref{sec:commitments}, and the
Schnorr protocol of \S\ref{signatures:subsec:schnorrid}, which will
turn out to have been this section's canonical example all along.

Conventions are those of \S\ref{sec:model} throughout: \(\lambda\) is
the security parameter, adversaries are PPT (non-uniform where it
matters), and \(\negl\) denotes a negligible function. Proof systems
introduce one re-indexing: the objects of study are indexed by an
\emph{instance} \(x\) rather than by \(1^{\lambda}\), and asymptotic
statements are measured in the instance length \(\abs{x}\)---for the
instances this volume cares about, \(\abs{x}\) and \(\lambda\) are
polynomially related, so the two scales agree up to the polynomial
slop the definitions already absorb.

\subsection{Languages, relations, and witnesses}
\label{zk:subsec:relations}

What is a ``statement'', and what does it mean to hold the ``why'' of
one? The vocabulary is that of languages and relations over bit
strings.

\begin{definition}[Language, relation, witness]\label{zk:def:relation}
Fix the alphabet \(\{0,1\}\). A \emph{language} is a set
\(L \subseteq \{0,1\}^{*}\). A \emph{binary relation} is a set
\(R \subseteq \{0,1\}^{*} \times \{0,1\}^{*}\) of pairs \((x, w)\),
where \(x\) is called the \emph{instance} (or \emph{statement}) and
\(w\) a \emph{witness} for \(x\). The relation \emph{induces} the
language
\[
L_R \;:=\; \{\, x \;:\; \exists\, w \text{ with } (x, w) \in R \,\},
\]
and for each instance \(x\) the (possibly empty) \emph{witness set}
is \(R(x) := \{\, w : (x, w) \in R \,\}\). The relation \(R\) is an
\emph{NP-relation} if it is
\begin{enumerate}
\item \emph{polynomially balanced}: there is a polynomial \(p\) with
  \(\abs{w} \le p(\abs{x})\) for every \((x, w) \in R\); and
\item \emph{polynomial-time decidable}: a deterministic algorithm
  decides, in time polynomial in \(\abs{x}\), whether a given pair
  \((x, w)\) belongs to \(R\).
\end{enumerate}
\end{definition}

The class \(\mathsf{NP}\) is exactly the class of languages induced
this way: a language \(L\) lies in \(\mathsf{NP}\) if and only if
\(L = L_R\) for some NP-relation \(R\). This certificate
characterisation of \(\mathsf{NP}\) we take from complexity theory
without proof; it is the bridge over which every ``statement with a
short checkable reason'' enters the theory. Membership in \(L_R\) is
existentially quantified---\emph{some} witness exists---while holding
an element of \(R(x)\) is possession of the reason itself. The gap
between the two quantifiers is where this section lives.

\begin{example}[Three running relations]\label{zk:exp:relations}
Three relations recur through the section, in increasing order of
expressive power.
\begin{enumerate}
\item \emph{Discrete logarithm.} In the setting of
  \S\ref{signatures:subsec:schnorrid}---a cyclic group
  \(\G = \langle G \rangle\) of prime order \(q\), written
  additively---the relation is
  \[
  R_{\mathrm{dl}}
  \;=\; \bigl\{\, (X, x) \in \G \times \Fq \;:\; X = [x]G \,\bigr\}.
  \]
  Because \(x \mapsto [x]G\) is a bijection \(\Fq \to \G\), every
  instance \(X \in \G\) has exactly one witness: the induced language
  is \emph{trivial}---all of \(\G\)---and membership is vacuous. What
  is hard, under the DLP assumption of
  \S\ref{assumptions:subsec:dlp}, is to \emph{produce} the witness.
  This is the prototype of why proving membership differs from
  proving knowledge.
\item \emph{Graph three-colouring.} With instances encoding finite
  graphs \(G = (V, E)\),
  \[
  R_{3\mathrm{col}}
  \;=\; \bigl\{\, (G, c) \;:\; c \colon V \to \{1,2,3\}
    \text{ and } c(u) \ne c(v) \text{ for all } \{u,v\} \in E
  \,\bigr\},
  \]
  the relation of \emph{proper} three-colourings. The induced
  language \(L_{3\mathrm{col}}\) is \(\mathsf{NP}\)-complete---every
  \(\mathsf{NP}\) language reduces to it in polynomial time, a
  notion and fact we take from complexity theory without proof---so
  a zero-knowledge protocol for it yields one for \emph{every}
  \(\mathsf{NP}\) language by reduction
  (\S\ref{zk:subsec:snark}).
\item \emph{Circuit satisfiability.} With instances encoding
  circuits,
  \[
  R_{\mathrm{sat}}
  \;=\; \bigl\{\, (C, w) \;:\; C \text{ a Boolean (or arithmetic)
    circuit and } C(w) = 1 \,\bigr\}.
  \]
  This is the universal target of practical proof systems: a
  computation is first \emph{arithmetised} into the satisfiability
  of an arithmetic circuit over a finite field \(\F\), and the
  prover proves knowledge of a satisfying assignment
  (\S\ref{zk:subsec:snark}).
\end{enumerate}
The discrete-logarithm example already shows that ``\(x \in L\)'' is
too weak a target for cryptography: every group element is some
multiple of \(G\), so a protocol establishing mere membership
establishes nothing. The content of the Schnorr protocol is that the
prover \emph{knows} the scalar---and formalising ``knows'' for an
algorithm, an object with no mental states, is the central conceptual
move of this section (\S\ref{zk:subsec:knowledge}).
\end{example}

\subsection{Interactive protocols, proofs, and arguments}
\label{zk:subsec:proofs}

\begin{definition}[Interactive protocol]\label{zk:def:protocol}
An \emph{interactive protocol} is a pair \((P, V)\) of algorithms,
the \emph{prover} and the \emph{verifier}. Both receive a common
input \(x\); the prover additionally receives a private input \(w\)
(a witness, or an arbitrary auxiliary string); each party holds its
own private random tape. The parties alternate sending messages
\(a_1, a_2, \dots\) over a number of rounds polynomial in
\(\abs{x}\); after the last message the verifier outputs a single
bit, written
\[
\langle P(w), V \rangle(x) \;\in\; \{0, 1\},
\]
with \(1\) for \emph{accept}---a random variable over both parties'
coins. The \emph{transcript} of an execution is the full sequence of
exchanged messages (together with \(x\), and with the verifier's
coins when relevant). The protocol is \emph{public-coin} if every
message the verifier sends is a fresh uniformly random string that it
also reveals---the verifier's messages \emph{are} its coins---and its
final decision is a deterministic function of the transcript. The
verifier is required to run in probabilistic polynomial time,
always.
\end{definition}

The last clause is the asymmetry on which everything turns. The
verifier is the party being \emph{protected}, so its power is fixed
at PPT once and for all: a guarantee that only holds against verifiers
with super-polynomial resources would protect nobody. The prover is
the party whose power we \emph{vary}, and that variation is exactly
the distinction between proofs and arguments.

\begin{definition}[Interactive proof and interactive argument]
\label{zk:def:proofargument}
Let \(L \subseteq \{0,1\}^{*}\) be a language, and let
\(c, s \colon \N \to [0, 1]\). The protocol \((P, V)\) is an
\emph{interactive proof system} for \(L\) with \emph{completeness
error} \(c\) and \emph{soundness error} \(s\) if:
\begin{enumerate}
\item \emph{Completeness.} For every \(x \in L\) there is a prover
  input \(w\) with
  \[
  \Pr\bigl[\langle P(w), V \rangle(x) = 1\bigr]
  \;\ge\; 1 - c(\abs{x}).
  \]
\item \emph{Soundness.} For every \(x \notin L\) and \emph{every}
  prover strategy \(P^{*}\)---of \emph{unbounded} computational
  power, with arbitrary auxiliary input \(z\)---
  \[
  \Pr\bigl[\langle P^{*}(z), V \rangle(x) = 1\bigr]
  \;\le\; s(\abs{x}).
  \]
\end{enumerate}
The soundness clause is a game in the sense of
\S\ref{model:subsec:games}: the honest verifier is the challenger,
the prover strategy is the adversary, and the win condition is an
accept on a false instance. The two thresholds must be bounded apart,
\[
c(\abs{x}) + s(\abs{x}) \;\le\; 1 - \frac{1}{\poly(\abs{x})},
\]
and in the protocols of this volume both errors are negligible (often
\(c = 0\)). The class of languages admitting such proof systems---
with the honest prover unbounded---is \(\mathsf{IP}\). Every
protocol this volume constructs is for a language \(L = L_R\)
induced by an NP-relation \(R\), with the honest prover required to
run in polynomial time given a witness \(w \in R(x)\) as its private
input; the unbounded honest prover matters only for the
complexity-theoretic statements below.

An \emph{interactive argument system} (or \emph{computationally
sound} proof) relaxes only the soundness quantifier: the guarantee
holds against every \emph{PPT} cheating prover \(P^{*}\), that is,
for \(x \notin L\) and PPT \(P^{*}\) with arbitrary auxiliary input
\(z\), \(\Pr[\langle P^{*}(z), V \rangle(x) = 1] \le s(\abs{x})\),
where \(s\) is required to be negligible. Soundness of an argument
typically rests on a computational hardness assumption.
\end{definition}

\begin{remark}[Proofs versus arguments: neither dominates]
\label{zk:rem:proofvsarg}
A proof resists an unbounded cheating prover: its soundness is
information-theoretic and unconditional. An argument resists only
efficient adversaries, and an unbounded prover can typically break
one outright---by breaking the binding of a commitment the protocol
relies on (Theorem~\ref{commitments:thm:impossibility} showed
perfectly hiding commitments have this exposure), or simply by
brute-forcing a discrete logarithm. In exchange, arguments achieve
properties provably impossible for proofs, chief among them
\emph{succinctness}: a sound interactive proof for an
\(\mathsf{NP}\)-complete language whose communication is much smaller
than the witness would let an unbounded prover compress
\(\mathsf{NP}\) certificates below what complexity theory is believed
to allow, collapsing classes believed distinct, whereas succinct
\emph{arguments} exist under standard assumptions
(\S\ref{zk:subsec:snark}). For cryptographic applications---all of
the SNARK literature included---the prover is a real machine with
real resource bounds, and the argument notion is exactly right:
bounded-prover soundness is the price of succinctness, and its
gateway.
\end{remark}

Why interaction, and why randomness? Strip both away and see what is
left. With a deterministic verifier and a single prover message, the
system proves exactly the class \(\mathsf{NP}\): the one message is a
certificate, the verifier's check is the relation's decision
procedure---and there is no hope of zero knowledge, since the
certificate is, in spirit, the witness itself, and the verifier could
replay the prover's analysis at leisure. Verifier randomness changes
the game: a cheating prover must commit to its messages \emph{before}
knowing the challenge, so a random spot-check catches what a
predictable one never would. Interaction lets the verifier probe
adaptively, round after round. Together they buy real power: the
equality \(\mathsf{IP} = \mathsf{PSPACE}\)---interactive proofs
capture exactly the languages decidable with polynomial memory and
unrestricted running time, a class believed far larger than
\(\mathsf{NP}\)---is a landmark of complexity theory (Lund, Fortnow,
Karloff, and Nisan; Shamir) that we cite without proof. For this
section the sharper point is qualitative: interaction and randomness
are what make zero knowledge and knowledge extraction possible at
all. A single static certificate can be neither hidden nor rewound.

\begin{proposition}[Soundness amplification by repetition]
\label{zk:prop:amplify}
Let \((P, V)\) be an interactive proof (or argument) for \(L\) with
completeness error \(c\) and soundness error \(s\), and let
\(k = k(\abs{x})\) be polynomially bounded. The protocol that runs
\(k\) independent instances and accepts only if all \(k\) accept has
completeness error at most \(k \cdot c\) and soundness error at most
\(s^{k}\). Sequential repetition achieves this always; parallel
repetition achieves it for proofs and, with care, for many arguments.
In particular, if \(s \le 1 - \delta\) for a noticeable function
\(\delta\), then choosing \(k\) with
\(\delta(\abs{x}) \cdot k(\abs{x}) = \omega(\log \abs{x})\)---always
possible with \(k\) polynomial---drives the soundness error below
every inverse polynomial, so a base protocol needs only a noticeable
soundness \emph{gap}.
\end{proposition}

\begin{proof}
For completeness, the honest prover fails at least one of the \(k\)
runs with probability at most \(k \cdot c\) by the union bound
(\emph{Math Guide}, \S``The union bound and a birthday
calculation''). For soundness under sequential repetition, fix
\(x \notin L\) and any cheating strategy; conditioned on the
transcripts of the first \(i - 1\) runs, the \(i\)-th run is an
execution against a fresh independently-random verifier, so it
accepts with probability at most \(s\) whatever state the prover
carries forward; multiplying the conditional bounds gives \(s^{k}\).
(For parallel repetition of proofs the same bound holds by a more
careful argument, which we omit; for arguments parallel repetition
can fail to amplify in general, and we use only the sequential
form.) Finally
\(s^{k} \le (1 - \delta)^{k} \le e^{-\delta k}\), which is negligible
once \(\delta k = \omega(\log \abs{x})\).
\end{proof}

\subsection{Knowledge soundness and extractors}
\label{zk:subsec:knowledge}

Soundness, as defined, is silent exactly where cryptography speaks.
For the relation \(R_{\mathrm{dl}}\) every instance \(X \in \G\) lies
in the induced language, so the soundness clause of
Definition~\ref{zk:def:proofargument}---quantified over
\(x \notin L\)---holds \emph{vacuously}: a verifier that always
accepts is ``sound'' for the trivial language. The assurance actually
wanted is that a convincing prover \emph{possesses} the discrete
logarithm. But possession is not a property an algorithm wears on its
sleeve; the definition reads it operationally. A prover knows a
witness if an efficient procedure, granted full control over the
prover---the ability to run it, feed it messages, and \emph{rewind}
it to an earlier state to try different continuations---can compute
the witness from it. The procedure is called the \emph{extractor},
and what it extracts is the meaning of ``knows''.

% FIGURE (planned, dedicated figures pass): fig:extractor --- the
% extractor cloning and rewinding a prover across challenges. A boxed
% extractor E containing two (or more) parallel copies of the same
% prover P* run from a cloned state: a shared prefix band up to the
% commitment a, then divergent challenge arrows e and e' into the two
% copies, responses z and z' flowing out, and both transcripts feeding
% a witness-computation box that emits w. A caption note that E holds
% P* as a subroutine --- rewinding is legitimate because the prover is
% code, not a live counterparty.

\begin{definition}[Knowledge soundness; proof of knowledge]
\label{zk:def:knowledge}
Let \((P, V)\) be an interactive protocol and \(R\) an NP-relation.
The protocol is a \emph{proof of knowledge} for \(R\) with
\emph{knowledge error} \(\kappa \colon \N \to [0,1]\) if there exist
a probabilistic oracle algorithm \(\mathcal E\) (the
\emph{extractor}) and a polynomial \(q\) such that for every prover
strategy \(P^{*}\)---possibly cheating, possibly computationally
unbounded---and every instance \(x\), writing
\[
\varepsilon(x) \;:=\; \Pr\bigl[\langle P^{*}, V \rangle(x) = 1\bigr]
\]
for the probability that \(P^{*}\) convinces the honest verifier, the
following holds whenever \(\varepsilon(x) > \kappa(x)\): the
extractor \(\mathcal E\), given \(x\) and oracle access to
\(P^{*}\) \emph{with the ability to rewind}---to run \(P^{*}\) from
any point with the same coins, feeding verifier messages of
\(\mathcal E\)'s choosing---outputs a witness \(w \in R(x)\) within
expected running time
\[
\frac{q(\abs{x})}{\varepsilon(x) - \kappa(x)}.
\]
If the guarantee is required only for \emph{PPT} prover strategies
\(P^{*}\), the protocol is an \emph{argument of knowledge}.
\end{definition}

Three consequences can be read straight off the definition. The
extraction cost may blow up as \(\varepsilon\) approaches
\(\kappa\): the bound \(q(\abs{x})/(\varepsilon - \kappa)\) is
useful precisely when the prover's success is noticeably above the
knowledge error. A prover succeeding with probability at most
\(\kappa\) need yield nothing at all---\(\kappa\) is the ``free''
success rate, typically the probability of blind guessing, below
which conviction carries no evidence of knowledge. And any prover
noticeably above \(\kappa\) surrenders the witness: whatever code and
coins produce conviction can be mined, efficiently, for the very
object the conviction asserts. Knowledge is thereby defined as
\emph{efficient computability relative to the prover's own code and
coins}---an algorithm knows what can be efficiently computed from
it. Note the expected-time clause: extractors are the standing
reason this volume's adversary class admits expected polynomial time
(Remark~\ref{model:rem:variants}).

\begin{remark}[The rewinding justification]\label{zk:rem:rewinding}
No live counterparty could rewind the prover, and none ever does:
the extractor is a \emph{thought experiment}, never executed during
an honest run of the protocol. It exists to give the security claim
its meaning---and, in security proofs, its teeth. In a reduction the
extractor \emph{is} the reduction: the cheating prover is a
subroutine the reduction runs through its input/output interface
alone (Definition~\ref{model:def:reduction}), and black-box access
already includes choosing the subroutine's random tape and
restarting it from the top with the same coins---so replaying it
under different challenges is not science fiction but the ordinary
right of a machine over its own subroutines, no reading of its code
required. The same
device runs in both directions across this section: the
\emph{simulator} of \S\ref{zk:subsec:zk} rewinds a verifier to fake
transcripts without a witness, and the extractor rewinds a prover to
mine real witnesses from conviction. Rewinding is the single most
important technical device in the theory.
\end{remark}

\begin{proposition}[Knowledge soundness implies soundness]
\label{zk:prop:ks-sound}
A proof of knowledge for \(R\) with knowledge error \(\kappa\) is an
interactive proof for \(L_R\) with soundness error
\(s(x) \le \kappa(x)\). The converse fails: a protocol can be sound
without being a proof of knowledge.
\end{proposition}

\begin{proof}
Fix \(x \notin L_R\), so that \(R(x) = \emptyset\). If some prover
strategy \(P^{*}\) achieved
\(\varepsilon(x) > \kappa(x)\), the extractor run on \(P^{*}\) would
halt in finite expected time and output some \(w \in R(x)\)---
impossible, as the witness set is empty. Hence
\(\varepsilon(x) \le \kappa(x)\) for every \(P^{*}\), which is the
soundness bound. For the failure of the converse: soundness
constrains only false instances, so a protocol for a trivial
language---such as \(R_{\mathrm{dl}}\)'s, where no false instances
exist---is automatically sound however it behaves, while nothing
forces a witness to be extractable from it. One may be convinced of
a true statement by a prover who does not know why it is true;
knowledge soundness is strictly stronger, and it---not mere
soundness---is the notion cryptographic applications actually rely
on (\S\ref{zk:subsec:snark}).
\end{proof}

\subsection{The simulation paradigm and zero knowledge}
\label{zk:subsec:zk}

What should ``the verifier learned nothing'' mean? The definitional
move---often called the most influential idea in modern
cryptography---is the \emph{simulation paradigm}: the verifier
learned nothing beyond the truth of the statement if it could have
produced, entirely on its own and without ever talking to the prover,
a transcript indistinguishable from the real one. Whatever a party
can manufacture for itself it did not learn from the conversation;
if the whole conversation is reproducible without the witness, the
conversation taught nothing about the witness. The three grades of
indistinguishability---perfect, statistical, computational---carry
over verbatim from \S\ref{model:subsec:indist}
(Definitions~\ref{model:def:statdist} and~\ref{model:def:indist},
with the unbounded-distinguisher bound of
Proposition~\ref{model:prop:statdist} and the hierarchy of
Proposition~\ref{model:prop:hierarchy}) under one re-indexing:
ensembles are indexed by instances \(x\) (and auxiliary inputs)
rather than by the security parameter alone, and negligibility is
measured in \(\abs{x}\).

\begin{definition}[Zero knowledge]\label{zk:def:zk}
Let \((P, V)\) be an interactive proof or argument for \(L = L_R\).
For a verifier strategy \(V^{*}\), an instance \(x\), a witness
\(w \in R(x)\), and an auxiliary input \(z\), the \emph{view}
\[
\mathrm{View}_{V^{*}}\bigl(P(w), V^{*}(z)\bigr)(x)
\]
is the random variable consisting of \(x\), the random tape of
\(V^{*}\), and every message \(V^{*}\) receives---everything the
verifier sees. The protocol is \emph{zero knowledge} if there exists
a PPT simulator \(\mathrm{Sim}\) such that for every \(x \in L\),
every witness \(w \in R(x)\), and every auxiliary input \(z\), the
ensembles
\[
\bigl\{ \mathrm{Sim}(x, z) \bigr\}
\quad\text{and}\quad
\bigl\{ \mathrm{View}_{V^{*}}\bigl(P(w), V^{*}(z)\bigr)(x) \bigr\}
\]
are indistinguishable; the grade of indistinguishability---perfect,
statistical, computational---gives \emph{perfect} (PZK),
\emph{statistical} (SZK), or \emph{computational} (CZK) zero
knowledge. Two regimes for the quantifier over \(V^{*}\):
\begin{enumerate}
\item \emph{Honest-verifier zero knowledge} (HVZK): the guarantee is
  required only for \(V^{*} = V\), the verifier that follows the
  protocol, with \(z\) empty; the simulator must reproduce the
  distribution of honest-execution transcripts.
\item \emph{Malicious-verifier (auxiliary-input) zero knowledge}:
  the guarantee holds for \emph{every} PPT verifier strategy
  \(V^{*}\), deviating arbitrarily, with arbitrary auxiliary input
  \(z\). Here \(\mathrm{Sim}\) may use \(V^{*}\) as a subroutine
  \emph{and rewind it}, and is permitted expected polynomial running
  time (Remark~\ref{model:rem:variants}).
\end{enumerate}
\end{definition}

Why does the simulator's existence suffice? The simulator receives
\(x\)---and \(z\)---but \emph{not} the witness \(w\). If its output
is nevertheless indistinguishable from the real view, then the real
view carries no efficiently extractable information about \(w\)
beyond the fact \(x \in L\): any property of \(w\) an efficient
verifier could compute from the interaction, it could compute
equally well from the simulator's output, which was manufactured
without \(w\)---otherwise the computation would itself be a
distinguisher. The simulator wields a power the prover lacks:
it may rewind \(V^{*}\), retrying until the verifier's challenges
happen to align with a transcript preparable without the witness.
That power is legitimate for exactly the reason rewinding always is
(Remark~\ref{zk:rem:rewinding}): the simulator is the analyst's
tool, holding the verifier as a subroutine, not a participant in any
real execution.

The auxiliary input \(z\) models the prior knowledge a malicious
verifier may bring to the conversation---the history of earlier
protocol runs, partial information about the witness, anything.
Quantifying over all \(z\) is what makes zero knowledge
\emph{compose}: a protocol that is zero knowledge for every
auxiliary input remains zero knowledge when run as a subroutine of a
larger system, whose surrounding state simply becomes the \(z\) of
the analysis. Practice therefore adopts the auxiliary-input
formulation, and so does this volume.

\subsection{Sigma-protocols}\label{zk:subsec:sigma}

The general definitions are now in place, and one protocol shape
realises them with remarkable economy: three moves, public coins,
and algebra doing all the work. The shape is named for the capital
\(\Sigma\) its message flow traces on the page.

% FIGURE (planned, dedicated figures pass): fig:sigma --- the
% three-move Sigma-protocol message flow and its Fiat--Shamir
% collapse. Left panel: prover and verifier columns, arrows a (down
% to V), e (back to P, stamped uniform from C), z (down to V), with
% Verify(x,a,e,z) at the bottom; the a-e-z zigzag traces the capital
% sigma. Right panel: the same protocol collapsed by Fiat--Shamir ---
% a single proof string pi = (a, z) with e = H(x, a) computed locally
% on both sides, the hash box replacing the verifier's arrow.
% Caption ties to Definition zk:def:fs.

\begin{definition}[Sigma-protocol]\label{zk:def:sigma}
Let \(R\) be an NP-relation. A \emph{Sigma-protocol} for \(R\) with
\emph{challenge space} \(C\) is a three-move public-coin protocol on
common input \(x\), the prover holding \(w\) with \((x, w) \in R\):
\begin{enumerate}
\item \emph{Commitment.} The prover sends a first message \(a\)
  (the \emph{commitment}, or announcement), computed from
  \((x, w)\) and fresh randomness.
\item \emph{Challenge.} The verifier sends
  \(e \xleftarrow{\$} C\), uniformly random.
\item \emph{Response.} The prover sends \(z\), computed from
  \((x, w, a, e)\) and its randomness.
\end{enumerate}
The verifier decides by a deterministic predicate
\(\mathrm{Verify}(x, a, e, z) \in \{0, 1\}\). Three properties are
required.
\begin{itemize}
\item \emph{Completeness.} An honest prover holding
  \((x, w) \in R\) convinces the honest verifier with probability
  \(1\).
\item \emph{Special soundness.} There is a PPT algorithm that, given
  \(x\) and any two accepting transcripts \((a, e, z)\) and
  \((a, e', z')\) with the \emph{same} commitment \(a\) but
  \emph{distinct} challenges \(e \ne e'\), outputs a witness
  \(w \in R(x)\).
\item \emph{Special honest-verifier zero knowledge} (sHVZK). There
  is a PPT simulator that, on input \(x\) and \emph{any} challenge
  \(e \in C\), outputs a pair \((a, z)\) such that
  \(\mathrm{Verify}(x, a, e, z) = 1\) and the triple \((a, e, z)\)
  is distributed exactly---or statistically, or computationally,
  close to---a real honest transcript conditioned on the challenge
  being \(e\). The simulator chooses \(a\) \emph{after} fixing
  \(e\); reversing the honest order is what makes simulation
  possible without \(w\).
\end{itemize}
Two accepting transcripts sharing a commitment but differing in
challenge are called a \emph{collision}. Special soundness says a
collision \emph{is} a witness: it is the local, two-transcript
reformulation of knowledge extraction, and the engine of the theorem
that follows.
\end{definition}

\begin{theorem}[Special soundness yields knowledge soundness]
\label{zk:thm:special}
Let \((P, V)\) be a Sigma-protocol for \(R\) with challenge space
\(C\) of size \(N := \abs{C}\). Then \((P, V)\) is a proof of
knowledge for \(R\) with knowledge error \(\kappa = 1/N\).
\end{theorem}

\begin{proof}
We construct an extractor establishing the slightly weaker knowledge
error \(2/N\), with expected running time
\(O(1/(\varepsilon - 2/N))\) prover invocations whenever the
prover's success probability satisfies \(\varepsilon > 2/N\); the
closing paragraph addresses the gap to the optimal \(1/N\).

\emph{The prover as a matrix.} Fix \(x\) and a prover strategy
\(P^{*}\); fold any auxiliary input into its code. The behaviour of
\(P^{*}\) is a deterministic function of its random tape \(\rho\):
fixing \(\rho\) fixes the commitment \(a(\rho)\), and then each
challenge \(e\) determines the response \(z(\rho, e)\). Form the
Boolean matrix
\[
H[\rho, e] \;:=\;
\begin{cases}
1 & \text{if } \mathrm{Verify}\bigl(x,\, a(\rho),\, e,\,
      z(\rho, e)\bigr) = 1,\\
0 & \text{otherwise},
\end{cases}
\]
with one row per random tape and one column per challenge; then
\(\varepsilon = \Pr_{\rho, e}[H[\rho, e] = 1]\), and a collision is
two \(1\)-entries in a single row.

\emph{The extractor.} \emph{Step 1.} Run \(P^{*}\) with fresh
\(\rho\) and a uniform challenge \(e\), repeatedly, until an
accepting transcript \((a, e, z)\) appears. The number of trials is
geometric with success probability \(\varepsilon\), so this costs
\(1/\varepsilon\) invocations in expectation---the geometric waiting
time (\emph{Math Guide}, \S``Beyond finite spaces: countable
additivity, limits, and densities''). \emph{Step 2.} Alternate two
kinds of trial: \emph{(i)} rewind \(P^{*}\) to the moment just after
it emitted \(a\)---the tape \(\rho\) stays fixed---and feed one
fresh uniform challenge \(e' \ne e\); \emph{(ii)} make one entirely
fresh probe, as in Step 1. If a kind-(i) trial accepts, a collision
\((a, e, z), (a, e', z')\) is in hand: go to Step 3. If a kind-(ii)
trial accepts, \emph{abandon} the current row and continue Step 2
from the new accepting transcript (a fresh \emph{phase}).
\emph{Step 3.} Apply the special-soundness algorithm of
Definition~\ref{zk:def:sigma} to the collision and output the
witness.

\emph{Analysis.} Call a row \(\rho\) \emph{heavy} if its acceptance
density \(\delta_\rho := \Pr_e\bigl[H[\rho, e] = 1\bigr]\) is at
least \(\varepsilon/2\). The heavy-row lemma---a Markov-style
averaging---states that an accepting entry lies in a heavy row with
probability at least \(1/2\): the total accepting mass contributed
by light rows is less than \(\varepsilon/2\) (each light row
contributes less than \(\varepsilon/2\) of its own row measure), so
conditioned on acceptance the light rows carry less than half the
mass. Every phase opens with a fresh accepting transcript
distributed exactly as Step 1's (the opening transcript is either
Step 1's or a kind-(ii) success, both conditioned draws from the
accepting entries), so the lemma applies at every phase, not only
the first.

Each alternation of Step 2 contains a kind-(ii) trial, which accepts
with probability \(\varepsilon\); the number of alternations in a
phase is therefore dominated by a geometric variable with mean
\(1/\varepsilon\), and a phase costs \(O(1/\varepsilon)\)
invocations in expectation. This interleaving is not an
optimisation but a necessity: a row whose only accepting challenge
is the already-used \(e\)---any row with
\(\delta_\rho \le 1/N\)---contains no collision at all, and an
extractor that only rewound would sit in such a row for unbounded
expected time. The fresh probes cap every phase.

Now condition on the phase's row being heavy. A kind-(i) trial draws
\(e'\) uniformly from \(C \setminus \{e\}\) and accepts with
probability at least
\[
\frac{\delta_\rho N - 1}{N - 1}
\;\ge\; \delta_\rho - \frac{1}{N}
\;\ge\; \frac{\varepsilon}{2} - \frac{1}{N},
\]
which is positive precisely when \(\varepsilon > 2/N\). The phase
ends at its first success of either kind, and the chance that this
first success is a collision rather than a row-switch is at least
\[
\frac{\varepsilon/2 - 1/N}{(\varepsilon/2 - 1/N) + \varepsilon}
\;=\; \Omega\!\left(\frac{\varepsilon - 2/N}{\varepsilon}\right).
\]
Un-conditioning costs the heavy-row factor \(1/2\), so each phase
yields a collision with probability
\(\Omega\bigl((\varepsilon - 2/N)/\varepsilon\bigr)\), and the
number of phases is dominated by a geometric variable with mean
\(O\bigl(\varepsilon/(\varepsilon - 2/N)\bigr)\). By Wald's
identity---the expected total is the expected number of phases times
the expected cost per phase (\emph{Math Guide}, same section)---the
extractor's expected number of prover invocations is
\[
O\!\left(\frac{\varepsilon}{\varepsilon - 2/N}\right) \cdot
O\!\left(\frac{1}{\varepsilon}\right)
\;=\; O\!\left(\frac{1}{\varepsilon - 2/N}\right),
\]
of the form \(q(\abs{x})/(\varepsilon - 2/N)\) demanded by
Definition~\ref{zk:def:knowledge} with \(\kappa = 2/N\).

\emph{Closing the factor of two.} The threshold
\(\delta_\rho \ge \varepsilon/2\) defining heaviness is an artefact
of this exposition, and it is exactly what costs the extra \(1/N\)
in the knowledge error. Sharper accounting removes it:
Damg\aa rd's refined heavy-row analysis in his Sigma-protocol
lecture notes, and, in full generality, the
\(k\)-special-soundness extractor of Attema, Cramer, and Kohl, which
achieves knowledge error exactly \((k-1)/N\)---here \(k = 2\),
giving the optimal \(\kappa = 1/N\) that matches the soundness
error. We cite these rather than reproduce them; every use in this
volume tolerates the factor of two.
\end{proof}

\begin{remark}[Knowledge error equals soundness error, and how to
size the challenge space]\label{zk:rem:challengesize}
The knowledge error \(1/N\) of Theorem~\ref{zk:thm:special} is no
accident of the analysis: an explicit attack meets it. A prover with no
witness guesses the challenge in advance: it picks \(e^{*} \in C\),
runs the sHVZK simulator on \((x, e^{*})\) to obtain an accepting
pair \((a, z)\)---for a protocol such as Schnorr's the simulator's
algebra works for \emph{every} instance---and sends \(a\). If the
verifier's challenge happens to equal \(e^{*}\), which occurs with
probability exactly \(1/N\), the prepared \(z\) is accepted. One
challenge is thus always answerable without a witness, so \(1/N\) is
also the protocol's soundness error, and
Theorem~\ref{zk:thm:special} says the two coincide. One caveat
keeps the lower bound honest: under
Definition~\ref{zk:def:knowledge} an extractor facing this guessing
prover (\(\varepsilon = 1/N\)) is granted expected time
\(q(\abs{x})/(1/N - \kappa)\), so for exponential \(N\) a claimed
\(\kappa\) below \(1/N\) is not formally refuted---the definition
would hand the extractor a super-polynomial budget, enough to
brute-force a witness unaided. The attack does rule out any
noticeably smaller \(\kappa\) once the extractor is required to run
in strict polynomial time and witnesses are genuinely hard to
compute: such an extractor, run on this witness-free prover, would
otherwise mine a witness from nothing---for \(R_{\mathrm{dl}}\),
solving the discrete-logarithm problem itself. To make the
error negligible there are two routes: take \(C\) exponentially
large---for instance \(C = \Fq\) for a \(\lambda\)-bit prime \(q\),
so \(N = q\)---or repeat a small-challenge protocol
\(\omega(\log \lambda)\) times
(Proposition~\ref{zk:prop:amplify}). Large challenge spaces are why
a \emph{single} round of Schnorr already has negligible knowledge
error, while a protocol with a polynomially small challenge space
needs many repetitions.
\end{remark}

Multi-round protocols generalise the collision to a tree. The
succinct arguments of \S\ref{zk:subsec:snark} run many challenge
rounds, and their extractors need more than two transcripts.

\begin{definition}[Tree special soundness]\label{zk:def:treess}
Let \((P, V)\) be a \((2\mu + 1)\)-move public-coin protocol for
\(R\) in which the prover speaks first and the verifier issues
\(\mu\) uniform challenges from a space \(C\), and let
\(k_1, \dots, k_\mu\) be positive integers. A
\emph{\((k_1, \dots, k_\mu)\)-tree of accepting transcripts} for
\(x\) is a tree of depth \(\mu\) in which every node at level \(i\)
has \(k_i\) children labelled by \emph{pairwise distinct} challenges
for round \(i\), every root-to-leaf path carries the prover messages
of one accepting transcript, and transcripts passing through a
common node share their prefix up to that node. The protocol is
\emph{\((k_1, \dots, k_\mu)\)-special-sound} if a PPT algorithm
computes a witness \(w \in R(x)\) from any such tree. A
Sigma-protocol is the case \(\mu = 1\), \(k_1 = 2\).
\end{definition}

Extraction generalises phase by phase: under the hypotheses of a
tree-extraction theorem---the Attema--Cramer--Kohl analysis cited in
the proof of Theorem~\ref{zk:thm:special} is the general
form---rewinding the prover round by round assembles the required
tree of \(\prod_i k_i\) accepting transcripts: the protocol is a
proof of knowledge with knowledge error growing from \(1/\abs{C}\)
to the order of \(\bigl(\sum_i (k_i - 1)\bigr)/\abs{C}\), and with
extraction cost polynomial provided the tree size \(\prod_i k_i\)
is. Both degradations are benign in the regime the compiled
protocols of this section occupy: \(C\) exponentially large,
\(\mu\) logarithmically bounded, and the \(k_i\) constant, so the
error stays negligible and \(\prod_i k_i\) polynomial. This is
the notion under which the inner-product argument's extractor
operates (Proposition~\ref{commitments:prop:ipa-props}): each of
the \(\log_2 d\) folding rounds contributes one level of branching
to the tree, three accepting children at every node. Each round's
accepting relation is linear in three unknown parent terms with
coefficients that are powers of the challenge, so three suitably
distinct challenges make that system invertible. The
\emph{Halo~2 Guide} carries out
that extraction for the deployed protocol, whose IPA normalisation
involves each round's challenge through the coefficients
\(u^{-1}\), \(1\), \(u\): the extractor accordingly needs pairwise
distinct \emph{nonzero} challenges at every node---a
negligible-fraction exclusion absorbed into the knowledge
error---and not the pairwise distinct squares (\(u \ne \pm u'\))
that the alternative \(u^{-2}, 1, u^{2}\) normalisation would
demand.

\begin{proposition}[sHVZK implies HVZK]\label{zk:prop:shvzk}
A Sigma-protocol satisfying special honest-verifier zero knowledge
is honest-verifier zero knowledge in the sense of
Definition~\ref{zk:def:zk}, with the same grade (perfect,
statistical, or computational).
\end{proposition}

\begin{proof}
The honest view is a triple \((a, e, z)\) in which \(e\) is uniform
on \(C\). The HVZK simulator draws \(e \xleftarrow{\$} C\) itself,
invokes the sHVZK simulator on \((x, e)\) to obtain \((a, z)\), and
outputs \((a, e, z)\). For each fixed \(e\), the simulated pair
\((a, z)\) is distributed as a real run conditioned on that
challenge---exactly, or up to the stated grade of closeness---and
the simulator's \(e\) has the same uniform marginal as the honest
verifier's; averaging the conditional distributions over uniform
\(e\) therefore reproduces the honest view, with any statistical or
computational slack preserved under the averaging.
\end{proof}

\begin{example}[Schnorr: the canonical Sigma-protocol]
\label{zk:exp:schnorr}
The Schnorr identification protocol of
\S\ref{signatures:subsec:schnorrid} is a Sigma-protocol for the
discrete-logarithm relation \(R_{\mathrm{dl}}\) of
Example~\ref{zk:exp:relations}, and the properties proved there,
self-contained, are exactly the clauses of
Definition~\ref{zk:def:sigma}. The commitment is \(R = [r]G\); the
challenge is a uniform \(c \in \Fq\); the response is
\(s = r + cx\); the predicate is \([s]G = R + [c]X\). Perfect
completeness is Proposition~\ref{signatures:prop:completeness};
\(2\)-special soundness is
Theorem~\ref{signatures:thm:specialsound}, whose collision algebra
yields the witness \(x = (s - s')(c - c')^{-1}\); special HVZK is
Theorem~\ref{signatures:thm:hvzk}, whose simulator samples \(s\) and
solves \(R = [s]G - [c]X\)---choosing the commitment after the
challenge, as the definition prescribes. The challenge space is
\(\Fq\), so \(N = q\) and the knowledge error \(1/q\) is negligible
in a single round (Remark~\ref{zk:rem:challengesize}). By
Theorem~\ref{zk:thm:special} and Proposition~\ref{zk:prop:shvzk},
Schnorr's protocol is a perfect-HVZK proof of knowledge of a
discrete logarithm.
\end{example}

\begin{remark}[From honest to malicious verifiers: three remedies]
\label{zk:rem:malicious}
Sigma-protocols are honest-verifier zero knowledge only. A malicious
\(V^{*}\) may choose \(e\) as a function of the commitment \(a\),
and the sHVZK simulator---which fixes \(e\) before producing
\(a\)---does not obviously cope: nothing lets it predict which
challenge the verifier's function will return for a commitment it
has yet to manufacture. Three standard remedies close the gap.
\emph{(1) Guess and rewind}, for small challenge spaces: the
simulator guesses \(e^{*} \xleftarrow{\$} C\), runs the sHVZK
simulator on \((x, e^{*})\), and feeds the resulting \(a\) to
\(V^{*}\); if the verifier's challenge equals \(e^{*}\), the round
is simulated, and otherwise the simulator rewinds \(V^{*}\) and
guesses afresh. The honest prover fixes \(a\) before any challenge
exists, so the simulated \(a\) carries (almost) no information
about the guess and each attempt succeeds with probability close
to \(1/N\): for polynomially large \(N\) the expected \(O(N)\)
attempts stay polynomial, and sequential repetition of the
small-challenge protocol---simulated round by round, the
auxiliary input carrying the history---is full malicious-verifier
zero knowledge, at the price of many rounds. For exponentially
large \(N\) the guess never lands, and repetition does not
generically upgrade honest-verifier to malicious-verifier zero
knowledge. \emph{(2) Challenge commitment}: the verifier commits to
\(e\) before seeing \(a\) (a coin-tossing step using the
commitments of \S\ref{sec:commitments}), so \(e\) cannot depend on
\(a\); under the additional conditions such a compiler imposes for
simulation, the simulator rewinds past the opening. \emph{(3) The
Fiat--Shamir transform}: remove the verifier's choice altogether by
deriving \(e\) from a hash of the transcript---under the hypotheses
of Theorem~\ref{zk:thm:fs}, in the programmable random oracle
model, the resulting \emph{non-interactive} protocol is zero
knowledge against \emph{all} verifiers, since there is no live
challenge left to skew. For the SNARK pipeline route (3) is the one
that matters, and honest-verifier zero knowledge is exactly the
property the transform requires. It is next.
\end{remark}

\subsection{The Fiat--Shamir transform: from interactive to
non-interactive}\label{zk:subsec:fiatshamir}

In a public-coin protocol the verifier contributes nothing but fresh
randomness: its messages are coin tosses, its decision a
deterministic function of the transcript
(Definition~\ref{zk:def:protocol}). Every such message can therefore
be replaced by a value that \emph{anyone} can compute---a hash of
everything sent so far---and the interaction collapses into a single
proof string, verifiable by recomputing the hashes. The prover no
longer talks to a live, possibly malicious verifier but to a fixed
public function, which is why honest-verifier zero knowledge
upgrades to full non-interactive zero knowledge. The cost is that
the hash must be idealised: the analysis lives in the random oracle
model, with \(H \colon \{0,1\}^{*} \to C\) a random oracle whose
range is the challenge space
(Definition~\ref{signatures:def:rom-range}), and it exploits both
powers established in Proposition~\ref{hash:prop:rom-powers}:
\emph{observability}---the reduction sees every query the adversary
makes---and \emph{programmability}---the reduction may fix the
oracle's output at a freshly queried point to any value, provided
answers stay consistent and uniformly distributed.

\begin{definition}[The Fiat--Shamir transform]\label{zk:def:fs}
Let \((P, V)\) be a public-coin protocol for \(R\) with challenge
space \(C\), and let \(H\) be a random oracle with range \(C\). The
\emph{Fiat--Shamir transform} \(\mathrm{FS}[(P, V), H]\) is the
non-interactive system defined, in the Sigma case, by:
\begin{itemize}
\item \emph{Prover.} Compute the commitment \(a\) as in \((P, V)\);
  set \(e := H(x, a)\); compute the response \(z\); output the proof
  \(\pi := (a, z)\) (equivalently \((a, e, z)\)).
\item \emph{Verifier.} Given \(x\) and \(\pi = (a, z)\), recompute
  \(e := H(x, a)\) and accept iff
  \(\mathrm{Verify}(x, a, e, z) = 1\).
\end{itemize}
A multi-round public-coin protocol applies the rule per round: the
\(i\)-th challenge is
\[
e_i \;:=\; H(x,\, a_1,\, e_1,\, \dots,\, a_i),
\]
each challenge hashing the \emph{entire} transcript prefix. To bind
a message \(m\) to the proof---turning it into a signature---one
includes \(m\) in the hash: \(e := H(x, a, m)\).
\end{definition}

\begin{remark}[The pitfall: weak Fiat--Shamir]\label{zk:rem:weakfs}
The definition's insistence that each challenge hash \emph{all}
prior messages, the instance \(x\) included, is load-bearing, and
omitting any part of it has broken deployed systems. Hashing
\(e := H(a)\) without \(x\)---``weak Fiat--Shamir''---lets an
attacker fix \((a, e, z)\) first and choose the instance
\emph{afterwards}, searching for an \(x\) that the frozen challenge
happens to validate: the interactive protocol's soundness quantified
over provers who commit before the challenge, and an unhashed \(x\)
re-opens exactly that order of quantifiers, yielding proofs of false
statements. In multi-round protocols, hashing only the latest
message rather than the running prefix breaks soundness the same
way, one round at a time. The maxim: \emph{the challenge must be a
hash of everything the verifier would have seen up to the moment it
speaks.}
\end{remark}

\begin{theorem}[Fiat--Shamir security in the ROM]\label{zk:thm:fs}
Let \((P, V)\) be a Sigma-protocol for \(R\) with challenge space
\(C\) of size \(N\), special soundness, and special HVZK; assume
further that, conditioned on \((x, e)\), the first message output
by the sHVZK simulator has min-entropy \(\omega(\log \lambda)\)
(min-entropy \(H_{\infty}\):
Definition~\ref{symmetric:def-minentropy}).
In the random oracle model,
\(\Pi = \mathrm{FS}[(P, V), H]\) is a non-interactive argument for
\(R\) that is:
\begin{enumerate}
\item \emph{perfectly complete}, inheriting the protocol's
  completeness;
\item \emph{knowledge sound}: from any PPT prover \(P^{*}\) that
  makes at most \(Q\) random-oracle queries and outputs an accepting
  proof for \(x\) with probability \(\varepsilon\), an extractor
  obtains a witness \(w \in R(x)\) with probability at least
  \(\varepsilon \bigl(\varepsilon/Q - 1/N\bigr)\) up to constant
  factors---non-negligible whenever \(\varepsilon\) is
  non-negligible and \(N\) is super-polynomial; and
\item \emph{non-interactive zero knowledge}: a PPT simulator that
  controls the random oracle produces, without the witness, proofs
  indistinguishable from real ones---for all verifiers, including
  malicious ones.
\end{enumerate}
\end{theorem}

\begin{proof}
\emph{Completeness.} The honest prover computes \(e = H(x, a)\) and
an honest response; the verifier recomputes the very same \(e\) from
the same inputs and runs the same predicate, so the protocol's
probability-\(1\) completeness transfers verbatim.

\emph{Zero knowledge, by programming.} The simulator, on input
\(x\), draws \(e \xleftarrow{\$} C\), runs the sHVZK simulator on
\((x, e)\) to obtain an accepting pair \((a, z)\), and then
\emph{programs} the oracle:
\[
H(x, a) \;:=\; e .
\]
Programming fails only if \((x, a)\) was queried before this
moment; by hypothesis the simulated first message has, conditioned
on \((x, e)\), min-entropy \(\omega(\log \lambda)\), so a given
prior query collides with \(a\) with probability
\(2^{-H_\infty(a)}\), and a
union bound over the adversary's prior queries makes the failure
probability negligible. Conditioned on no failure, the programmed
answer is a fresh uniform value---exactly what an unprogrammed
oracle would have returned---so the oracle's joint distribution is
untouched, and by sHVZK the output \((a, z)\) is distributed as a
real proof. The programming is invisible to every verifier, and
there is no live challenge for a malicious verifier to skew: its
malice is moot, and honest-verifier simulation delivers
full non-interactive zero knowledge.

\emph{Knowledge soundness, by forking.} Let \(P^{*}\) output an
accepting proof \((a, z)\) for \(x\) with probability
\(\varepsilon\). With overwhelming probability a successful
\(P^{*}\) actually queried \(H\) at \((x, a)\): otherwise the
challenge \(e = H(x, a)\) used by the verifier is a uniform value
the prover never saw, and the proof verifies with probability at
most \(1/N\). Say the critical query is the \(j\)-th of the at most
\(Q\) queries. The extractor runs \(P^{*}\) once, recording the
transcript of oracle answers and the index \(j\) of the critical
query; it then \emph{rewinds} \(P^{*}\) to just before the \(j\)-th
query and re-runs it with the same coins, answering the first
\(j - 1\) queries identically---so the commitment \(a\) is
unchanged---and the \(j\)-th and subsequent queries with fresh
uniform values, the \(j\)-th being \(e' \xleftarrow{\$} C\). This is
precisely the experiment of the general forking lemma
(Theorem~\ref{signatures:thm:forking}, with \(h = N\)): if \(P^{*}\)
succeeds with probability \(\varepsilon\) over \(Q\) query slots,
then with probability at least
\(\varepsilon(\varepsilon/Q - 1/N)\) both runs succeed at the same
index \emph{and} \(e \ne e'\). Conditioned on that event, the two
runs yield accepting transcripts \((a, e, z)\) and \((a, e', z')\)
with a common commitment and distinct challenges---a collision---and
the special-soundness algorithm of Definition~\ref{zk:def:sigma}
outputs \(w \in R(x)\).
\end{proof}

\begin{remark}[Caveats, and the refinements SNARKs need]
\label{zk:rem:rom-caveats}
The random oracle is an idealisation: a real hash function is fixed,
public, and not random, and there exist contrived protocols secure
in the ROM yet insecure under \emph{every} concrete instantiation.
The transform is nonetheless trusted in practice, with the working
stance of Remark~\ref{model:rem:rom-caveat}, and it underlies
essentially every
deployed non-interactive proof and signature---Fiat--Shamir-compiled
SNARKs included. Two refinements matter for the compiled arguments
of \S\ref{zk:subsec:snark}. First, Theorem~\ref{zk:thm:fs} covers
the one-challenge Sigma case only, and does not by itself prove
security for a multi-round protocol such as Halo~2's: there the
extractor must fork the transcript at \emph{each} round, its
success probability degrading with the round count and with the
protocol's soundness structure, and special soundness must
generalise to tree special soundness
(Definition~\ref{zk:def:treess})---a separate multi-round theorem,
whose hypotheses must match the protocol at hand. Second, when the
underlying interactive protocol is itself only computationally
sound---an argument---the ROM extraction \emph{layers on top of}
the computational soundness reduction rather than replacing it.
There is also a gap between proofs that merely observe the oracle
and proofs that program it (the non-programmable ROM is strictly
weaker as a proof technique); deployed systems adopt the
programmable model---the zero-knowledge clause above, which
programs the oracle outright, already shows why---and an analysis
must state which programmable-oracle (and, where used, algebraic)
model it works in.
\end{remark}

\begin{example}[Schnorr signatures, recovered]
\label{zk:exp:schnorrsig}
Applying Definition~\ref{zk:def:fs} to Schnorr's Sigma-protocol
(Example~\ref{zk:exp:schnorr}) with a message \(m\) folded into the
challenge hash yields, symbol for symbol, the Schnorr signature
scheme of Construction~\ref{signatures:con:schnorrsig}; and its
EUF-CMA security in the ROM
(Theorem~\ref{signatures:thm:schnorr-euf}) is established by exactly
the forking argument in the proof of Theorem~\ref{zk:thm:fs},
augmented with a signing-oracle simulation. The signatures section
proved the instance; this section supplies the general statement it
instantiates. The slogan: \emph{proof of knowledge + Fiat--Shamir =
signature}.
\end{example}

\subsection{SNARKs: succinct non-interactive arguments of knowledge}
\label{zk:subsec:snark}

The Sigma-protocol paradigm reaches all of \(\mathsf{NP}\): a
commit--challenge--reveal proof of knowledge for the
\(\mathsf{NP}\)-complete relation \(R_{3\mathrm{col}}\) of
Example~\ref{zk:exp:relations}---zero knowledge under one-way
functions alone, a result of Goldreich, Micali, and Wigderson taken
without proof---extends to every \(\mathsf{NP}\) language by
reduction. The generic reduction is hopelessly inefficient, however,
its proofs polynomially \emph{longer} than the computation they
certify, so practical systems abandon it and arithmetise the
computation directly.

A blockchain verifier operates under constraints the classical
theory never contemplated. It cannot interact---a transaction is
broadcast once and verified by parties who will never speak to its
author, years later, from a proof string sitting in a block. It
cannot afford to re-execute---every full node checks every proof,
so verification must cost less than the computation being certified,
by orders of magnitude. And it must resist a prover with every
incentive to cheat: an accepting proof \emph{is} money. The object
meeting all three constraints at once is the SNARK.

\begin{definition}[Non-interactive argument]\label{zk:def:nizk}
A \emph{non-interactive argument} for an NP-relation \(R\) is a
triple of PPT algorithms:
\begin{itemize}
\item \(\mathrm{Setup}(1^{\lambda}, R) \to srs\), producing a
  \emph{structured reference string}---the public parameters. In
  the random oracle model all three algorithms additionally access
  one globally sampled oracle; the oracle is not part of the
  \(srs\)---a PPT setup cannot output an infinite object---though
  public parameters may be derived through it. A setup holding no
  secret trapdoor is \emph{transparent}: no secret is ever held, so
  none must be trusted to have been discarded
  (Remark~\ref{commitments:rem:setup-trust}).
\item \(\mathrm{Prove}(srs, x, w) \to \pi\), for
  \((x, w) \in R\).
\item \(\mathrm{Verify}(srs, x, \pi) \to \{0, 1\}\).
\end{itemize}
The system is \emph{complete} if honestly generated proofs always
verify. A \emph{preprocessing} argument lets \(\mathrm{Setup}\)
depend on the circuit or relation and output, alongside the proving
parameters, a short \emph{verification key}, amortising the
verifier's per-statement work across all statements of that shape.
\end{definition}

\begin{definition}[SNARK]\label{zk:def:snark}
A \emph{succinct non-interactive argument of knowledge} (SNARK) for
an NP-relation \(R\) is a non-interactive argument that is
additionally:
\begin{enumerate}
\item \emph{Knowledge sound.} For every PPT prover \(P^{*}\) there
  is a PPT extractor \(\mathcal E\)---in the ROM, observing
  \(P^{*}\)'s oracle queries and rewinding or forking it---such
  that, on the same setup,
  \[
  \Pr\bigl[\mathrm{Verify}(srs, x, \pi) = 1
    \;\wedge\; (x, \mathcal E^{P^{*}}) \notin R\bigr]
  \;\le\; \negl(\lambda),
  \]
  the probability over the setup and both algorithms' coins: an
  accepting proof implies an extractable witness, except with
  negligible probability.
\item \emph{Succinct.} The proof size \(\abs{\pi}\) and the
  verifier's running time are \emph{sublinear} in the size of the
  computation proved---typically
  \(\abs{\pi} = O(\poly(\lambda))\) or
  \(O(\poly(\lambda) \log \abs{C})\), and verification
  \(O(\poly(\lambda) \log \abs{C})\) for a circuit \(C\),
  independent of, or polylogarithmic in, the witness and circuit
  sizes. (Some of the literature reserves ``succinct'' for
  polylogarithmic proofs; the constant-size proofs of pairing-based
  schemes are the strongest form.)
\end{enumerate}
With zero knowledge---the non-interactive, ROM form of
Definition~\ref{zk:def:zk}, as delivered in the Sigma case by
Theorem~\ref{zk:thm:fs}, clause 3---it is a \emph{zk-SNARK}.
\end{definition}

\begin{remark}[Why these systems are arguments]
\label{zk:rem:succinctness-args}
The ``AR'' in SNARK is earned by necessity, not by preference.
Remark~\ref{zk:rem:proofvsarg} showed that a succinct \emph{proof}
for an \(\mathsf{NP}\)-complete language is essentially ruled out:
a transcript shorter than the witness cannot
information-theoretically pin the witness down, and an unbounded
prover could otherwise compress certificates below what complexity
theory is believed to allow. Succinct \emph{arguments} escape
because a computationally bounded prover cannot \emph{find} a
convincing-but-false short proof even though one may
exist---succinctness and computational soundness arise together,
by necessity. The constructions below realise the necessity at a
specific spot: their polynomial commitments are binding only
against \emph{efficient} adversaries, under cryptographic
assumptions, and an unbounded prover could equivocate on a
commitment and defeat extraction. It is at that commitment layer
that the information-theoretic core turns into an argument.
\end{remark}

What does \emph{knowledge} soundness buy an application? Consider a
private transaction whose statement reads: \emph{there exists a note
I am authorised to spend, and input and output values balance}.
Plain soundness guarantees only that \emph{some} satisfying witness
exists---and for such existential statements about a large anonymity
set, witnesses exist in abundance: other people's notes satisfy the
existential quantifier perfectly well. Soundness alone would not
prevent a prover holding no spendable note from proving the
existential statement, should she somehow come by a proof. Knowledge
soundness closes the gap: whoever produced the accepting proof
\emph{actually holds} a concrete witness---a real spendable note, a
real authorisation---that the extractor could recover from her. The
application may then act---release funds, grant access---exactly as
though the witness had been presented in the clear, while it stays
hidden. Soundness protects against false statements; knowledge
soundness protects against provers who do not know what they claim.
Almost every cryptographic use, and certainly every monetary one,
requires the latter (Proposition~\ref{zk:prop:ks-sound}).

\subsubsection*{The pairing-based route}

The first SNARKs to reach deployment, and the ones the deployed
Sapling protocol still uses, verify with a bilinear pairing.

\begin{definition}[Bilinear pairing]\label{zk:def:pairing}
Let \(\G_1, \G_2, \G_T\) be cyclic groups of prime order \(r\), the
first two written additively with generators \(G\) and \(\hat G\).
A \emph{bilinear pairing} is a map
\(e \colon \G_1 \times \G_2 \to \G_T\) that is
\begin{enumerate}
\item \emph{bilinear}: \(e([a]P, [b]Q) = e(P, Q)^{ab}\) for all
  \(P \in \G_1\), \(Q \in \G_2\), and \(a, b \in \Fr\);
\item \emph{non-degenerate}: \(e(G, \hat G) \ne 1\), so that
  \(e(G, \hat G)\) generates \(\G_T\); and
\item \emph{efficiently computable}.
\end{enumerate}
Bilinearity lets a verifier test one multiplication of hidden
scalars---whether \(c = ab\), given \([a]G\), \([b]\hat G\), and
\([c]G\)---without learning any of them
(Example~\ref{assumptions:exa:pairing}); the KZG check of
Construction~\ref{commitments:def:kzg-eval} is such a test. The
groups come from \emph{pairing-friendly} elliptic curves, chosen
with a small embedding degree (the quantity of
Definition~\ref{assumptions:def:secure-curve}, there required
large) so that \(\G_T\) lives in a manageable extension field; the
curve of the deployed Sapling protocol is BLS12-381, a
Barreto--Lynn--Scott curve of embedding degree twelve (protocol
specification \href{https://zips.z.cash/protocol/protocol.pdf\#blspairing}{\S~5.4.9.2}). The deployed system treats the pairing
as a black box, and so does this volume.
\end{definition}

The lineage is short. The BCTV14 family (Ben-Sasson, Chiesa, Tromer,
and Virza) is a preprocessing argument in the sense of
Definition~\ref{zk:def:nizk}: the setup encodes one fixed arithmetic
circuit into a structured reference string of group elements, the
prover's work is a handful of multiscalar multiplications over that
string, and the proof is a constant number of group
elements---eight, in the deployed encoding (protocol specification
\href{https://zips.z.cash/protocol/protocol.pdf\#bctv}{\S~5.4.10.1})---checked by a constant number of pairing equations,
whatever the circuit's size. Groth's
refinement (Groth16) cuts the proof to \emph{three} group elements,
\(\pi_A, \pi_C \in \G_1\) and \(\pi_B \in \G_2\), and verification
to a single pairing-product equation: three pairings---\(\pi_A\)
against \(\pi_B\), \(\pi_C\) against a fixed verification-key
element, and a public-input-weighted combination of verification-key
elements against another---compared with one pairing precomputed
once per circuit. Deployed Zcash uses Groth16
over BLS12-381 for every Sapling spend and output proof (protocol
specification \href{https://zips.z.cash/protocol/protocol.pdf\#groth}{\S~5.4.10.2}; proofs produced and verified by the
\texttt{bellman} crate's \texttt{groth16} module, consumed through the
\texttt{zcash\_proofs} crate). PLONK, the
third generation, is the recipe of Theorem~\ref{zk:thm:recipe}
below with the KZG scheme of Construction~\ref{commitments:def:kzg}
as its commitment: the gate-and-permutation arithmetisation Halo~2
also uses, each polynomial opening checked by two pairings, so proof
size and verification cost stay constant. Constant-size proofs and
constant-time verification, independent of the circuit, are the
family's signature---the strongest form of the succinctness clause
of Definition~\ref{zk:def:snark}.

The price is the setup. The reference string is produced from secret
scalars---the evaluation point \(\tau\) of
Construction~\ref{commitments:def:kzg} and, in Groth16, further
secrets that tie the proof elements to one circuit---which must be
\emph{destroyed} once the parameters are published: they are the
\emph{toxic waste}. Whoever retains them forges accepting proofs of
false statements at will, exactly as the trapdoor \(\tau\) forges
KZG openings (Remark~\ref{commitments:rem:kzg-trade}), and no
verifier can tell: knowledge soundness collapses silently. For
Groth16 the waste is \emph{per circuit}---every new statement shape
needs a fresh setup---whereas PLONK's KZG string is universal, one
powers-of-\(\tau\) string serving every circuit up to a size bound,
but a trusted one still. The mitigation is a \emph{multiparty
ceremony}: participants contribute secret randomness in turn, each
re-randomising the running parameters and destroying its own share,
so that the waste stays unknown as long as \emph{one} participant
was honest. Zcash ran such ceremonies for the parameters of its
earlier proof systems---the BCTV14 Sprout circuit and the Groth16
Sapling circuits (protocol specification \href{https://zips.z.cash/protocol/protocol.pdf\#bctvparameters}{\S\S~5.7}--\href{https://zips.z.cash/protocol/protocol.pdf\#grothparameters}{5.8})---and each
remains sound only on that one-honest-participant assumption.

The other route through the design space keeps the discrete
logarithm and nothing else. The inner-product commitment of
Construction~\ref{commitments:def:ipa} needs no pairing and no
trusted setup---its generators are hashed into existence, with no
secret behind them---at the price of logarithmic proofs and a
linear-time opening check
(Proposition~\ref{commitments:prop:ipa-props}); it is the road the
\emph{Halo~2 Guide} takes, and the one the rest of this subsection
paves.

\subsubsection*{Polynomial interactive oracle proofs}

Modern SNARKs are built in two layers---one information-theoretic,
one cryptographic---and the information-theoretic layer speaks the
language of polynomials.

\begin{definition}[Polynomial interactive oracle proof]
\label{zk:def:polyiop}
Fix a finite field \(\F\). A \emph{polynomial interactive oracle
proof} (Poly-IOP) for a relation \(R\) over \(\F\) is a public-coin
interactive protocol in which, in place of ordinary messages, the
prover in each round sends an \emph{oracle} for a polynomial
\(p_i \in \F[X]\) of some bounded degree; the verifier sends
uniformly random field-element challenges, and may \emph{query} each
polynomial oracle at points of its choice, receiving the evaluation
\(p_i(\zeta)\). After the interaction the verifier accepts or
rejects as a function of its challenges and the few evaluations it
queried. Completeness and (knowledge) soundness are as for
interactive proofs, with the soundness analysis purely
algebraic---resting, invariably, on the Schwartz--Zippel lemma: a
nonzero polynomial of degree at most \(d\) over \(\F\) vanishes at
a uniformly random point with probability at most \(d/\abs{\F}\)
(\emph{Math Guide}, \S``The Schwartz--Zippel lemma''). A good
Poly-IOP verifier makes only \(O(1)\) or \(O(\log)\) queries, each
a single evaluation.
\end{definition}

The oracle is an idealisation---no real prover can hand over an
opaque box answering evaluation queries honestly---but a productive
one: it isolates the algebra from the cryptography. Soundness of the
Poly-IOP layer is \emph{unconditional}, resting on no assumption
beyond the rigidity of low-degree polynomials; the cryptography
enters only when the oracles are instantiated.

\emph{Arithmetisation} is how a computation becomes a Poly-IOP
statement. The claim ``I know \(w\) with \(C(w) = 1\)''
(the relation \(R_{\mathrm{sat}}\) of
Example~\ref{zk:exp:relations}) is encoded as a small set of
polynomial identities over \(\F\) that hold if and only if the
computation is correct. The computation is first expressed as a
\emph{constraint system}---the customisable
\emph{PLONKish} gate-and-permutation form used by Halo~2---whose
satisfaction is equivalent to a handful of polynomials satisfying
identities on a chosen evaluation domain, a fixed finite set of
points with one point per constraint row. In the PLONKish form,
polynomials encoding the wire values \(a, b, c\) and fixed selector
polynomials \(q_L, q_R, q_M, q_O, q_C\) chosen at circuit-definition
time must satisfy the gate identity
\[
q_L \cdot a \;+\; q_R \cdot b \;+\; q_M \cdot a b
\;+\; q_O \cdot c \;+\; q_C \;=\; 0
\]
at every point of the domain---selectors switching each row among
addition, multiplication, and custom behaviours---while a
\emph{permutation argument} enforces the copy constraints wiring one
gate's output to another's input. One more step turns the
on-domain requirement into a spot-checkable one: the identity is
demanded only on the domain \(H\), and a polynomial vanishes on all
of \(H\) exactly when it is divisible by the vanishing polynomial
\(Z_H := \prod_{h \in H}(X - h)\) (\emph{Math Guide}, \S``Roots and
the factor theorem''), so the prover also sends an oracle for the
quotient \(t\) with
\[
q_L \cdot a + q_R \cdot b + q_M \cdot ab + q_O \cdot c + q_C
\;=\; t \cdot Z_H
\]
as an identity of polynomials---one that must now hold at
\emph{every} point of \(\F\). Checking it at a single random point
is the Poly-IOP: the verifier queries the oracles at a random
\(\zeta\), checks the displayed equation there, and by
Schwartz--Zippel a false statement survives with probability at
most \(d/\abs{\F}\), with \(d\) bounding the degree of either
side.

\emph{The polynomial commitment scheme is the cryptographic half.}
The PCS of Definition~\ref{commitments:def:pcs} supplies exactly
what the idealised oracle promised: a short, binding commitment to a
polynomial, with short proofs---verifiable against the commitment
alone---that \(p(\zeta) = v\); evaluation binding forbids two
accepted values at one point, and the \emph{extractability} the
definition names is what SNARK compilation consumes: a prover whose
evaluation proofs verify must ``know'' an underlying polynomial.
Three instantiations mark out the design space: KZG
(Construction~\ref{commitments:def:kzg})---pairing-based,
constant-size, trusted setup; the inner-product-argument commitment
over a single elliptic curve
(Construction~\ref{commitments:def:ipa})---transparent, no trusted
setup, logarithmic-size openings, the choice made by Halo~2; and
the hash-based transparent schemes, outside this volume's deployed
path and mentioned only to bound the territory.

\begin{theorem}[The compilation recipe, informal]\label{zk:thm:recipe}
Let \(\mathrm{PIOP}\) be a public-coin polynomial IOP for \(R\),
complete and knowledge-sound, and let \(\mathrm{PC}\) be a
polynomial commitment scheme that is correct, evaluation-binding,
and extractable. Compile as follows: wherever the prover would send
a polynomial oracle \(p_i\), it sends
\(\mathrm{Commit}(p_i)\) instead; wherever the verifier would query
\(p_i(\zeta)\), the prover sends the claimed value together with a
\(\mathrm{PC}\) evaluation proof, which the verifier checks against
the commitment. Under compatible composition hypotheses the result
is a public-coin interactive argument of knowledge for \(R\), and a
suitable multi-round Fiat--Shamir theorem in the ROM
(Definition~\ref{zk:def:fs} gives the transform: every challenge a
hash of the entire transcript so far) makes it non-interactive. Its
proof is succinct; under the strict Definition~\ref{zk:def:snark}
it is a SNARK only if verification is also sublinear in the proved
computation. If, additionally, the Poly-IOP is properly
blinded---random masking terms on the prover's polynomials, hiding
commitments (Remark~\ref{commitments:rem:pcs-hiding})---and the
corresponding simulation theorem applies, the result is zero
knowledge.
\end{theorem}

\begin{proof}[Proof sketch: layered extraction and the succinctness
account]
\emph{Completeness} transfers verbatim: honest commitments open to
honest evaluations, and the honest Poly-IOP verifier accepts.

\emph{Knowledge soundness} is \emph{layered}, one extractor feeding
another. Given a successful prover against the compiled argument,
the SNARK extractor first invokes \(\mathrm{PC}\)
\emph{extractability} on each commitment, recovering actual
polynomials \(\hat p_i\) of the right degrees consistent with every
evaluation the prover got accepted; this step is where the
cryptographic assumption lives, and why the compiled object is an
\emph{argument}---a prover that could equivocate on its commitments
would defeat it, and only computational binding rules that out. The
extractor then hands the \(\hat p_i\) to the Poly-IOP's
\emph{unconditional} knowledge extractor, which is sound against
\emph{any} polynomials---in particular the extracted ones---and
outputs \(w \in R(x)\); evaluation binding guarantees the prover
could not have answered queries inconsistently with the committed
\(\hat p_i\), so the algebraic soundness analysis applies to the
extracted polynomials as run. Fiat--Shamir then converts the
public-coin interactive argument to a non-interactive one in the
ROM by the multi-round generalisation of Theorem~\ref{zk:thm:fs}:
challenges become transcript hashes
(Remark~\ref{zk:rem:weakfs} governing what they must absorb),
observability and forking supply the rewinding that the layered
extractor needs---tree-style, per
Definition~\ref{zk:def:treess}---and the honest-verifier zero
knowledge of the blinded protocol becomes full non-interactive zero
knowledge by oracle programming.

\emph{Succinctness} is an accounting exercise, and the account is
paid entirely by \(\mathrm{PC}\). The proof consists of a constant
or logarithmic number of commitments and evaluation proofs---short
by construction, sublinear in the circuit, independent of the
witness---so the compiled \emph{proof} is succinct whenever the
commitment scheme's objects are. The \emph{verifier} is sublinear
only when the evaluation checks are themselves succinct, as for
KZG's two pairings; with the inner-product commitment each
evaluation check costs \(O(d)\) group operations
(Proposition~\ref{commitments:prop:ipa-props}), so the compiled
verifier runs in time linear in the circuit while the proof stays
succinct. The compiled object therefore satisfies every clause of
Definition~\ref{zk:def:snark} when the \(\mathrm{PC}\) openings
verify succinctly, and otherwise yields succinct proofs with
linear-time verification---the trade Halo~2 accepts and then
recovers through batching.
\end{proof}

\begin{remark}[Halo 2 is the recipe, instantiated]
\label{zk:rem:halo2}
The proving system deployed in Orchard (protocol specification
\href{https://zips.z.cash/protocol/protocol.pdf\#halo2}{\S~5.4.10.3}) instantiates the recipe of Theorem~\ref{zk:thm:recipe},
and each ingredient can be pointed to in the deployed source. The
Poly-IOP is a PLONKish arithmetisation---custom gates declared as
polynomial expressions, lookup arguments asserting that witness
values lie in a declared table, and a permutation argument enforcing
copy constraints. The polynomial commitment is the inner-product-argument
scheme over a single elliptic curve, \emph{transparent}---the
generators are derived by hashing to the curve from the fixed
domain string \texttt{Halo2-Parameters}, so the setup holds no
trapdoor---with logarithmic-size openings. Fiat--Shamir over a
concrete hash makes the argument non-interactive: the transcript is
a running BLAKE2b state into which every prover message is
absorbed, domain-separated by type, and out of which every
challenge is squeezed---the entire-prefix discipline of
Definition~\ref{zk:def:fs} enforced by construction.

The linear cost of each IPA opening check
(Proposition~\ref{commitments:prop:ipa-props}) makes the deployed
verifier's work linear in the commitment dimension: Halo~2 has
succinct \emph{proofs} but does not meet the strict
sublinear-verifier clause of Definition~\ref{zk:def:snark}, and it
is called a zk-SNARK under the broader convention that asks
succinctness of the proof alone. The system's signature
innovation---the technique that named
\emph{Halo}---is \emph{accumulation}: the one expensive step, the
\(O(d)\) final check of each IPA opening, is deferred and passed
along rather than paid inside a circuit, one proof attesting to
the verification of another---which is what makes \emph{recursive}
composition possible without a trusted setup. Deployed Zcash does
not use recursion: each Orchard proof is verified directly, and
what the deployed verifier takes from the Halo idea is the deferral
alone, exploited for \emph{batching}---many proofs' deferred final
checks collapsing into one multiscalar multiplication
(consumed for Orchard bundles as described in
Remark~\ref{commitments:rem:kzg-vs-ipa}). These claims rest on model and
composition hypotheses beyond the generic recipe
(Remark~\ref{zk:rem:rom-caveats}).
\end{remark}

Three load-bearing ideas from this section carry forward into
everything the upper volumes build. \emph{Rewinding}
(Remark~\ref{zk:rem:rewinding}) is the universal device: the same
cloning of an algorithm's state powers the extractor that defines
knowledge and the simulator that defines zero knowledge.
\emph{The layered split} of Theorem~\ref{zk:thm:recipe}---an
unconditionally sound information-theoretic core, wrapped in a
cryptographically enforced commitment layer---is the architecture
of every modern proof system, and the right frame for auditing one:
algebraic bugs and cryptographic bugs live in different layers.
And \emph{knowledge soundness}, not mere soundness, is what an
application stakes its security on: a shielded protocol does not
care that a spendable note exists; it cares that the prover has
one.

\section{Merkle trees and commitments to sets}\label{sec:merkle}

The adversary of this final section tells lies about a set. Her first
lie is \emph{inclusion}: she claims that a value sits in a committed
collection it never entered---a payment out of thin air, spent against
a ledger that never recorded it. Her second lie is
\emph{equivocation}: she publishes one short commitment and later
produces two different collections that both explain it, so that the
commitment pinned down nothing. Her third role is not a lie but a
gaze: as an \emph{observer} she watches an honest member prove that
it belongs, and learns \emph{which} member spoke---enough, in a
payment system, to trace every coin. This section defeats all three
with a single construction, built from nothing beyond the
collision-resistant hash of \S\ref{sec:hash}, and closes the volume by
spending everything the earlier sections banked: the commitments of
\S\ref{sec:commitments} become the leaves, the algebraic hashes of
\S\ref{sec:hash} and \S\ref{sec:commitments} become the internal
nodes, and the zero-knowledge arguments of \S\ref{sec:zk} hide the
path from leaf to root.

The plan runs from public to private. We first fix the two statements
a verifier might want proved about a large collection, and recall the
compression function and its collision game. The binary Merkle tree
follows, with its root as a binding commitment to the whole list; a
worked depth-two tree keeps every index honest. Authentication paths
then give logarithmic-size membership proofs, complete and sound; the
soundness proof introduces the one technique this section reuses
until the end---extracting a hash collision from the first level at
which two disagreeing computations agree. Layer and domain separation
keep that extraction honest against structural confusion between
leaves and nodes, with a deployed failure as the cautionary tale.
Append-only trees and the frontier make the structure grow in
logarithmic time; the vector-commitment abstraction names exactly what
the root achieves; and zero-knowledge paths convert the public proof
into a private one. The capstone assembles an Orchard shielded spend
end to end.

\subsection{Two statements about a set}\label{merkle:sec-statements}

The setting is asymmetric. A prover holds a large collection of
values---transactions in a block, public keys in a directory, records
in a log; for this volume, the ever-growing set of note commitments in
a shielded payment system. A verifier wants conviction about that
collection without storing it, indeed without ever seeing most of it.
Two statements cover the uses that matter here.

\begin{itemize}
\item \emph{Membership.} ``Value \(v\) occurs at position \(i\) of the
  list''---or, forgetting the position, ``\(v\) belongs to the set.''
\item \emph{Integrity.} ``The collection is exactly the one committed
  to earlier, unaltered.''
\end{itemize}

What is wanted is a single short \emph{commitment} that determines the
whole collection, together with short \emph{membership proofs}:
``short'' meaning of size independent of the number of elements, or at
worst logarithmic in it. A commitment of that shape makes integrity a
one-comparison check, and membership a small certificate that travels
with the value it certifies. The construction achieving this---the
\emph{Merkle tree}, due to Ralph Merkle---needs nothing more than a
collision-resistant hash function, and it is the membership backbone
of Orchard: every shielded spend proves, against a published root,
that some note commitment lies in the tree.

Notation is as fixed in \S\ref{hash:subsec:syntax} and carried
throughout the volume: \(\{0,1\}^{*}\) is the set of finite binary
strings, \(\{0,1\}^{n}\) the strings of length exactly \(n\),
\(x \,\|\, y\) concatenation, and \(\lambda\) the security parameter.
All algorithms are PPT in \(\lambda\), and \(\negl(\lambda)\) denotes
a negligible function (\emph{Math Guide}, \S``Polynomial, exponential,
and negligible functions'').

\subsection{The compression function and collision resistance}
\label{merkle:sec-compression}

The tree consumes one ingredient, which we recall in the keyed form
that makes its security a meaningful asymptotic statement.

\begin{definition}[Compression function]
\label{merkle:def-compression}
Let \(H \colon \mathcal{K} \times \{0,1\}^{*} \to \{0,1\}^{\ell}\) be
a keyed hash family exactly as in Definition~\ref{hash:def:keyed}: a
key \(s \in \mathcal{K}\) is sampled publicly at random and fixed,
the digest length is \(\ell = \ell(\lambda)\), and
\(H_{s} := H(s, \cdot)\) is the family member in force. When the
input length is restricted to a fixed multiple of the output length,
the family is a \emph{compression function}; a \emph{two-to-one}
compression function has signature
\[
h_{s} \colon \{0,1\}^{2\ell} \longrightarrow \{0,1\}^{\ell},
\]
taking two digests to one.
\end{definition}

The key \(s\) models the choice of a concrete function from a family,
which is what rescues collision resistance from the hardcoded-collision
triviality discussed in \S\ref{hash:subsec:syntax}: for a single
fixed function a colliding pair exists and some small program outputs
it, so only the keyed formulation supports a clean quantifier. In
practice one deploys a fixed standardised function---or, inside
arithmetic circuits, an algebraic hash in the sense of
\S\ref{hash:subsec:poseidon} and
\S\ref{commitments:subsec:sinsemilla}---and we suppress \(s\) from the
notation, writing \(h\) and \(H\) plain.

The security property the tree needs is collision resistance, exactly
as in Definition~\ref{hash:def:col} (\S\ref{hash:subsec:games}): the
game \(\mathsf{COL}^{\mathcal{A}}_{H}\) hands the adversary the key
\(s\), the adversary outputs a pair \((x, x')\), and it wins if
\(x \neq x'\) while \(H_{s}(x) = H_{s}(x')\); the family is
collision resistant if every PPT adversary wins with probability
\(\negl(\lambda)\). This is the strongest of the three classical hash
notions (Proposition~\ref{hash:prop:hierarchy}), and the generic
birthday attack caps an \(\ell\)-bit digest at \(\ell/2\) bits of
collision security (Lemma~\ref{hash:lem:birthday})---whence the
\(256\)-bit outputs used at the \(128\)-bit level here. Merkle trees
require \emph{full} collision resistance, not merely second-preimage
resistance, for a reason the soundness theorem
(Theorem~\ref{merkle:thm-soundness}) makes exact: the party who
assembles the tree may itself be the adversary, choosing every leaf,
which is precisely the both-inputs-adversarial regime that
Remark~\ref{hash:rem:which-notion} assigns to collision resistance.

\subsection{Binary Merkle trees}\label{merkle:sec-tree}

The construction hashes the collection pairwise, level by level, until
one digest remains.

\begin{definition}[Binary Merkle tree]\label{merkle:def-tree}
Fix a two-to-one compression function
\(h \colon \{0,1\}^{2\ell} \to \{0,1\}^{\ell}\), a depth \(d \in \N\),
and an ordered list \(L = (L_0, \dots, L_{2^d - 1})\) of \(2^d\)
leaf values in \(\{0,1\}^{\ell}\). The \emph{Merkle tree of depth
\(d\) over \(L\)} is the complete binary tree with \(2^d\) leaves
whose nodes carry labels in \(\{0,1\}^{\ell}\), defined level by
level from the leaves (level \(0\)) to the root (level \(d\)) by
\begin{equation}\label{merkle:eq-recurrence}
N_{0,j} = L_j \quad (0 \le j < 2^d),
\qquad
N_{t+1,j} = h\bigl(N_{t,2j} \,\|\, N_{t,2j+1}\bigr)
\quad (0 \le t < d,\; 0 \le j < 2^{d-t-1}).
\end{equation}
Level \(t\) has \(2^{d-t}\) nodes, and \(N_{t,j}\) is the \(j\)-th
node from the left at level \(t\). The unique level-\(d\) node is the
\emph{root}:
\(\mathsf{root}(L) := N_{d,0}\).
\end{definition}

\begin{example}[A depth-two tree]\label{merkle:exp-depth2}
With \(d = 2\) and leaves \(L_0, L_1, L_2, L_3\), the tree has seven
nodes: the four leaves, the two level-one nodes
\[
N_{1,0} = h(L_0 \,\|\, L_1), \qquad N_{1,1} = h(L_2 \,\|\, L_3),
\]
and the root
\[
\mathsf{root}(L) = N_{2,0}
= h\bigl(\, h(L_0 \,\|\, L_1) \,\|\, h(L_2 \,\|\, L_3) \,\bigr).
\]
The single \(\ell\)-bit root summarises the whole list of four
\(\ell\)-bit strings; the example carries through the next two
subsections.
\end{example}

A list whose length \(m\) is not a power of two is handled by fixing
the depth \(d\) with \(2^d \ge m\) and padding the remaining
\(2^d - m\) leaf slots with a canonical value---the all-zero string,
say, or the hash of the empty string. Padding must not create
ambiguity: two different lists must never yield the same padded leaf
sequence, else the root fails to determine the list before any
cryptography is consulted. Binding the true length \(m\) into the
root's computation---one further compression of the top node with an
encoding of \(m\)---removes any such ambiguity. An
alternative is the \emph{unbalanced} tree, in which a node with a
single child passes its child's label upward unchanged; it is
convenient but subtle, and we return to it alongside second-preimage
attacks in Remark~\ref{merkle:rem-unbalanced}.

The root is short; the claim that it \emph{commits} to the list is a
security property, stated as a game in the tradition of
Definition~\ref{model:def:game}.

\begin{definition}[Merkle commitment; binding]\label{merkle:def-commit}
The \emph{Merkle commitment} to a list \(L\) of \(2^d\) leaves is
\(\mathsf{root}(L) \in \{0,1\}^{\ell}\). The game
\(\mathsf{BIND}^{\mathcal{A}}_{H,d}(\lambda)\) runs as follows.
\begin{enumerate}
\item The challenger samples \(s \xleftarrow{\$} \mathcal{K}\),
  fixing the compression function \(h := H_{s}\), and sends \(s\) to
  the adversary.
\item The adversary outputs two lists \(L, L'\), each of \(2^d\)
  leaves in \(\{0,1\}^{\ell}\).
\item The game outputs \(1\) iff \(L \neq L'\) and
  \(\mathsf{root}(L) = \mathsf{root}(L')\).
\end{enumerate}
The commitment is \emph{binding} if every PPT adversary wins with
probability \(\negl(\lambda)\).
\end{definition}

\begin{proposition}[Binding of the root]\label{merkle:prop-binding}
If \(h\) is collision resistant, the Merkle commitment is binding.
Concretely, there is a deterministic procedure that, from any two
distinct equal-length lists \(L \neq L'\) with
\(\mathsf{root}(L) = \mathsf{root}(L')\), extracts a collision for
\(h\) in time linear in the tree size---one root-to-leaf walk of
\(O(d)\) hash comparisons. The reduction is tight: an adversary
winning \(\mathsf{BIND}\) with probability \(\varepsilon\) yields a
collision finder succeeding with the same probability
\(\varepsilon\).
\end{proposition}

\begin{proof}
Build both trees. The lists differ at some leaf index \(j_0\), so
\(N_{0,j_0} \neq N'_{0,j_0}\), writing \(N\) and \(N'\) for the labels
of the two trees. Follow the leaf-to-root path
\[
N_{0,j_0},\; N_{1,\floor{j_0/2}},\; N_{2,\floor{j_0/4}},\; \dots,\;
N_{d,0}
\]
in both trees. The labels differ at level \(0\) and agree at level
\(d\)---the common root---so there is a \emph{smallest} level \(t\) at
which they first agree: with \(j = \floor{j_0/2^{t}}\),
\[
N_{t,j} = N'_{t,j}, \qquad\text{yet}\qquad
N_{t-1,2j} \neq N'_{t-1,2j} \;\text{ or }\;
N_{t-1,2j+1} \neq N'_{t-1,2j+1},
\]
since one of the two children at level \(t-1\) is the path node, at
which the labels still disagree by minimality of \(t\). By the
recurrence \eqref{merkle:eq-recurrence},
\[
h\bigl(N_{t-1,2j} \,\|\, N_{t-1,2j+1}\bigr)
= N_{t,j} = N'_{t,j}
= h\bigl(N'_{t-1,2j} \,\|\, N'_{t-1,2j+1}\bigr),
\]
while the two \(2\ell\)-bit inputs differ in at least one half: a
collision for \(h\), located after at most \(d\) comparisons along
one path. The reduction wrapping this extractor runs the
\(\mathsf{BIND}\) adversary once and extracts whenever it wins, so
collision resistance forces the winning probability to be negligible.
\end{proof}

The proof's engine deserves its name now, because the section reuses
it almost verbatim twice more: \emph{a discrepancy at the leaves and
agreement at the root must, by this pigeonhole-like descent, manifest
a collision somewhere on the path between them.} There is a smallest
level at which two disagreeing computations agree; the compression
call producing that first agreement has two differing preimages. All
three security theorems of this section are this one sentence wearing
different inputs.

\subsection{Authentication paths and membership proofs}
\label{merkle:sec-auth}

% TikZ figure placeholder --- fig:merkle-path: the depth-2 tree of
% Example~\ref{merkle:exp-depth2} with the authentication path for
% leaf index 2 highlighted: path nodes N_{0,2}, N_{1,1}, N_{2,0}
% emphasised, siblings s_0 = N_{0,3} and s_1 = N_{1,0} in the accent
% colour, direction bits b_0 = 0, b_1 = 1 annotated on the edges.
% Dedicated figures pass supplies the drawing.

The root alone certifies integrity; membership needs a certificate
relating one leaf to the root. The certificate is the list of
siblings along the leaf's path upward---everything the verifier needs
to recompute the root from the leaf, and nothing more.

\begin{definition}[Authentication path]\label{merkle:def-path}
Let a depth-\(d\) tree over \(L\) be given, and let \(i\) be a leaf
index, \(0 \le i < 2^d\), with binary expansion
\(i = b_{d-1} b_{d-2} \cdots b_0\), so that bit \(b_t\) records
whether the level-\(t\) node on the leaf-to-root path is a left child
(\(b_t = 0\)) or a right child (\(b_t = 1\)). The
\emph{authentication path} (also \emph{Merkle proof} or
\emph{membership witness}) for index \(i\) is
\[
\pi_i = \bigl((s_0, b_0), (s_1, b_1), \dots, (s_{d-1}, b_{d-1})\bigr),
\]
where \(s_t\) is the label of the \emph{sibling} of the level-\(t\)
node on the path. Explicitly, with \(j_t = \floor{i/2^{t}}\) the
path node's index at level \(t\),
\[
s_t =
\begin{cases}
N_{t,\,j_t + 1} & \text{if } b_t = 0
  \quad\text{(sibling on the right)},\\
N_{t,\,j_t - 1} & \text{if } b_t = 1
  \quad\text{(sibling on the left)}.
\end{cases}
\]
The path has exactly \(d\) entries.
\end{definition}

\begin{definition}[Path verification]\label{merkle:def-verify}
The verifier \(\mathsf{Verify}(\mathit{rt}, i, v, \pi_i)\), on a root
\(\mathit{rt}\), an index \(i\), a claimed leaf value
\(v \in \{0,1\}^{\ell}\), and a path
\(\pi_i = ((s_0, b_0), \dots, (s_{d-1}, b_{d-1}))\), proceeds in two
stages. First the \emph{index-consistency check}: reject unless, for
every \(t\), the bit \(b_t\) carried in \(\pi_i\) equals bit \(t\) of
the binary expansion of \(i\). Then set \(a_0 := v\) and fold the path
upward: for \(t = 0, \dots, d-1\),
\begin{equation}\label{merkle:eq-fold}
a_{t+1} :=
\begin{cases}
h(a_t \,\|\, s_t) & \text{if } b_t = 0,\\
h(s_t \,\|\, a_t) & \text{if } b_t = 1.
\end{cases}
\end{equation}
Accept iff \(a_d = \mathit{rt}\). We call \(a_d\) \emph{the root
recomputed from \((i, v, \pi_i)\)}. The bit \(b_t\) places the
running value on the side matching the argument order of the
tree-building recurrence: a left child is hashed on the left.
\end{definition}

\begin{proposition}[Completeness]\label{merkle:prop-completeness}
For every list \(L\), every index \(i\), and \(\pi_i\) the genuine
authentication path for \(i\) in the tree over \(L\),
\[
\mathsf{Verify}\bigl(\mathsf{root}(L),\, i,\, L_i,\, \pi_i\bigr)
= \mathsf{accept}.
\]
\end{proposition}

\begin{proof}
The index-consistency check passes by construction, the genuine path
carrying exactly the bits of \(i\). We show by induction on \(t\)
that \(a_t = N_{t, \floor{i/2^{t}}}\), the label of the level-\(t\)
node on the leaf-to-root path. The base is
\(a_0 = L_i = N_{0,i}\). For the step, the level-\(t\) path node and
its sibling \(s_t\) are precisely the two children of the
level-\((t+1)\) path node, and the fold \eqref{merkle:eq-fold} places
\(a_t\) on the side dictated by \(b_t\) with \(s_t\) on the other
side---reproducing exactly the two arguments of \(h\) in the
recurrence \eqref{merkle:eq-recurrence}, so
\(a_{t+1} = N_{t+1, \floor{i/2^{t+1}}}\). At \(t = d\) this reads
\(a_d = N_{d,0} = \mathsf{root}(L)\), and the verifier accepts.
\end{proof}

\begin{proposition}[Logarithmic proof size]\label{merkle:prop-size}
An authentication path in a depth-\(d\) tree consists of \(d\)
sibling labels and \(d\) direction bits, \(d\ell + d\) bits in all,
and verification performs exactly \(d\) evaluations of \(h\).
Choosing the minimal sufficient depth \(d = \ceil{\log_2 m}\) for a
list of \(m \ge 2\) leaves (any \(2^d \ge m\) suffices) yields proofs
of \(\ceil{\log_2 m}\,(\ell + 1)\) bits---logarithmic in the length
of the list.
\end{proposition}

\begin{proof}
One sibling label, one direction bit, and one compression per level,
over the \(d\) levels above the leaves.
\end{proof}

The concrete scale is worth pausing on. For a set of a billion
elements---\(m = 2^{30} = 1\,073\,741\,824\)---a membership proof
consists of \(30\) hash values and \(30\) direction bits. With a
\(256\)-bit hash the sibling labels occupy
\(30 \cdot 256 = 7680\) bits, which is \(960\) bytes; direction bits
included, \(964\) bytes---under a kilobyte, however the billion
elements are arranged. (The byte counts are the output of the
volume's companion script.) This logarithmic succinctness is the central
objective: the verifier's storage is one digest, the certificate fits
in a packet, and neither grows meaningfully as the set does.

\begin{example}[A depth-two membership proof]\label{merkle:exp-proof}
Continue Example~\ref{merkle:exp-depth2}, and prove membership of
\(L_2\) at index \(i = 2 = (10)_2\), so \(b_0 = 0\) and \(b_1 = 1\).
At level \(0\) the path node \(N_{0,2} = L_2\) is a left child, with
sibling \(s_0 = N_{0,3} = L_3\). At level \(1\) the path node
\(N_{1,1} = h(L_2 \,\|\, L_3)\) is a right child, with sibling
\(s_1 = N_{1,0} = h(L_0 \,\|\, L_1)\). The path is
\[
\pi_2 = \bigl((L_3,\, 0),\; (h(L_0 \,\|\, L_1),\, 1)\bigr).
\]
Verification folds upward:
\(a_1 = h(L_2 \,\|\, L_3)\), then
\(a_2 = h\bigl(h(L_0 \,\|\, L_1) \,\|\, a_1\bigr) = \mathsf{root}(L)\),
and the verifier accepts. Note what the verifier never saw: neither
\(L_0\) nor \(L_1\) themselves, only their digest
\(h(L_0 \,\|\, L_1)\)---yet it is convinced that \(L_2\) sits at
index \(2\) of the committed list.
\end{example}

\subsection{Soundness: forging a path implies a collision}
\label{merkle:sec-soundness}

Completeness says honest paths verify; the security question is
whether dishonest ones can. The threat is the inclusion lie of the
section's opening: an adversary who convinces the verifier, against
the genuine root, that position \(i\) holds a value it does not.

\begin{definition}[Path forgery]\label{merkle:def-forgery}
Fix a list \(L\) of \(2^d\) leaves with root
\(\mathit{rt} = \mathsf{root}(L)\). The game
\(\mathsf{PATHFORGE}^{\mathcal{A}}_{H,L}(\lambda)\) runs as follows.
\begin{enumerate}
\item The challenger samples \(s \xleftarrow{\$} \mathcal{K}\) to fix
  \(h := H_{s}\), builds the tree over \(L\), and sends \(s\)
  together with \(L\)---equivalently, the whole tree---to the
  adversary.
\item The adversary outputs an index \(i\) and a pair
  \((v^{*}, \pi^{*})\).
\item The game outputs \(1\) iff
  \(\mathsf{Verify}(\mathit{rt}, i, v^{*}, \pi^{*}) =
  \mathsf{accept}\) and \(v^{*} \neq L_i\).
\end{enumerate}
A winning output is a \emph{path forgery at index \(i\)}: the
verifier accepts, against the honestly computed root, a leaf value
different from the genuine one at position \(i\).
\end{definition}

\begin{theorem}[Soundness of Merkle membership proofs]
\label{merkle:thm-soundness}
Let \(h\) be collision resistant. Then no PPT adversary, given the
list \(L\) and the whole tree, produces a path forgery except with
probability \(\negl(\lambda)\). In constructive form: there is a
deterministic extractor that from any list \(L\), index \(i\), and
accepting forgery \((v^{*}, \pi^{*})\) with \(v^{*} \neq L_i\)
outputs a collision for \(h\), in time linear in the tree
size---recomputing the genuine leaf-to-root labels from \(L\)
dominates; given the built tree, one verifier run plus a
level-by-level comparison, \(O(d)\) work in all, completes the
extraction.
\end{theorem}

\begin{proof}
Run the verifier's fold \eqref{merkle:eq-fold} on
\((i, v^{*}, \pi^{*})\), obtaining the chain
\(a^{*}_0 = v^{*},\, a^{*}_1,\, \dots,\, a^{*}_d = \mathit{rt}\)
(the final equality because the forgery is accepting). Alongside it
place the genuine chain
\(a_t = N_{t, \floor{i/2^{t}}}\), the labels of the leaf-to-root path
in the honest tree; by the induction of
Proposition~\ref{merkle:prop-completeness} this is the chain the fold
produces on the genuine inputs, and \(a_d = \mathit{rt}\).

The two chains disagree at the bottom,
\(a^{*}_0 = v^{*} \neq L_i = a_0\), and agree at the top,
\(a^{*}_d = a_d = \mathit{rt}\). Hence
\[
t^{*} := \min\,\{\, t : a^{*}_t = a_t \,\} \;\in\; \{1, \dots, d\}
\]
exists, and \(a^{*}_{t^{*}-1} \neq a_{t^{*}-1}\) by minimality.
Because the forgery is accepting, the index-consistency check of
Definition~\ref{merkle:def-verify} forces the direction bits carried
in \(\pi^{*}\) to equal the bits of \(i\)---the same bits the genuine
path carries---so at level \(t^{*}-1\) \emph{both} chains place their
running value on the same side of the compression call: both compute
\(h(\cdot \,\|\, \text{sibling})\) if \(b_{t^{*}-1} = 0\), both
\(h(\text{sibling} \,\|\, \cdot)\) if \(b_{t^{*}-1} = 1\). Form the
two \(2\ell\)-bit preimages accordingly,
\[
X^{*} =
\begin{cases}
a^{*}_{t^{*}-1} \,\|\, s^{*}_{t^{*}-1} & \text{if } b_{t^{*}-1} = 0,\\
s^{*}_{t^{*}-1} \,\|\, a^{*}_{t^{*}-1} & \text{if } b_{t^{*}-1} = 1,
\end{cases}
\qquad
X =
\begin{cases}
a_{t^{*}-1} \,\|\, s_{t^{*}-1} & \text{if } b_{t^{*}-1} = 0,\\
s_{t^{*}-1} \,\|\, a_{t^{*}-1} & \text{if } b_{t^{*}-1} = 1,
\end{cases}
\]
where \(s^{*}_{t^{*}-1}\) is the sibling claimed in \(\pi^{*}\) and
\(s_{t^{*}-1}\) the genuine sibling. Then
\[
h(X^{*}) = a^{*}_{t^{*}} = a_{t^{*}} = h(X),
\qquad X^{*} \neq X,
\]
the inequality because the two preimages differ in the half occupied
by the running value. The pair \((X^{*}, X)\) is a collision for
\(h\), found deterministically in \(O(d)\) work once the genuine
chain is at hand---a read-off from the built tree, or \(2^{d} - 1\)
hashes from the bare list \(L\).
\end{proof}

The reduction wrapping the extractor is immediate. If a PPT adversary
\(\mathcal{A}\) produces path forgeries with probability
\(\varepsilon(\lambda)\), the collision finder \(\mathcal{B}\) that
receives the key, samples or receives \(L\), runs \(\mathcal{A}\),
and applies the extractor outputs a collision whenever
\(\mathcal{A}\) succeeds---hence with probability
\(\varepsilon(\lambda)\), at an overhead of building the tree over
\(L\), which \(\mathcal{B}\) does in any case to play the challenger,
plus one verifier run and \(O(d)\) comparisons. Collision resistance
of \(h\) therefore forces
\(\varepsilon(\lambda) \in \negl(\lambda)\), and the reduction is
tight: no advantage is lost in the translation.

\begin{remark}[Scope: a fixed honest root]
\label{merkle:rem-soundness-scope}
Theorem~\ref{merkle:thm-soundness} is a statement about a \emph{fixed
honest root}: relative to a root that genuinely is
\(\mathsf{root}(L)\), the only leaf value admitting an accepting path
at position \(i\) is \(L_i\), up to a collision in \(h\). It does not
by itself prevent an adversary who is free to \emph{choose} the root
from committing to a maliciously crafted tree; that threat is what
binding (Proposition~\ref{merkle:prop-binding}) rules out. With an
untrusted party supplying the root, both properties are necessary and
they divide the labour cleanly: binding ensures the root determines
at most one list, and path soundness ensures openings are faithful to
it. Together they make the root a sound, binding commitment with
succinct openings---the abstraction \S\ref{merkle:sec-vc} names, where
a single argument will deliver both faces at once.
\end{remark}

\subsection{Layer and domain separation}\label{merkle:sec-domain}

The soundness theorem has a structural blind spot, inherited from an
innocent-looking convenience. Suppose the leaves are arbitrary-length
data, hashed into \(\{0,1\}^{\ell}\) by the \emph{same} function used
internally. Then an internal label \(h(A \,\|\, B)\) is itself a
well-formed leaf value---it has the right length---and one string can
play both roles. An adversary exploits the ambiguity directly: she
presents the internal node \(N_{1,0} = h(L_0 \,\|\, L_1)\) as if it
were a \emph{leaf} near the top of a shallower tree, equipped with
the correspondingly shorter authentication path, and a verifier who
does not know the depth accepts. The root now has a second list
explaining it---a genuine second-preimage attack, with no collision
of \(h\) anywhere in sight. Absent a way for a label to assert ``I am
a leaf'' or ``I am an internal node'', the leaf-to-root structure
that Theorem~\ref{merkle:thm-soundness} relies on collapses.

The repair is the domain-separation discipline of
\S\ref{hash:subsec:domainsep}, applied \emph{within} a single
structure: distinct structural positions get distinct tags.

\begin{definition}[Layer and domain separation]\label{merkle:def-domain}
A Merkle construction has \emph{domain separation} if the function
applied at distinct structural positions is forced to differ, by
prepending a position-dependent tag to every hash input.
Concretely, from a variable-input-length hash \(H\) define
\[
\mathsf{LeafHash}(v) := H(\mathtt{0x00} \,\|\, v),
\qquad
\mathsf{NodeHash}_t(A, B) :=
H(\mathtt{0x01} \,\|\, \langle t \rangle \,\|\, A \,\|\, B),
\]
where the byte tags \(\mathtt{0x00}\) and \(\mathtt{0x01}\)
distinguish the \emph{layer kind}---leaf versus internal node---and
\(\langle t \rangle\) is a fixed-length encoding of the level index,
distinguishing internal nodes by depth. The tree recurrence becomes
\[
N_{0,j} = \mathsf{LeafHash}(L_j),
\qquad
N_{t+1,j} = \mathsf{NodeHash}_{t+1}\bigl(N_{t,2j},\, N_{t,2j+1}\bigr),
\]
and the verifier's fold applies the same tagged functions at the same
positions.
\end{definition}

\begin{proposition}[Separation forces within-domain collisions]
\label{merkle:prop-domain}
With the tagging of Definition~\ref{merkle:def-domain}, any collision
extracted by the binding proof
(Proposition~\ref{merkle:prop-binding}) or the soundness proof
(Theorem~\ref{merkle:thm-soundness}) is a collision for \(H\) between
two inputs \emph{sharing the same leading tag}: a leaf preimage
against a leaf preimage, or a level-\(t\) node preimage against a
level-\(t\) node preimage. In particular, no accepting forgery can
substitute an internal node for a leaf, or a node at one level for a
node at another, without exhibiting a within-domain collision of
\(H\)---which collision resistance forbids.
\end{proposition}

\begin{proof}[Proof sketch]
The extractors of both proofs locate the smallest level at which two
disagreeing computations agree, and collide the two preimages of the
compression call at that position. Under the tagged recurrence, the
verifier's recomputation applies the position's prescribed function
exactly as the honest tree does: at the leaves both preimages pass
through \(\mathsf{LeafHash}\), at level \(t \ge 1\) both pass through
\(\mathsf{NodeHash}_t\). Both colliding preimages therefore begin
with the same tag---\(\mathtt{0x00}\), or
\(\mathtt{0x01} \,\|\, \langle t \rangle\)---and the collision lies
within a single domain. A cross-domain confusion, by contrast, would
require two inputs with \emph{different} leading tags and equal
\(H\)-images, which is again an \(H\)-collision; either way the
adversary has broken collision resistance, and the shallower-tree
attack above is foreclosed because a leaf preimage and a node
preimage can no longer be the same string.
\end{proof}

\begin{remark}[Unbalanced trees, padding, and a deployed failure]
\label{merkle:rem-unbalanced}
The subtleties are not hypothetical. Bitcoin's transaction Merkle
tree handles an odd level by \emph{duplicating} its last node before
hashing pairwise, and the duplication rule is not injective: a list
of transactions and the same list with its tail duplicated can yield
the same root. This is CVE-2012-2459, a consensus-splitting
vulnerability---an attacker could present two \emph{different}
transaction lists sharing a root, one valid and one invalid, and
nodes could be made to cache the invalid one as permanently rejected
(documented at length in Zcash's inherited validator:
\texttt{zebra-chain} crate, \texttt{src/block/\allowbreak merkle.rs},
which reproduces the upstream warning that the flawed algorithm must
not be reused; the deployed defence---rejecting any block whose
transaction hashes are not unique---lives in
\texttt{zebra-consensus}, \texttt{src/block/\allowbreak check.rs}).
The lesson generalises to three rules. First, layer-separate leaves
from internal nodes, as Definition~\ref{merkle:def-domain}
prescribes. Second, make every
padding or duplication rule \emph{injective} on admissible lists---for
instance by binding the true leaf count into the root's computation.
Third, feed \(\mathsf{NodeHash}\) only fixed-length inputs, so that
no length-extension structure of the underlying hash
(\S\ref{hash:subsec:domainsep}) can intervene between layers.
\end{remark}

\begin{remark}[Algebraic hashes in the tree; Orchard's
\(\mathsf{MerkleCRH}\)]\label{merkle:rem-poseidon}
When the Merkle path is verified inside a zero-knowledge circuit
(\S\ref{merkle:sec-zk}), arithmetic gates over \(\Fp\) evaluate the
compression function, and the choice of \(h\) dominates the cost. A
bit-oriented hash such as SHA-256 costs tens of thousands of
constraints per call; one therefore uses an \emph{algebraic} hash
with compression \(h \colon \F_p^2 \to \F_p\). Poseidon
(\S\ref{hash:subsec:poseidon}) is a low-degree permutation-based map
computable in few constraints; Sinsemilla
(\S\ref{commitments:subsec:sinsemilla}) is cheap in constraints
through lookups and incomplete curve additions, \emph{not} through
low degree. For Poseidon, layer and domain separation come not from
byte tags but from distinct field constants---initial capacity
values or round constants---fed into the permutation; the tree's
security argument is identical, with ``collision'' read as a
collision of the field-valued compression function.

Orchard's Merkle hash \(\mathsf{MerkleCRH}\) retains the tag
mechanism of Definition~\ref{merkle:def-domain} in encoded form. Each
internal call prepends the \(10\)-bit encoded layer index to the two
\(255\)-bit child encodings---a single fixed input length of
\(520\) bits (a total the volume's companion script checks), exactly
the fixed-length-per-personalisation regime in
which the deployed Sinsemilla hash is collision resistant
(\S\ref{commitments:subsec:sinsemilla})---with \emph{all} layers
sharing one Sinsemilla instance (\texttt{MerkleHashOrchard::\allowbreak
combine}, prepending \texttt{i2lebsp\_k(level)} to two
\texttt{L\_ORCHARD\_MERKLE} \(= 255\)-bit child encodings;
protocol specification \href{https://zips.z.cash/protocol/protocol.pdf\#merklecrh}{\S~5.4.1.3}). Separation between
\emph{distinct} Sinsemilla instances comes from the personalisation
string---here \texttt{z.cash:\allowbreak Orchard-MerkleCRH}---fed to
hash-to-curve to derive the instance-specific base point \(Q\), while
the generator table \(S\) is shared across all instances
(\texttt{HashDomain::\allowbreak new} derives \(Q\) from the
personalisation via the \texttt{z.cash:\allowbreak SinsemillaQ}
hash-to-curve domain; the table derives from the fixed
\texttt{z.cash:\allowbreak SinsemillaS}). Sinsemilla is collision
resistant under a discrete-logarithm assumption
(Proposition~\ref{commitments:prop:sinsemilla-cr}), so the tree's
hashing and its commitment role fuse into one Pedersen-like map: the
same structure that compresses the children also binds them.
\end{remark}

\subsection{Append-only and incrementally updatable trees}
\label{merkle:sec-frontier}

A shielded pool's set of note commitments only ever grows, and it
grows one leaf at a time. Rebuilding a depth-\(32\) tree from scratch
on every insertion would cost \(2^{32}\) hashes; the structure of the
tree permits the update in \(d\) hashes and \(O(d)\) state, and the
state has a name.

\begin{definition}[Append-only Merkle tree]\label{merkle:def-append}
Fix a maximum depth \(d\), giving capacity \(2^d\), and the
level-indexed \(\mathsf{NodeHash}\) functions of
Definition~\ref{merkle:def-domain}. Leaf values here enter as
level-\(0\) labels directly, with no \(\mathsf{LeafHash}\): the
leaves are already digests, and the fixed depth together with the
level tags (Remark~\ref{merkle:rem-poseidon}) carries the layer
separation---the convention of deployed Orchard, whose leaf is the
extracted coordinate itself (Example~\ref{merkle:exp-orchard}). An
\emph{append-only} (or \emph{incremental}) \emph{Merkle tree}
maintains a position counter \(m\)---the number of leaves inserted so
far, \(0 \le m \le 2^d\)---and supports one mutating operation
\(\mathsf{Append}(v)\), which places \(v\) at leaf index \(m\) and
increments \(m\). Empty leaf slots \(m, \dots, 2^d - 1\) carry a
fixed \emph{empty-leaf} value \(\bot_0\), and the labels of entirely
empty subtrees follow by precomputation:
\[
\bot_{t+1} := \mathsf{NodeHash}_{t+1}(\bot_t, \bot_t),
\qquad 0 \le t < d.
\]
The \emph{root after \(m\) appends}, written \(\mathsf{root}_m\), is
the root of the depth-\(d\) tree, under this recurrence, whose first
\(m\) level-\(0\) labels are the inserted values and whose remaining
ones are \(\bot_0\).
\end{definition}

\begin{definition}[Frontier]\label{merkle:def-frontier}
The \emph{frontier} of an append-only tree holding \(m\) leaves is
the minimal state needed to (i) compute the current root and (ii)
update incrementally on the next \(\mathsf{Append}\). For each level
\(t \in \{0, \dots, d-1\}\) it holds at most one \emph{carry}: the
label of the left child of the next node to be completed at level
\(t+1\), retained precisely when leaf index \(m\) has a \(1\) bit at
position \(t\)---a completed left sibling awaiting its right
neighbour. Equivalently, writing \(m\) in binary as
\(b_{d-1} \cdots b_0\), the frontier stores, for each \(t\) with
\(b_t = 1\), the already-finalised left-sibling label \(f_t\) at
level \(t\) along the rightmost filled path. Its size is at most
\(d\) labels---one per set bit of \(m\), hence \(O(\log m)\) state.
\end{definition}

The right picture is arithmetic. The binary representation of \(m\)
decomposes the inserted leaves into completed perfect subtrees of
distinct heights---one of height \(t\) for each set bit
\(b_t\)---and the frontier holds exactly the roots of those subtrees:
the pieces that are ``complete and waiting for a right neighbour''.
An append is then the Merkle analogue of incrementing a binary
counter: a carry at level \(t\) corresponds to combining two
height-\(t\) subtrees into one of height \(t+1\), and the chain of
carries is as long as the run of trailing \(1\)-bits of \(m\).

% TikZ figure placeholder --- fig:frontier: the frontier as a binary
% counter across three appends at depth 3. Three panels (m = 1, 2, 3)
% of the same depth-3 tree outline; filled leaves v_0, v_1, v_2;
% carried labels f_0/f_1 highlighted at their levels; below each
% panel the binary counter 001 -> 010 -> 011 with set bits aligned to
% the occupied frontier levels. Dedicated figures pass supplies the
% drawing.

\begin{proposition}[Incremental append in logarithmic time]
\label{merkle:prop-incremental}
From the frontier of an append-only tree holding \(m < 2^d\) leaves,
\(\mathsf{Append}(v)\) computes the new frontier and the new root
\(\mathsf{root}_{m+1}\) using at most \(d\) evaluations of
\(\mathsf{NodeHash}\) and \(O(d)\) additional work, while storing
only \(O(d)\) labels throughout.
\end{proposition}

\begin{proof}
The algorithm is carry propagation followed by an upward fold. Set
the running label \(c \gets v\) at level \(t \gets 0\). While bit
\(t\) of \(m\) is \(1\), the new path node at level \(t\) is a right
child whose left sibling is the stored carry \(f_t\): set
\(c \gets \mathsf{NodeHash}_{t+1}(f_t, c)\), clear the carry
\(f_t\), and increment \(t\). When bit \(t\) of \(m\) is \(0\), store
\(c\) as the new carry \(f_t\) and stop. Writing \(z\) for the number
of trailing \(1\)-bits of \(m\), the carry phase costs exactly \(z\)
hashes and deposits \(f_z\) at level \(z\)---the root of the
just-completed height-\(z\) subtree containing the new leaf, matching
the binary increment \(m \mapsto m + 1\), which clears \(z\) trailing
ones and sets bit \(z\).

To obtain \(\mathsf{root}_{m+1}\), fold upward from level \(z\) with
the just-deposited \(f_z\) as the running value---it is the path
label at level \(z\), not a sibling. For \(t = z, \dots, d - 1\),
combine the running value with the stored carry \(f_t\) as left
sibling when bit \(t\) of \(m\) is \(1\)---those carries lie above
the cleared run and survive the append untouched---and with the
precomputed empty-subtree label \(\bot_t\) as right sibling
otherwise; at level \(z\) itself bit \(z\) of \(m\) is \(0\), so the
first combination is with \(\bot_z\), never with the deposited carry.
One hash per level makes \(d - z\) hashes in all (zero if \(z = d\),
when the tree has just filled). The total is exactly
\(z + (d - z) = d\) evaluations of \(\mathsf{NodeHash}\), and the
retained state is the at-most-\(d\) carries: precisely the frontier
for \(m + 1\) leaves.
\end{proof}

\begin{example}[The frontier as a binary counter]
\label{merkle:exp-frontier}
Take depth \(d = 3\) (capacity \(8\)) and insert \(v_0, v_1, v_2\) in
turn; the frontier evolves exactly as the binary counter of the leaf
count \(m\).
\begin{itemize}
\item After \(v_0\), \(m = 1 = (001)_2\): one completed height-\(0\)
  subtree; the frontier is \(\{\, f_0 = v_0 \,\}\).
\item After \(v_1\), \(m = 2 = (010)_2\): leaf \(v_1\) at index \(1\)
  is a right child, triggering one carry
  \(c = \mathsf{NodeHash}_1(v_0, v_1)\)---a completed height-\(1\)
  subtree; the frontier is \(\{\, f_1 = c \,\}\), with \(f_0\)
  cleared.
\item After \(v_2\), \(m = 3 = (011)_2\): leaf \(v_2\) at index \(2\)
  is a left child and becomes a new height-\(0\) carry; the frontier
  is \(\{\, f_1 = \mathsf{NodeHash}_1(v_0, v_1),\; f_0 = v_2 \,\}\)
  ---exactly the two set bits of \(3\).
\end{itemize}
Each append touched only the carry chain to the new leaf, never the
whole tree. (The trace involves no numeric hash values and follows
mechanically from the algorithm of
Proposition~\ref{merkle:prop-incremental}; the volume's companion
script replays it symbolically and checks the occupied frontier
levels against the set bits of \(m\) at every step.)
\end{example}

Appending changes the root, which raises a question the definition of
soundness so far does not answer: which root is a proof valid
against?

\begin{definition}[History binding]\label{merkle:def-append-sound}
An append-only tree is \emph{history-binding} when the surrounding
protocol treats every historical root as an acceptable commitment in
its own right: each membership proof verifies against the root it was
formed for, and appending further leaves never invalidates a
previously valid \((\mathit{root}, \mathit{leaf}, \mathit{path})\)
triple for that root. The soundness of each individual proof is
exactly Theorem~\ref{merkle:thm-soundness} applied to the
corresponding fixed root.
\end{definition}

Deployed Orchard is history-binding by construction. The chain
records each block's note-commitment-tree root---its \emph{anchor}
---and a spend may prove membership against \emph{any} recorded
historical anchor: consensus requires the anchor named by a
transaction to be the root of the Orchard note-commitment tree as it
stood at the end of some earlier block, checked by membership in the
set of recorded anchors (protocol specification \href{https://zips.z.cash/protocol/protocol.pdf\#actions}{\S~3.7}, `Action Transfers and their
Descriptions').
Per-proof soundness is then the fixed-root soundness theorem, once
per anchor. The frontier is likewise the deployed representation of
the growing tree (\texttt{incrementalmerkletree} crate,
types \texttt{Frontier} and
\texttt{NonEmptyFrontier}; the empty-subtree labels are the
\texttt{EMPTY\_ROOTS} table of \texttt{orchard}).

\begin{remark}[Why append-only, and what it costs]
\label{merkle:rem-append-only}
Deletion or in-place update could be supported at \(O(\log m)\) per
operation, given the relevant siblings; the restriction to appends is
a design choice, made for a monotonicity it buys. Once a leaf is
inserted it is never removed, so a membership proof valid against the
root at insertion time remains valid against \emph{that historical
root} forever. This is exactly what a shielded pool requires---a
note, once added, can always be spent later by proving against some
past anchor---and it spares the spender from tracking a moving
target as the tree grows under everyone else's insertions. The cost
falls on the verifier, who must accept a window of historical roots
rather than a single current one; the previous paragraph is that
window, deployed.
\end{remark}

\subsection{Vector commitments}\label{merkle:sec-vc}

The section has been proving properties of one construction; it pays
to name the abstraction the construction inhabits, both to state
precisely what the root achieves and to mark the boundary with its
relatives. The abstraction is a commitment to an \emph{ordered
vector} that opens one position at a time. (Its unordered
sibling---the \emph{cryptographic accumulator}, which certifies set
membership with no notion of position---exists in the literature, but
Orchard never needs it: notes occupy positions, and the position is
part of what the spend circuit witnesses. We do not develop it.)

\begin{definition}[Vector commitment]\label{merkle:def-vc}
A \emph{vector commitment} is a tuple of PPT algorithms
\((\mathsf{Setup}, \mathsf{Commit}, \mathsf{Open}, \mathsf{Vf})\):
\begin{itemize}
\item \(\mathsf{Setup}(1^{\lambda}, n) \to \mathit{pp}\) produces
  public parameters for vectors of length \(n\);
\item \(\mathsf{Commit}(\mathit{pp}, (v_0, \dots, v_{n-1})) \to
  (C, \mathit{aux})\) outputs a short commitment and auxiliary state;
\item \(\mathsf{Open}(\mathit{pp}, i, \mathit{aux}) \to \pi_i\)
  produces a proof that position \(i\) holds \(v_i\);
\item \(\mathsf{Vf}(\mathit{pp}, C, i, v, \pi) \to
  \{\mathsf{accept}, \mathsf{reject}\}\) verifies an opening.
\end{itemize}
The scheme is \emph{correct} if honest openings always verify, and
\emph{succinct} if \(\abs{C}\) and \(\abs{\pi_i}\) are bounded by a
fixed polynomial in \(\lambda\) and \(\log n\). It is
\emph{position-binding} if every PPT adversary wins the following
game with probability \(\negl(\lambda)\): given \(\mathit{pp}\), the
adversary outputs \((C, i, v, v', \pi, \pi')\) and wins iff
\(v \neq v'\) while both
\(\mathsf{Vf}(\mathit{pp}, C, i, v, \pi)\) and
\(\mathsf{Vf}(\mathit{pp}, C, i, v', \pi')\) accept. Note that the
adversary chooses \(C\) itself: binding must hold even for
maliciously formed commitments, with no honest committer anywhere in
the game.
\end{definition}

\begin{theorem}[Merkle trees are succinct position-binding vector
commitments]\label{merkle:thm-merkle-vc}
Instantiate the syntax by: \(\mathsf{Setup}\) samples \(h\) (and
fixes \(d\) with \(n = 2^d\)); \(\mathsf{Commit}\) computes the tree
and outputs \(C = \mathsf{root}\) with \(\mathit{aux}\) the full
tree; \(\mathsf{Open}\) returns the authentication path; and
\(\mathsf{Vf} = \mathsf{Verify}\). If \(h\) is collision resistant,
this is a correct, succinct, position-binding vector commitment, with
commitment size \(\ell\) and opening size \(O(\ell \log n)\).
\end{theorem}

\begin{proof}
Correctness is completeness
(Proposition~\ref{merkle:prop-completeness}), and the size bounds are
Proposition~\ref{merkle:prop-size}. For position-binding, suppose an
adversary outputs two accepting openings \((v, \pi)\) and
\((v', \pi')\) at the same index \(i\) under the same root \(C\),
with \(v \neq v'\). Run the verifier's fold \eqref{merkle:eq-fold} on
each, obtaining two recomputed chains that disagree at the bottom
(\(v \neq v'\)) and agree at the top (both reach \(C\), since both
openings accept). The index-consistency check forces both paths to
carry the direction bits of \(i\), so at every level both chains
place their running value on the same side of the compression call.
The first-level-of-agreement extraction of
Theorem~\ref{merkle:thm-soundness} now applies verbatim to this pair
of chains and yields two distinct \(2\ell\)-bit inputs with the same
\(h\)-image: a collision. The argument needed \emph{no honest
reference tree}---it pits the two forged openings against each
other---so it binds even a maliciously chosen root, which is the
strong form of position-binding the definition demands.
\end{proof}

\begin{remark}[Position-binding versus path soundness]
\label{merkle:rem-binding-vs-sound}
The two theorems answer different verifiers. Path soundness
(Theorem~\ref{merkle:thm-soundness}) fixes an \emph{honest} root and
forbids opening a position to anything but its true value; it speaks
to a verifier who trusts where the root came from. The
vector-commitment theorem (Theorem~\ref{merkle:thm-merkle-vc})
forbids opening a single position two \emph{different} ways even
under an \emph{adversarial} root; it speaks to a verifier who
received the root from an untrusted source, and it is what
``commitment'' should mean---\(C\) determines at most one value per
position, whoever made \(C\). Both reduce to the same
first-level-of-agreement collision extraction; the only difference is
which pair of chains is collided, forged-against-genuine in the one
case and forged-against-forged in the other.
\end{remark}

\subsection{Privacy-preserving membership via zero-knowledge paths}
\label{merkle:sec-zk}

The Merkle membership proof is succinct but not private---and for a
shielded payment system, not private is fatal. Revealing the
authentication path of \S\ref{merkle:sec-auth} discloses the leaf
value itself (here, the note commitment being spent) and the index
\(i\) at which it sits, and even the sibling labels can leak
structural information about the neighbourhood of the leaf. An
observer correlating spends with the appends that created their
leaves links payments end to end: spending a note must not reveal
\emph{which} note is spent. The remedy is to prove the
\emph{existence} of a valid path in zero knowledge, treating the
entire path as a private witness and publishing only the root.

\begin{definition}[The membership relation]\label{merkle:def-zk-relation}
For the domain-separated Merkle construction of depth \(d\), with
leaves entering as level-\(0\) labels under the convention of
Definition~\ref{merkle:def-append}, define
\[
R_{\mathrm{mem}} :=
\bigl\{\, \bigl(\mathit{rt} \,;\, (v, i, \pi)\bigr) :
\mathsf{Verify}(\mathit{rt}, i, v, \pi) = \mathsf{accept} \,\bigr\},
\]
with the \emph{statement} the public root \(\mathit{rt}\) and the
\emph{witness} comprising the leaf \(v\), its index \(i\), and the
authentication path \(\pi\). A zero-knowledge argument for
\(R_{\mathrm{mem}}\) lets a prover convince a verifier that
\emph{some} leaf lies in the tree with root \(\mathit{rt}\) while
revealing none of \(v\), \(i\), or \(\pi\). We assume a zk-SNARK as
\S\ref{sec:zk} develops it (Definition~\ref{zk:def:snark}):
\emph{completeness} (honest provers convince), \emph{knowledge
soundness} (an accepting prover ``knows'' a witness, formalised by an
extractor), and \emph{zero knowledge} (the proof reveals nothing
beyond the truth of the statement, formalised by a simulator).
\end{definition}

% TikZ figure placeholder --- fig:zkpath: the statement/witness split
% of R_mem drawn on the depth-2 tree. The root rt above a horizontal
% public/private divider; below it, greyed-out leaf v at index i with
% its authentication path and direction bits, boxed and labelled
% "witness (v, i, pi)"; the root alone labelled "statement rt"; an
% arrow from witness to root through the fold (dagger). Dedicated
% figures pass supplies the drawing.

\begin{construction}[In-circuit path verification]
\label{merkle:def-circuit}
The Merkle-path circuit \(C_{\mathrm{mem}}\) of depth \(d\) takes as
public input the root \(\mathit{rt}\), and as private (witness)
inputs the leaf \(v\), the direction bits \(b_0, \dots, b_{d-1}\),
and the siblings \(s_0, \dots, s_{d-1}\). It implements the \(d\)
compression calls of the fold \eqref{merkle:eq-fold}---the level-\(t\)
call being the position's prescribed function,
\(\mathsf{NodeHash}_{t+1}\) in the domain-separated construction,
with no \(\mathsf{LeafHash}\) call since the leaf enters as the
level-\(0\) label (Definition~\ref{merkle:def-append})---each
preceded by the conditional swap that selects the argument order
\((a_t, s_t)\) or \((s_t, a_t)\) according to \(b_t\), and enforces
the single output constraint \(a_d = \mathit{rt}\). Its size is
\(d\) times the cost of one compression call plus \(O(d)\)
constraints for the swaps. The position enters \emph{only} through the bits
\(b_0, \dots, b_{d-1}\)---the binary decomposition of \(i\)---so the
index-consistency check of Definition~\ref{merkle:def-verify} is
enforced structurally: there is no separate index to disagree with
its bits. Deployed Orchard matches this shape exactly: the circuit's
Merkle gadget decomposes the witnessed position into little-endian
bits and applies a conditional swap per layer
(protocol specification \href{https://zips.z.cash/protocol/protocol.pdf\#merklepath}{\S~4.9}, `Merkle Path
Validity').
\end{construction}

\begin{theorem}[Privacy-preserving membership]
\label{merkle:thm-zk-membership}
Let \(h\) be collision resistant and let \((P, V)\) be a zk-SNARK for
\(R_{\mathrm{mem}}\)---equivalently, for satisfiability of
\(C_{\mathrm{mem}}\). Then the protocol ``the prover sends a SNARK
proof for the statement \(\mathit{rt}\); the verifier checks it'' is
a membership argument that is:
\begin{enumerate}
\item \emph{complete}: a prover holding a genuine leaf and its path
  produces an accepting proof;
\item \emph{sound}: any prover producing an accepting proof for an
  honest root \(\mathit{rt} = \mathsf{root}(L)\) knows---via the
  SNARK extractor---a witness \((v, i, \pi)\) with
  \(\mathsf{Verify}(\mathit{rt}, i, v, \pi) = \mathsf{accept}\), and
  \(v\) equals the genuine leaf \(L_i\) unless a collision in \(h\)
  is thereby found: the proven leaf is a real member of the
  committed list;
\item \emph{private}: the proof reveals nothing about \(v\), \(i\),
  or \(\pi\) beyond the fact that some valid leaf exists under
  \(\mathit{rt}\).
\end{enumerate}
\end{theorem}

\begin{proof}[Proof idea]
Completeness composes SNARK completeness with tree completeness
(Proposition~\ref{merkle:prop-completeness}): the genuine
\((L_i, i, \pi_i)\) satisfies \(C_{\mathrm{mem}}\), and the honest
prover proves it. Soundness composes the SNARK knowledge extractor
with the path-forgery-to-collision extractor of
Theorem~\ref{merkle:thm-soundness}: from an accepting prover, the
SNARK extractor produces \((v, i, \pi)\) accepted by
\(\mathsf{Verify}\); if \(v \neq L_i\), that output is precisely a
path forgery, hence yields a collision for \(h\), hence occurs with
probability \(\negl(\lambda)\). Zero knowledge carries over verbatim
from the SNARK simulator, which produces from \(\mathit{rt}\) alone a
proof indistinguishable from the real one---so the real proof cannot
be conveying \(v\), \(i\), or \(\pi\).
\end{proof}

\begin{remark}[Public root, private path; deployment trade-offs]
\label{merkle:rem-zk-design}
The asymmetry in \(R_{\mathrm{mem}}\) is the essential design point.
The root must be \emph{public} so that the verifier checks the proof
against a state everyone agrees on---a previously published
note-commitment-tree anchor---and soundness is relative to that
public root, exactly as Remark~\ref{merkle:rem-soundness-scope}
delimits. The path must be \emph{private} so that observers cannot
link the spend to a particular leaf. SNARK succinctness adds the
final ingredient: the verifier checks a proof whose size is
essentially independent of \(d\)---polylogarithmic---even though the
witness contains \(d\) siblings, so the construction \emph{compresses}
the membership proof while simultaneously \emph{hiding} it. This is
the precise sense in which a Merkle root over note commitments,
opened in zero knowledge, is a privacy-preserving membership
primitive.

Deployment chooses where to pay for that verification. Pairing-based
schemes such as KZG (\S\ref{commitments:subsec:poly}) achieve
constant-size proofs with fast verification, at the price of a
trusted setup. Halo~2's IPA-based proofs grow logarithmically in
size, but their verification is \emph{linear} in the circuit size;
the deployment amortises it by batch-verifying many proofs in a
single multiexponentiation. The accumulation technique of the Halo
lineage---which would make per-proof verifier work sublinear by
deferring the linear check across recursively composed proofs---is
not deployed: each Orchard proof is verified, in batch, directly.
\end{remark}

\subsection{A shielded spend, end to end}\label{merkle:sec-spend}

The volume closes by assembling the pieces. Every primitive below has
been constructed in a previous section or in this one; the shielded
spend is their composition, and the description follows the deployed
Orchard sources.

\begin{example}[An Orchard shielded spend]\label{merkle:exp-orchard}
Each shielded output of a transaction appends a \emph{note
commitment} \(\mathsf{cm}\)---itself a hiding, binding Sinsemilla
commitment to the note's value, recipient, and randomness
(\S\ref{commitments:subsec:sinsemilla})---as a leaf of an append-only
Merkle tree built with the Sinsemilla compression function of
Remark~\ref{merkle:rem-poseidon}. Strictly, the leaf is the extracted
\(x\)-coordinate
\[
\mathsf{cm}_x = \mathsf{ExtractP}(\mathsf{cm}),
\]
a \(255\)-bit Pallas base-field element (\emph{Math Guide},
\S``Pallas and Vesta assembled'') matching the child encodings of
Remark~\ref{merkle:rem-poseidon}. The protocol
periodically publishes the current root---the \emph{anchor}---on
chain, one per block.

To spend a note, its owner:
\begin{enumerate}
\item locates the leaf \(\mathsf{cm}_x\) in the public tree and reads
  off its authentication path \(\pi\) and position \(i\)
  (\texttt{MerklePath::\allowbreak root}
  recomputes the anchor from exactly these data);
\item chooses a recent anchor \(\mathit{rt} = \mathsf{root}_m\) whose
  tree already contains the leaf---any recorded historical anchor
  serves, per \S\ref{merkle:sec-frontier};
\item proves in zero knowledge, per
  Theorem~\ref{merkle:thm-zk-membership}, the statement
  \(\mathit{rt}\) with witness \((\mathsf{cm}_x, i, \pi)\): membership
  of the note commitment under the anchor, well-formedness of the
  note opening, and correct derivation of a \emph{nullifier}---a
  deterministic tag whose on-chain uniqueness prevents
  double-spending, computed from the note and the spender's nullifier
  key on \(\mathsf{PRF}^{\mathsf{nfOrchard}}\) of
  \S\ref{hash:subsec:poseidon} by a construction of the
  \emph{Ironwood Guide}, \S``The nullifier, and the loop back to
  minting'' (the recomputed root is constrained to
  equal the public anchor, and the instance carries the anchor and
  the nullifier; protocol specification
  \href{https://zips.z.cash/protocol/protocol.pdf\#merklepath}{\S~4.9}, `Merkle Path Validity', and \href{https://zips.z.cash/protocol/protocol.pdf\#actionstatement}{\S~4.18.4}, `Action
  Statement');
\item publishes the proof and the nullifier---but neither
  \(\mathsf{cm}\), nor \(i\), nor \(\pi\).
\end{enumerate}
By Theorem~\ref{merkle:thm-zk-membership}, the spend convinces the
network that the spender owns \emph{some} note genuinely added to the
tree---soundness resting on the collision resistance of Sinsemilla
(Proposition~\ref{commitments:prop:sinsemilla-cr}) and the knowledge
soundness of the proof system---while the network learns nothing
about \emph{which} note. (One refinement of the deployed circuit is
worth recording: the root-equals-anchor constraint is enforced for
notes of nonzero value, while zero-valued dummy spends may witness an
arbitrary path---by design, so that every action carries a spend
whether or not a real note backs it.) The Merkle root is thus a
succinct, binding, privacy-preserving commitment to the set of all
notes ever created.
\end{example}

The spend consumes the whole volume. The note commitment is the
hiding and binding envelope of \S\ref{sec:commitments}; the tree
hash and its personalised domains are the discipline of
\S\ref{sec:hash} instantiated by Sinsemilla; the nullifier is derived
under the spender's nullifier key from a PRF of \S\ref{sec:prf},
instantiated by Poseidon (\S\ref{hash:subsec:poseidon}), by the
construction of the \emph{Ironwood Guide}, \S``The nullifier, and the
loop back to minting''; the proof is a
zk-SNARK of \S\ref{sec:zk} for the relation
\(R_{\mathrm{mem}}\) extended with the note's opening; and the
authorising signatures over the spend---re-randomised RedDSA for
spend authority, the binding signature for value balance---are those
of \S\ref{sec:signatures}. How the pieces sit inside a transaction,
per Action, with the key hierarchy and note-delivery encryption
around them, is the business of the \emph{Ironwood Guide}.

\end{document}
